\documentclass[12pt,twoside,openany]{book}
% BEGIN ../tex/zh.tex
\usepackage[a4paper,inner=26mm,outer=24mm,top=23mm,bottom=25mm,marginparwidth=15mm,marginparsep=6mm]{geometry}
\usepackage{fontspec}
\defaultfontfeatures{Ligatures=TeX}
\setmainfont{TeX Gyre Pagella}
\setsansfont{TeX Gyre Heros}
\setmonofont{TeX Gyre Cursor}
\usepackage{xeCJK}
\setCJKmainfont{SimSun}[AutoFakeBold=2.2]
\setCJKsansfont{Microsoft YaHei}[AutoFakeBold=2.2]
\setCJKmonofont{Microsoft YaHei}
\xeCJKsetup{PunctStyle=kaiming}
\usepackage{microtype}
\usepackage{graphicx}
\usepackage{mathtools}
\usepackage{amsthm}
\usepackage{unicode-math}
\setmathfont{Latin Modern Math}
\usepackage{bookmark}
\usepackage{hyperref}
\hypersetup{hidelinks,unicode=true}

\setcounter{secnumdepth}{3}
\setcounter{tocdepth}{2}
\numberwithin{equation}{section}
\setlength{\parindent}{2em}
\setlength{\parskip}{0pt}
\linespread{1.20}
\renewcommand{\contentsname}{目录}
\renewcommand{\proofname}{证明}

\newcommand{\fgaunitid}{unbound}
\newcommand{\src}[1]{%
  \phantomsection
  \hypertarget{src:\fgaunitid:#1}{}%
  \hypertarget{fga.\fgaunitid.p#1}{}%
  \ifcsname fgasrcseen#1\endcsname\else
    \expandafter\gdef\csname fgasrcseen#1\endcsname{1}%
    \hypertarget{src:#1}{}%
  \fi
  \marginpar{\raggedleft\footnotesize\sffamily #1}%
}
\newcommand{\sid}[1]{%
  \gdef\fgaunitid{#1}%
  \phantomsection\hypertarget{fga:#1}{}\hypertarget{fga.#1}{}%
}
\newcommand{\Osh}{\underline{O}}
\newcommand{\Zsh}{\underline{\mathbf Z}}
\newcommand{\Ccat}{\mathcal C}
\newcommand{\Homsh}{\underline{\operatorname{Hom}}}
\newcommand{\Extsh}{\underline{\operatorname{Ext}}}
\newcommand{\hsh}{\underline{h}}
\newcommand{\fquote}[1]{“#1”}
\newcommand{\zhterm}[1]{\textbf{#1}}
\newcommand{\seminarzh}[1]{%
  \begin{flushleft}\sffamily 布尔巴基讨论班\\#1\end{flushleft}%
}
\newcommand{\paperhead}[2]{%
  \begin{center}
    {\Large\bfseries #1\par}
    \vspace{0.6\baselineskip}
    #2
  \end{center}
  \vspace{0.8\baselineskip}
}

\theoremstyle{plain}
\newtheorem{proposition}{命题}[section]
\newtheorem{theoreme}[proposition]{定理}
\newtheorem{theorem}[proposition]{定理}
\newtheorem{corollaire}[proposition]{推论}
\newtheorem{corollary}[proposition]{推论}
\newtheorem{lemme}[proposition]{引理}
\newtheorem{lemma}[proposition]{引理}

\theoremstyle{definition}
\newtheorem{definition}[proposition]{定义}
\newtheorem{definitionen}[proposition]{定义}

\theoremstyle{remark}
\newtheorem*{remarque}{注}
\newtheorem*{remark}{注}

\newcommand{\unit}[1]{%
  \chapter*{#1}%
  \addcontentsline{toc}{chapter}{#1}%
  \markboth{#1}{#1}%
}
\newcommand{\fsection}[2]{%
  \section*{#1. #2}%
  \addcontentsline{toc}{section}{#1. #2}%
}
\newcommand{\result}[2]{%
  \par\medskip\noindent\textbf{#1%
  \if\relax\detokenize{#2}\relax\else\space#2\fi。——}\enspace
}

% END ../tex/zh.tex

\setlength{\emergencystretch}{1.5em}
\hypersetup{
  pdftitle={代数几何基础——第149、182、190、195、212、221、232、236次报告、1962年评注及勘误（简体中文版）},
  pdfauthor={Alexandre Grothendieck}
}

\begin{document}
\frontmatter
\begin{titlepage}
  \centering
  \vspace*{0.16\textheight}
  {\Huge\bfseries 代数几何基础\par}
  \vspace{1.2\baselineskip}
  {\Large 亚历山大·格罗滕迪克\par}
  \vfill
  {\large 简体中文版（中国大陆用语）\par}
  \vspace{0.5\baselineskip}
  {\normalsize 第149、182、190、195、212、221、232、236次报告、1962年评注及勘误 · 累积工作版\par}
\end{titlepage}
\tableofcontents
\mainmatter
% BEGIN 149.tex
\unit{第149次报告——凝聚代数层的对偶定理}
\sid{149}
\src{169}
\seminarzh{（1957年5月）}
\paperhead{凝聚代数层的对偶定理}%
  {亚历山大·格罗滕迪克 著}

下述结果受塞尔的\fquote{代数对偶定理}启发，发现于1955年冬和1956年冬。借助
关于射影空间上同调的若干相当初等的结果[3]，并大量采用[2]所述形式的
Cartan--Eilenberg同调代数，便可极其简明地证明这些结果。

\fsection{1}{模层的 Ext（[2]，第3、4章）}

设 $X$ 是一个拓扑空间，并赋有一个带幺环（不必交换）的层 $\Osh$。考虑由
$\Osh$-模层构成的阿贝尔范畴 $\Ccat^{\Osh}$；这些对象也称为 $\Osh$-模。已知
该范畴的每个对象都具有内射消解，因而可在其中定义具有熟知形式性质的
$\operatorname{Ext}$ 函子。为避免混淆，更确切地说，我们以
$\operatorname{Hom}_{\Osh}(X;A,B)$，或简记为 $\operatorname{Hom}(X;A,B)$，
表示从 $A$ 到 $B$ 的 $\Osh$-同态所成的阿贝尔群；而以
$\Homsh_{\Osh}(A,B)$ 表示从 $A$ 到 $B$ 的同态芽所组成的层
（$A,B\in\Ccat^{\Osh}$）。固定 $A\in\Ccat^{\Osh}$，定义函子 $h_A$ 和
$\hsh_A$；它们分别取值于阿贝尔群范畴 $\Ccat$ 和 $X$ 上阿贝尔群层的范畴
$\Ccat^{\Zsh}$，其定义式为：
\begin{equation}\tag{1.1}
  h_A(B)=\operatorname{Hom}_{\Osh}(X;A,B),
  \qquad
  \hsh_A(B)=\Homsh_{\Osh}(A,B).
\end{equation}

函子 $h_A$ 和 $\hsh_A$ 都是协变左正合函子。考察它们的右导出函子，分别记为
$\operatorname{Ext}_{\Osh}^{p}(X;A,B)$ 和 $\Extsh_{\Osh}^{p}(A,B)$。因此按定义有
\begin{equation}\tag{1.2}
\left\{
\begin{aligned}
\operatorname{Ext}_{\Osh}^{p}(X;A,B)
  &= (R^{p}h_A)(B)
   = H^{p}\bigl(\operatorname{Hom}_{\Osh}(X;A,C(B))\bigr),\\
\Extsh_{\Osh}^{p}(A,B)
  &= (R^{p}\hsh_A)(B)
   = H^{p}\bigl(\Homsh_{\Osh}(A,C(B))\bigr),
\end{aligned}
\right.
\end{equation}

其中符号 $R^{p}$ 表示取右导出函子，$C(B)$ 表示 $B$ 在 $\Ccat^{\Osh}$ 中的
任意内射消解。以
\[
  \Gamma:\Ccat^{\Zsh}\longrightarrow\Ccat
\]
表示\fquote{截面}函子。回顾其右导出函子记为
\[
  B\longmapsto H^{p}(X,B).
\]

\src{170}
\begin{equation}\tag{1.3}
  H^{p}(X,B)=(R^{p}\Gamma)(B)=H^{p}\bigl(\Gamma(C(B))\bigr).
\end{equation}
显然 $h_A=\Gamma\hsh_A$；另一方面，容易验证 $\hsh_A$ 将内射对象变为
$\Gamma$-无环对象。由熟知的结论可得：

\result{命题}{1} 对每个 $\Osh$-模 $A$，在 $\Ccat^{\Osh}$ 上存在一个上同调
谱函子，它收敛于分次函子
$\bigl(\operatorname{Ext}_{\Osh}^{*}(X;A,B)\bigr)$，且其初始项为
\begin{equation}\tag{1.4}
  E_{2}^{p,q}(A,B)=H^{p}\bigl(X,\Extsh_{\Osh}^{q}(A,B)\bigr).
\end{equation}
由此得到边缘同态和一条五项正合列，兹不写出。

\result{推论}{1} 若 $A$ 局部同构于 $\Osh^{n}$，则有典范同构
\begin{equation}\tag{1.5}
  \operatorname{Ext}_{\Osh}^{p}(X;A,B)
  \overset{\sim}{\longleftarrow}
  H^{p}\bigl(X,\Homsh_{\Osh}(A,B)\bigr)
\end{equation}
（由该谱序列的边缘同态给出）。特别地，有典范同构
\begin{equation}\tag{1.6}
  \operatorname{Ext}_{\Osh}^{p}(X;\Osh,B)=H^{p}(X,B).
\end{equation}

为应用这些结果，必须能够明确描述 $\Extsh_{\Osh}^{p}(A,B)$。这些函子可以
局部计算：若 $U$ 是 $X$ 的开子集，则
\[
  \Extsh_{\Osh}^{p}(A,B)|U
  =\Extsh_{\Osh|U}^{p}(A|U,B|U),
\]
这是因为内射 $\Osh$-模限制到 $U$ 上后仍是内射 $\Osh|U$-模。此外，固定
$x\in X$，有函子性同态
\begin{equation}\tag{1.7}
  \Homsh_{\Osh}(A,B)_x
  \longrightarrow
  \operatorname{Hom}_{\Osh_x}(A_x,B_x),
\end{equation}
它唯一延拓为关于 $B$ 的上同调函子之间的一个同态：
\begin{equation}\tag{1.8}
  \Extsh_{\Osh}^{p}(A,B)_x
  \longrightarrow
  \operatorname{Ext}_{\Osh_x}^{p}(A_x,B_x).
\end{equation}

\result{命题}{2} 若在 $x$ 的某个邻域内，$A$ 同构于某个同态
$\Osh^{m}\longrightarrow\Osh^{n}$ 的余核，则对每个 $p$，(1.8) 都是同构。%
\footnote{印刷文本作“(1.7)”；$p$ 的出现及前一句均要求此处为“(1.8)”。}
特别地，若 $A$ 是凝聚 $\Osh$-模 [3]，上述结论成立。

\src{171}
\result{命题}{3} 设 $L_{*}=(L_i)$ 是 $\Osh$-模 $A$ 的一个左消解，其中每个
$\Osh$-模都局部同构于某个 $\Osh^{n}$。则 $\Extsh_{\Osh}^{*}(A,B)$ 可与
$H^{*}\bigl(\Homsh_{\Osh}(L_{*},B)\bigr)$ 等同，而
$\operatorname{Ext}_{\Osh}^{*}(X;A,B)$ 可与复形
$\Homsh_{\Osh}(L_{*},B)$ 关于 $X$ 的超上同调等同。

证明是标准的。考虑双复形 $\Homsh_{\Osh}(L_{*},C(B))$，其中 $C(B)$ 是 $B$
的一个内射消解，再考虑从 $\Homsh_{\Osh}(L_{*},B)$ 和
$\Homsh_{\Osh}(A,C(B))$ 到此双复形的自然同态即可。

最后指出，(1.2) 中引入的两个 $\operatorname{Ext}$ 函子不仅是关于 $B$ 的
上同调函子，而且是上同调双函子：关于 $B$ 协变，关于 $A$ 反变。

\fsection{2}{Ext 中的合成律}

本节结果由 CARTIER 和 YONEDA 各自独立得到；详情见 CARTIER 的报告 [1]。
设 $\Ccat$ 为阿贝尔范畴，$K,L$ 为 $\Ccat$ 中两个分次对象。以
$\operatorname{Hom}(K,L)$ 表示如下分次阿贝尔群：其 $n$ 次分量由从 $K$ 到
$L$ 的 $n$ 次齐次同态组成（即同态 $K^{i}\longrightarrow L^{i+n}$ 所成的
系统 $(u_i)$）。

若 $K,L$ 是复形（为明确起见，设其微分算子的次数为 $+1$），则在
$\operatorname{Hom}(K,L)$ 上引入微分算子
\begin{equation}\tag{2.1}
  \delta(u)=du+(-1)^{n+1}ud,\quad
  \mbox{\upshape\fontspec{SimSun}其中 }n=\deg(u),
\end{equation}
从而使它成为微分次数为 $+1$ 的复形。$n$ 次闭元就是与 $d$ 反交换的 $n$ 次
算子（视为齐次算子）。\footnote{印刷文本作“与 $u$”；方程 (2.1) 要求此处
为“与 $d$”。} 因而可以考虑 $H^{*}\bigl(\operatorname{Hom}(K,L)\bigr)$；
它是 $K,L$ 同伦型的不变量，可记为 $H^{*}(K,L)$。

若还有第三个复形 $M$，则同态的合成定义一个配对
\[
  \operatorname{Hom}(K,L)\times\operatorname{Hom}(L,M)
  \longrightarrow\operatorname{Hom}(K,M)
\]
且与微分算子相容，因此取上同调后得到配对
\begin{equation}\tag{2.2}
  H^{*}(K,L)\times H^{*}(L,M)\longrightarrow H^{*}(K,M),
\end{equation}
记作 $(u,v)\longmapsto vu$。这些配对显然满足结合律。特别地，
$H^{*}(K,K)$ 是带幺结合分次环，而 $H^{*}(K,L)$（相应地，$H^{*}(L,K)$）
是该环上的分次右模（相应地，分次左模），等等。在 $0$ 次时，(2.2) 化为
复形同态的合成。最后，若有复形的正合列
\[
  0\longrightarrow K'\longrightarrow K\longrightarrow K''\longrightarrow 0,
\]

\src{172}
并且每个 $K'^{i}$ 都可与 $K^{i}$ 的一个直和因子等同，则它诱导群复形
$\operatorname{Hom}(K'',L)$ 等的正合列，进而给出连接算子
\[
  H^{i}(K',L)\longrightarrow H^{i+1}(K'',L).
\]
相对于 $L$ 中一条正合列的连接算子也用同样方式定义。配对 (2.2) 按通常意义
与这些连接算子相容。

现在设 $\Ccat$ 是一个范畴，其中每个对象 $A$ 都有内射消解 $C(A)$。运用
双复形定理的诸多变式之一可知
\[
  H^{*}\bigl(C(A),C(B)\bigr)
  =H^{*}\bigl(\operatorname{Hom}(C(A),C(B))\bigr)
\]
典范同构于
\[
  H^{*}\bigl(\operatorname{Hom}(A,C(B))\bigr)
  =\operatorname{Ext}^{*}(A,B).
\]
上面考虑的连接算子给出 $\operatorname{Ext}$ 的连接同态。此外，配对 (2.2)
在此给出与连接同态相容的结合配对：
\begin{equation}\tag{2.3}
  \operatorname{Ext}^{*}(A,B)\times\operatorname{Ext}^{*}(B,C)
  \longrightarrow\operatorname{Ext}^{*}(A,C).
\end{equation}
特别地，$\operatorname{Ext}^{*}(A,A)$ 是带幺结合分次环，等等。（类似地可证，
$\operatorname{Ext}$ 函子作用于任意函子的导出函子；此处不使用这一事实。）

若所考虑的范畴是 $X$ 上 $\Osh$-模的范畴 $\Ccat^{\Osh}$，便得到配对
\begin{equation}\tag{2.4}
  \operatorname{Ext}_{\Osh}^{p}(X;A,B)
  \times\operatorname{Ext}_{\Osh}^{q}(X;B,C)
  \longrightarrow\operatorname{Ext}_{\Osh}^{p+q}(X;A,C),
\end{equation}
它们可以按前述方式明确写出。另一方面，作同样的推导，但把阿贝尔群范畴替换为
$X$ 上阿贝尔群层的范畴，并把 $\operatorname{Hom}$ 函子替换为 $\Homsh$
函子，还可定义具有相同形式性质、而且这一次具有\fquote{局部性质}的配对：
\begin{equation}\tag{2.5}
  \Extsh_{\Osh}^{p}(A,B)\times\Extsh_{\Osh}^{q}(B,C)
  \longrightarrow\Extsh_{\Osh}^{p+q}(A,C).
\end{equation}
再注意同态 (1.8) 与 $\operatorname{Ext}$ 之间的配对相容，便可进一步明确
这些配对。

\src{173}
最后回顾，函子 $H^{p}(X,A)$ 之间也有乘法结构（杯积）。由此可见，命题 1
中的谱序列与乘法结构相容。更确切地说，谱序列 $E(A,B)$ 与谱序列 $E(B,C)$
有一个取值于谱序列 $E(A,C)$ 的配对；在收敛项上，它化为整体
$\operatorname{Ext}$ 群之间的配对，而在初始项上，它化为 (1.4) 右端中由
杯积和局部 $\operatorname{Ext}$ 群配对所导出的配对。特别地，由此可知边缘同态
\begin{equation}\tag{2.6}
  \operatorname{Ext}_{\Osh}^{n}(X;A,B)
  \longrightarrow H^{0}\bigl(X,\Extsh_{\Osh}^{n}(A,B)\bigr)
\end{equation}
和
\begin{equation}\tag{2.7}
  H^{n}\bigl(X,\Homsh_{\Osh}(A,B)\bigr)
  \longrightarrow\operatorname{Ext}_{\Osh}^{n}(X;A,B)
\end{equation}
均与乘法结构相容。因此，若只考虑局部同构于某个 $\Osh^{m}$ 的层，则结合
同构 (1.5)，即可完全用杯积明确描述整体 $\operatorname{Ext}$ 之间的合成。

\fsection{3}{局部上同调的结果}

设 $A$ 是一个带幺交换环，并赋有一个理想 $J$。对变动的 $A$-模 $M$，我们将
定义函子性同态
\begin{equation}\tag{3.1}
\left\{
\begin{aligned}
\operatorname{Ext}_{A}^{p}(A/J,M)
  &\longrightarrow
  \operatorname{Hom}_{A}\!\left(\bigwedge^{p}J/J^{2},M\otimes A/J\right),\\
\operatorname{Tor}_{p}^{A}(A/J,M)
  &\longleftarrow
  \left(\bigwedge^{p}J/J^{2}\right)
  \otimes\operatorname{Hom}_{A}(A/J,M),
\end{aligned}
\right.
\end{equation}
其中张量积和外积都在环 $A$ 上取。还应注意，$J/J^{2}$ 实际上是一个
$A/J$-模，而它作为 $A$-模或作为 $A/J$-模所取的外幂相同。定义同态 (3.1)，
等价于对 $J$ 中元素的每个系统 $\underline{x}=(x_1,\ldots,x_p)$ 定义同态
$\varphi_{\underline{x}}$：
\begin{equation}\tag{3.2}
\left\{
\begin{aligned}
\varphi_{\underline{x}}:
  \operatorname{Ext}_{A}^{p}(A/J,M)&\longrightarrow M\otimes A/J,\\
\varphi_{\underline{x}}:
  \operatorname{Hom}_{A}(A/J,M)&\longrightarrow
  \operatorname{Tor}_{p}^{A}(A/J,M),
\end{aligned}
\right.
\end{equation}
并使其满足下列条件：

\src{174}
\begin{enumerate}
\item[\textup{i.}] $\varphi_{x_1,\ldots,x_p}$ 按交替 $A$-多线性的方式依赖于
$x_i\in J$ 所成的系统。
\item[\textup{ii.}] 只要有一个 $x_i$ 属于 $J^{2}$，就有
$\varphi_{x_1,\ldots,x_p}=0$。
\end{enumerate}
事实上，ii 由 i 推出，因为对 $a\in J$ 有 $a\varphi_{\underline{x}}=0$；只需
注意 (3.2) 中所有的模都被 $J$ 零化。

为定义 $\varphi_{\underline{x}}$，考虑复形 $K_{\underline{x}}$：其底层
$A$-模是 $\bigwedge A^{p}$，其微分算子是内乘算子
$i_{\underline{x}}$，这里将 $\underline{x}$ 视为 $A^{p}$ 上分量为
$x_1,\ldots,x_p$ 的线性形式。微分算子的次数为 $-1$，各次数非负，而该复形
在 $0$ 维的同调为 $A/(x_1A+\cdots+x_pA)$。由于 $x_i\in J$，便得到增广
$K_{\underline{x},0}\longrightarrow A/J$。因此 $K_{\underline{x}}$ 是一个
以 $A/J$ 为增广模的增广自由复形。按熟知的方式，由此得到同态
\[
\operatorname{Ext}_{A}^{*}\bigl(H_0(K_{\underline{x}}),M\bigr)
\longrightarrow
H^{*}\bigl(\operatorname{Hom}_{A}(K_{\underline{x}},M)\bigr)
\]
以及
\[
\operatorname{Tor}_{*}^{A}\bigl(H_0(K_{\underline{x}}),M\bigr)
\longleftarrow
H_{*}\bigl(K_{\underline{x}}\otimes_A M\bigr).
\]
再与增广同态 $H_0(K_{\underline{x}})\longrightarrow A/J$ 所诱导的
$\operatorname{Ext}$ 和 $\operatorname{Tor}$ 上的同态合成，得到同态
\begin{equation}\tag{3.3}
\left\{
\begin{aligned}
\psi_{\underline{x}}:
\operatorname{Ext}_{A}^{*}(A/J,M)
&\longrightarrow H^{*}\bigl(\operatorname{Hom}_{A}(K_{\underline{x}},M)\bigr),\\
\psi_{\underline{x}}:
\operatorname{Tor}_{*}^{A}(A/J,M)
&\longleftarrow H_{*}\bigl(K_{\underline{x}}\otimes_A M\bigr).
\end{aligned}
\right.
\end{equation}
立即可见，在最高维数 $p$，右端的同调是
$M/(x_1M+\cdots+x_pM)$（相应地，是 $M$ 中被每个 $x_i$ 零化的元素所成的
集合）。由于 $x_i\in J$，便得到同态
\begin{equation}\tag{3.4}
\left\{
\begin{aligned}
H^{p}\bigl(\operatorname{Hom}_{A}(K_{\underline{x}},M)\bigr)
&\longrightarrow M\otimes A/J,\\
H_{p}\bigl(K_{\underline{x}}\otimes_A M\bigr)
&\longleftarrow\operatorname{Hom}_{A}(A/J,M).
\end{aligned}
\right.
\end{equation}
将同态 (3.3) 与 (3.4) 合成，便得到所要定义的同态 (3.2)。条件 i 的验证
虽然繁琐，却没有困难。

\src{175}
\result{命题}{4} 设 $A$ 是带幺交换环，$(x_1,\ldots,x_p)$ 是 $A$ 中的一个
元素序列，并且对 $1\leq i\leq p$，$x_i$ 在 $A$ 对由
$(x_1,\ldots,x_{i-1})$ 生成的理想所得商环中的像不是零因子。令 $J$ 为
$x_i$ 生成的理想。则 $J/J^{2}$ 是自由 $A/J$-模，其基为 $x_i$ 的典范像；
复形 $K_{\underline{x}}$ 是 $A/J$ 的一个自由消解；而且对任意 $A$-模 $M$，
(3.1) 中次数为 $p$ 的同态都是双射。只要 $J M=0$，在任意次数 $i$ 定义的
类似同态也具有相同结论。

（关键之处是 $K_{\underline{x}}$ 的无环性；其余结论都由此推出。在上述条件下，
这是熟知的事实。）

\result{推论}{1} 设 $A,J$ 如上，并进一步假设 $A$ 是完美域 $k$ 上维数为
$n$ 的正则仿射代数，且 $A/J$ 也是正则仿射代数。以 $\Omega^{1}(A)$、
$\Omega^{1}(A/J)$ 表示 Kähler 微分模。则有与局部化相容的典范同构
\begin{equation}\tag{3.5}
  \operatorname{Ext}_{A}^{p}\bigl(\Omega^{n-p}(A/J),\Omega^{n}(A)\bigr)=A/J.
\end{equation}
事实上，$\Omega^{n-p}(A/J)$ 是秩为 $1$ 的自由 $A/J$-模，同样，
$\Omega^{n}(A)$ 是秩为 $1$ 的自由 $A$-模，\footnote{印刷文本作“秩 $n$”；
两个最高次形式模的秩均为 $1$，下文与 $\bigwedge^{p}(J/J^{2})$ 的等同也要求
如此。} 因而左端等于
\[
  \operatorname{Ext}_{A}^{p}(A/J,A)
  \otimes\Omega^{n-p}(A/J)'\otimes\Omega^{n}(A)
\]
（其中符号“$'$”表示对偶 $A/J$-模）。后两个因子的张量积可与
$\bigwedge^{p}(J/J^{2})$ 等同，故整个表达式可与
\[
  \operatorname{Ext}_{A}^{p}\!\left(A/J,\bigwedge^{p}(J/J^{2})\right)
\]
等同，再由命题与
\[
  \operatorname{Hom}_{A}\!\left(\bigwedge^{p}J/J^{2},
  \bigwedge^{p}J/J^{2}\right)
\]
等同，也就是与 $A/J$ 等同。特别地，
$\operatorname{Ext}_{A}^{p}\bigl(\Omega^{n-p}(A/J),\Omega^{n}(A)\bigr)$ 中
有一个对应于 $A/J$ 单位元的特出元素，称为\zhterm{理想 $J$ 在 $A$ 中的基本类}。
（事实上，它可以在宽泛得多的条件下定义。）推论 1 可以表为更几何、更整体的
形式：

\result{推论}{2} 设 $X$ 是代数闭域 $k$ 上的非奇异簇，$Y$ 是 $X$ 的闭
非奇异子簇；$\Osh_X$ 是 $X$ 的局部环层，$\Osh_Y$ 是 $Y$ 的局部环层，并将
它视为 $\Osh_X$ 的商层。设 $X$ 的维数为 $n$，$Y$ 的维数为 $n-p$。

\src{176}
令 $\Omega_X$（相应地，$\Omega_Y$）为 $X$（相应地，$Y$）上正则微分形式
的芽层。则有典范同构
\begin{equation}\tag{3.6}
  \Extsh_{\Osh_X}^{p}\bigl(\Omega_Y^{n-p},\Omega_X^{n}\bigr)=\Osh_Y,
\end{equation}
或者等价地，
\begin{equation}\tag{3.6 bis}
  \Extsh_{\Osh_X}^{p}\bigl(\Osh_Y,\Omega_X^{n}\bigr)=\Omega_Y^{n-p}.
\end{equation}
当 $Y$ 是奇异簇时，公式 (3.6 bis) 可用作 $\Omega_Y^{n-p}$ 的定义。更确切地说：

\result{命题}{5} 设 $Y$ 是维数为 $q=n-p$ 的代数子集，包含在维数为 $n$ 的
非奇异代数簇 $X$ 中。设 $F$ 是 $X$ 上支集包含于 $Y$ 的凝聚代数层，$L$ 是
$X$ 上的局部自由代数层。则当 $i<p$ 时，层 $\Extsh_{\Osh_X}^{i}(F,L)$ 为零；
而当 $i=p$ 时，有典范同构
\begin{equation}\tag{3.7}
  \Extsh_{\Osh_X}^{p}(F,L)
  =\Homsh_{\Osh_X}\!\left(F,
  \Extsh_{\Osh_X}^{p}(\Osh_X/J,L)\right),
\end{equation}
其中 $J$ 表示 $X$ 上任意一个零化 $F$、且零点集为 $Y$ 的理想层。特别地，
若 $F$ 是 $Y$ 上的凝聚代数层，则
\begin{equation}\tag{3.7 bis}
  \Extsh_{\Osh_X}^{p}(F,L)
  =\Homsh_{\Osh_Y}\!\left(F,\Extsh_{\Osh_X}^{p}(\Osh_Y,L)\right).
\end{equation}
最后，仍设 $F$ 是 $Y$ 上的凝聚代数层，则诸层
\[
  E^{i}(F)=\Extsh_{\Osh_X}^{p+i}(F,\Omega_X^{n})
\]
不依赖于把代数空间 $Y$ 嵌入非奇异簇 $X$ 的方式。

问题是局部的，故可设 $X$ 仿射且 $L=\Osh_X$。这把问题归结为交换代数问题，
甚至局部代数问题：若 $A$ 是正则局部环，$M$ 是支集维数
$\leq q=n-p$ 的 $A$-模，则要证明当 $i<p$ 时
$\operatorname{Ext}_{A}^{i}(M,A)=0$，并且
\[
  \operatorname{Ext}_{A}^{p}(M,A)
  =\operatorname{Hom}_{A}\!\left(M,\operatorname{Ext}_{A}^{p}(A/J,A)\right),
\]
其中 $J$ 是任意一个零化 $M$ 且\fquote{维数} $\leq q$ 的理想。对于第一点，
对 $q$ 作归纳：一次直接的分解化简把问题归约到 $M=A/J$ 的情形；随后把 $J$
换成一个更小的理想，并运用归纳假设和 $\operatorname{Ext}$ 的正合列，又可归约
到 $J$ 由命题 4 中那样的一个\fquote{参数系}生成的情形，此时结论立即成立。
前一结果蕴含：若 $J$ 是一个理想

\src{177}
固定且\fquote{维数} $\leq q$，则 $A/J$-模范畴上的反变函子
$E(M)=\operatorname{Ext}_{A}^{p}(M,A)$ 是左正合的；它还把直和变为直积，
由此容易得到
\[
  E(M)=\operatorname{Hom}_{A}\bigl(M,E(A/J)\bigr).
\]
\footnote{印刷文本作 $E(A)$；在所考虑的范畴中，表示对象是 $A/J$，而 (3.7)
所需的等式为 $E(A/J)$。} 最后，命题 5 的最后一个断言更为微妙；它由
$E^{i}(F)$ 的一个内蕴刻画推出，而该刻画使用了此处无法陈述的局部对偶定理。

\result{推论}{} 以 $\omega_Y^{q}$ 表示层
$\Extsh_{\Osh_X}^{p}(\Osh_Y,\Omega_X^{n})$。则对 $Y$ 上的凝聚代数层 $F$，
有函子性同构
\begin{equation}\tag{3.8}
  \Extsh_{\Osh_X}^{p}(F,\Omega_X^{n})
  =\Homsh_{\Osh_X}(F,\omega_Y^{q}).
\end{equation}

\fsection{4}{与子簇相伴的上同调类}

下文始终以 $X$ 表示定义在域 $k$ 上的维数为 $n$ 的代数集；为简便起见，假设
$k$ 代数闭。除第 6 节外，均假设 $X$ 非奇异。以 $\Osh_X$ 表示 $X$ 上的局部
环层，以
\[
  \Omega_X^{*}=\bigoplus_p\Omega_X^{p}
\]
表示 $X$ 上微分形式的芽层。\footnote{印刷文本使用并集号；这里的分次结构是
诸层 $\Omega_X^p$ 的直和。} 若 $Y$ 是 $X$ 的闭子集，就把 $Y$ 上的凝聚代数层
视同为 $X$ 上在 $Y$ 外为零的凝聚层；$\Osh_Y$ 和 $\Omega_Y$ 尤其如此。

\result{引理}{1} 设 $F$ 是 $X$ 上支集维数 $\leq n-p$ 的凝聚代数层，$L$
是 $X$ 上的局部自由凝聚代数层。则当 $i<p$ 时，
$\operatorname{Ext}_{\Osh_X}^{i}(X;F,L)$ 为零，并有典范同构
\begin{equation}\tag{4.1}
  \operatorname{Ext}_{\Osh_X}^{p}(X;F,L)
  =H^{0}\bigl(X,\Extsh_{\Osh_X}^{p}(F,L)\bigr).
\end{equation}
若 $F$ 是 $X$ 的闭子集 $Y$ 上的凝聚代数层，且 $Y$ 的维数
$\leq n-p$，则有典范同构
\begin{equation}\tag{4.1 bis}
  \Extsh_{\Osh_X}^{p}(F,L)
  =\Homsh_{\Osh_X}\!\left(
    F\otimes L'\otimes\Omega_X^{n},\omega_Y^{n-p}
  \right),
\end{equation}
其中 $\omega_Y^{n-p}$ 是命题 5 之后的推论所定义的 $Y$ 上的层（若 $Y$ 非奇异，
它可与 $\Omega_Y^{n-p}$ 等同）。

公式 (4.1) 直接由命题 1 的谱序列和命题 5 推出；根据公式 (3.8)，可以写成

\src{178}
\begin{align*}
\Extsh_{\Osh_X}^{p}(F,L)
&=L\otimes(\Omega_X^{n})'\otimes
  \Extsh_{\Osh_X}^{p}(F,\Omega_X^{n})\\
&=L\otimes(\Omega_X^{n})'\otimes
  \Homsh_{\Osh_X}(F,\omega_Y^{q})\\
&=\Homsh_{\Osh_X}\!\left(
  F\otimes L'\otimes\Omega_X^{n},\omega_Y^{q}\right),
\end{align*}
其中 $q=n-p$，由此得到公式 (4.1 bis)。

特别取 $F=\Osh_Y$、$L=\Omega_X^{p}$，并注意到
$\Omega_X^{n}\otimes(\Omega_X^{p})'=\Omega_X^{n-p}$，便得到典范同构
\begin{equation}\tag{4.2}
  \operatorname{Ext}_{\Osh_X}^{p}(X;\Osh_Y,\Omega_X^{p})
  =\operatorname{Hom}_{\Osh_X}
  (X;\Omega_X^{n-p},\omega_Y^{n-p}).
\end{equation}
为简便起见，现在假设 $Y$ 非奇异，因而
$\omega_Y^{n-p}=\Omega_Y^{n-p}$。从 $\Omega_X^{n-p}$ 到
$\Omega_Y^{n-p}$ 有自然同态，故得到层
$\Extsh_{\Osh_X}^{p}(\Osh_Y,\Omega_X^{p})$ 的典范截面 $s_Y$。若 $Y$ 的
所有分支维数均为 $n-p$，就称它为层
$\Extsh_{\Osh_X}^{p}(\Osh_Y,\Omega_X^{p})$ 的\zhterm{基本截面}。由 (4.1)，
该截面定义 $\operatorname{Ext}_{\Osh_X}^{p}(X;\Osh_Y,\Omega_X^{p})$ 中的
一个元素。另一方面，自然同态 $\Osh_X\longrightarrow\Osh_Y$ 定义同态
\[
  \operatorname{Ext}_{\Osh_X}^{p}(X;\Osh_Y,\Omega_X^{p})
  \longrightarrow
  \operatorname{Ext}_{\Osh_X}^{p}(X;\Osh_X,\Omega_X^{p})
  =H^{p}(X,\Omega_X^{p}).
\]
这样便得到 $H^{p}(X,\Omega_X^{p})$ 中的一个元素，记作 $P_X(Y)$，称为
\zhterm{$Y$ 在 $X$ 中的上同调类}。因此，它由
$\Extsh_{\Osh_X}^{p}(\Osh_Y,\Omega_X^{p})$ 的截面 $s_Y$ 通过下列同态图式
得到：
\begin{equation}\tag{4.3}
\begin{aligned}
H^{p}(X,\Omega_X^{p})
&=\operatorname{Ext}_{\Osh_X}^{p}(X;\Osh_X,\Omega_X^{p})\\
&\longleftarrow\operatorname{Ext}_{\Osh_X}^{p}(X;\Osh_Y,\Omega_X^{p})
\xrightarrow{\ \sim\ }
H^{0}\!\left(X,\Extsh_{\Osh_X}^{p}(\Osh_Y,\Omega_X^{p})\right).
\end{aligned}
\end{equation}

所谓\zhterm{维数为 $n-p$ 的非奇异循环}，是指由 $X$ 中维数为 $n-p$ 的不可约
非奇异子簇生成的自由阿贝尔群中的元素。于是函数 $Y\longmapsto P(Y)$ 延拓为
从 $X$ 上维数为 $n-p$ 的非奇异循环群到群 $H^{p}(X,\Omega_X^{p})$ 的同态。

设 $Z^{n-p}$ 和 $Z'^{n-p'}$ 是维数分别为 $n-p$ 和 $n-p'$ 的非奇异循环。
若 $Z$ 的每个分支都与 $Z'$ 的每个分支横截相交，就称它们\zhterm{横截相交}。
此时循环 $Z.Z'$ 有定义，并且是维数为 $n-p-p'$ 的非奇异循环。在这些约定下，
有如下结论。

\src{179}
\result{定理}{1} 若 $Z^{n-p}$ 和 $Z'^{n-p'}$ 是横截相交的非奇异循环，则
\begin{equation}\tag{4.4}
  P_X(Z.Z')=P_X(Z).P_X(Z'),
\end{equation}
其中右端的乘法是杯积
\[
  H^{p}(X,\Omega_X^{p})\times H^{p'}(X,\Omega_X^{p'})
  \longrightarrow H^{p+p'}(X,\Omega_X^{p+p'}).
\]
（假设 $X$ 同构于某个射影空间的一个局部闭子集。）

最后这个假设仅用于保证 $X$ 上每个凝聚代数层都是某个局部自由凝聚代数层的
商（SERRE），因而具有由局部自由层组成的左消解。为证明定理 1，可以假设
$Z,Z'$ 是横截相交的不可约非奇异子簇 $Y,Y'$。设 $L_{*}$ 是 $\Osh_Y$ 的一个
由局部自由层组成的左消解；则根据命题 3，同态图式 (4.3) 可与下图等同：
\[
R^{p}\Gamma(\Omega_X^{p})
\xleftarrow{\alpha}
\mathbb R^{p}\Gamma\!\left(\Homsh_{\Osh_X}(L_{*},\Omega_X^{p})\right)
\xrightarrow[\sim]{\beta}
\Gamma\!\left(H^{p}\bigl(\Homsh_{\Osh_X}(L_{*},\Omega_X^{p})\bigr)\right).
\]
这里 $\Gamma$ 是 $X$ 上阿贝尔群层范畴的\fquote{截面群}函子，
$\mathbb R^{p}\Gamma$ 表示该函子的 $p$ 维超上同调，而 $R^{p}\Gamma$ 表示其
第 $p$ 个导出函子。为简便起见，可设 $L_0=\Osh_X$，并设增广
$L_0\longrightarrow\Osh_Y$ 为自然同态（这是允许的）。此时 $\alpha$ 由复形
同态 $\Osh_X\longrightarrow L_{*}$ 诱导（把 $\Osh_X$ 视为集中在次数 $0$ 的
复形）；这里用到：若 $K$ 是集中在次数 $0$ 的层复形，则
$\mathbb R^{p}\Gamma(K)=R^{p}\Gamma(K_0)$。同态 $\beta$ 是熟知的边缘同态。
对 $\Osh_{Y'}$ 的一个局部自由消解 $L'_{*}$ 考虑类似图式，再考虑下列配对的
交换图：
\begin{equation}\tag{4.5}
\resizebox{\textwidth}{!}{$
\begin{array}{ccc}
R^{p}\Gamma(\Omega_X^{p})
&\longleftarrow
\mathbb R^{p}\Gamma\!\left(\Homsh_{\Osh_X}(L_{*},\Omega_X^{p})\right)
&\xrightarrow{\sim}
\Gamma H^{p}\!\left(\Homsh_{\Osh_X}(L_{*},\Omega_X^{p})\right)
\\[3pt]
\times&\times&\times\\[3pt]
R^{p'}\Gamma(\Omega_X^{p'})
&\longleftarrow
\mathbb R^{p'}\Gamma\!\left(\Homsh_{\Osh_X}(L'_{*},\Omega_X^{p'})\right)
&\xrightarrow{\sim}
\Gamma H^{p'}\!\left(\Homsh_{\Osh_X}(L'_{*},\Omega_X^{p'})\right)
\\[3pt]
\downarrow&\downarrow&\downarrow\\[3pt]
R^{p+p'}\Gamma(\Omega_X^{p+p'})
&\longleftarrow
\mathbb R^{p+p'}\Gamma\!\left(
\Homsh_{\Osh_X}(L_{*}\otimes L'_{*},\Omega_X^{p+p'})\right)
&\xrightarrow{\sim}
\Gamma H^{p+p'}\!\left(
\Homsh_{\Osh_X}(L_{*}\otimes L'_{*},\Omega_X^{p+p'})\right).
\end{array}
$}
\end{equation}
\src{180}
(4.5) 右侧两列中的配对由下列层复形的配对诱导：
\[
\Homsh_{\Osh_X}(L_{*},\Omega_X^{p})
\times\Homsh_{\Osh_X}(L'_{*},\Omega_X^{p'})
\longrightarrow
\Homsh_{\Osh_X}(L_{*}\otimes L'_{*},\Omega_X^{p+p'}),
\]
该配对借助外积
$\Omega_X^{p}\times\Omega_X^{p'}\longrightarrow\Omega_X^{p+p'}$ 定义。
第一列中的配对是（相对于此外积的）杯积。我们断言，(4.5) 的最后一行可与
类似于 (4.3) 的同构图式等同，其中以 $Y\cap Y'$ 代替 $Y$，以 $p+p'$ 代替
$p$。为此只需证明 $L\otimes L'$ 是 $\Osh_{Y\cap Y'}$ 的一个消解（它显然
局部自由）。事实上，
\[
  H_0(L\otimes L')=\Osh_Y\otimes\Osh_{Y'}=\Osh_{Y\cap Y'}
\]
且
\[
  H_i(L\otimes L')=\operatorname{Tor}_{i}^{\Osh_X}(\Osh_Y,\Osh_{Y'})=0
  \qquad(i>0),
\]
因为 $Y,Y'$ 横截相交。现在定理 1 由下式推出：
\begin{equation}\tag{4.6}
  s_Y.s_{Y'}=s_{Y.Y'}
\end{equation}
（左端的乘法是 (4.5) 最后一列中的乘法）。公式 (4.6) 具有纯局部性质；对
$L_{*},L'_{*}$ 取命题 4 中所考虑的消解，即可毫无困难地验证它。类似地（而且
证明更容易）可证 $Z\longmapsto P_X(Z)$ 与笛卡儿积相容：
\begin{equation}\tag{4.7}
  P_{X\times X'}(Z\times Z')=P_X(Z)\otimes P_{X'}(Z')
\end{equation}
（当 $Z$、相应地 $Z'$ 是非奇异簇 $X$、相应地 $X'$ 上的非奇异循环时，该式
成立；这里把 $Z\times Z'$ 视为 $X\times X'$ 上的非奇异循环）。由 (4.4)
和 (4.7) 可知，$P_X(Z)$ 还与非奇异簇之间的态射
$f:X\longrightarrow X'$ 所给出的\fquote{逆像}运算相容：
\begin{equation}\tag{4.8}
  P_X\bigl(f^{-1}(Z')\bigr)=f^{*}\bigl(P_{X'}(Z')\bigr),
\end{equation}
当 $Z'$ 是 $X'$ 上的非奇异循环且 $f$ 与 $Z'$\fquote{横截}时，该式成立；
也就是说，$f$ 的图像在 $X\times X'$ 中与循环 $X\times Z'$ 横截。

\src{181}
\result{推论}{1} 设 $X,X'$ 是某个射影空间中的两个局部闭非奇异簇，并设
$X'$ 完备。设 $U$ 是 $X\times X'$ 上的非奇异循环，$a,b$ 是 $X'$ 中的两个
点，并且 $U$ 与循环 $X\times(a)$ 和 $X\times(b)$ 都横截相交。设 $Z,Z'$ 是
$X$ 上满足
\[
  Z\times(a)=(X\times(a)).U,
  \qquad
  Z'\times(b)=(X\times(b)).U
\]
的非奇异循环。则
\[
  P_X(Z)=P_X(Z').
\]

事实上，令 $f_a:X\longrightarrow X\times X'$ 由 $f_a(x)=(x,a)$ 定义，
根据 (4.8) 有 $P(Z)=f_a^{*}(P(U))$；同理，$P(Z')=f_b^{*}(P(U))$。利用
Künneth 公式
\[
  H^{*}(X\times X',\Omega_{X\times X'}^{*})
  =H^{*}(X,\Omega_X^{*})\otimes H^{*}(X',\Omega_{X'}^{*})
\]
以及 $H^{0}(X',\Osh_{X'})$ 仅由标量组成这一事实，容易得到
$f_a^{*}=f_b^{*}$，从而得证。

对每个 $x\in X$，点 $(x)$ 都是 $X$ 中余维为 $n$ 的非奇异子簇，因而定义
$H^{n}(X,\Omega_X^{n})$ 中的一个元素 $\varepsilon_x$。若 $X$ 是非奇异射影簇，
由推论 1 可知 $\varepsilon_x$ 不依赖于所选的点 $x$；将其记为
$\varepsilon_X$，称为 $H^{n}(X,\Omega_X^{n})$ 的\zhterm{基本类}。

\result{注}{} 为得到令人满意的理论，应当对任意循环 $Z$ 定义 $P_X(Z)$，并对
循环的正常相交证明定理 1。（撰写本报告时，这一点尚未在完全一般的情形中完成。）
假定这些工作已经完成，推论 1 就变成：若 $Z,Z'$ 是两个代数等价的循环，则
$P_X(Z)=P_X(Z')$（即使 $Z,Z'$ 非奇异，这一陈述似乎也不能由前述结果推出）。

\fsection{5}{对偶定理}

本节中，$X$ 表示维数为 $n$ 的非奇异射影簇。

\result{定理}{2} $H^{n}(X,\Omega_X^{n})$ 的基本类 $\varepsilon_X$ 是该向量
空间的一组基。

（证明将在下文给出。）借助前一定理，可以把 $H^{n}(X,\Omega_X^{n})$ 与域
$k$ 等同。现在考虑第 2 节中引入的配对；它特别给出配对

\src{182}
\[
\operatorname{Ext}_{\Osh_X}^{p}(X;\Osh_X,F)
\times\operatorname{Ext}_{\Osh_X}^{n-p}(X;F,\Omega_X^{n})
\longrightarrow
\operatorname{Ext}_{\Osh_X}^{n}(X;\Osh_X,\Omega_X^{n}),
\]
即
\begin{equation}\tag{5.1}
  H^{p}(X,F)
  \times\operatorname{Ext}_{\Osh_X}^{n-p}(X;F,\Omega_X^{n})
  \longrightarrow H^{n}(X,\Omega_X^{n}).
\end{equation}
考虑到定理 2，该配对定义同态
\begin{equation}\tag{5.2}
  \operatorname{Ext}_{\Osh_X}^{n-p}(X;F,\Omega_X^{n})
  \longrightarrow\bigl(H^{p}(X,F)\bigr)'.
\end{equation}
该同态关于 $F$ 具有函子性，并与正合列
$0\longrightarrow F'\longrightarrow F\longrightarrow F''\longrightarrow0$
所对应的连接同态可交换。

\result{定理}{3} 同态 (5.2) 是同构。

特别地，由此重新得到 Serre 的结果：

\result{推论}{} 设 $E$ 是 $X$ 上的代数向量丛，$\Osh_X(E)$ 是 $E$ 的正则
截面芽层。则有典范同构
\begin{equation}\tag{5.3}
  \bigl(H^{p}(X,\Osh_X(E))\bigr)'
  =H^{n-p}\bigl(X,\Omega_X^{n}\otimes\Osh_X(E')\bigr).
\end{equation}
只需把定理 3 和命题 1 的推论 1 结合使用。

定理 2 和定理 3 将由下述命题推出：

\noindent\textbf{(D) 下列同态}
\begin{equation}\tag{5.2 bis}
  \operatorname{Ext}_{\Osh_X}^{n-p}(X;F,\Omega_X^{n})
  \longrightarrow\bigl(H^{p}(X,F)\bigr)'\otimes L,
  \qquad L=H^{n}(X,\Omega_X^{n}),
\end{equation}
由配对 (5.1) 诱导，并且是同构。

先证明 (D) 蕴含定理 2。令 $k_x=\Osh_{(x)}$ 为缩约到点 $x\in X$ 的簇的
局部环层。考虑典范同态 $\Osh_X\longrightarrow k_x$ 以及相伴的同态
\begin{equation}\tag{5.3 bis}
  H^{0}(X,\Osh_X)\longrightarrow H^{0}(X,k_x),
\end{equation}
其转置可与下列同态等同：
\begin{equation}\tag{5.4}
  \operatorname{Ext}_{\Osh_X}^{n}(X;k_x,\Omega_X^{n})\otimes L'
  \longrightarrow
  \operatorname{Ext}_{\Osh_X}^{n}(X;\Osh_X,\Omega_X^{n})\otimes L'
\end{equation}
该同态由 $\Osh_X\longrightarrow k_x$ 所对应的 $\operatorname{Ext}^{n}$ 上的
同态
\begin{equation}\tag{5.5}
  \operatorname{Ext}_{\Osh_X}^{n}(X;k_x,\Omega_X^{n})
  \longrightarrow
  \operatorname{Ext}_{\Osh_X}^{n}(X;\Osh_X,\Omega_X^{n})
  =H^{n}(X,\Omega_X^{n}).
\end{equation}
诱导。

\src{183}
由于 (5.3 bis) 是同构，\footnote{印刷文本在第 182 页给出了第二个公式
“(5.3)”，并在此处引用它；标识“(5.3 bis)”把该公式与第一个 (5.3) 区分开。}
(5.4) 也是同构，因而 (5.5) 也是同构。根据 (4.2)，$s_{(x)}$ 是
$\operatorname{Ext}_{\Osh_X}^{n}(X;k_x,\Omega_X^{n})$ 的一组基，故其像
$\varepsilon_X$ 确实是 $H^{n}(X,\Omega_X^{n})$ 的一组基。

还需证明命题 (D)。它将由以下诸引理所概括的初等事实纯形式地推出。在这些引理
中，假设 $X$ 是维数为 $r$ 的射影空间 $P$ 的闭子集（可以奇异，也可以非奇异）。
以 $\Osh_P(m)$ 表示 [3] 中记作 $\Osh(m)$ 的 $P$ 上的层，以 $\Osh_X(m)$
表示 $X$ 上类似的层。

\result{引理}{2} 若 $X=P$ 且 $F=\Osh_P(m)$，则命题 (D) 成立。

该引理可由直接计算验证。$H^{i}(P,\Osh_P(m))$ 的显式计算见 [3]，不过也可用
更初等的方法完成。为明确写出配对 (5.1) 而需要计算的杯积
\[
  H^{i}(P,\Osh_P(m))\times H^{j}(P,\Osh_P(m'))
  \longrightarrow H^{i+j}(P,\Osh_P(m+m'))
\]
没有困难。\footnote{印刷文本作 $\Osh_P(r+r')$；两个因子 $\Osh_P(m)$ 和
$\Osh_P(m')$ 要求此处为 $\Osh_P(m+m')$。}

\result{引理}{3} $X$ 上每个凝聚代数层 $F$ 都同构于某个层
$\Osh_X(-m)^{k}$ 的商层，其中 $m$ 可以取得任意大。

这是因为当 $m$ 足够大时，$F\otimes\Osh_X(m)$\fquote{由其截面生成}；
参见 [3]。

\result{引理}{4} 设 $i>0$。当 $m$ 足够大时，
$H^{r-i}(P,\Osh_P(-m))=0$；并且对 $X$ 上任意凝聚代数层 $B$，当 $m$
足够大时有
\[
  \operatorname{Ext}_{\Osh_X}^{i}(X;\Osh_X(-m),B)=0.
\]

第一个断言由上面提到的显式计算推出。对第二个断言，注意有同构
\[
  \operatorname{Ext}_{\Osh_X}^{i}(X;\Osh_X(-m),B)
  =H^{i}(X,B\otimes\Osh_X(m))
\]
（命题 1，推论 1），再由 [3] 中一个熟知的结果即得结论。结合前两个引理，
得到以下推论。

\result{推论}{} 设 $i>0$。则 $P$ 上凝聚代数层范畴中的函子
$F\longmapsto H^{r-i}(P,F)$ 是余可抹去的；$X$ 上凝聚代数层范畴中的函子
$\operatorname{Ext}_{\Osh_X}^{i}(X;F,B)$ 也是余可抹去的。

\src{184}
\result{引理}{5} 设 $A,B$ 是 $X$ 上两个凝聚代数层，并令
$A(m)=A\otimes\Osh_X(m)$。则当 $m$ 足够大时，典范同态
\[
  \operatorname{Ext}_{\Osh_X}^{1}(X;A(-m),B)
  \longrightarrow H^{0}\bigl(X,\Extsh_{\Osh_X}^{1}(A(-m),B)\bigr)
  =H^{0}\bigl(X,\Extsh_{\Osh_X}^{1}(A,B)(m)\bigr)
\]
是同构。

把命题 1 的谱序列用于 $A(-m),B$，便立即得到此结论，因为此时
\[
  E_{2}^{p,q}(A(-m),B)
  =H^{p}\bigl(X,\Extsh_{\Osh_X}^{q}(A(-m),B)\bigr)
  =H^{p}\bigl(X,\Extsh_{\Osh_X}^{q}(A,B)(m)\bigr),
\]
而当 $p>0$ 且 $m$ 足够大时，此式为零。

现在证明 $X=P$ 时的 (D)。先证当 $p=n$ 时 (5.2 bis) 为同构；此时两端都是
左正合函子（因为 $H^{r+1}(P,F)=0$），故由引理 3，只需对
$F=\Osh_P(-m)$ 证明该断言，而这正包含在引理 2 中。由于同态 (5.2 bis)
具有函子性并与连接算子相容，且当 $p<n$ 时 (5.2 bis) 两端都是关于 $F$ 的
余可抹去函子（引理 4 的推论），由标准论证可知 (5.2 bis) 对每个 $p$ 都是
同构。这就证明了射影空间的对偶定理。

现在设 $X$ 任意但非奇异。由 $P$ 的对偶定理，有同构
\[
  H^{n}(X,F)'=H^{n}(P,F)'
  =\operatorname{Ext}_{\Osh_P}^{r-n}(P;F,\Omega_P^{r}).
\]
根据第 4 节引理 1，最后一端可与
\[
  \operatorname{Hom}_{\Osh_P}(P;F,\omega_X^{n})
  =\operatorname{Hom}_{\Osh_X}(X;F,\Omega_X^{n})
  =\operatorname{Ext}_{\Osh_X}^{0}(X;F,\Omega_X^{n})
\]
等同，因而有同构
\begin{equation}\tag{5.6}
  H^{n}(X,F)'=\operatorname{Hom}_{\Osh_X}(X;F,\Omega_X^{n})
  =\operatorname{Ext}_{\Osh_X}^{0}(X;F,\Omega_X^{n}).
\end{equation}
取 $F=\Omega_X^{n}$，得到同构
\begin{equation}\tag{5.7}
  \gamma:H^{n}(X,\Omega_X^{n})\xrightarrow{\sim}k.
\end{equation}

\src{185}
可以验证，借助 (5.7) 取 $L=k$ 后，同构 (5.6) 正是 $p=n$ 时的
(5.2 bis)。因此当 $p=n$ 时，(5.2 bis) 是同构。为证明它对每个 $p$ 都是
同构，只需再证明：当 $p<n$ 时，(5.2 bis) 两端都是关于 $F$ 的余可抹去函子；
进一步说（考虑到引理 3），只需证明取 $F=\Osh_X(-m)$ 且 $m$ 足够大时两端
都为零。对第一端，这由引理 4 成立；对第二端，利用 $P$ 上的对偶定理写成
\[
  H^{p}(X,\Osh_X(-m))'
  =\operatorname{Ext}_{\Osh_P}^{r-p}
    (P;\Osh_X(-m),\Omega_P^{r}).
\]
当 $p<n$ 且 $m$ 足够大时，第二端为零。这由引理 5（其中取 $X=P$）以及下述
事实推出：$\Osh_X$ 作为 $P$ 上的凝聚代数层，其上同调维数
$\leq r-n$（因为 $X$ 非奇异），所以
\[
  \Extsh_{\Osh_P}^{r-p}(\Osh_X,\Omega_P^{r})=0
  \qquad\mbox{\upshape\fontspec{SimSun}当 }p<n.
\]

\fsection{6}{奇异簇的对偶定理}

设 $X$ 是维数为 $r$ 的射影空间 $P$ 中维数为 $n$ 的闭子集。此处公式 (5.6)
应写成
\begin{equation}\tag{6.1}
  H^{n}(X,F)'
  \simeq\operatorname{Hom}_{\Osh_X}(X;F,\omega_X^{n})
  =\operatorname{Ext}_{\Osh_X}^{0}(X;F,\omega_X^{n}),
\end{equation}
其中定义
\begin{equation}\tag{6.2}
  \omega_X^{n}=E^{0}(\Osh_X)
  =\Extsh_{\Osh_P}^{r-n}(\Osh_X,\Omega_P^{r}).
\end{equation}
如命题 5 所述，这样引入的层实际上不依赖于所选的把 $X$ 嵌入非奇异簇 $P$
的方式。在 (6.1) 中取 $F=\omega_X^{n}$，得到
\begin{equation}\tag{6.2 bis}
  H^{n}(X,\omega_X^{n})'
  \simeq\operatorname{Hom}_{\Osh_X}(X;\omega_X^{n},\omega_X^{n}),
\end{equation}
\footnote{印刷文本也把该公式编号为“(6.2)”；后缀“bis”用于区分两次出现。}
因而 $H^{n}(X,\omega_X^{n})'$ 中有一个对应于 $\omega_X^{n}$ 恒等同态的
特出元素：
\begin{equation}\tag{6.3}
  \gamma:H^{n}(X,\omega_X^{n})\longrightarrow k.
\end{equation}

\src{186}
现考察由 $\operatorname{Ext}$ 的复合所定义的配对：
\begin{equation}\tag{6.4}
  H^{p}(X,F)\times
  \operatorname{Ext}_{\Osh_X}^{n-p}(X;F,\omega_X^{n})
  \longrightarrow H^{n}(X,\omega_X^{n}),
\end{equation}
再把它们与 (6.3) 中的同态 $\gamma$ 复合，便得到与连接算子相容的函子性同态
\begin{equation}\tag{6.5}
  \operatorname{Ext}_{\Osh_X}^{n-p}(X;F,\omega_X^{n})
  \longrightarrow H^{p}(X,F)'
\end{equation}
（它推广了 (5.2)）。可以验证，当 $p=n$ 时，由此得到同构 (6.1)。在这些约定下，
有如下结论。

\result{定理}{3 bis} 对给定整数 $k\geq0$，下列关于 $X$ 的四个条件等价：
\begin{enumerate}
  \item[\textit{i.}] 当 $n-k\leq p\leq n$ 时，函子性同态 (6.5) 是同构。
  \item[\textit{ii.}] 当 $m$ 足够大且 $n-k\leq p<n$ 时，
    $H^{p}(X,\Osh_X(-m))=0$。
  \item[\textit{iii.}] 当 $n-k\leq p<n$ 时，$X$ 上凝聚代数层范畴中的函子
    $H^{p}(X,F)$ 是余可抹去的。
  \item[\textit{iv.}] 当 $0<i\leq k$ 时，
    $E^{i}(\Osh_X)=\Extsh_{\Osh_P}^{r-n+i}(\Osh_X,\Omega_P^{r})=0$。
\end{enumerate}

\noindent\textbf{证明。}——由引理 4 得
\textit{i.}$\Rightarrow$\textit{ii.}，由引理 3 得
\textit{ii.}$\Rightarrow$\textit{iii.}。又因为当 $n-k\leq p<n$ 时，
(6.5) 两端都是余可抹去函子（前者由引理 4），故由一个熟知的标准论证得到
\textit{iii.}$\Rightarrow$\textit{i.}。最后证明
\textit{ii.}$\Leftrightarrow$\textit{iv.}；这由下一命题的推论推出。

\result{命题}{6} 设 $F$ 是 $X$ 上的凝聚代数层，$i$ 是整数。则当 $m$ 足够大时，
有同构
\begin{equation}\tag{6.6}
  H^{i}(X,F(-m))'
  \simeq H^{0}\bigl(X,E^{n-i}(F)(m)\bigr),
\end{equation}
其中规定（比较第 3 节命题 5）
\begin{equation}\tag{6.7}
  E^{j}(F)=\Extsh_{\Osh_P}^{r-n+j}(F,\Omega_P^{r}).
\end{equation}
事实上，由 $P$ 上的对偶定理，(6.6) 的左端同构于
$\operatorname{Ext}_{\Osh_P}^{r-i}(P;F(-m),\Omega_P^{r})$，故 (6.6)
由引理 5 推出。

\src{187}
\result{推论}{} 当 $m$ 足够大时 $H^{i}(X,F(-m))=0$，当且仅当
$E^{n-i}(F)=0$。

回顾 $E^{j}(F)$ 实际上不依赖于所考察的射影嵌入。该推论中的条件完全是局部的；
因此，若它对 $F$ 成立，就对任一局部同构于层 $F^{n}$ 的层成立。特别地，若该
条件对 $\Osh_X$ 成立，就对每个局部自由凝聚代数层成立。例如：若 $X$ 非奇异，
则对每个 $i<n$ 都成立；若 $X$ 的任何分支都不退化为一个点，则对 $i=0$
成立；若 $X$ 正规且其所有分支的维数都 $>1$，则对 $i=0,1$ 成立（参见
[3]）。\footnote{印刷文本此处写作“$S$ 正规”；本段只涉及代数簇 $X$。}
按照定义，要使它对每个 $i<k$ 都成立，充要条件是局部环
$\Osh_x$（$x\in X$）的\fquote{同调余维数} $\geq k$（关于此概念参见
[4]）。若 $k=n$，则由定理 3 bis，这意味着对偶定理对 $X$ 成立，即对每个
$p$ 和每个 $F$，(6.5) 都是同构。关于局部环 $\Osh_x$，可以给出许多与此
等价的条件（NAGATA），例如要求它们满足 Cohen--Macaulay 等维定理。比如，
若 $X$ 局部地是某个非奇异环境代数簇中的\fquote{完全交}，这一条件便成立。

\fsection{7}{庞加莱对偶}

设 $X$ 是维数为 $n$ 的非奇异射影代数簇。则
\[
  H^{**}(X)=H^{*}(X,\Omega_X^{*})
\]
是一个有限维反交换双分次代数；我们按总次数给它分次，因而
$H^{p,q}(X)=H^{p}(X,\Omega_X^{q})$ 的次数为 $p+q$。$H^{*}(X)$ 的次数介于
$0$ 与 $2n$ 之间。由定理 2 和定理 3 的推论，$H^{*}(X)$ 是维数为 $2n$ 的
\fquote{庞加莱代数}：也就是说，$H^{2n}(X)$ 具有到基域 $k$ 的一个同构，且乘法
\[
  H^{m}(X)\times H^{2n-m}(X)\longrightarrow H^{2n}(X)=k
\]
给出 $H^{m}(X)$ 与 $H^{2n-m}(X)$ 之间的对偶。此外，若 $Y$ 是另一个非奇异
射影代数簇，则凝聚代数层的 Künneth 公式给出
\begin{equation}\tag{7.1}
  H^{*}(X\times Y)=H^{*}(X)\otimes H^{*}(Y),
\end{equation}
该同构与庞加莱代数结构相容。再者，作为交换代数，$H^{*}(X)$ 是 $X$ 的反变
函子；态射 $f:Y\longrightarrow X$ 以显然方式定义分次代数同态

\src{188}
\begin{equation}\tag{7.2}
  f^{*}:H^{*}(X)\longrightarrow H^{*}(Y).
\end{equation}
由于这里涉及的是庞加莱代数，转置便给出向量空间同态
\begin{equation}\tag{7.3}
  f_{*}:H^{*}(Y)\longrightarrow H^{*}(X).
\end{equation}
我们在第 4 节已经看到，$f^{*}$ 对非奇异循环所对应上同调类的作用，可在几何上
解释为取这些循环的逆像所对应的上同调类。在当前情形，重要的是证明 (7.3)
同样对应于循环的\fquote{正像}运算。至少在适当的非奇异性条件下，这由下述特殊
情形推出。

\result{定理}{4} 若 $f$ 是非奇异子簇 $Y^{m}\subset X^{n}$ 到 $X^{n}$ 的
自然包含映射，并以 $1_Y$ 表示 $H(Y)$ 的单位元，则
\begin{equation}\tag{7.4}
  f_{*}(1_Y)=P_X(Y),
\end{equation}
其中右端是 $Y$ 在 $X$ 中所对应的上同调类。

该公式等价于
\begin{equation}\tag{7.4 bis}
  \bigl\langle\xi^{m,m},P_X(Y)\bigr\rangle\varepsilon_Y
  =f^{*}(\xi^{m,m})
  \qquad
  \bigl(\xi^{m,m}\in H^{m}(X,\Omega_X^{m})\bigr),
\end{equation}
\footnote{印刷文本写作 $f_{*}(\xi^{m,m})$。由于
$\xi^{m,m}\in H^{m}(X,\Omega_X^{m})$，且等式两端均属于
$H^{m}(Y,\Omega_Y^{m})$，这里需要的映射是 $f^{*}$。}
其中 $\varepsilon_Y$ 是 $H^{m}(Y,\Omega_Y^{m})$ 中的基本元；对于非奇异
射影代数簇，该式给出了 $Y$ 所对应上同调类的一个新定义。为证明定理 4，
借助定理 3，把同态
\[
  H^{m}(X,\Omega_X^{m})
  \longrightarrow H^{m}(Y,\Omega_Y^{m})
  =H^{m}(X,\Omega_Y^{m})
\]
的转置解释为同态
\begin{equation}\tag{7.5}
\begin{split}
  H^{n-m}(X,\Omega_X^{n-m})
  &\simeq
  \operatorname{Ext}_{\Osh_X}^{n-m}(X;\Omega_X^{m},\Omega_X^{n})\\
  &\longleftarrow
  \operatorname{Ext}_{\Osh_X}^{n-m}(X;\Omega_Y^{m},\Omega_X^{n})
  \simeq
  \operatorname{Hom}_{\Osh_X}(X;\Omega_Y^{m},\Omega_Y^{m}).
\end{split}
\end{equation}
可以验证，$H^{m}(Y,\Omega_Y^{m})$ 的对偶中的元素 $1_Y$，与右端中对应于
$\Omega_Y^{m}$ 恒等自同态的元素等同；另一方面，该元素在
$H^{n-m}(X,\Omega_X^{n-m})$ 中的像正是 $P_X(Y)$。这些相容性本可在第 4 节
就给出；事实上，它们可以对任意非奇异代数簇表述，并且确实成立。例如，第二个
相容性来自下列典范自同态图的交换性：

\src{189}
\begin{equation}\tag{7.6}
\resizebox{\textwidth}{!}{$
\left\{
\begin{array}{ccccc}
\operatorname{Ext}_{\Osh_X}^{n-m}(X;\Omega_X^{m},\Omega_X^{n})
&\longleftarrow&
\operatorname{Ext}_{\Osh_X}^{n-m}(X;\Omega_Y^{m},\Omega_X^{n})
&\simeq&
\operatorname{Hom}_{\Osh_X}(X;\Omega_Y^{m},\Omega_Y^{m})\\[3pt]
\Big\downarrow&&&&\Big\downarrow\\[3pt]
\operatorname{Ext}_{\Osh_X}^{n-m}(X;\Osh_X,\Omega_X^{n-m})
&\longleftarrow&
\operatorname{Ext}_{\Osh_X}^{n-m}(X;\Osh_Y,\Omega_X^{n-m})
&\simeq&
\operatorname{Hom}_{\Osh_X}(X;\Omega_X^{m},\Omega_Y^{m})
\end{array}
\right.
$}
\end{equation}

由此得到紧致定向流形上庞加莱对偶形式体系的一个完全对应物。特别地，定理 4
可以确定 $X\times X$ 的对角线所对应的上同调类。再由一个熟知的论证，例如可得
如下 Lefschetz 公式。

\result{定理}{5} 设 $f$ 是非奇异射影代数簇 $X$ 的一个自同态，并且 $f$ 的
所有不动点的重数均为 1。则这些不动点的数目模 $k$ 的特征，同余于 $f$ 在各
$H^{i}(X)$ 上所诱导自同态之迹的交错和。

对 $f$ 所作的限制源于第 4 节注记中提到的困难。不过应当注意，若 $f$ 与一个
所有不动点重数均为 1 的自同态\fquote{同伦}，Lefschetz 公式仍然成立。

\fsection{8}{对偶定理的推广}

设 $X$ 是一个非奇异代数簇，并且 $X$ 上每个凝聚代数层 $F$ 都同构于某个局部
自由凝聚代数层的商（当 $X$ 是射影空间中的局部闭子集时即属此情形）。于是
$X$ 上每个凝聚代数层 $F$ 都具有一个由局部自由层组成的有限消解 $L$；对于
任意两个这样的消解，总能找到第三个消解以及从它到前两个消解、且与增广相容的
同态。同样，若 $L$ 是 $F$ 的局部自由有限消解，并给定同态
$F'\longrightarrow F$，则存在 $F'$ 的局部自由有限消解 $L'$ 以及与
$F'\longrightarrow F$ 相容的同态 $L'\longrightarrow L$；若
$F'\longrightarrow F$ 是满射，还可以令后一同态为满射。由此，对于给定的两个
整数 $r,s\geq0$，可以定义两个关于诸变元
$A_1,\ldots,A_r$；$B_1,\ldots,B_s$ 的上同调多重函子；

\src{190}
这些变元在 $X$ 上凝聚代数层范畴中变化，而两个函子分别取值于 $X$ 上凝聚
代数层范畴和 $H^{0}(X,\Osh_X)$-模范畴；其定义式为
\begin{equation}\tag{8.1}
\resizebox{\textwidth}{!}{$
\left\{
\begin{aligned}
\underline{T}_{r}^{s*}(A_1,\ldots,A_r;B_1,\ldots,B_s)
  &={}H^{*}\!\left(
    \Homsh_{\Osh_X}\!\left(
      L(A_1)\otimes\cdots\otimes L(A_r),
      L(B_1)\otimes\cdots\otimes L(B_s)
    \right)\right),\\[3pt]
T_{r}^{s*}(A_1,\ldots,A_r;B_1,\ldots,B_s)
  &={}\mathbf R^{*}\Gamma\!\left(
    \Homsh_{\Osh_X}\!\left(
      L(A_1)\otimes\cdots\otimes L(A_r),
      L(B_1)\otimes\cdots\otimes L(B_s)
    \right)\right).
\end{aligned}
\right.
$}
\end{equation}
在该公式中，$L(F)$ 表示凝聚代数层 $F$ 的一个局部自由有限消解，
$\mathbf R^{*}\Gamma(K)$ 表示空间 $X$ 关于层复形 $K$ 的超上同调。若 $r$ 或
$s$ 为零，就分别用 $\Osh_X$ 代替诸 $L(A_i)$ 或诸 $L(B_j)$ 的张量积。特别地，
$\underline{T}_{0}^{0}$ 与 $T_{0}^{0}$ 是零变元的分次函子：
$\underline{T}_{0}^{0}$ 只在次数 0 非零，此时它是层 $\Osh_X$；而
$T_{0}^{0}$ 等同于 $H^{*}(X,\Osh_X)$。(8.1) 右端不依赖所选消解：对第一行
这无论如何是显然的（因为问题是局部的）；对第二行，则由此前的一般性说明以及
层复形
\[
  K=\Homsh_{\Osh_X}\bigl(L(A_1)\otimes\cdots,
      L(B_1)\otimes\cdots\bigr)
\]
的超上同调谱序列推出。该谱序列收敛于 $X$ 关于 $K$ 的超上同调，其初始项是
$H^{p}(X,H^{q}(K))$，即
\begin{equation}\tag{8.2}
  E_{2}^{p,q}
  =H^{p}\!\left(X,
    (\underline{T}_{r}^{s})^{(q)}
      (A_1,\ldots,A_r;B_1,\ldots,B_s)\right).
\end{equation}
由此可见，该谱序列本身也不依赖所选消解，而其收敛项就是 (8.1) 中相应的
左端。又因每个正合列
$0\longrightarrow F'\longrightarrow F\longrightarrow F''\longrightarrow0$
都能由局部自由有限复形的一个正合列来消解，故容易定义关于各变元
$A_i,B_j$ 的连接算子。

在函子系 $\underline{T}_{r}^{s*}$ 以及 $T_{r}^{s*}$ 中，可以定义类似于
张量演算的运算；它们由定义式 (8.1) 立即给出。例如，有一个复合运算
（它推广了第 2 节所考察的复合）
\begin{equation}\tag{8.3}
\resizebox{\textwidth}{!}{$
\begin{aligned}
&T_{r}^{s*}(A_1,\ldots,A_r;B_1,\ldots,B_s)
\times T_{r'}^{s'*}(A'_1,\ldots,A'_{r'};B'_1,\ldots,B'_{s'})\\
&\qquad\longrightarrow
T_{r+r'}^{s+s'*}
(A_1,\ldots,A_r,A'_1,\ldots,A'_{r'};
 B_1,\ldots,B_s,B'_1,\ldots,B'_{s'}),
\end{aligned}
$}
\end{equation}
它满足显然的结合律，并与函子性同态、连接同态和谱序列相容。同样，还有一些
对称运算，其具体写法留给读者。再者，每当某个变元 $A_i$

\src{191}
等于某个变元 $B_j$ 时，就有一个缩并运算：
\begin{equation}\tag{8.4}
\resizebox{\textwidth}{!}{$
\begin{aligned}
&T_{r}^{s*}(A_1,\ldots,A_{i-1},C,A_{i+1},\ldots,A_r;
 B_1,\ldots,B_{j-1},C,B_{j+1},\ldots,B_s)\\
&\qquad\longrightarrow
T_{r-1}^{s-1*}(A_1,\ldots,\widehat{A_i},\ldots,A_r;
 B_1,\ldots,\widehat{B_j},\ldots,B_s).
\end{aligned}
$}
\end{equation}
此外，若变元 $A_i$ 是局部自由层，则可以消去它，只需把另一个 $A_j$
（$j\neq i$）换成 $A_j\otimes A_i$，或把某个 $B_k$ 换成
$B_k\otimes A_i'$（其中 $A_i'=\Homsh_{\Osh_X}(A_i,\Osh_X)$）；当某个变元
$B_i$ 局部自由时，也有类似的计算规则。特别地，总可以消去等于 $\Osh_X$ 的
变元。若除至多一个 $B_i$ 外所有变元都局部自由，刚才的规则给出函子性同构
\begin{equation}\tag{8.5}
  T_{r}^{s*}(A_1,\ldots,A_r;B_1,\ldots,B_s)
  =H^{*}\bigl(X,A_1'\otimes\cdots\otimes A_r'
    \otimes B_1\otimes\cdots\otimes B_s\bigr)
\end{equation}
（因为这可归结到 $r=0,s=1$ 的情形，而该情形是立即可见的；也可以直接使用
初始项为 (8.2) 的谱序列）。对 $\underline{T}_{r}^{s}$ 也可定义与前述运算
相应的运算。如此引入的各运算之间的关系，与张量演算中类似运算之间的关系相同。

设 $n$ 是 $X$ 的维数。依次应用张量复合 (8.3)，再对重复出现的变元施行缩并
(8.4)，便得到配对
\begin{equation}\tag{8.6}
\resizebox{\textwidth}{!}{$
  (T_{r}^{s})^{p}(A_1,\ldots;B_1,\ldots)
  \times
  (T_{s}^{r})^{n-p}(B_1,\ldots;A_1,\ldots,A_r\otimes\Omega_X^{n})
  \longrightarrow H^{n}(X,\Omega_X^{n}).
$}
\end{equation}

\result{定理}{6} 若 $X$ 是非奇异射影代数簇，则配对 (8.6) 都是对偶。

这完全形式地由定理 3 的推论得出。事实上，从该推论容易推出：若 $K$ 是由
局部自由凝聚代数层组成的复形，则 $X$ 关于 $K$ 的超上同调，与 $X$ 关于
$K'\otimes\Omega_X^{n}$ 的超上同调，通过自然配对
\begin{equation}\tag{8.7}
  \mathbf R^{p}\Gamma(K)
  \times\mathbf R^{n-p}\Gamma(K'\otimes\Omega_X^{n})
  \longrightarrow\mathbf R^{n}\Gamma(\Omega_X^{n})
  =H^{n}(X,\Omega_X^{n}).
\end{equation}
互为对偶。这可由初始项为
$H^{p}(X,H^{q}(K))$\footnote{印刷文本写作
$H^{p}(H^{q}(X,K))$；初始项应为谱序列 (8.2) 的初始项，即
$H^{p}(X,H^{q}(K))$。}
的谱序列，以及 $K'\otimes\Omega_X^{n}$ 的类似谱序列看出。利用定义 (8.1)，
定理 6 即由上述结果推出。

\src{192}
\subsection*{注记}

\noindent 1\textsuperscript{o} 在定理 6 之前的定义中，并不需要假设 $X$
非奇异，因为只使用有限消解并非不可缺少。但是，当 $X$ 奇异时，不能再先验地
断言
$(\underline{T}_{r}^{s})^{p}(A_1,\ldots;B_1,\ldots)$ 是凝聚层，因为在层复形
\[
  \Homsh_{\Osh_X}\bigl(L(A_1)\otimes\cdots,L(B_1)\otimes\cdots\bigr)
\]
中，每个给定的总次数都会有无穷多个分量。

\medskip
\noindent 2\textsuperscript{o} 容易验证，在公式 (8.1) 中，可以把某个
$L(B_i)$ 换成 $B_i$。结合命题 3，这就表明
\begin{equation}\tag{8.8}
\left\{
\begin{aligned}
  \underline{T}_{1}^{1*}(A;B)&=\Extsh_{\Osh_X}^{*}(A,B),\\
  T_{1}^{1*}(A;B)&=\operatorname{Ext}_{\Osh_X}^{*}(X;A,B).
\end{aligned}
\right.
\end{equation}
特别地，在 (8.6) 中取 $r=s=1$、$A_1=\Osh_X$，便重新得到定理 3。公式
(8.8) 还蕴含
\[
  T_{0}^{1*}(B)=H^{*}(X,B),
  \qquad
  T_{1}^{0*}(A)=\operatorname{Ext}_{\Osh_X}^{*}(X;A,\Osh_X).
\]

\medskip
\noindent 3\textsuperscript{o} 由公式 (8.1) 可见，一般而言，函子
$\underline{T}_{r}^{s*}$ 和 $T_{r}^{s*}$ 同时具有正次数和负次数的分量。
利用注记 2 可见：若 $X$ 的维数为 $n$，则当 $s>0$ 时，
$\underline{T}_{r}^{s*}$ 的非零分量介于 $-(s-1)n$ 与 $rn$ 之间；当
$s=0$ 时，介于 $0$ 与 $rn$ 之间。对于 $T_{r}^{s*}$，当 $s>0$ 时，其
非零分量介于 $-(s-1)n$ 与 $(r+1)n$ 之间；当 $s=0$ 时，介于 $0$ 与
$(r+1)n$ 之间（而且，若无误，当 $r>0$ 时，甚至分别介于 $-(s-1)n$
与 $rn$ 之间以及 $0$ 与 $rn$ 之间）。

\section*{参考文献}

\noindent[1] Cartier (Pierre). --- Les groupes
$\operatorname{Ext}^{s}(A,B)$, Séminaire A. Grothendieck : Algèbre homologique,
t. 1, 1957, n\textsuperscript{o} 3.

\noindent[2] Grothendieck (Alexandre). --- Sur quelques points d'algèbre
homologique, \emph{Tohoku Math. J.}, t. 9, 1957, p. 119--183.

\noindent[3] Serre (Jean-Pierre). --- Faisceaux algébriques cohérents,
\emph{Annals of Math.}, t. 61, 1955, p. 197--278.

\noindent[4] Serre (Jean-Pierre). --- Sur la dimension homologique des anneaux
et des modules noethériens, \emph{Proc. Intern. Symp. on Alg. Number Theory}
[1955, Tokyo et Nikko]. --- Tokyo, Science Council of Japan, 1956 ;
p. 175--189.

\clearpage
\src{193}
\section*{补记}

第 13 页注记中指出的困难目前已经完全解决，这是由于对偶定理已推广到任意
代数簇（允许任意奇点）。关于这一理论——它可以在\fquote{概形理论}的一般框架
中表述——我们请读者参阅 J. Dieudonné 与 A. Grothendieck 当时正在编写的
《Éléments de Géométrie algébrique》。下列文献给出了一些说明：

\begin{quote}
Grothendieck (A.). --- Cohomology theory of algebraic varieties,
\emph{Congrès int. des Mathématiciens}, 1958, Edinburgh (à paraître).
\end{quote}

\begin{flushright}
[1959年4月]
\end{flushright}

\begin{center}\rule{36mm}{0.5pt}\end{center}

% END 149.tex

% BEGIN 182.tex
\unit{第182次报告——形式几何与代数几何}
\sid{182}
\src{193}
\seminarzh{（1959年5月）}
\paperhead{形式几何与代数几何}%
  {亚历山大·格罗滕迪克 著}

\fsection{1}{概形}

众所周知，定义在域 $k$ 上的仿射代数空间本质上由其仿射代数 $A$
（定义在 $k$ 上的正则函数环）所决定；代数空间的态射
$X\longrightarrow Y$ 与 $k$-代数同态
$A(Y)\longrightarrow A(X)$ 一一对应。与一个代数空间相应的仿射代数是
有限生成的 $k$-代数，并且从\fquote{经典}观点来看，它没有幂零元；反过来，
每个这类代数都可由一个定义在 $k$ 上的代数空间作为其仿射代数而得到。
于是有一套熟知的词典，可将有关仿射代数空间的情形解释为交换代数的
情形。人们早已发现，由此得到的命题其实更为一般：通常不再需要假设所涉
各环都属于刚才考察的类型，诺特性假设大多已经足够。特别地，无论是否
指定了基域，都没有理由排除这些环含有幂零元的情形。以往几何学家一直
拒绝正视这些启示，执意把研究限制在无幂零元的仿射代数上，也就是说，
限制在结构层中不含幂零元的代数空间上（而且往往还要限制在
\fquote{绝对不可约}的代数空间上）。报告人认为，这种观念严重阻碍了
代数几何中真正自然的方法的发展。

设 $A$ 是交换环。众所周知，$A$ 的素理想所成的集合
$X=\operatorname{Spec}(A)$ 带有一个自然拓扑，称为\fquote{Zariski 拓扑}
或谱拓扑。另一方面，$X$ 上有一个交换环层 $\Osh_X$，其在
$\mathfrak p\in X$ 处的茎是局部化环 $A_{\mathfrak p}$，其截面环可与 $A$
等同。这样，$X$ 就成为一个\emph{赋环空间}，称为 $A$ 的\emph{素谱}。
环同态 $f:A\longrightarrow B$ 定义了赋环空间的一个态射
$f':\operatorname{Spec}(B)\longrightarrow\operatorname{Spec}(A)$，
其底层集合映射正是
$\mathfrak p\longmapsto f^{-1}(\mathfrak p)$。由这种方式得到的

\src{194}
从 $\operatorname{Spec}(B)$ 到 $\operatorname{Spec}(A)$ 的赋环空间态射，
恰好是使所有同态 $\Osh_x\longrightarrow\Osh_y$（其中 $x=f'(y)$）均为
局部同态的那些态射（即极大理想的逆像仍为极大理想）。

与某个 $\operatorname{Spec}(A)$ 同构的赋环空间称为\emph{仿射概形}；
局部仿射的赋环空间称为\emph{预概形}，也就是说，它的每一点都有一个
开邻域，在诱导结构下该邻域是仿射概形。预概形的\emph{态射}按显然方式
定义；局部地，它们与环同态相对应。

固定一个预概形 $S$ 并考察预概形态射 $X\longrightarrow S$ 时，$S$ 所起的
作用类似于基域或基环（更确切地说，是一个纤维化中的基空间）。这时称
$X$ 为一个 $S$-预概形；若 $S=\operatorname{Spec}(A)$，这也意味着
$\Osh_X$ 是一个 $A$-代数层。于是每个预概形都能以唯一方式看成一个
$\mathbf Z$-预概形。当然，$S$-预概形构成一个范畴；进一步可以证明，
在这个范畴中任意两个对象 $X,Y$ 的乘积总是存在，记作 $X\times_S Y$。
借助乘积的概念，可以定义与一个态射 $S'\longrightarrow S$ 相应的
$S$-预概形的基变换：事实上，$X\times_S S'$ 可看成一个 $S'$-预概形。

若 $X$ 在 $X\times_S X$ 中的对角像是闭的，就称 $X$ 在 $S$ 上
\emph{分离}。在 $\mathbf Z$ 上分离的预概形称为\emph{概形}；于是它在
任意基底上都是分离的。为简单起见，下文只谈概形，并且还总假设它们是
\emph{诺特的}，也就是说，可由有限个仿射开集覆盖，且这些开集都是诺特环
的谱。若对 $S$ 的每个仿射开集 $U$，其在 $X$ 中的逆像都能由有限个仿射
开集覆盖，并且这些开集的环都是 $U$ 的环上的有限型代数，就称 $X$ 在
$S$ 上\emph{有限型}。正是这类 $S$-概形适合进行真正的几何研究。特别地，
对每个 $s\in S$，$X$ 在 $s$ 上的纤维 $f^{-1}(s)$ 是定义在 $s$ 的局部环
$\Osh_s$ 的剩余域 $\kappa(s)$ 上的代数概形。因此，在某种程度上可以把
$X$ 看成一族\fquote{代数空间} $f^{-1}(s)$，其中参数 $s$ 遍历 $S$
（也就是说，从局部观点来看，遍历某个给定环的素理想集合）。当然，各个
$\kappa(s)$ 可以具有不同的特征。若 $S=\operatorname{Spec}(k)$，其中
$k$ 是一个域，便基本恢复了通常的\fquote{代数空间}概念，唯一的区别是
现在结构层可以含有幂零元。

\src{195}
仿照熟知的概念，可以定义\emph{射影态射}，以及更一般的\emph{固有态射}。
这类态射是有限型的；此外，它将闭子集映为闭子集，并且在任意基变换下
仍保持这一性质。

设 $X$ 是一个概形（照例是诺特的），则层 $\Osh_X$ 是[2]意义下的
\emph{凝聚环层}。因此，$X$ 上的凝聚模层也正是那些局部同构于某个同态
$\Osh_X^{m}\longrightarrow\Osh_X^{n}$ 的余核的层。

\fsection{2}{形式概形}

设 $X$ 是一个概形，$X'$ 是 $X$ 的一个闭子集。于是存在 $\Osh_X$ 的一个
凝聚子层 $J$，使得 $X'=\operatorname{Supp}(\Osh_X/J)$（而且还存在其中
最大的一个）。赋予 $\Osh_X/J$ 后，$X'$ 成为一个概形，记作 $X_0$；
这样的概形称为\fquote{$X$ 的闭子概形}。对每个 $n$，还可以考察赋有
$\Osh_X/J^{n+1}$ 的 $X'$，记作 $X_n$；它是 $X$ 的一个闭子预概形，
其底层集合仍为 $X'$，但结构层已经不同，即
$\Osh_{X_n}=\Osh_X/J^{n+1}$。显然，各 $\Osh_{X_n}$ 构成 $X'$ 上环层的
一个逆系统；其逆极限 $\overline{\Osh}_X$ 称为 $\Osh_X$ 沿 $X'$ 的
\emph{形式完备化}。赋予这个环层后，$X'$ 称为 $X$ 沿 $X'$ 的
\emph{形式完备化}；所以它是一个赋环空间，但一般不是概形。对 $X$ 上的
每个凝聚层 $F$，同样可以考察 $F$ 沿 $X'$ 的形式完备化
$\overline F=\varprojlim_n F_n$（其中
$F_n=F\otimes_{\Osh_X}\Osh_X/J^{n+1}$）；它是 $\overline X$ 上的模层。
它的截面称为\fquote{$F$ 沿 $X'$ 的形式截面}，并可与
$\varprojlim_n\Gamma(X',F_n)$ 的元素等同。取 $F=\Osh_X$ 时，得到
Zariski 意义下 $X$ 沿 $X'$ 的\fquote{全纯函数}（由于这一称法与经典术语
相冲突，本文不沿用它）。

所谓\emph{形式概形}（默认是诺特的），是指一个拓扑空间 $\mathfrak X$，
连同其上的一个拓扑环层 $\Osh_{\mathfrak X}$，并满足以下条件：存在拓扑
环层的同构
$\Osh_{\mathfrak X}=\varprojlim_n\Osh_n$，其中各 $\Osh_n$ 构成
$\mathfrak X$ 上环层的一个逆系统，并且每个 $\Osh_n$ 都使
$\mathfrak X$ 成为一个概形 $\mathfrak X_n$；当 $m\geq n$ 时，同态
$\Osh_m\longrightarrow\Osh_n$ 是满射，其核为 $J_m^{n+1}$，其中 $J_m$ 是

\src{196}
$\Osh_m\longrightarrow\Osh_0$ 的核。由此可以证明，
$\Osh_{\mathfrak X}$ 是一个凝聚的诺特局部环层。

根据定义，上述形式完备化 $\overline X$ 是一个形式概形；反过来，每个形式
概形\emph{局部地}都具有这种形式。事实上，给定一个\emph{仿射}形式概形
（即 $\mathfrak X_0$ 为仿射概形；这蕴含所有 $\mathfrak X_n$ 都是仿射
概形），等价于给定一个在 $J$-进拓扑下分离且完备的诺特拓扑环。

普通概形的各项通常定义——态射、有限型态射、固有态射等等——都能毫无
困难地推广到形式概形。

\fsection{3}{三个基本定理}

设 $f:X\longrightarrow Y$ 是概形（照例是诺特概形）的固有态射，$Y'$ 是
$Y$ 的一个闭子集，$X'$ 是它在 $X$ 中的逆像；考察相应的形式完备化
$\overline Y$ 和 $\overline X$。于是 $f$ 诱导形式概形的态射
$\overline f:\overline X\longrightarrow\overline Y$，而且这个态射仍是
固有的。设 $F$ 是 $X$ 上的凝聚层，则 $\overline F$ 是 $\overline X$ 上的
凝聚层。在定理1中，我们不再提 $X,Y,F$，只考察形式概形的固有态射
$\overline f$ 和 $\overline X$ 上的凝聚层 $\overline F$。（不过，报告人
只在从 $X,Y,f,F$ 出发的情形中写出了完整证明。）

\result{定理}{1（有限性定理）}
\begin{enumerate}
  \item[\textit{i.}] 各 $R^{q}\overline f_{*}(\overline F)$ 都是
    $\overline Y$ 上的凝聚层。
  \item[\textit{ii.}] 自然同态
    \[
      R^{q}\overline f_{*}(\overline F)
      \longrightarrow \varprojlim_n R^{q}f_{n*}(F_n)
    \]
    都是同构。
\end{enumerate}

在这个命题中，假定已经选取了 $\Osh_Y$ 的一个定义 $Y'$ 的凝聚子层 $J$；
由逆像得到 $\Osh_X$ 的一个定义 $X'$ 的凝聚子层，进而按第2节定义
$F_n,X_n,Y_n$ 以及 $f_n:X_n\longrightarrow Y_n$。若从任意两个形式概形
之间的固有态射出发，为明确这些记号而需要作出的细微调整也是显然的。

定理1只涉及\fquote{形式上同调}。下一个定理

\src{197}
把它与\fquote{代数上同调}联系起来，并且类似于 Serre 关于代数上同调与
解析上同调之比较的一个熟知定理[4]。

\result{定理}{2（第一比较定理）} 各 $R^{q}f_{*}(F)$ 都是 $Y$ 上的凝聚层
（这是定理1的一个特例），并且自然同态
\[
  \overline{R^{q}f_{*}(F)}
  \longrightarrow \varprojlim_n R^{q}f_{n*}(F_n)
\]
都是同构。

\result{推论}{1} 有典范同构
\[
  \overline{R^{q}f_{*}(F)}=R^{q}\overline f_{*}(\overline F).
\]

当 $q=0$ 时，这个推论是 Zariski 的\fquote{全纯函数基本定理}的推广；我们
将由它推出 Zariski 的\fquote{连通性定理}的推广。还应注意，定理1
（\textit{ii}）在 $q=0$ 时虽然是平凡的，定理2或它的等价表述（推论1）
却绝非如此。事实上，证明对 $q$ 作降归纳（当 $q$ 很大时两端都为零，
所以结论平凡），于是 $q=0$ 的情形正是归纳的最后一步，因而可以说是
\fquote{最困难}的情形。

\result{推论}{2} 设 $Y=\operatorname{Spec}(A)$，而 $Y'$ 由 $A$ 的理想
$J$ 定义。则对 $X$ 上的每个凝聚层 $F$，各 $H^{q}(X,F)$ 都是有限生成的
$A$-模，其 $J$-进完备化就是 $H^{q}(\overline X,\overline F)$。

最后将这个推论应用于
$H=\Homsh_{\Osh_X}(F,G)$，得到：

\result{推论}{3} 设 $Y=\operatorname{Spec}(A)$，而 $Y'$ 由 $A$ 的理想
$J$ 定义。若 $F,G$ 是 $X$ 上的两个凝聚层，则
$\operatorname{Hom}(F,G)$ 是有限生成的 $A$-模，其 $J$-进完备化可与
$\operatorname{Hom}(\overline F,\overline G)$ 等同。

当然，自然映射
$\operatorname{Hom}(F,G)\longrightarrow
\operatorname{Hom}(\overline F,\overline G)$ 把同态
$u:F\longrightarrow G$ 映为它\fquote{依连续性}得到的延拓
$\overline u:\overline F\longrightarrow\overline G$（采用这一约定后，
$\overline F$ 成为关于 $F$ 的函子）。

现在假设 $A$ 在其 $J$-进拓扑下分离且完备，则上述推论2和推论3给出
\[
  H^{q}(X,F)=H^{q}(\overline X,\overline F),
  \qquad
  \operatorname{Hom}(F,G)
  =\operatorname{Hom}(\overline F,\overline G).
\]

\src{198}
最后这个等式表明，$X$ 上凝聚层的范畴可与 $\overline X$ 上凝聚层范畴的
一个\emph{子范畴}（其态射为所诱导的态射）等同。事实上还有更强的结论：

\result{定理}{3} $\overline X$ 上的一个模层是凝聚的，当且仅当它同构于
形如 $\overline F$ 的层，其中 $F$ 是 $X$ 上的凝聚层（由定理2的推论3，
$F$ 在典范同构的意义下唯一确定）。[回顾此时
$Y=\operatorname{Spec}(A)$，其中 $A$ 是一个在 $J$-进拓扑下分离且完备的
诺特拓扑环。]

\result{推论}{1} $X$ 的闭子概形与 $\overline X$ 的闭形式子概形一一对应。

事实上，它们分别对应于 $\Osh_X$ 和 $\overline{\Osh}_X$ 的凝聚子层。
把态射的图看成闭子概形，由推论1得到：

\result{推论}{2} 设 $X,Z$ 是 $A$ 上的两个固有概形（$A$ 是在 $J$-进拓扑
下分离且完备的诺特环）。则映射 $g\longmapsto\overline g$ 在从 $X$ 到
$Z$ 的 $Y$-态射与从 $\overline X$ 到 $\overline Z$ 的 $\overline Y$-态射
之间给出一一对应。

换言之，$A$ 上的固有代数概形构成 $\overline Y$ 上固有形式概形范畴的一个
\emph{子范畴}（其态射为所诱导的态射）。但须注意，$\overline Y$ 上存在
不能\fquote{代数化}的固有形式概形，也就是说，它们并不同构于任何
$\overline X$，其中 $X$ 在 $A$ 上固有（正如存在并非来自复数域上代数簇的
紧复解析簇一样）。这类形式概形在\fquote{模理论}中会自然出现。不过，
下面这个形式概形可代数化的特殊情形值得注意：

\result{定理}{4} 设 $A$ 是剩余域为 $k$ 的完备诺特局部环，$\mathfrak X$
是 $A$ 上的固有形式概形（赋有其 $\mathfrak m(A)$-进拓扑）。假设：
\begin{enumerate}
  \item[\textit{i.}] $\Osh_{\mathfrak X}$ 的各局部环都是平坦 $A$-模；
    或者等价地，若给 $\Osh_{\mathfrak X}$ 和 $A$ 赋予由 $A$ 的极大理想
    的各次幂定义的滤过，则相伴分次对象满足关系
    \[
      \operatorname{gr}(\Osh_{\mathfrak X})
      \simeq
      \operatorname{gr}^{0}(\Osh_{\mathfrak X})
      \otimes_k\operatorname{gr}(A).
    \]
  \item[\textit{ii.}] 把
    $\mathfrak X_0=\mathfrak X\otimes_A k$ 看成 $k$ 上的代数概形时，有
    \[
      H^2(\mathfrak X_0,\Osh_{\mathfrak X_0})=0.
    \]
  \item[\textit{iii.}] $\mathfrak X_0$ 是射影的。
\end{enumerate}
在这些条件下，$\mathfrak X$ 可以代数化；更确切地说，它同构于
$\overline X$，其中 $X$ 是一个射影 $A$-概形。

当 $\mathfrak X_0$ 是 $k$ 上的一条非奇异曲线时，条件\textit{(ii)}和
\textit{(iii)}特别成立；因此定理4尤其可以应用于给定亏格曲线的
\fquote{模理论}\ldots{} 下面说明如何证明定理4：可以证明（参见下文命题3）
条件\textit{(i)}和\textit{(ii)}蕴含，$\mathfrak X_0$ 上每个局部同构于
结构层的凝聚层，都可以从 $\mathfrak X$ 上的一个同类层经约化得到。
从 $\mathfrak X_0$ 上的一个\fquote{丰沛}层出发（由\textit{(iii)}知这样的
层存在），把它提升为 $\mathfrak X$ 上的一个可逆层；再用定理1证明，
该层的某个张量幂定义了 $\mathfrak X$ 到某个概形 $\mathbf P_A^r$ 沿闭集的
形式完备化中的一个浸入（$\mathbf P_A^r$ 是 $A$ 上维数为 $r$ 的
\fquote{标准射影}概形）。

定理1至定理3的证明见[1]。

\src{199}
\fsection{4}{连通性定理与 Zariski\fquote{主定理}的应用}

设 $f:X\longrightarrow Y$ 是概形的固有态射。由有限性定理，
$f_*(\Osh_X)=\mathcal A$ 是 $Y$ 上的凝聚层；它还是一个交换代数层，
因而对应于一个在 $Y$ 上有限的 $Y$-概形
$g:Y'\longrightarrow Y$（其定义条件是在 $Y$ 上仿射，即一个仿射开集的
逆像仍为仿射，并且 $g_*(\Osh_{Y'})=\mathcal A$）。立即可见，$f$ 典范地
分解为 $f=gf'$，其中 $f':X\longrightarrow Y'$ 是从 $X$ 到 $Y'$ 的态射，
并且这一次满足 $f'_*(\Osh_X)=\Osh_{Y'}$。$f$ 的这个分解称为 $f$ 的
\emph{Stein 分解}。把第一比较定理的推论1应用于 $f'$，并取 $Y'$ 中仅含
一点 $y'$ 的部分，便知 $f'^{-1}(y')=X'$ 是连通的（否则，$X$ 沿 $X'$ 的
形式截面不会构成局部环，而 $f'_*(\Osh_X)_{y'}=\Osh_{Y',y'}$ 的完备化
却是局部的！）。我们已经证明：

\result{定理}{5（Zariski\fquote{连通性定理}）} 设
$f:X\longrightarrow Y$ 是固有态射，则 $f$ 可唯一地（不计同构）分解为
$f=gf'$，其中 $g:Y'\longrightarrow Y$ 是有限态射，并且
$f':X\longrightarrow Y'$ 满足 $f'_*(\Osh_X)=\Osh_{Y'}$（从而
$g_*(\Osh_{Y'})=f_*(\Osh_X)$）。$f'$ 的各纤维都是连通的；也就是说，
$f$ 的纤维 $f^{-1}(y)$ 的连通分支集合，与 $Y'$ 中位于 $y$ 上方的点集
一一对应，也即与 $f_*(\Osh_X)_y$ 中的极大理想集合一一对应。

连通性定理的各个通常变体立即随之得到。这里只记下下面的推论。

\src{200}
\result{推论}{1} $X$ 的点 $x$ 在其纤维 $f^{-1}(y)$ 中孤立，当且仅当纤维
$f'^{-1}(y')$（其中 $y'=f'(x)$）只有点 $x$；或者等价地，$f'$ 在 $x$ 的
某个邻域与 $y'$ 的某个邻域之间诱导同构。这类点的集合是一个开集 $U$，
且 $f'$ 把 $U$ 同构地映到 $Y'$ 的一个开集上。

为证明 $f'$ 在 $x$ 处是局部同构，注意 $f'$ 诱导同构
$\Osh_{y'}\longrightarrow\Osh_x$；这是由 $f'_*(\Osh_X)=\Osh_{Y'}$，以及
下述事实看出的：$f'$ 是闭映射且在 $y'$ 上的纤维只有 $x$，所以当 $V$
遍历 $y'$ 的一个基本邻域系时，各 $f'^{-1}(V)$ 遍历 $x$ 的一个基本邻域系。
由此立即得到下面的结论；在\fquote{几何}情形中，这个结论归功于 Chevalley。

\result{推论}{2} 态射 $f$ 是有限态射，当且仅当它是固有态射并且其纤维
都是有限的。

在这种情形中，由上述结果，$f'$ 的确是同构。

设 $f:X\longrightarrow Y$ 是一个不必固有的态射，但假设 $X$ 作为开集包含
在一个固有 $Y$-概形 $\widetilde f:\widetilde X\longrightarrow Y$ 中
（特别地，当 $f$ 拟射影时便是如此）。应用推论1可知，$\widetilde f'$
把 $X$ 中在各自纤维内孤立的点所成的集合 $U$ 同构地映到 $Y'$ 的一个
开集上（并且 $U$ 的确是开集）。由此得到 Zariski\fquote{主定理}的如下
整体形式。

\result{定理}{6（Zariski\fquote{主定理}）} 设
$f:X\longrightarrow Y$ 是有限型态射。则 $X$ 中在各自纤维内孤立的点所成
的集合 $U$ 是开集；若 $f$ 分离，\footnote{1962年的勘误在此把印出的假设
\fquote{$f$ 拟射影}改为较弱的假设\fquote{$f$ 分离}；依据 SGA VIII，6.2，
每个拟有限且分离的态射都是拟射影的。} 则 $U$ 在 $Y$ 上同构于某个在
$Y$ 上有限的概形 $Y'$ 的一个开部分。

由于有限型态射局部仿射，因而\textit{更不消说}局部拟射影，Zariski 主定理
的各个通常局部形式立即由定理6推出。

\fsection{5}{固有平坦态射的上同调研究}

设 $f:X\longrightarrow Y$ 是固有态射，$F$ 是 $X$ 上的凝聚层，并假设
$F$ 在 $Y$ 上平坦，即各 $F_x$ 都是环 $\Osh_y$（$y=f(x)$）上的平坦模。
这也意味着，对每个 $y\in Y$，若沿纤维 $f^{-1}(y)$ 用
$\mathfrak m_y^nF$ 对 $F$ 作滤过（$\mathfrak m_y$ 是 $\Osh_y$ 的极大
理想），则相伴分次对象同构于
$(F/\mathfrak m_yF)\otimes_{\kappa(y)}\operatorname{gr}(\Osh_y)$；换言之，
对每个整数 $n$ 都有
\[
  \mathfrak m_y^nF/\mathfrak m_y^{n+1}F
  =F_y\otimes_{\kappa(y)}
  (\mathfrak m_y^n/\mathfrak m_y^{n+1})
\]
这里 $X_y$ 表示把纤维 $f^{-1}(y)$ 看成 $y$ 的剩余域 $\kappa(y)$ 上的
固有概形，而 $F_y$ 表示 $F$ 在 $X_y$ 上诱导的层
$F/\mathfrak m_yF$。结合这个同构和定理2，由 $X_y$ 关于系数层 $F_y$ 的
上同调，可以得到 $y$ 附近各 $R^qf_*(F)$ 的估计，有时还可作出计算。
定理2在这里具有形式
\[
  \widehat{R^qf_*(F)_y}
  =\varprojlim_n H^q(X_y,F/\mathfrak m_y^nF).
\]
这里只指出下面这个推论：

\src{201}
\result{命题}{1} 设 $f:X\longrightarrow Y$ 是固有态射，$F$ 是 $X$ 上在
$Y$ 上平坦的凝聚层。设 $y\in Y$，$q$ 是一个整数，并假设
$H^q(X_y,F_y)=0$。则 $R^qf_*(F)$ 在 $y$ 的一个邻域内为零，而且对每个
$n$，自然同态
\[
  R^{q-1}f_*(F)_y\longrightarrow
  H^{q-1}(X_y,F/\mathfrak m_y^nF)
\]
都是满射。

特别地，若 $f$ 是平坦态射（即 $\Osh_X$ 在 $Y$ 上平坦），则 $X$ 上每个
局部自由凝聚层 $F$ 都在 $Y$ 上平坦。设 $F,G$ 是两个这样的层，把命题1
应用于 $\Homsh_{\Osh_X}(F,G)$ 并取 $q=1$，得到：

\result{定理}{7} 设 $f$ 是固有平坦态射，$F,G$ 是 $X$ 上的两个局部自由
凝聚层，$y\in Y$，并假设
\[
  H^1\bigl(X_y,\Homsh_{\Osh_{X_y}}(F_y,G_y)\bigr)=0,
\]
则每个同态 $u_0:F_y\longrightarrow G_y$ 都由一个同态
\[
  u:F|V\longrightarrow G|V,
  \qquad V=f^{-1}(U),
\]
诱导，其中 $U$ 是 $y$ 的一个邻域。

\result{推论}{} 若 $u_0$ 是同构（相应地，是单态射或满态射），则当 $U$
充分小时，$u$ 也具有同一性质。

特别地：

\src{202}
\result{推论}{2} 设 $E_0$ 是 $X_y$ 上的局部自由凝聚层，并且
\[
  H^1\bigl(X_y;\Homsh_{\Osh_{X_y}}(E_0,E_0)\bigr)=0,
\]
则两个局部自由层如果在 $X_y$ 上的限制都同构于 $E_0$，它们就在 $X_y$
的一个邻域内同构。

于是：

\result{推论}{3} 假设 $H^1(X_y,\Osh_{X_y})=0$。则 $X$ 上两个可逆层
（即局部同构于 $\Osh_X$ 的层），只要它们在 $X_y$ 上的限制同构，它们
自身就同构。

由此得到：

\result{命题}{2} 设 $Y$ 是连通概形，$E$ 是 $Y$ 上的局部自由凝聚层。
考察与 $E$ 相伴的射影空间丛 $X=\mathbf P(E)$，并赋以其熟知的可逆层
$\Osh_X(1)$。则 $X$ 上的每个可逆层 $L$ 都同构于形如
\[
  f^*(L')\otimes\Osh_X(n),
\]
的层，其中 $L'$ 是 $Y$ 上的可逆层，$n$ 是整数。这个整数 $n$ 唯一确定，
而 $L'$ 在同构意义下唯一确定。

上述推论3证明，在每条纤维的一个邻域内，$L$ 都同构于某个
$\Osh_X(n)$。其余部分几乎是形式的。

借助命题2，可以确定从 $X=\mathbf P(E)$ 到另一个射影丛的 $Y$-态射。
特别地得到：

\result{推论}{} 设 $u$ 是 $X=\mathbf P(E)$ 的一个自同构。则存在 $Y$
上的一个可逆层 $L'$ 和一个同构 $v:E\longrightarrow E\otimes L'$，使得
$u$ 是相应的同构
\[
  \mathbf P(E)\xrightarrow{\sim}\mathbf P(E\otimes L')=\mathbf P(E);
\]
二元组 $(v,L')$ 在同构意义下唯一确定。

设 $\Gamma$ 是 $Y$ 上满足 $E\otimes L'\simeq E$ 的可逆丛 $L'$ 的同构类
所成的集合。它的元素都是挠元：事实上，若 $n$ 是 $E$ 的秩，取第 $n$
外幂便知必有 $L'^{\otimes n}\simeq\Osh_Y$。因此，上述推论也可以表述为
存在群的正合列
\[
  e\longrightarrow
  \operatorname{Aut}(E)/\Gamma(Y,\Osh_Y^*)
  \longrightarrow \operatorname{Aut}_Y(X)
  \longrightarrow \Gamma\longrightarrow e
\]
（此外，它也可以由下面这个群层正合列所诱导的上同调正合列推出：
\[
  e\longrightarrow \Osh_Y^*
  \longrightarrow \underline{\operatorname{Aut}}(E)
  \longrightarrow \underline{\operatorname{Aut}}_Y(X)
  \longrightarrow e,
\]
其中 $\Osh_Y^*$ 是 $\Osh_Y$ 的\fquote{可逆元}层，并等同于
$\underline{\operatorname{Aut}}(E)$ 的中心）。

\src{203}
\fsection{6}{完备 $J$-进环上层与概形的存在性及唯一性定理的应用}

定理7利用定理1和定理2给出了局部自由凝聚层的一个唯一性结论。现在利用
定理3，我们将推出层、概形态射或概形的存在性定理。下文中，$A$ 表示一个
分离且完备的诺特局部环。一般方法始终是先作形式构造；这实质上就是在
阿廷环上作代数几何，然后借助三个基本定理推出具有\fquote{代数}性质的
结论。

\result{命题}{3} 设 $\mathfrak X$ 是 $A$ 上的固有平坦形式概形，$F_0$
是 $X_0$ 上的局部自由层，并满足
\[
  H^2\bigl(X_0,\Homsh_{\Osh_{X_0}}(F_0,F_0)\bigr)=0.
\]
则 $\mathfrak X$ 上存在一个局部自由层 $F$，它在 $X_0$ 上诱导的层同构于
$F_0$。（此外，若
$H^1(X_0,\Homsh_{\Osh_{X_0}}(F_0,F_0))=0$，则 $F$ 在同构意义下唯一。）

逐次在各 $X_n$ 上构造相互诱导的局部自由层 $F_n$。构造 $F_n$ 时遇到的
障碍位于
\[
  H^2\bigl(X_0,\Homsh_{\Osh_{X_0}}(F_0,F_0)\bigr)
  \otimes_{A/J}(J^n/J^{n+1}),
\]
所以按假设该障碍为零。现在使用定理3，得到：

\result{推论}{1} 设 $X$ 是 $A$ 上的固有平坦概形，$F_0$ 如上。则 $X$
上存在一个局部自由层 $F$，它在 $X_0$ 上诱导的层同构于 $F_0$。此外，
若 $H^1(X_0,\Homsh_{\Osh_{X_0}}(F_0,F_0))=0$，则 $F$ 在同构意义下唯一。

设 $X_0$ 是域 $k$ 上的有限型概形；假设 $X_0$ 在 $k$ 上非奇异（这里指
绝对非奇异），但不必在 $k$ 上固有。设 $A$ 是剩余域为 $k$ 的阿廷局部环。
我们要寻找在 $A$ 上平坦并且满足 $X\otimes_A k=X_0$ 的概形 $X$
（这是 $X_0$ 的\fquote{模理论}或\fquote{结构变形}的出发点）。给出这样
一个 $X$，等价于在拓扑空间 $X_0$ 上给出一个环层 $\Osh_X$，并赋予如下
结构：
\begin{enumerate}
  \item[\textit{i.}] 它是一个 $A$-代数层。
  \item[\textit{ii.}] 它带有一个增广同态
    $\Osh_X\longrightarrow\Osh_{X_0}$（与 $A$-代数结构相容）；这些数据
    还须满足如下条件：该增广诱导同构
    $\Osh_X\otimes_A k\xrightarrow{\sim}\Osh_{X_0}$；$\Osh_X$ 在 $A$
    上平坦，也就是说，用 $A$ 的极大理想 $\mathfrak m$ 的各次幂对
    $\Osh_X$ 作滤过后所得的相伴分次对象，同构于
    $\operatorname{gr}^0(\Osh_X)\otimes_k\operatorname{gr}(A)$；亦即有同构
    \[
      \mathfrak m^n\Osh_X/\mathfrak m^{n+1}\Osh_X
      =\Osh_{X_0}\otimes_k(\mathfrak m^n/\mathfrak m^{n+1}).
    \]
\end{enumerate}
基本事实如下：

\src{204}
\result{定理}{8} 设 $X_0$ 是域 $k$ 上有限型且非奇异的概形，并假设
$X_0$ 仿射。设 $A$ 是剩余域为 $k$ 的阿廷局部环。则存在一个在 $A$
上平坦的 $A$-概形 $X$，使得 $X\otimes_A k=X_0$；任意两个这样的概形
必定同构。

这里的同构并非典范的，因为一般而言，$X$ 会有在 $X_0$ 上诱导恒等的
非平凡 $A$-自同构。另一方面，一般不存在满足所给条件的 $X$ 的
\fquote{典范}选择；唯一的例外是 $A$ 为 $k$-代数（等特征情形），此时
可以取 $X=X_0\otimes_k A$，即
$\Osh_X=\Osh_{X_0}\otimes_k A$（事实上无论 $X_0$ 是否仿射都可以）。
在混合特征情形中，当 $X_0$ 不仿射时，我一般不知道能否把 $X_0$
\fquote{提升}为定义在 $A$ 上的 $X$。不过，设 $n>0$ 是整数，令
$A_{n-1}=A/\mathfrak m^n$，并假设已经把 $X_0$ 提升为一个平坦
$A_{n-1}$-概形 $X_{n-1}$；现在要把 $X_{n-1}$ 提升为一个平坦
$A_n$-概形 $X_n$。由定理8已知这在局部是可能的；另一方面容易验证，
若 $U_n$ 提升了 $X_{n-1}$ 的开集 $U_{n-1}$，则 $U_n$ 的自同构群层
（在 $U_{n-1}$ 上诱导恒等的那些自同构）典范同构于
\[
  \Theta_{X_0/k}\otimes_k(\mathfrak m^n/\mathfrak m^{n+1})
\]
在 $U_{n-1}$ 上的限制，因而特别是交换的（$\Theta_{X_0/k}$ 表示 $X_0$
上的 $k$-导子芽层）。由此容易看出，构造提升 $X_{n-1}$ 的 $X_n$ 时有
一个障碍，它位于
\[
  H^2(X_0,\Theta_{X_0/k})\otimes_k
  \mathfrak m^n/\mathfrak m^{n+1}.
\]
因此：

\result{推论}{1} 设 $X_0$ 是 $k$ 上有限型且非奇异的概形，并假设
$H^2(X_0,\Theta_{X_0/k})=0$。则对每个剩余域为 $k$ 的阿廷局部环 $A$，
都存在一个在 $A$ 上平坦的 $A$-概形 $X$，使得
$X\otimes_A k=X_0$。

此外，如果已经找到一个在 $A$ 上平坦并提升 $X_0$ 的 $X$，那么由定理8，
提升 $X_0$ 的平坦 $A$-概形的同构类集合可与
\[
  H^1\bigl(X_0,\underline{\operatorname{Aut}}(X)\bigr),
\]
等同；这里 $\underline{\operatorname{Aut}}(X)$ 当然表示 $A$-代数层
$\Osh_X$ 的、与增广相容的自同构芽层。$\Osh_X$ 的滤过定义了
$\underline{\operatorname{Aut}}(X)$ 的一个滤过；这个层对第 $n$ 个滤过
子群的商为 $\underline{\operatorname{Aut}}(X_n)$。该滤过的相伴分次对象
是交换的，并可与 $\Theta_{X_0/k}\otimes_k\operatorname{gr}(A)$ 等同。
特别地，若 $\mathfrak m^{n+1}$ 是 $\mathfrak m$ 的第一个为零的幂，则
$F^n(\underline{\operatorname{Aut}}(X))$（滤过的最后一项）包含在
$\underline{\operatorname{Aut}}(X)$ 的中心中，并同构于
$\Theta_{X_0/k}\otimes_k\mathfrak m^n$；它也正是 $X$ 的、在
$X_{n-1}=X\otimes_A A/\mathfrak m^n$ 上诱导恒等的自同构芽层。利用这些
结果，立即得到以下命题：

\src{205}
\result{推论}{2} 设 $X_0$ 是 $k$ 上有限型且非奇异的概形，$A$ 是剩余域
为 $k$、极大理想为 $\mathfrak m$ 的阿廷局部环；假设
$\mathfrak m^{n+1}=0$，令 $A_{n-1}=A/\mathfrak m^n$；再设 $X_{n-1}$
是一个满足 $X_{n-1}\otimes_A k=X_0$ 的平坦 $A_{n-1}$-概形。则所有满足
$X_n\otimes_A A_{n-1}=X_{n-1}$ 的平坦 $A$-概形 $X_n$ 的同构类
（所取同构须在 $X_{n-1}$ 上诱导恒等）所成的集合，或者为空，或者是群
\[
  H^1(X_0,\Theta_{X_0/k})\otimes_k\mathfrak m^n.
\]
作用下的主齐性空间。

（注意，一般而言，这个空间没有特选的原点，因为没有特选的方式把
$X_{n-1}$ 提升为 $X_n$。）

\result{推论}{3} 设 $X_0$ 是 $k$ 上有限型且非奇异的概形，并假设
$H^1(X_0,\Theta_{X_0/k})=0$。则对每个剩余域为 $k$ 的阿廷局部环 $A$，
至多存在一个（不计同构）满足 $X\otimes_A k=X_0$ 的平坦 $A$-概形 $X$。

推论1和推论3立即蕴含表面上更一般的结论：只须假设 $A$ 是剩余域为 $k$
的完备诺特局部环，并把 $X$ 作为 $A$ 上的形式概形引入即可：

\result{定理}{9} 设 $k$ 是域，$X_0$ 是 $k$ 上有限型且非奇异的概形。
对每个剩余域为 $k$ 的完备诺特局部环 $A$，令 $F(A)$ 表示所有满足
$\mathfrak X\otimes_A k=X_0$、在 $A$ 上有限型且平坦的形式概形
$\mathfrak X$ 的同构类集合（所取同构须在 $X_0$ 上诱导恒等）。采用这些
记号：
\begin{enumerate}
  \item[\textit{i.}] 若 $H^1(X_0,\Theta_{X_0/k})=0$，则对每个 $A$，
    $F(A)$ 至多有一个元素。
  \item[\textit{ii.}] 若 $H^2(X_0,\Theta_{X_0/k})=0$，则对每个 $A$，
    $F(A)$ 至少有一个元素。
\end{enumerate}

\src{206}
\result{推论}{1} 假设 $X_0$ 在 $k$ 上固有。在条件\textit{(i)}下，对每个
$A$，至多存在一个（不计在 $X_0$ 上诱导恒等的同构）满足
$X\otimes_A k=X_0$、在 $A$ 上固有且平坦的概形 $X$。

这里使用定理3的推论2。例如：

\result{推论}{2} 若 $X$ 是 $A$ 上的固有平坦概形，并且
$X\otimes_A k$ 同构于 $k$ 上维数为 $r$ 的标准射影概形
$\mathbf P_k^r$，则 $X$ 同构于 $\mathbf P_A^r$。

（也可由命题3的推论1推出这个结果。）

\result{推论}{3} 设 $X_0$ 是 $k$ 上的非奇异射影概形，并假设
\[
  H^2(X_0,\Osh_{X_0})=H^2(X_0,\Theta_{X_0/k})=0.
\]
则对每个 $A$，都存在一个在 $A$ 上射影且平坦的概形 $X$，使得
$X\otimes_A k=X_0$。

结合定理9\textit{(ii)}和定理4即可。特别地：

\result{推论}{4} 设 $X_0$ 是 $k$ 上一条完备非奇异代数曲线的概形。
则对每个剩余域为 $k$ 的完备诺特局部环 $A$，都存在一个 $A$ 上的
\fquote{光滑曲线概形} $X$，使得 $X\otimes_A k=X_0$。

\result{注记}{}

1）推论3和推论4主要在 $k$ 的特征为 $p\ne0$ 时有意义：可取 $A$ 为特征
为 $0$、剩余域为 $k$ 的离散赋值环，例如取\fquote{最小的可能 $A$}，
即其极大理想由 $p$ 生成的环。（事实上，由 Cohen 定理，只须对一个这样的
环 $A$ 知道推论3和推论4即可。）在这方面应当指出，据专家所知，目前尚
不知道域 $k$ 上是否存在不能由定义在这类环 $A$ 上的平坦概形模 $p$ 约化
得到的概形。至少，本节的结果提供了系统研究这一问题的方法。首先应当
考察 $H^2(X_0,\Theta_{X_0/k})$ 中的第一个障碍是否必然为零。

\emph{1962年补记。}---J.-P. Serre 构造了特征 $p>0$ 的代数闭域 $k$
上的一个三维非奇异射影簇；它并非某个固有概形的约化，而该固有概形定义
在一个剩余域为 $k$、分式域特征为零的整局部环上。据说 Mumford 找到了
一个以非奇异射影曲面为例的类似结果。

2）应注意，定理3及相应的方法只对完备基环（为具体起见，可设为局部环）
有效。

要从局部环的完备化上已知的结果过渡到该局部环自身的相应结果，还需要
第四个\fquote{基本定理}，而它的最终表述仍有待找到。

\src{207}
3）应把本节的结果（特别是前面的推论1和推论2）以及下一节的结果，与
Kodaira--Spencer 关于复结构变形的结果相比较。利用刚才提到的猜想性定理，
在推论1的条件下但不再假设 $A$ 完备，应当可以断言：存在一个包含 $A$、
在 $A$ 上有限且非分歧的环 $A'$，使
$X\otimes_A A'$ 与 $X'\otimes_A A'$ 在 $A'$ 上同构（这里 $X,X'$ 是
两个给定的固有平坦 $A$-概形，并满足
$X\otimes_A k=X'\otimes_A k=X_0$）。至少当 $X_0=\mathbf P_k^r$ 时，
利用命题2的推论可以证明这一点。无论如何，当
$H^1(X_0,\Theta_{X_0/k})=0$ 时，可以证明 $X,X'$ 在
$Y=\operatorname{Spec}(A)$ 的每一点 $y$ 上的纤维同构，至少在把剩余域
$\kappa(y)$ 换成其代数闭包以后如此。（还有一个表面上更强的局部结果，
其中不再需要假设 $A$ 是局部环。）关于射影空间的\fquote{结构变形}，
再指出一个由 Kodaira--Spencer 的相应问题所启发的问题。设 $X$ 是整局部
环 $A$ 上的固有平坦概形，$A$ 的分式域为 $K$、剩余域为 $k$；假设
$X\otimes_A K$ 同构于 $\mathbf P_K^r$，那么
$X\otimes_A k=X_0$ 是否同构于 $\mathbf P_k^r$（至少在 $k$ 的代数闭包
上如此）？在这个问题中，可以假设 $A$ 是完备离散赋值环。当 $X_0$ 是
阿贝尔簇时也有类似的问题。

\emph{1962年补记。}---此时报告人对上述猜想性结果已不那么乐观。不过，
射影空间结构变形的问题已由 Hironaka 肯定地解决，阿贝尔簇的相应问题
则由 Koizumi 解决。

\fsection{7}{模理论的应用}

报告人本人也才刚开始研究这一理论，因此这里只能给出一些说明。为简单
起见，取基底为域 $k$，也就是说只讨论等特征情形；不过定理8看来也能在
没有本质变化的情况下处理一般情形。目前我们还没有越过\fquote{形式}
阶段；但报告人希望在某些情形中能够由此构造真正的模概形，特别是对每个
整数 $g$，在整数环上构造一个概形，使其扮演亏格为 $g$ 的非奇异曲线的
万有模概形。

\emph{1962年补记。}---Mumford 已构造了亏格为 $g$ 的曲线的模概形
（参见 SHMI）。此外，定理10还证明，模理论中的\fquote{$n$ 级 Jacobi
概形}是非奇异的，甚至在 $\mathbf Z$ 上光滑。

\src{208}
回到定理9的情形和记号，现在假设 $A$ 是 $k$ 上的有限维局部代数；为简单
起见，还假设 $k$ 代数闭。于是可把 $F(A)$ 看成关于 $A$ 的、取值于集合
范畴的协变函子：$k$-代数同态 $A\longrightarrow B$ 定义映射
$F(A)\longrightarrow F(B)$，因为每个满足 $X\otimes_A k=X_0$ 的平坦
$A$-概形 $X$ 都产生具有同样性质的 $B$-概形 $X\otimes_A B$。假设可以
找到一个完备诺特局部 $k$-代数 $\mathcal O$ 和一个函子性同构
\begin{equation}
  \tag{*}
  \operatorname{Hom}(\mathcal O,A)\xrightarrow{\sim}F(A)
\end{equation}
（左端表示 $k$-代数同态）。容易看出，这样的 $\mathcal O$ 在典范同构的
意义下唯一确定；这时称 $\mathcal O$ 的形式谱 $\mathfrak Y$（即只含一个
点的拓扑空间，并赋以值为 $\mathcal O$ 的拓扑环层）为 $X_0$ 的
\emph{形式模概形}。（它不一定存在。）设 $\mathfrak r$ 是 $\mathcal O$
的极大理想；对每个 $n$ 令
$\mathcal O_n=\mathcal O/\mathfrak r^{n+1}$（所以 $\mathcal O_0=k$）。
典范同态 $\mathcal O\longrightarrow\mathcal O_n$ 是
$\operatorname{Hom}(\mathcal O,\mathcal O_n)$ 的一个元素，因而定义
$F(\mathcal O_n)$ 的一个元素，也就是一个平坦 $\mathcal O_n$-概形
$X_n$，其模 $\mathfrak r$ 约化为 $X_0$。这些 $X_n$ 通过标量扩张
（在这里即约化）相互诱导，因此它们来自形式模概形 $\mathfrak Y$ 上一个
唯一确定的形式概形 $\mathfrak X$；$\mathfrak X$ 在 $\mathfrak Y$ 上
平坦，并且 $\mathfrak X_0=X_0$。立即可见，同构(*)是这样给出的：对每个
$k$-代数同态 $\mathcal O\longrightarrow A$，对应 $A$-概形
$\mathfrak X\otimes_{\mathcal O}A$ 的同构类（也就是说，对每个 $k$-概形
态射 $Y'=\operatorname{Spec}(A)\longrightarrow\mathfrak Y$，对应由基变换
得到的 $Y'$-概形 $\mathfrak X\otimes_{\mathfrak Y}Y'$）。进一步，即使
只假设 $A$ 是完备诺特局部 $k$-代数（而不必是阿廷的），同构(*)及上述
描述仍然成立。当然，照例 $\mathcal O$\textit{先验地}完全可能含有幂零元；
看来很可能存在 $\mathcal O$ 本身是阿廷环而又不等于 $k$ 的情形。这说明
Kodaira--Spencer 的观点（只取正则环 $A$）\textit{先验地}多么不适合一般
情形。

\src{209}
现在还须给出在 $X_0$ 于 $k$ 上固有时，其形式模概形存在的充分条件。一般
而言，对一个从有限维局部 $k$-代数到集合的函子 $A\longmapsto F(A)$，
容易给出简单的充要条件，使它能够写成适当 $\mathcal O$ 的
$\operatorname{Hom}(\mathcal O,A)$。这里不详述这些条件。只指出，在我们
所考虑的情形中，这些条件对 $X_0$ 施加了非平凡的上同调条件；虽然报告人
尚未构造出反例，但它们看来不太可能总是满足。另一方面，
$H^0(X_0,\Theta_{X_0/k})=0$ 看来很可能是形式模概形存在的一个充分条件
（虽然绝非必要条件）。这里只陈述一个特别简单情形中的定理（它在解析
空间理论中的类似结果是熟知的，参见 Kodaira--Spencer）；借助上一节的
结果便可毫无困难地证明它：

\result{定理}{10} 设 $X_0$ 是域 $k$ 上的固有非奇异概形，并且
\[
  H^0(X_0,\Theta_{X_0/k})=H^2(X_0,\Theta_{X_0/k})=0.
\]
则 $X_0$ 存在形式模概形，对应于一个正则局部环 $\mathcal O$（即 $k$
上的形式幂级数代数）。

如前所述，$\mathcal O$ 上的形式概形 $\mathfrak X$ 一般并不一定能够
代数化；但已知当 $X_0$ 射影且 $H^2(X_0,\Osh_{X_0})=0$ 时确能代数化
（定理4），例如当 $X_0$ 的维数为1时。由此，构造给定亏格曲线的整数环上
模概形的希望多少得到了一些支持\ldots{}

还应注意，本节所述方法也可以用于 Picard 簇的构造和研究，以及许多其他
构造。我们不久将再讨论这些问题。

\fsection{8}{基本群的应用}

上述方法使我们能够仿照拓扑理论，对基本群展开系统研究。本节陈述的前
两个定理推广了 Lang--Serre 一篇新近论文中的结果。

设 $X$ 是概形。若一个 $X$-概形 $X'$ 满足：
\begin{enumerate}
  \item[\textit{i.}] $X'$ 在 $X$ 上有限，也就是说，它由 $X$ 上的一个
    凝聚代数层 $\mathcal A=\mathcal A(X')$ 定义；
  \item[\textit{ii.}] $\mathcal A$ 是 $X$ 上的局部自由层；
  \item[\textit{iii.}] 对每个 $x\in X$，商
    \[
      \mathcal A_x/\mathfrak m_x\mathcal A_x
      =\mathcal A_x\otimes_{\Osh_x}\kappa(x)
    \]
    是 $\kappa(x)$ 上的可分代数，
\end{enumerate}
则称 $X'$ 是 $X$ 的一个\emph{非分歧覆盖}。

\src{210}
这个非分歧覆盖概念（归功于 Serre 和报告人）具有一切可以合理期待的
初等性质，这里不逐一列举。只须指出，它产生了一套仿照经典 Galois 理论
（并包含后者；其证明还比后者通常采用的证明更为简单）以及拓扑覆盖的
Galois 理论而建立的 Galois 理论。更确切地说，概形 $X$ 的一个
\emph{几何点}，是指从代数闭域 $\Omega$ 的谱 $\xi$ 到 $X$ 的态射 $a$；
也就是说，给定 $X$ 的一点 $x=|a|$ 的剩余域 $\kappa(x)$ 的一个代数闭
扩张（称 $x$ 为几何点 $a$ 的\emph{所在点}）。若 $X'$ 是 $X$ 的非分歧
覆盖，就可为它配上\fquote{$X'$ 中位于 $a$ 上方的几何点}所成的集合
$E_a(X')$，也就是由如下二元组组成的集合：位于 $x$ 上方的一点
$x'\in X'$，以及一个 $\kappa(x)$-同态
$\kappa(x')\longrightarrow\Omega$。这样，在固定 $(X,a)$ 后，便得到一个
函子 $F(X,a)$，它从 $X$ 的非分歧覆盖 $X'$ 所成的范畴 $R(X)$ 到有限集
范畴。若 $X$ 连通，则 $R(X)$ 与 $F(X,a)$ 所成的二元组具有所需的形式
性质，因而同构于由某个适当的紧全不连通拓扑群 $\pi$（即有限群的逆极限）
定义的类似二元组：取 $\pi$ 连续作用的有限集 $E$ 所成的范畴 $C(\pi)$，
以及从该范畴到有限集范畴的遗忘函子 $F(\pi)(E)=E$。此外，由
$(C(\pi),F(\pi))$ 与给定二元组同构这一条件，$\pi$ 在典范同构的意义下
唯一确定。在当前情形中，$\pi$ 称为连通概形 $X$ 在几何点 $a$ 处的
\emph{基本群}，记作 $\pi_1(X,a)$。若 $X$ 不连通，则以 $x=|a|$ 所在的
连通分支替代 $X$。若 $X$ 连通，则对 $X$ 的两个几何点 $a',a''$，群
$\pi_1(X,a')$ 与 $\pi_1(X,a'')$ 同构（该同构只在相差内自同构的意义下
确定）；因此照例可以按问题需要选择 $a$，例如当 $X$ 不可约时取其一般点。
当然，$\pi_1(X,a)$ 是关于带点概形 $(X,a)$ 的协变函子。关于非分歧覆盖
分类的每个命题\footnote{印刷文本作\fquote{不可分覆盖}；但本段讨论的是
非分歧覆盖范畴 $R(X)$ 以及它同带连续 $\pi_1(X,a)$-作用的有限集范畴之间
的等价。}，都可以依照熟知的词典翻译成群论语言（只是这里还须顾及所涉
群为拓扑群这一事实）。

\src{211}
我们的目标是对固有态射 $f:X\longrightarrow Y$ 建立纤维空间同伦正合列的
类似物。显然，由于尚不知道高阶同伦群应当是什么，所得结果必然不完整。
为应用第3节的基本定理，首先必须明确说明关于域或阿廷环上概形的几个
初等引理（这与一般方法完全一致！）。

\result{引理}{1} 设 $(X',a')$ 是一个带点非分歧覆盖，它对应于
$\pi_1(X,a)$ 在有限集 $E$（带有标记点 $e$）上的一个带点表示。则典范
同态 $\pi_1(X',a')\longrightarrow\pi_1(X,a)$ 把前一个群等同于
$\pi_1(X,a)$ 中 $e$ 的稳定子（因而它是单射）。

\result{引理}{2} 设 $X$ 是域 $k$ 上的代数概形，$k'$ 是 $k$ 的纯不可分
扩张。则 $X\otimes_k k'$ 的每个非分歧覆盖都由 $X$ 的某个非分歧覆盖取
逆像（即扩张标量）而来，且后者在典范同构的意义下唯一确定。

特别地，由这两个引理可知，对 $k$ 的每个代数扩张 $K$，以及
$X'=X\otimes_k K$ 中映到 $X$ 的几何点 $a$ 的每个几何点 $a'$，函子性
同态 $\pi_1(X',a')\longrightarrow\pi_1(X,a)$ 都是单射。

\result{引理}{3} 设 $X$ 是阿廷局部环 $A$ 上的完备概形，并满足
$H^0(X,\Osh_X)=A$；设 $X'$ 是 $X$ 的非分歧覆盖，并令
$A'=H^0(X',\Osh_{X'})$，所以 $A'$ 是 $A$ 上的有限环（它
\textit{先验地}可能在 $A$ 上分歧）。设 $X_0,X'_0$ 分别为与 $X,X'$
相伴的约化子概形（即分别把 $\Osh_X,\Osh_{X'}$ 除以其幂零元层而得）；
设 $k$ 是 $A/\mathfrak r(A)$ 的一个子域，使 $A/\mathfrak r(A)$ 在 $k$
上有限（于是 $X_0$ 是 $k$ 上的完备代数概形，而 $X'_0$ 是它的非分歧
覆盖）。最后，设 $\Omega$ 是 $k$ 的一个代数闭扩张，并考察
$X_0\otimes_k\Omega$ 的非分歧覆盖 $X'_0\otimes_k\Omega$。则以下两个
条件等价：
\begin{enumerate}
  \item[\textit{i.}] $X'_0\otimes_k\Omega$ 在
    $X_0\otimes_k\Omega$ 上完全分裂。
  \item[\textit{ii.}] 自然态射
    $X'\longrightarrow X\otimes_A A'$ 是同构。
\end{enumerate}
在这些条件下，$A'$ 是 $A$ 的非分歧扩张。最后，若 $X'$ 连通，则条件
\textit{(i)}等价于下面这个表面上较弱的条件：
\begin{enumerate}
  \item[\textit{i bis.}] $X'_0\otimes_k\Omega$ 在
    $X_0\otimes_k\Omega$ 上有一个正则截面。
\end{enumerate}
当条件\textit{(ii)}成立时，称 $X$ 的非分歧覆盖 $X'$ 是
\emph{几何平凡的}。

\src{212}
\result{引理}{4} 设 $f:X\longrightarrow Y$ 是固有态射，并满足
$f_*(\Osh_X)=\Osh_Y$。设 $a$ 是 $X$ 的一个几何点，$b$ 是它在 $Y$ 中
的像。则 $\pi_1(X,a)\longrightarrow\pi_1(Y,b)$ 是满射。

事实上，只须证明如下命题：若 $Y$ 的一个非分歧覆盖 $Y'$（对应于一个
局部自由代数层 $\mathcal A$）使得 $X\otimes_Y Y'$ 不连通，则 $Y'$ 也
不连通。的确，$f^*(\mathcal A)$ 将是两个非零环层的直和，因而其正像也
如此；但后一正像正是
$\mathcal A\otimes f_*(\Osh_X)=\mathcal A$。

\result{引理}{5} 设 $X$ 是域 $k$ 上的完备概形。假设
$H^0(X,\Osh_X)$ 是局部环 $A$，并且 $A/\mathfrak r(A)$ 在 $k$ 上纯不可分。
设 $\Omega$ 是 $k$ 的一个代数闭包，并令
$\overline X=X\otimes_k\Omega$（它是连通的）。选取 $\overline X$ 的一个
几何点 $\overline a$，使其映到 $X$ 的几何点 $a$，则有正合列
\[
  e\longrightarrow\pi_1(\overline X,\overline a)
  \longrightarrow\pi_1(X,a)
  \longrightarrow\pi_1(k,b)\longrightarrow e
\]
（其中 $\pi_1(k,b)$ 是 $\Omega$ 在 $k$ 上的 Galois 群）。

第一个同态的单射性已由引理1和引理2得到；中间的正合性来自引理3；最后
一个同态的满射性来自引理4（只有这一步用到
$A/\mathfrak r(A)$ 在 $k$ 上纯不可分）。

\result{命题}{4} 设 $f:X\longrightarrow Y$ 是固有平坦态射，并且对每个
$y\in Y$，$H^0(f^{-1}(y),\Osh_{f^{-1}(y)})$ 都是剩余域 $\kappa(y)$ 上的
可分代数（例如，当 $f^{-1}(y)$ 是 $\kappa(y)$ 上的可分概形时便是如此，
也就是说，它是约化的，并且与其各不可约分支相应的域都是 $\kappa(y)$ 的
可分扩张）。则与 $f_*(\Osh_X)$ 相伴的 $Y$ 的覆盖 $Y'$ 是非分歧的。

证明由定理2即可容易得到。

\src{213}
结合引理1，这个命题实际上把固有平坦态射（纤维可分）的同伦研究归结到
$f_*(\Osh_X)=\Osh_Y$ 的情形（因为利用 Stein 分解，可以用 $Y'$ 代替
$Y$）。

\result{注}{} 有限型平坦态射如果其纤维都是可分概形（相应地，光滑概形），
就称为\emph{可分的}（相应地，\emph{光滑的}）。可以证明，若 $f$ 平坦，
且 $f^{-1}(y)$ 可分（相应地，光滑），则存在 $f^{-1}(y)$ 的一个邻域，
$f$ 在其上可分（相应地，光滑）。对\fquote{绝对正规}也有同样的结果
（\fquote{Bertini 定理}）。

设 $f:X\longrightarrow Y$ 是固有态射，并满足
\begin{equation}
  \tag{i}
  f_*(\Osh_X)=\Osh_Y,
\end{equation}
且设 $X'$ 是 $X$ 上的有限概形。设 $Y'$ 是与
$X'\longrightarrow Y$ 的 Stein 分解相应的 $Y$ 的覆盖（参见定理5）。
取 $y\in Y$；于是 $X'$ 在 $y$ 上的纤维 $F'$ 的连通分支集合，可与 $Y'$
中位于 $y$ 上方的点 $y'$ 的集合等同（定理5）。考察由自然态射
$X'\longrightarrow X$ 和 $X'\longrightarrow Y'$ 所诱导的显然态射
\begin{equation}
  \tag{*}
  X'\longrightarrow X\times_Y Y'
\end{equation}
\footnote{印刷文本在第二个 $X'$ 上漏掉了撇号；到 $Y'$ 的态射是
$X'\to Y$ 的 Stein 分解中的态射。} 每当 $X'$ 形如 $X\times_Y Y''$，
其中 $Y''$ 是 $Y$ 的非分歧覆盖时，这个态射就是同构；此时 $Y'$ 正是
$Y''$，而(*)是恒等。我们要给出使 $X'$ 具有上述形式的条件，也就是使
$Y'$ 非分歧且(*)为同构的条件。为此，引入 $X$ 在 $y$ 上的纤维 $F$；
它是 $\kappa(y)$ 上的固有概形，而 $F'$ 是它的一个覆盖（若 $X'$ 非分歧，
则这个覆盖也非分歧）。设 $F'_1$ 是 $F'$ 的一个连通分支，它对应于 $Y'$
中位于 $y$ 上方的一点 $y'_1$。再假设：
\begin{enumerate}
  \item[\textit{(ii)}] $X'$ 在 $F'_1$ 的各点处在 $X$ 上非分歧（所以
    $F'_1$ 是 $F$ 的非分歧覆盖）；
  \item[\textit{(iii)}] $F'_1$ 是 $F$ 的几何平凡覆盖（参见引理3）。
\end{enumerate}

\result{定理}{11} 在这些条件下，$Y'$ 中存在 $y'_1$ 的一个开邻域 $U'$，
使(*)在 $U'$ 上是同构。此外，$Y'$ 在 $y'_1$ 处在 $Y$ 上非分歧（但在
$Y'$ 中位于 $y$ 上方的其他点 $y'$ 处可能分歧）。

当然，条件\textit{(ii)}和\textit{(iii)}也是定理结论成立的必要条件。
利用引理3和定理2即可容易地证明本定理。

\src{214}
\result{推论}{1} 仍假设\textit{(i)}成立。$X$ 的一个非分歧覆盖 $X'$
同构于 $Y$ 的某个非分歧覆盖 $Y'$ 的逆像，当且仅当 $X'$ 在每条纤维
$f^{-1}(y)$ 上诱导几何平凡覆盖。

由定理11，使这一条件成立的 $Y$ 中各点所成的集合是开集，所以只须在闭点
$y$ 上验证即可\ldots{} 还可注意到与推论1等价的如下表述：同态
$\pi_1(X)\longrightarrow\pi_1(Y)$（由引理4知它是满射）的核，是由各群
$\pi_1(\overline{f^{-1}(y)})$ 在 $\pi_1(X)$ 中的像所生成的闭正规子群，
其中
\[
  \overline{f^{-1}(y)}
  =f^{-1}(y)\otimes_{\kappa(y)}\overline{\kappa(y)},
\]
$\overline{\kappa(y)}$ 表示 $\kappa(y)$ 的一个代数闭包。应注意，由于不能
为所有纤维选取同一个基点，一旦为 $X$（从而也为 $Y$）选定基点，同态
$\pi_1(\overline{f^{-1}(y)})\longrightarrow\pi_1(X)$ 无论如何也只能在
复合 $\pi_1(X)$ 的内自同构的意义下确定。

\result{推论}{2} 在定理11的一般条件下，再假设 $Y,X,X'$ 都是整概形，
其函数域分别为 $K,L,L'$。则 $L'/K$ 中存在一个可分中间扩张 $K'$，它与
$L$ 线性无关，并且 $L'=LK'$（从而 $L'=L\otimes_K K'$）。

（把引理3的最后部分应用于 $X$ 的几何纤维。）定理11最有意义的应用情形
是 $f$ 为可分态射。此时 $X'$ 在 $Y$ 上也可分，所以由命题4，$Y'$ 在
$Y$ 上非分歧；因而(*)右端 $X\times_Y Y'$ 在 $X$ 上非分歧。由此容易
推出：

\result{推论}{3} 除\textit{(i)}外，再假设 $f$ 可分。设 $X'$ 是 $X$
的一个连通非分歧覆盖；则 $X'$\footnote{印刷文本此处作 $X$；这个判据的
主语是前一句引入的覆盖 $X'$。} 是 $Y$ 的某个非分歧覆盖 $Y'$ 的逆像，
当且仅当它在某条几何纤维 $\overline F=f^{-1}(y)$ 上诱导的覆盖
$\overline F'$ 有一个正则截面。

应注意，不必假设 $\overline F'$ 在 $\overline F$ 上几何平凡（这将在
\textit{事后}成立，尽管\textit{先验地}这个条件强得多）。推论3事实上
等价于以下命题：

\src{215}
\result{推论}{4} 设 $f:X\longrightarrow Y$ 是满足
$f_*(\Osh_X)=\Osh_Y$ 的固有可分态射。设 $\overline F$ 是 $Y$ 中一点
$y$ 的几何纤维，并在 $\overline F$ 中选取一个几何点；由态射
$\overline F\longrightarrow X\longrightarrow Y$ 得到 $X,Y$ 中的几何点，
并取作 $\overline F,X,Y$ 的基本群的基点。在这些条件下，有正合列
\[
  \boxed{
    \pi_1(\overline F)\longrightarrow\pi_1(X)
    \longrightarrow\pi_1(Y)\longrightarrow e
  }.
\]

由此容易推出 Serre--Lang 的下面两个命题；这里已完全除去了正规性假设：

\result{推论}{5} 设 $X,Y$ 是域 $k$ 上的连通概形。假设约化概形
$X_{\mathrm{red}}$ 在 $k$ 上可分（若 $k$ 完美，这一点自动成立），并且
$X$ 或 $Y$ 在 $k$ 上固有。\footnote{此条件由1962年的勘误加入。} 在
$X$（相应地，$Y$）中选取几何点 $a$（相应地，$b$），从而在
$X\times_kY$ 中得到几何点 $c=(a,b)$ 以及自然态射
\[
  \pi_1(X\times_kY,c)\longrightarrow
  \pi_1(X,a)\times\pi_1(Y,b)
\]
（它由从 $\pi_1(X\times Y,c)$ 到 $\pi_1(X,a)$ 和 $\pi_1(Y,b)$ 的函子性
态射诱导）。这个态射是单射；若 $k$ 代数闭，它还是双射。

（后一情形中的满射性几乎是平凡的。）与 Serre--Lang 一起，由此推出：

\result{推论}{6} 设 $X$ 是代数闭域 $k$ 上的连通代数概形，$K$ 是 $k$
的一个代数闭扩张。则 $X$ 与 $X\otimes_kK$ 的基本群相同；也就是说，
后一个概形的每个非分歧覆盖都由 $X$ 的一个非分歧覆盖扩张标量而来，且
后者在同构意义下唯一。

\result{注记}{}

1）利用命题4可以看出，推论4中的假设 $f_*(\Osh_X)=\Osh_Y$ 并非必要。
在一般情形中，不应在 $\pi_1(Y)$ 后接单位群 $e$，而应继续为
\[
  \pi_0(\overline F)\longrightarrow\pi_0(X)
  \longrightarrow\pi_0(Y)\longrightarrow e.
\]
如同代数拓扑中一样。

\src{216}
2）一般说来，对于
$\pi_1(\overline F)\longrightarrow\pi_1(X)$ 的核，目前还无法作出
结论；这里看来应当涉及一个 $\pi_2(Y)$。不过，若 $Y$ 是局部环 $A$
的谱，似乎应当可以借助下文定理12证明
$\pi_1(\overline F)\longrightarrow\pi_1(X)$ 是单射；定理12在 $A$
完备时断言的正是这一点。

定理11只用到了定理1和定理2。现在我们将应用定理3，并使用如下初等引理：

\result{引理}{6} 设 $X$ 是概形，$X_0$ 是相应的约化概形（也就是说，
把幂零元杀掉所得的概形）。则 $X_0$ 的每个非分歧覆盖 $X'_0$ 都由
$X$ 的一个非分歧覆盖 $X'$ 诱导，且 $X'$ 在典范同构的意义下唯一。

这个纯局部性质的引理，在这里所起的作用类似于模理论中的定理8。把它与
存在性定理（定理3）结合起来，就得到：

\result{定理}{12} 设 $A$ 是剩余域为 $k$ 的完备诺特局部环，$X$ 是
$A$ 上的固有概形。则 $X_0=X\otimes_Ak$ 的每个非分歧覆盖 $X'_0$
都由 $X$ 的一个非分歧覆盖 $X'$ 诱导，且后者在同构意义下唯一。

换言之：

\result{推论}{1} 在 $X_0$ 中选取一个几何点，作为 $X_0$ 和 $X$ 的
基本群的基点。则典范同态
$\pi_1(X_0)\longrightarrow\pi_1(X)$ 是同构。

现在把引理5应用于 $X_0$（为简单起见，假设
$H^0(X_0,\Osh_{X_0})=k$），并注意到由于 $A$ 完备，它的非分歧扩张
与其剩余域的非分歧扩张相对应；也就是说，
$\pi_1(Y)=\pi_1(k)$，其中 $Y=\operatorname{Spec}(A)$。于是得到正合列
\[
  e\longrightarrow\pi_1(\overline X_0)
  \longrightarrow\pi_1(X)\longrightarrow\pi_1(Y)\longrightarrow e.
\]

\result{推论}{2} 设 $f:X\longrightarrow Y$ 是固有平坦态射，$y_1$
是 $Y$ 中一点，$y_0$ 是 $y_1$ 的一个特殊化，并考虑相应的
\fquote{几何}纤维 $\overline X_1$ 和 $\overline X_0$。假设
$\overline X_0$ 可分且连通（这蕴含 $\overline X_1$ 也满足同样的条件）。
则存在群同态
$\pi_1(\overline X_1)\longrightarrow\pi_1(\overline X_0)$；它在复合
一个内自同构的意义下确定，并且是满射。

\src{217}
人们也许会期望这个同态总是双射；遗憾的是，当 $\kappa(y_0)$ 的特征
$>0$ 时，一般并非如此。不过我们将在下文给出这个同态的核的一个上界
（至少在 $\overline X_0$ 非奇异时），并由此推出：若 $\kappa(y_0)$
的特征为 $0$，则上述同态是双射（这一结果也可用超越方法证明）。无论如何，
我们已经能够借助一般纤维的 $\pi_1$ 来控制特殊纤维的 $\pi_1$。例如，
利用特征 $p$ 的代数曲线可提升为特征零曲线这一事实（定理9，推论4），
由超越方法可得：

\result{推论}{3} 设 $X_0$ 是任意特征的代数闭域上一条完备非奇异曲线
的概形，$g$ 是 $X_0$ 的亏格。则 $\pi_1(X_0)$ 有 $2g$ 个拓扑生成元，
它们满足熟知的关系。

再由熟知的超平面截面技巧推出：

\result{推论}{4} 设 $X$ 是任意特征的代数闭域上的非奇异射影概形。
则 $\pi_1(X)$ 有有限个拓扑生成元。

下面来考察同态
$\pi_1(\overline X_1)\longrightarrow\pi_1(\overline X_0)$ 的核。为此，
可以假设 $Y$ 是完备离散赋值环 $V=\Osh_y$（$y=y_0$）的谱。问题等价于：
给定 $\overline X_1$ 的一个非分歧覆盖 $\overline X'_1$（若愿意，还可
假设它是 Galois 覆盖），判定它在什么条件下来自 $X_0$ 的一个非分歧
覆盖。\textit{先验地}，这个给定覆盖由 $X_1\otimes_KK'$ 的一个非分歧
覆盖 $X'_1$ 扩张标量而来，其中 $K'$ 是 $V$ 的分式域 $K$ 的代数闭包
$\overline K$ 中的一个有限子扩张；若 $\overline X'_1$ 是 Galois 覆盖，
其群为 $G$，则可选取 $X'_1$ 也为群 $G$ 的 Galois 覆盖。于是，要使
$X'_1\otimes_{K'}\overline K=\overline X'_1$ 来自 $X_0$ 的一个非分歧
覆盖，当且仅当能够在 $\overline K$ 中找到
$K'$ 的一个有限扩张 $K''$，使
$X''_1=X'_1\otimes_{K'}K''$ 形如 $X''\otimes_{V''}K''$；这里 $V''$
是 $V$ 在 $K''$ 中的整闭包，而 $X''$ 是 $X\otimes_VV''$ 的一个
非分歧覆盖。例如，假设 $X_0$ 绝对正规；则 $X\otimes_VV''$ 正规
（因为它在 $V''$ 上平坦且特殊纤维正规），其函数域就是下列各概形共有
的函数域 $K''(X_1)$：
\[
  X_1\otimes_KK''
  =(X\otimes_VK)\otimes_KK''
  =X\otimes_VK''
  =(X\otimes_VV'')\otimes_{V''}K''.
\]
令 $L''=K''(X'_1)$ 为 $X'_1\otimes_{K'}K''$ 的函数域。它是
$K''(X_1)$ 的有限可分扩张，而上述条件也就是说，$L''$ 是
$X\otimes_VV''$ 的函数域的非分歧扩张（也就是说，
$X\otimes_VV''$ 在 $L''$ 中的正规化在 $X\otimes_VV''$ 上非分歧）。
只须检验 $L''$ 在 $X\otimes_VV''$ 的特殊纤维各点处非分歧即可
（因为它在一般纤维 $X_1\otimes_KK''$ 上方已经非分歧）。现在若
$X_0$ 非奇异，则由 Nagata--Zariski 的\fquote{纯性定理}，甚至只须检验
$L''$ 在 $X\otimes_VV''$ 的特殊纤维一般点处的局部环 $O''$ 上方
非分歧即可。这个环是离散赋值环，也是特殊纤维一般点处的局部环
$O\subset K(X_1)$ 在 $K''(X_1)$ 中的正规化。因此我们得到：

\src{218}
\result{推论}{5} 在上述条件和记号下，$\overline X_1=
X_1\otimes_K\overline K$ 的非分歧覆盖
$X'_1\otimes_{K'}\overline K$ 来自 $\overline X_0$ 的一个非分歧覆盖，
当且仅当存在 $\overline K/K'$ 的一个有限子扩张 $K''$，使
$K''(X'_1)$ 在离散赋值环 $O''\subset K''(X_1)$ 上方非分歧。

现在注意，$O''$ 是离散赋值环 $O'\subset K'(X_1)$ 在 $K''(X_1)$
中的正规化；这里 $O'$ 是 $O$ 在 $K'(X_1)$ 中的正规化。环 $O'$ 包含
$V$ 在 $K'$ 中的整闭包 $V'$，而 $V'$ 的一个素元 $u$ 也是 $O'$ 的
素元。再假设 $X'_1$ 是 Galois 覆盖，Galois 群 $G$ 的阶 $n$ 与
$\kappa(y_0)$ 的特征 $p$ 互素（$p$ 也是 $O'$ 的剩余域的特征）。于是
$K'(X'_1)$ 在 $O'$ 上方驯分歧。由 Abhyankar 引理容易推出：向它添加
$O'$ 的一个素元的 $n$ 次根 $v$ 后，所得扩张在 $O'$ 于
$K'(X_1)(v)$ 中的正规化上方非分歧。而 $v$ 可以取为 $V'$ 的一个素元
的 $n$ 次根，这就证明了推论5的条件成立。（利用 Abhyankar 引理和纯性
定理的这一办法，是 Serre 卖给我的。）为了表述所得结果，对每个紧致全
不连通群 $\pi$，令 $\widetilde\pi$ 表示 $\pi$ 关于其 Sylow $p$-子群
所生成的闭子群的商群；也就是说，它是 $\pi$ 的所有阶与 $p$ 互素的离散
商群的逆极限。采用这个记号，便有：

\result{定理}{13} 设 $f:X\longrightarrow Y$ 是固有平坦态射，$y_1$
是 $Y$ 中一点，$y_0$ 是 $y_1$ 的一个特殊化，并假设 $\overline X_0$
连通且非奇异。由定理12推论2中的满同态所诱导的同态
$\widetilde\pi_1(\overline X_1)\longrightarrow
\widetilde\pi_1(\overline X_0)$ 是同构。

换言之：

\src{219}
\result{推论}{1} Galois 群的阶与 $\kappa(y_0)$ 的特征 $p$ 互素的
非分歧 Galois 覆盖，在 $\overline X_0$ 和 $\overline X_1$ 上有相同的
分类。

特别地，若 $\kappa(y_0)$ 的特征为零，便由代数方法得出
$\pi_1(\overline X_1)\longrightarrow\pi_1(\overline X_0)$ 是双射。

最后指出，上述技巧还给出下面这个比定理13更一般的结果：

\result{定理}{14} 设 $f:X\longrightarrow Y$ 是固有光滑态射，$D$ 是
$X$ 的一个闭子概形，它在 $Y$ 上光滑，并且在每一点处的余维数都为1。
对 $f$ 的一条纤维 $Z=f^{-1}(z)$，令 $Z'=Z-Z\cap D$，并令
$\pi_1^t(Z')$ 为基本群 $\pi_1(Z')$ 的如下商群：它分类 $Z'$ 的那些
在 $\overline{Z\cap D}$ 上方驯分歧的非分歧覆盖。设 $y_0,y_1$ 如定理13。
则有一个满同态（在复合一个内自同构的意义下确定）
\[
  \pi_1^t(\overline X'_1)\longrightarrow\pi_1^t(\overline X'_0),
\]
并且相应的同态
\[
  \widetilde\pi_1^t(\overline X'_1)
  \longrightarrow\widetilde\pi_1^t(\overline X'_0)
\]
是同构。

由此可得定理13各推论以及定理12推论4的相应变体。同样，利用定理9
推论3，由超越方法可得：

\result{推论}{} 设 $X_0$ 是任意特征的代数闭域上一条完备非奇异曲线
的概形，$S=(s_i)_{1\le i\le n}$ 是 $X_0$ 中一个含 $n$ 个元素的有限
子集。则 $\pi_1^t(X_0-S)$ 有 $2g+n$ 个拓扑生成元
$x_i,y_i$（$1\le i\le g$）和 $\sigma_j$（$1\le j\le n$），它们
满足关系
\[
  \left(\prod_i x_i y_i x_i^{-1}y_i^{-1}\right)
  \sigma_1\cdots\sigma_n=1,
\]
其中 $\sigma_j$ 是与 $s_j$ 相应的惯性群的生成元。对于每个阶与该域
的特征互素的有限群 $G$，若它由满足上述关系的元素
$x_i,y_i,\sigma_j$ 生成，则存在 $X_0-S$ 的一个群为 $G$ 的非分歧
Galois 覆盖，而且它在点 $s_j$ 处的惯性群由 $\sigma_j$ 生成。

当 $X_0$ 的亏格为0且 $n=3$ 时，我们便得到\fquote{三点问题}的一个
解，至少对阶与该域特征互素的 Galois 覆盖如此。（这里定理9其实并不
需要；另一方面，看来还可以从三点问题所考察的特殊情形推出前一推论。）

\src{220}
\result{注记}{}

1）更完整的研究或许应当引入 $X,X_0,X_1$ 的广义 Galois 覆盖
（其 Galois 群可以是无穷小群），从而有可能恢复定理12推论2中的核。
相反，对具有非驯分歧的覆盖进行研究，看来要困难得多。

2）将引理6与 Grauert 关于沿超平面截面的非奇异射影概形的形式完备化
的一项结果结合起来（或者与定理11之后注记2中提到、目前尚未证明的定理
结合起来），也可以在\fquote{抽象}代数几何中证明经典的基本群 Lefschetz
定理。

\begin{center}
  \textbf{参考文献}
\end{center}

\begin{description}
  \item[{[1]}] Dieudonné (J.) 和 Grothendieck (A.). --- \emph{Éléments de
    géométrie algébrique}, Publications mathématiques de l'Institut des Hautes
    Études Scientifiques（即将出版）。
  \item[{[2]}] Grothendieck (Alexander). --- \emph{The cohomology theory of
    abstract algebraic varieties}. International Congress of Mathematicians
    [1958. Edinburgh]（即将出版）。
  \item[{[3]}] Serre (Jean-Pierre). --- \emph{Faisceaux algébriques
    cohérents}, Annals of Math., Series 2, t. 61, 1955, p. 197--278.
  \item[{[4]}] Serre (Jean-Pierre). --- \emph{Géométrie algébrique et
    géométrie analytique}, Ann. Institut Fourier Grenoble, t. 6, 1955--1956,
    p. 1--42.
\end{description}

% END 182.tex

% BEGIN 190.tex
\unit{第190次报告——下降技术 I}
\sid{190}
\src{299}
\seminarzh{第12年度，1959/60，第190号\\（1959年12月）}
\paperhead{下降技术与代数几何中的存在性定理\\
I. 一般理论：沿忠实平坦态射的下降}%
  {亚历山大·格罗滕迪克 著}

从技术观点看，本报告以及后续各篇都可视为著名的希尔伯特\fquote{定理90}
的种种变奏。把下降方法引入代数几何，似乎应归功于A.~Weil；他称之为
\fquote{基域下降}。不过，Weil 只处理域的有限可分扩张。随后，P.~Cartier
处理了高为1的纯不可分扩张。由于当时缺少概形语言，尤其因为所容许考察的
环中没有幂零元，Cartier 未能表述这两种情形在本质上的同一性。

目前看来，下面将要介绍的一般下降技术（必要时再与\fquote{形式几何}的
基本定理结合，参见[3]）构成了代数几何中大多数存在性定理的基础。
\footnote{1962年的勘误认为这一说法过强：对于 TDTE I--VI 系列前两篇
报告所用的非射影技术，这一说法在很大程度上成立；但对于射影技术
（TDTE IV、V、VI）则不成立。}还应指出，这项技术当然可以移植到
\fquote{解析几何}中；可以期待，在不久的将来，专家们能够证明报告II中
将给出的形式几何存在性定理的\fquote{解析}对应。无论如何，
Kodaira--Spencer 最近的工作已经相当清楚地表明，需要采用更接近概形理论的
方法：他们的方法似乎不适于在模簇的奇点附近定义并研究模簇（当然，概形论
方法还必须由超越方法补充）。

在本报告I中，我们将处理标题所示的最初等下降情形。不过，下面定理1、2、3
的应用已经极为丰富。我们只选取其中若干作为例子，并不试图在这里把它们
表述到尽可能一般的程度。

我们将自由使用概形语言；有关内容可参看前述报告及[2]。但须明确指出，本报告
所考察的预概形不必是诺特的，其中的

\src{300}
态射也不必是有限型的。因此，若 $A$ 是一个诺特局部环，其完备化为
$\widehat A$，我们便必须考察非诺特环
$\widehat A\otimes_A\widehat A$，以及由所涉各环之间的包含关系所对应的
仿射概形态射。

\fsection{A}{关于范畴的预备知识}

\subsection*{1. 纤维范畴、下降数据与 $\mathcal F$-下降态射}
\addcontentsline{toc}{subsection}{1. 纤维范畴与下降数据}

\result{a. 定义}{1.1} 一个以 $\mathcal C$ 为基的\emph{纤维范畴}
$\mathcal F$ 由以下数据组成：一个范畴 $\mathcal C$；对每个
$X\in\mathcal C$，一个范畴 $\mathcal F_X$；对每个
$\mathcal C$-态射 $f:X\longrightarrow Y$，一个函子
$f^*:\mathcal F_Y\longrightarrow\mathcal F_X$，对
$\xi\in\mathcal F_Y$ 也记作
\[
  f^*(\xi)=\xi\times_Y X
\]
（此时把 $X$ 看成\fquote{$Y$ 上的 $\mathcal C$-对象}，即配备了态射
$f$）；最后，对两个可复合态射
$X\xrightarrow{f}Y\xrightarrow{g}Z$，给定一个函子同构
\[
  c_{f,g}:(gf)^*\longrightarrow f^*g^*,
\]
这些数据须满足下列条件：

\begin{enumerate}
\item[(i)] $\operatorname{id}_X^*=\operatorname{id}$；
\item[(ii)] 若 $f$ 或 $g$ 是恒等同构，则 $c_{f,g}$ 是恒等同构；
\item[(iii)] 给定三个可复合态射
$X\xrightarrow{f}Y\xrightarrow{g}Z\xrightarrow{h}T$，由形如
$c_{u,v}$ 的同构组成的下图可交换：
\[
\begin{array}{ccc}
  (h(gf))^*=((hg)f)^* & \longrightarrow & f^*(hg)^* \\
  \big\downarrow && \big\downarrow \\
  (gf)^*h^* & \longrightarrow & (f^*g^*)h^*=f^*(g^*h^*).
\end{array}
\]
\end{enumerate}

\result{例}{1} 设 $\mathcal C$ 是一个存在纤维积的范畴。可如下定义一个以
$\mathcal C$ 为基的纤维范畴 $\mathcal F$：令 $\mathcal F_X$ 为
$X$ 上的 $\mathcal C$-对象所成的范畴；与态射
$f:X\longrightarrow Y$ 相应的函子
$f^*:\mathcal F_Y\longrightarrow\mathcal F_X$ 由纤维积
$Z\longmapsto Z\times_YX$ 定义。

\result{例}{2} 设 $\mathcal C$ 是预概形范畴；对
$X\in\mathcal C$，令 $\mathcal F_X$ 为 $X$ 上拟凝聚模层所成的范畴。
若 $f:X\longrightarrow Y$ 是预概形态射，则
$f^*:\mathcal F_Y\longrightarrow\mathcal F_X$ 是

\src{301}
模层的逆像函子。这样便得到一个以 $\mathcal C$ 为基的纤维范畴。

\result{b. 定义}{1.2} 集合映射图
\[
  E\xrightarrow{u}E'
  \overset{v_1}{\underset{v_2}{\rightrightarrows}}E''
\]
称为\emph{正合的}，如果 $u$ 把 $E$ 双射到 $E'$ 的如下子集：
其中各 $x'\in E'$ 满足 $v_1(x')=v_2(x')$。

\result{定义}{1.3} 设 $\mathcal F$ 是一个以 $\mathcal C$ 为基的纤维范畴；
考虑 $\mathcal C$ 中满足 $\alpha\beta_1=\alpha\beta_2$ 的态射图
\[
  \tag{+}
  S\xleftarrow{\alpha}S'
  \overset{\beta_1}{\underset{\beta_2}{\rightrightarrows}}S''
\]
若对 $\mathcal F_S$ 中任意一对对象 $\xi,\eta$，下列集合映射图都正合，
便称上图为 $\mathcal F$-\emph{正合的}：
\[
  \tag{+'}
  \operatorname{Hom}(\xi,\eta)
  \xrightarrow{\alpha^*}
  \operatorname{Hom}(\alpha^*\xi,\alpha^*\eta)
  \overset{\beta_1^*}{\underset{\beta_2^*}{\rightrightarrows}}
  \operatorname{Hom}(\gamma^*\xi,\gamma^*\eta),
\]
其中 $\gamma=\alpha\beta_1=\alpha\beta_2$。为简化记号，在后一图中我们
借助 $c_{\beta_i,\alpha}$，把 $\beta_i^*\alpha^*$ 与
$(\alpha\beta_i)^*=\gamma^*$ 视为相同。

\result{定义}{1.4} 设 $\mathcal F$ 是一个以 $\mathcal C$ 为基的纤维范畴；
考虑 $\mathcal C$ 中从 $S''$ 到 $S'$ 的两个态射
$\beta_1,\beta_2$。设 $\xi'\in\mathcal F_{S'}$。所谓 $\xi'$ 上的一个
\emph{粘合数据}（相对于 $(\beta_1,\beta_2)$），是一个从
$\beta_1^*(\xi')$ 到 $\beta_2^*(\xi')$ 的同构。若
$\xi',\eta'\in\mathcal F_{S'}$ 各自配备了粘合数据，则
$\mathcal F_{S'}$ 中的态射 $u:\xi'\longrightarrow\eta'$ 称为
\emph{与粘合数据相容的}，如果下图可交换：
\[
\begin{array}{ccc}
  \beta_1^*(\xi') & \longrightarrow & \beta_2^*(\xi') \\
  \big\downarrow && \big\downarrow \\
  \beta_1^*(\eta') & \longrightarrow & \beta_2^*(\eta').
\end{array}
\]
由此，$\mathcal F_{S'}$ 中配备粘合数据（相对于
$\beta_1,\beta_2$）的对象组成一个范畴。若
$\alpha:S'\longrightarrow S$ 满足
$\alpha\beta_1=\alpha\beta_2$，则对每个 $\xi\in\mathcal F_S$，对象
$\xi'=\alpha^*(\xi)$

\src{302}
在 $\mathcal F_{S'}$ 中带有一个典范粘合数据，因为
\[
  \beta_i^*\alpha^*(\xi)\simeq(\alpha\beta_i)^*(\xi)=\gamma^*(\xi),
\]
这里仍令 $\gamma=\alpha\beta_1=\alpha\beta_2$。此外，若
$u:\xi\longrightarrow\eta$ 是 $\mathcal F_S$ 中的一个态射，则
\[
  \alpha^*(u):\alpha^*(\xi)\longrightarrow\alpha^*(\eta)
\]
是 $\mathcal F_{S'}$ 中与典范粘合数据相容的态射。由此得到一个典范函子：
它从 $\mathcal F_S$ 映到 $\mathcal F_{S'}$ 中配备了相对于
$(\beta_1,\beta_2)$ 的粘合数据的对象所成的范畴。于是，定义1.3也可表述为：
该定义中的图（+）是 $\mathcal F$-正合的，当且仅当前述函子全忠实；也就是
说，它给出 $\mathcal F_S$ 与 $\mathcal F_{S'}$ 中配备了相对于
$(\beta_1,\beta_2)$ 的粘合数据的对象范畴的一个全子范畴之间的等价。

\result{定义}{1.5} 若 $\xi'\in\mathcal F_{S'}$ 配备的粘合数据使得
$\xi'$（连同此数据）同构于某个 $\alpha^*(\xi)$，其中
$\xi\in\mathcal F_S$，便称这个粘合数据（相对于 $\alpha$）是
\emph{有效的}。

当图（+）是 $\mathcal F$-正合的时，前一定义中的对象 $\xi$ 因而在唯一同构
意义下确定，而 $\mathcal F_S$ 等价于 $\mathcal F_{S'}$ 中配备有效粘合数据
的对象所成的范畴。

\medskip\noindent
c. 最重要的情形是
\[
  S''=S'\times_S S',
\]
其中 $\beta_i$ 是从 $S'\times_S S'$ 到两个因子的两个投影 $p_1,p_2$
（现假设 $\mathcal C$ 中存在纤维积）。这时，一个对象
$\xi'\in\mathcal F_{S'}$ 上的粘合数据
$\varphi:p_1^*(\xi')\longrightarrow p_2^*(\xi')$ 若要有效，立即有两个必要
条件。其一为
\[
  \tag{i}
  \Delta^*(\varphi)=\operatorname{id}_{\xi'},
\]
其中 $\Delta:S'\longrightarrow S'\times_S S'$ 是对角态射，并且我们把
$\Delta^*p_i^*(\xi')$ 与 $(p_i\Delta)^*(\xi')=\xi'$ 视为相同。其二为
\[
  \tag{ii}
  p_{31}^*(\varphi)=p_{32}^*(\varphi)\,p_{21}^*(\varphi),
\]

\src{303}
其中 $p_{ij}$ 表示从 $S'\times_S S'\times_S S'$ 到由其第 $i$、第 $j$
个因子构成的部分纤维积的典范投影。

\result{定义}{1.6} 相对于态射 $\alpha:S'\longrightarrow S$，所谓
$\xi'\in\mathcal F_{S'}$ 上的一个\emph{下降数据}，是 $\xi'$ 上相对于
典范投影对 $(p_1,p_2):S'\times_S S'\longrightarrow S'$ 的一个粘合数据，
并且满足上述条件（i）和（ii）。

\result{定义}{1.7} 若态射 $\alpha:S'\longrightarrow S$ 的态射图
\[
  S\xleftarrow{\alpha}S'
  \overset{p_1}{\underset{p_2}{\rightrightarrows}}S'\times_S S'
\]

是 $\mathcal F$-正合的（定义1.3），便称 $\alpha$ 为一个
\emph{$\mathcal F$-下降态射}；若此外 $\mathcal F_{S'}$ 中任一对象上的
每个下降数据（定义1.6）都有效，便称 $\alpha$ 为一个
\emph{严格 $\mathcal F$-下降态射}。

后一个条件也可更形象地说成：给定 $\mathcal F_S$ 中的一个对象，与给定
$\mathcal F_{S'}$ 中配备下降数据的一个对象，\fquote{是一回事}。

注意，若态射 $\alpha:S'\longrightarrow S$ 有一个\fquote{截面}
$s:S\longrightarrow S'$（即 $\alpha s=\operatorname{id}_S$），则它是一个
严格 $\mathcal F$-下降态射。\footnote{1962年的勘误指出，无须预先假设
$\alpha$ 是 $\mathcal F$-下降态射。}事实上，若
$\xi'\in\mathcal F_{S'}$ 配备了相对于 $\alpha$ 的下降数据，它便\fquote{来自}
$\xi=s^*(\xi')$。

\medskip\noindent
d. 还可用更直观的方式表述上述概念。对 $S$ 上的一个 $\mathcal C$-参数对象
$T$，引入集合
\[
  \operatorname{Hom}_S(T,S')=S'(T),
\]
其元素记作 $t,t'$ 等。给定 $\xi'\in\mathcal F_{S'}$ 后，每个
$t\in S'(T)$ 对应于 $\mathcal F_T$ 中的对象 $t^*(\xi')$，也记作
$\xi'_t$。于是，$\xi'$ 上相对于 $p_1,p_2$ 的一个粘合数据，就是对每个
$S$ 上的 $T$ 及每一对点 $t,t'\in S'(T)$，给定一个同构
\[
  \varphi_{t',t}:\xi'_t\longrightarrow\xi'_{t'}
\]
（当 $T$ 变化时还须满足显然的函子性条件）。1.c中的条件（i）、（ii）于是
分别写成
\[
  \tag{i bis}
  \varphi_{t,t}=\operatorname{id}_{\xi'_t}
  \qquad\text{对每个 $T$ 及每个 $t\in S'(T)$},
\]

\src{304}
\[
  \tag{ii bis}
  \varphi_{t,t''}=\varphi_{t,t'}\varphi_{t',t''}
  \qquad\text{对每个 $T$ 及所有 $t,t',t''\in S'(T)$}.
\]
此外，在（ii bis）中令 $t=t'=t''$，便得到
$\varphi_{t,t}=\varphi_{t,t}^2$；由于假设 $\varphi_{t,t}$ 是同构，遂得到
（i bis）。所以（i bis）其实是（ii bis）的推论（从而（i）是（ii）的
推论）。但若不再预先假设各 $\varphi_{t,t'}$ 为同构（即不再假设
$\varphi:p_1^*(\xi')\longrightarrow p_2^*(\xi')$ 是同构），则（ii bis）
未必推出（i bis）；不过，（ii bis）与（i bis）合在一起仍可推出各
$\varphi_{t,t'}$ 都是同构，因为此时
$\varphi_{t,t'}\varphi_{t',t}=\varphi_{t,t}=\operatorname{id}_{\xi'_t}$。

\subsection*{2. 正合图与严格满态射；下降态射。例}
\addcontentsline{toc}{subsection}{2. 正合图与严格满态射}

\result{a. 定义}{2.1} 设 $\mathcal C$ 是一个范畴。若对每个
$Z\in\mathcal C$，$\mathcal C$ 中的态射图
\[
  T\xrightarrow{\alpha}T'
  \overset{\beta_1}{\underset{\beta_2}{\rightrightarrows}}T''
\]
所对应的集合映射图
\[
  \operatorname{Hom}(Z,T)\longrightarrow
  \operatorname{Hom}(Z,T')
  \rightrightarrows\operatorname{Hom}(Z,T'')
\]
都正合（定义1.2），便称原图\emph{正合}。这时称 $(T,\alpha)$
（或略去严格说法而称 $T$）为态射对 $(\beta_1,\beta_2)$ 的一个\emph{核}。

显然，这个核在唯一同构意义下确定。若 $\mathcal C$ 是集合范畴，上述定义与
定义1.2相容。对偶地，定义 $\mathcal C$ 中态射图
\[
  S\xleftarrow{\alpha}S'
  \overset{\beta_1}{\underset{\beta_2}{\rightrightarrows}}S''
\]
的正合性；这时称 $(S,\alpha)$ 为态射对 $(\beta_1,\beta_2)$ 的一个
\emph{余核}。

\result{定义}{2.2} 若态射 $\alpha:S'\longrightarrow S$ 是满态射，并且对每个
$u:S'\longrightarrow Z$，下述必要条件也是 $u$ 可分解成
$S'\longrightarrow S\longrightarrow Z$ 的充分条件，便称 $\alpha$ 为一个
\emph{严格满态射}：对每个 $S''\in\mathcal C$ 以及每一对满足
$\alpha\beta_1=\alpha\beta_2$ 的态射
$\beta_1,\beta_2:S''\longrightarrow S'$，也有 $u\beta_1=u\beta_2$。

若纤维积 $S'\times_S S'$ 存在，这等价于说图
\[
  S\xleftarrow{\alpha}S'
  \overset{p_1}{\underset{p_2}{\rightrightarrows}}S'\times_S S'
\]

\src{305}
正合，也就是说，$S$ 是态射对 $(p_1,p_2)$ 的一个余核。无论如何，一个
余核态射总是严格满态射。还应注意，既是严格满态射又是单态射的态射必为
同构。其对偶概念\emph{严格单态射}留给读者自行展开。

为了精确说明 $\mathcal F$-下降态射与严格满态射之间的关系，再引入下列定义。

\result{定义}{2.3} 若对 $S$ 上的每个 $T$，纤维积
$T'=S'\times_ST$ 都存在，而且投影 $T'\longrightarrow T$ 是满态射
（相应地，是严格满态射），便称态射 $\alpha:S'\longrightarrow S$ 为
\emph{万有满态射}（相应地，\emph{万有严格满态射}）。

在性质很好的范畴中（例如集合范畴、一个拓扑空间上的集合层范畴、阿贝尔
范畴等），刚才引入的四种满态射性概念彼此重合；但在预概形范畴，或一个
给定非空预概形 $S$ 上的预概形范畴中，它们各不相同，即使只考虑在 $S$ 上
有限的 $S$-概形也是如此。

\result{定义}{2.4} 若态射 $\alpha:S'\longrightarrow S$ 是
$\mathcal F$-下降态射（相应地，严格 $\mathcal F$-下降态射；参见定义1.7），
其中 $\mathcal F$ 是以 $\mathcal C$ 为基、由 $\mathcal C$ 的对象上的
$\mathcal C$-对象组成的纤维范畴（第1节，例1），便称 $\alpha$ 为一个
\emph{下降态射}（相应地，\emph{严格下降态射}）。

\result{命题}{2.1} 若 $\mathcal C$ 中存在有限积和（有限）纤维积，则
$\mathcal C$ 中的下降态射恰好就是万有严格满态射。

\result{b. 例}{} 设 $\mathcal C$ 是预概形范畴，$S\in\mathcal C$；并设
$S',S''$ 是 $S$ 上的两个有限预概形，即分别对应于 $S$ 上的代数层
$\mathcal A',\mathcal A''$，它们作为模层是拟凝聚且有限型的（若 $S$
局部诺特，也就是凝聚的）。设 $\alpha:S'\longrightarrow S$ 是 $S'$ 的结构
态射，并设 $\beta_1,\beta_2$ 是从 $S''$ 到 $S'$ 的两个 $S$-态射；它们由
代数同态 $\mathcal A'\longrightarrow\mathcal A''$ 定义，后者仍记为
$\beta_1,\beta_2$。利用有限态射是闭态射这一事实（Cohen--Seidenberg
第一定理），很容易证明，$\mathcal C$ 中的图
\[
  \tag{+}
  S\xleftarrow{\alpha}S'
  \overset{\beta_1}{\underset{\beta_2}{\rightrightarrows}}S''
\]
正合，当且仅当 $S$ 上的层图

\src{306}
\[
  \mathcal O_S=\mathcal A\longrightarrow\mathcal A'
  \overset{\beta_1}{\underset{\beta_2}{\rightrightarrows}}\mathcal A''
\]
正合。特别地，若 $\alpha:S'\longrightarrow S$ 是对应于 $S$ 上代数层
$\mathcal A'$ 的有限态射，则 $\alpha$ 是严格满态射，当且仅当层图
\[
  \mathcal O_S=\mathcal A\longrightarrow\mathcal A'
  \overset{p_1}{\underset{p_2}{\rightrightarrows}}
  \mathcal A'\otimes_{\mathcal A}\mathcal A'
\]
正合（$\alpha$ 是满态射，当且仅当
$\mathcal A\longrightarrow\mathcal A'$ 是单射）。若 $S$ 是以 $A$ 为环的
仿射概形，从而 $S'$ 是以有限 $A$-代数 $A'$ 为环的仿射概形，则
$S'\longrightarrow S$ 是严格满态射，当且仅当 $A\longrightarrow A'$
诱导出 $A$ 与 $A'$ 的如下子环之间的同构：该子环由满足
\[
  1_{A'}\otimes_Ax'-x'\otimes_A1_{A'}=0
\]
的 $x'\in A'$ 组成的子环（它是满态射，当且仅当
$A\longrightarrow A'$ 是单射）。如前所述，即使 $S$ 是一个阿廷局部环的
谱，一个有限满态射 $S'\longrightarrow S$ 也未必是严格满态射。不过，
可以证明，若 $S$ 是诺特预概形，则每个有限满态射
$S'\longrightarrow S$ 都是有限多个（同样有限的）严格满态射的复合。这还
表明，两个严格满态射的复合未必是严格满态射。

\result{c}{} 若 $(+)$ 是由有限态射组成的正合图，则对预概形的每个平坦态射
$T\longrightarrow S$，经基变换 $T\longrightarrow S$ 得到的 $(+)$ 的变换图
仍然正合。因此，若 $X,Y$ 是两个 $S$-预概形，且 $X$ 在 $S$ 上平坦，则下列
集合映射图正合（其中 $X',Y'$ 是 $X,Y$ 在 $S'$ 上的逆像，$X'',Y''$ 是它们
在 $S''$ 上的逆像）：
\[
  \operatorname{Hom}_S(X,Y)\longrightarrow
  \operatorname{Hom}_{S'}(X',Y')
  \rightrightarrows
  \operatorname{Hom}_{S''}(X'',Y'')
\]
特别地，设 $\mathcal F$ 是以预概形范畴 $\mathcal C$ 为基的纤维范畴，
并且对 $X\in\mathcal C$，$\mathcal F_X$ 是平坦 $X$-预概形所成的范畴；
则图 $(+)$ 是 $\mathcal F$-正合的。（若不作平坦性假设，这个结论便不成立；
特别是，有限严格满态射未必是下降态射。）同样可见，当 $\mathcal F$ 是满足
如下条件的纤维范畴时，$(+)$ 也是 $\mathcal F$-正合的：$\mathcal F_X$ 是
预概形 $X$ 上拟凝聚平坦层所成的范畴（这里平坦性假设仍然

\src{307}
至关重要）。在这两种情形下，$S'$ 上的平坦对象所带粘合数据是否有效
（特别地，当 $S''=S'\times_SS'$ 时，下降数据是否有效）是一个微妙问题；
本系列报告的主要目标之一，正是回答若干特殊情形下的这一问题。报告人并不
知道：对每个有限严格满态射 $S'\longrightarrow S$，$S'$ 上拟凝聚平坦层的
每个下降数据是否都有效；即使假设 $S$ 是一个阿廷局部环的谱，并且只考察
秩为1的局部自由层，亦不知答案。

更一般地，设 $A$ 是环，$A'$ 是一个 $A$-代数（所涉各环均交换），并且映射图
\[
  A\longrightarrow A'\rightrightarrows A'\otimes_AA'
\]
正合。这也等价于：$A,A'$ 的谱所对应的态射 $S'\longrightarrow S$ 是一个
$\mathcal F$-下降态射，其中 $\mathcal F$ 是拟凝聚平坦层所成的纤维范畴。
设 $M'$ 是配备了一个相对于 $A\longrightarrow A'$ 的下降数据的平坦
$A'$-模；也就是说，配备了一个
$A'\otimes_AA'$-模同构
\[
  \varphi:M'\otimes_AA'\xrightarrow{\sim}A'\otimes_AM'
\]
并满足1.c中的条件（i）、（ii）（其模论表述留给读者写出）。这个数据是否
有效（相对于拟凝聚平坦层的纤维范畴）？令 $M$ 为 $M'$ 中满足
\[
  \varphi(x'\otimes_A1_{A'})=1_{A'}\otimes_Ax',
\]
的所有 $x'\in M'$ 所组成的子集；它是 $M'$ 的一个 $A$-子模。典范嵌入
$M\longrightarrow M'$ 定义了一个 $A'$-模同态
$M\otimes_AA'\longrightarrow M'$。于是，$\varphi$ 的有效性正是指：$M$
是平坦 $A$-模，并且上述同态是同构。

\result{注}{} 在此前的讨论中，我们没有对图 $(+)$ 中的态射作任何平坦性
假设；为了得到下降技术，因而不得不对所考察的 $S,S'$ 上对象作平坦性假设。
在第2节中，我们将对 $\alpha:S'\longrightarrow S$ 作平坦性假设；这会使我们
能够对 $S,S'$ 上不再受任何平坦性条件约束的对象采用下降技术。无论哪种
情形，总有一个平坦性假设参与其中。这正是平坦性概念在代数几何中如此重要
的主要原因之一（只要把基限制为域，这一作用便无法显现，因为在域上所讨论
的一切对象确实都是平坦的！）。

\src{308}
\subsection*{3. 平展态射的应用}
\addcontentsline{toc}{subsection}{3. 平展态射的应用}

设 $A$ 是局部环，$B$ 是 $A$ 上的局部代数，并且 $B$ 的极大理想在 $A$ 中
诱导出 $A$ 的极大理想。若 $B$ 满足下列条件，我们就称 $B$ 在 $A$ 上
\emph{平展}（不用别处所用的\fquote{非分歧}一词）：
\begin{enumerate}
\item[(i)] $B$ 在 $A$ 上平坦；
\item[(ii)] $B/\mathfrak mB$ 是 $A/\mathfrak m=k$ 的有限可分扩张（这里
$\mathfrak m$ 表示 $A$ 的极大理想）。
\end{enumerate}
当 $A,B$ 都是诺特环而 $k$ 代数闭时，这就是说，延拓
$A\longrightarrow B$ 的完备化同态
$\widehat A\longrightarrow\widehat B$ 是同构。有限型态射
$f:T\longrightarrow S$ 称为\emph{在 $x\in T$ 处平展}，或者称 $T$
\emph{在 $x$ 处平展于 $S$ 上}，若 $\mathcal O_x$ 在
$\mathcal O_{f(x)}$ 上平展；若 $f$ 在每个 $x\in T$ 处平展，则称 $f$
\emph{平展}，也称 $f$ 为一个\emph{平展态射}，或称 $T$
\emph{平展于 $S$ 上}。还要指出，若 $S$ 局部诺特，则 $T$ 中 $f$ 平展的点
组成开集；利用 Zariski \fquote{主定理}，还可以用一种熟知类型的方程，更
精确地描述这类点附近 $T/S$ 的结构。

当 $S$ 是复数域上的有限型概形时，按 Serre [5] 的意义，它对应于一个解析
空间 $\overline S$；这里只须容许 $\overline S$ 的结构层中可能含有幂零元，
而这并不改变 [5] 中任何本质内容。于是容易看出，$f$ 是平展态射，当且仅当
$\overline f:\overline T\longrightarrow\overline S$ 是平展态射；也就是说，
$\overline T$ 的每一点都有一个邻域，使 $\overline f$ 在其上诱导到
$\overline S$ 某个开集上的同构。特别地，$S$ 的每个\emph{平展覆盖} $T$
（即有限平展态射 $f:T\longrightarrow S$）对应于 $\overline S$ 的一个平展
覆盖 $\overline T$；后者连通，当且仅当 $T$ 连通 [5]。又容易看出，若
$T,T'$ 是两个平展于 $S$ 上的概形，则自然映射
\[
  \operatorname{Hom}_S(T,T')\longrightarrow
  \operatorname{Hom}_{\overline S}(\overline T,\overline{T'})
\]
是双射。换言之，从平展于 $S$ 上的概形范畴到平展于 $\overline S$ 上的
解析空间范畴的函子 $T\longmapsto\overline T$ 是\fquote{全忠实}的，因而
给出前一范畴与后一范畴的一个子范畴之间的等价。Grauert--Remmert [2] 的
一个定理蕴含：若 $S$ 正规，则由此得到 $S$ 的\emph{平展覆盖}范畴与
$\overline S$ 的（有限）平展覆盖范畴之间的等价；也就是说，$\overline S$
的每个平展覆盖都在 $\overline S$ 上同构于某个 $\overline T$，其中 $T$
是 $S$ 的平展覆盖。下面证明，\emph{即使不假设 $S$ 正规，
Grauert--Remmert 定理仍然成立}。先设 $S'\longrightarrow S$ 是有限严格
满态射，假定该定理已对 $S'$ 证明，我们来证明它对 $S$ 也成立。

\src{309}
事实上，设 $\mathfrak T$ 是 $\overline S$ 的一个平展覆盖，并考察它在
$\overline{S'}$ 上的逆像 $\mathfrak T'$。它对应于 $\overline{S'}$ 上的
凝聚解析代数层 $\mathcal A'$，后者是定义 $\mathfrak T$ 的
$\overline S$ 上代数层 $\mathcal A$ 的逆像。依假设，在
$\overline{S'}$ 上，$\mathfrak T'$ 来自 $S'$ 的一个平展覆盖 $T'$；也就是
说，$\mathcal A'$ 来自 $S'$ 上的凝聚代数层 $A'$。另一方面，
$\mathcal A'$ 具有相对于
$\overline{S'}\longrightarrow\overline S$ 的典范下降数据，即它在
\[
  \overline{S'}\times_{\overline S}\overline{S'}
  =\overline{(S'\times_SS')},
\]
上的两个逆像之间的一个同构（满足条件（i）、（ii））。根据上述讨论，
这个同构来自相应代数层之间的一个同构，也就是 $A'$ 相对于
$S'\longrightarrow S$ 的一个下降数据。容易验证，后一个下降数据是有效的
（因为 $\mathcal A'$ 上的数据有效，而上一小节已经具体说明，下降数据的
有效性可在所涉各模的\emph{完备化}上作局部判定）。由此得到 $S$ 上的一个
凝聚代数层 $A$，它定义了 $S$ 的一个平展覆盖 $T$，正是所求覆盖。显然，若
$S'\longrightarrow S$ 只是有限多个有限严格满态射的复合，上述结论仍然
成立；根据第2节指出的分解结果，这正是任意有限满态射的情形。因此，引入
$S$ 的正规化 $S'$，便看出当 $S$ 是\emph{约化}概形，即
$\mathcal O_S$ 不含幂零元时，Grauert--Remmert 定理仍然成立。从这里容易
过渡到一般情形。

完全类似的证明，再次利用有限严格满态射的分解结果以及下降数据有效性的
\fquote{形式}性质，可以证明如下结论：设 $S$ 是局部诺特预概形，并设
$S'\longrightarrow S$ 是\emph{有限且满射的纯不可分态射}（等价地，是一个满足
如下条件的有限型态射：对 $S$ 上每个 $T$，态射
$T'=S'\times_ST\longrightarrow T$ 都是同胚；这也表述为
$S'\longrightarrow S$ 是一个\emph{万有同胚}）。对每个平展于 $S$ 上的
$T$，考察它在 $S'$ 上的逆像 $T'=T\times_SS'$，它平展于 $S'$ 上。于是
函子 $T\longmapsto T'$ 给出\emph{平展于 $S$ 上的预概形 $T$ 所成范畴与
平展于 $S'$ 上的预概形 $T'$ 所成范畴之间的等价}。（这里用到：对两个
平展于 $S$ 上的预概形 $T_1,T_2$，映射
\[
  \operatorname{Hom}_S(T_1,T_2)\longrightarrow
  \operatorname{Hom}_{S'}(T'_1,T'_2)
\]
是双射，这一点很容易直接验证；还用到，当
$S'=(S,\mathcal O_S/\mathcal J)$ 时上述定理成立，其中

\src{310}
$\mathcal J$ 是 $\mathcal O_S$ 的凝聚幂零理想层（[4]，引理6）。）还应
注意，这里并没有假设所考察的 $T$ 在 $S$ 上有限。这个结论特别蕴含：态射
$S'\longrightarrow S$ 诱导 $S'$ 的基本群到 $S$ 的基本群的同构 $[\,]$
（\fquote{\emph{预概形基本群的拓扑不变性}}）。

\subsection*{4. 与一次上同调的关系}
\addcontentsline{toc}{subsection}{4. 与一次上同调的关系}

\medskip\noindent
a. 设 $\mathcal C$ 是任意两个对象的积都存在的范畴，$T\in\mathcal C$。
对每个有限集 $I\ne\varnothing$，可以考察 $T^I$；令 $I$ 变化，便得到从非空有限集
范畴到 $\mathcal C$ 的一个协变函子，也就是 $\mathcal C$ 的一个所谓
\emph{单纯对象}，记作 $K_T$。后者协变地依赖于 $T$；而且若 $u,v$ 是两个
态射 $T\longrightarrow T'$，则相应的态射
$K_T\longrightarrow K_{T'}$ 彼此同伦。若
$\operatorname{Hom}(T,T')\ne\varnothing$，就称 $T$ \emph{支配} $T'$；
这是 $\mathcal C$ 中一个向上滤过的预序关系。由前述事实可知，若 $T$ 支配
$T'$，就有一类在同伦意义下典范的单纯对象同态
$K_T\longrightarrow K_{T'}$；特别地，若 $K_T,K_{T'}$ 彼此支配，则它们
同伦等价。现在设 $F$ 是从 $\mathcal C$ 到某个阿贝尔范畴
$\mathcal C'$ 的函子（为确定起见，取作反变函子），那么
\[
  C^*(T,F)=F(K_T)
\]
是 $\mathcal C'$ 的一个余单纯对象，因而按熟知方式定义了
$\mathcal C'$ 中的一个（上链）复形；可以取它的上同调
\[
  H^*(T,F)=H^*(C^*(T,F))=H^*(F(K_T))
\]
（若可能混淆，可在 $H^*$ 的下标处添上 $\mathcal C$）。这是关于 $F$ 的
上同调函子；当 $T$ 变化时，它的变异性来自上面对 $K_T$ 所述的事实。更
确切地说，固定 $F$ 而令 $T$ 在 $\mathcal C$ 中变化（以支配关系为预序）
时，诸 $H^*(T,F)$ 构成 $\mathcal C'$ 中分次对象的一个归纳系；特别地，
若 $T,T'$ 彼此支配，则 $H^*(T,F)$ 与 $H^*(T',F)$ 典范同构。

假设 $\mathcal C$ 中存在纤维积。于是，对固定的 $S\in\mathcal C$，可将
上述构造应用于 $\mathcal C$ 中 $S$ 上对象所成的范畴 $\mathcal C_S$；若
要明确我们是在 $\mathcal C_S$ 中工作，就写 $C^*(T/S,F)$ 和
$H^*(T/S,F)$，而不写 $C^*(T,F)$ 和 $H^*(T,F)$。这样，

\src{311}
$C^*(T/S,F)$ 是 $\mathcal C'$ 中的一个上链复形，其 $n$ 次项等于
\[
  F(T\times_ST\times_S\cdots\times_ST)
\]
（括号中含 $n+1$ 个因子）。

还要照例指出，无须假设范畴 $\mathcal C'$ 是阿贝尔的，也可以定义
$H^0(T/S,F)$：它是态射对 $F(p_i)$（$i=1,2$）的核（定义2.1），若该核
存在；这里
\[
  F(T)\rightrightarrows F(T\times_ST)
\]
对应于两个投影 $p_1,p_2:T\times_ST\rightrightarrows T$。特别地，有自然
态射（称为\emph{增广}）
\[
  F(S)\longrightarrow H^0(T/S,F)
\]
在适当情形下它是同构（特别地，若 $T\longrightarrow S$ 是严格满态射而
$F$ \fquote{左正合}）。类似地，当 $F$ 取值于某个范畴 $\mathcal C''$
中的群所成范畴时，也可以定义 $H^1(T/S,F)$。当 $\mathcal C''$ 是集合范畴
时（也就是说，$F$ 取值于通常的、不必交换的群所成范畴），
$H^1(T/S,F)$ 是对 $C^1(T/S,F)=F(T\times_ST)$ 中由满足
\[
  F(p_{31})(g)=F(p_{32})(g)\,F(p_{21})(g)
\]
的 $g$ 所组成的子群 $Z^1(T/S,F)$ 作如下取商所得的集合：算子群 $F(T)$
作用于 $C^1(T/S,F)$，特别作用于其子集 $Z^1(T/S,F)$，作用公式为
\[
  \rho(g')\cdot g=F(p_2)(g')\,g\,F(p_1)(g')^{-1}.
\]

\medskip\noindent
b. 例如，设 $\mathcal F$ 是以 $\mathcal C$ 为基的纤维范畴。取
$\xi,\eta\in\mathcal F_S$，并对 $S$ 上每个 $S'$ 令
\[
  F_{\xi,\eta}(S')=
  \operatorname{Hom}(\xi\times_SS',\eta\times_SS').
\]
这样，$F_{\xi,\eta}$ 是从 $\mathcal C_S$ 到集合范畴的反变函子。在此
记号下，\emph{对每一对对象 $\xi,\eta\in\mathcal F_S$，增广态射}
\[
  F_{\xi,\eta}(S)\longrightarrow H^0(S'/S,F_{\xi,\eta})
\]
\emph{都是同构，正是说 $\alpha:S'\longrightarrow S$ 是
$\mathcal F$-下降态射}（定义1.7）。

\medskip\noindent
c. 类似地，对 $\xi\in\mathcal F_S$ 以及 $\mathcal C$ 中 $S$ 上的每个
对象 $S'$，令
\[
  G_\xi(S')=\operatorname{Aut}(\xi\times_SS') ;
\]
这样便定义了从 $\mathcal C_S$ 到群范畴的反变函子 $G_\xi$。

\src{312}
这时可见，$Z^1(S'/S,G_\xi)$ \emph{典范地等同于
$\xi'=\xi\times_SS'$ 相对于 $S'\longrightarrow S$ 的下降数据所成的集合}
（定义1.6），而 $H^1(S'/S,G_\xi)$ \emph{等同于如下对象的同构类集合：
$\mathcal F_{S'}$ 中配备了相对于 $\alpha:S'\longrightarrow S$ 的下降数据，
且作为 $\mathcal F_{S'}$ 中的对象同构于 $\xi'=\xi\times_SS'$ 的对象}。
因此，若 $\alpha:S'\longrightarrow S$ 是 $\mathcal F$-下降态射（参见（b）），
则 $H^1(S'/S,G_\xi)$ \emph{包含如下集合为其子集：$\mathcal F_S$ 中满足
$\eta\times_SS'$ 在 $\mathcal F_{S'}$ 中同构于 $\xi\times_SS'$ 的对象
$\eta$ 的同构类；这个包含是等号，当且仅当
$\xi'=\xi\times_SS'$ 相对于 $\alpha:S'\longrightarrow S$ 的每个下降数据
都有效}。（特别地，若 $\alpha:S'\longrightarrow S$ 是严格
$\mathcal F$-下降态射，便是这种情形。）

\result{注}{} $C^*(T/S,F)$ 这一类型的上链复形包含大多数熟知的标准复形
作为特殊情形（Čech 上同调、群上同调等等），并在代数几何中起重要作用，
尤其是在预概形的\fquote{Weil 上同调}中。

\medskip\noindent
d. \result{例}{1} 设 $S'$ 是 $S\in\mathcal C$ 上的一个对象，$\Gamma$ 是
$S'$ 的一个自同构群，使 $S'$ \fquote{在 $S$ 上形式地为主齐的，群为 $\Gamma$}；
也就是说，自然态射
\[
  \Gamma\times S'\longrightarrow S'\times_SS',
\]
是同构，其中 $\Gamma\times S'$ 表示 $\Gamma$ 份 $S'$ 的余积。（假设
$\mathcal C$ 中存在这里涉及的余积。）设 $F$ 是从 $\mathcal C$ 到阿贝尔群
范畴的反变函子。那么 $C^*(S'/S,F)$ 典范同构于标准齐次上链
$C^*(\Gamma,F(S'))$ 所成的单纯群，因而 $H^*(S'/S,F)$ 典范同构于
$H^*(\Gamma,F(S'))$。

\medskip\noindent
e. \result{例}{2} 设 $\mathcal C$ 是预概形范畴。以 $G_a$
（\fquote{加法群}）表示从 $\mathcal C$ 到阿贝尔群范畴的反变函子
\[
  G_a(X)=H^0(X,\mathcal O_X).
\]
类似地，定义函子 $G_m$（\fquote{乘法群}）为
\[
  G_m(X)=H^0(X,\mathcal O_X)^*
\]
（即环 $H^0(X,\mathcal O_X)$ 中可逆元所成的群），更一般地，定义函子
$\operatorname{Gl}(n)$（\fquote{$n$ 阶一般线性群}）为

\src{313}
\[
  \operatorname{Gl}(n)(X)=
  \operatorname{Gl}\bigl(n,H^0(X,\mathcal O_X)\bigr),
\]
这是从 $\mathcal C$ 到群范畴的函子（当 $n>1$ 时群不必交换；当 $n=1$
时便得到 $G_m$）。还可以把 $\operatorname{Gl}(n)$ 解释为一个自同构函子
（参见（c））：考察以 $\mathcal C$ 为基的纤维范畴 $\mathcal F$，其中对
$X\in\mathcal C$，$\mathcal F_X$ 是 $X$ 上局部自由层所成的范畴；事实上
\[
  \operatorname{Gl}(n)(X)=
  \operatorname{Aut}_{\mathcal F_X}(\mathcal O_X^n).
\]
由（b）可知，若 $\alpha:S'\longrightarrow S$ 是 $\mathcal F$-下降态射
（参见第2.c节），则 $H^1(S'/S,\operatorname{Gl}(n))$ 包含 $S$ 上如下局部
自由层的同构类集合：其在 $S'$ 上的逆像同构于 $\mathcal O_{S'}^n$；这个
包含是等号，当且仅当 $\mathcal O_{S'}^n$ 上相对于
$\alpha:S'\longrightarrow S$ 的每个下降数据都有效。当 $S$ 是局部环的谱
时，每个 $S$ 上的局部自由层都平凡，所以这就是说
$H^1(S'/S,\operatorname{Gl}(n))=(e)$。

指出态射 $\alpha:S'\longrightarrow S$ 的下列条件等价：
\begin{enumerate}
\item[(i)] 增广同态
\[
  H^0(S,\mathcal O_S)=G_a(S)\longrightarrow H^0(S'/S,G_a)
\]
是同构；
\item[(ii)] $S'\longrightarrow S$ 是 $\mathcal F$-下降态射（$\mathcal F$
是上面考察的以 $\mathcal C$ 为基的纤维范畴）。
\end{enumerate}
若 $S'\longrightarrow S$ 有限，这些条件还等价于：
\begin{enumerate}
\item[(iii)] $S'\longrightarrow S$ 是严格满态射（参见第2.c节）。
\end{enumerate}

现在假设 $S=\operatorname{Spec}(A)$、$S'=\operatorname{Spec}(A')$；则
\[
  C^n(S'/S,G_a)=C^n(A'/A,G_a)=\bigotimes_A^{n+1}A',
\]
上边缘算子
$C^n(A'/A,G_a)\longrightarrow C^{n+1}(A'/A,G_a)$ 是诸面算子的交错和，
其中
\[
  \partial_i(x_0\otimes x_1\otimes\cdots\otimes x_n)
  =x_0\otimes\cdots\otimes x_{i-1}\otimes1_{A'}\otimes
   x_i\otimes\cdots\otimes x_n.
\]
同样，
\[
  C^n(S'/S,G_m)=C^n(A'/A,G_m)
  \simeq\left(\bigotimes_A^{n+1}A'\right)^*,
\]
$C^*(A'/A,G_m)$ 中的单纯运算由 $C^*(S'/S,G_a)$ 中的运算诱导；
$C^*(A'/A,\operatorname{Gl}(n))$ 中的单纯运算也可类似地具体写出。
\emph{在报告人所知的一切情形中，都有}
\[
  H^i(A'/A,G_a)=0\qquad\text{对 $i>0$},
\]
\emph{并且若 $A$ 是局部环，则有}
\[
  H^1(A'/A,G_m)=0
  \quad\text{并且更一般地}\quad
  H^1(A'/A,\operatorname{Gl}(n))=(e)
\]
（这里 $S'\longrightarrow S$ 是 $\mathcal F$-下降态射，也就是说图
$A\longrightarrow A'\rightrightarrows A'\otimes_AA'$ 正合；比较第2.c节）。
还应注意，\emph{Hilbert 的\fquote{定理90}无非就是关系}
\src{314}
\[
  H^1(S'/S,G_m)=0
\]
\emph{在如下情形下的成立：$A$ 是域，$A'$ 是它的有限 Galois 扩张（参见
例1）；还可以说，在这种情形下，对秩为1的局部自由层所成纤维范畴，
$S'\longrightarrow S$ 是严格下降态射。}应当正是在后一形式下推广 Hilbert
定理：既改变对态射 $S'\longrightarrow S$ 的假设，也改变对所考察拟凝聚层
的假设。

最后，设 $A$ 是极大理想为 $\mathfrak m$ 的阿廷局部环，$A'$ 是
$A$-代数；对每个整数 $k>0$，用 $A_k$（相应地 $A'_k$）表示环
$A/\mathfrak m^{k+1}$（相应地 $A'/\mathfrak m^{k+1}A'$）。指出下列性质
等价：
\begin{enumerate}
\item[(i)] 对每个 $k$，$H^1(A'_k/A_k,G_a)=0$。
\item[(ii)] 对每个 $k$，$H^1(A'_k/A_k,G_m)=0$。
\item[(iii)] 对每个 $k$ 和每个 $n$，
$H^1(A'_k/A_k,\operatorname{Gl}(n))=(e)$。
\end{enumerate}
若 $S'\longrightarrow S$ 是严格满态射，则上述条件甚至蕴含：对
$A'$ 上的自由模（无论是否有限型），它是严格下降态射。

\result{注}{} 当 $S,S'$ 分别是域 $A,A'$ 的概形时，群
$H^i(S'/S,G_m)$ 的定义归功于 Amitsur。群 $H^2(S'/S,G_m)$ 尤其有趣，
它是 Brauer 群的一种\fquote{整体}变体；关于这种变体，可参阅 [1] 第VII章。

\fsection{B}{沿忠实平坦态射的下降}

\subsection*{1. 下降定理的陈述}
\addcontentsline{toc}{subsection}{1. 下降定理的陈述}

\result{定义}{1.1} 预概形态射 $\alpha:S'\longrightarrow S$ 称为
\emph{平坦}，若对每个 $x'\in S'$，$\mathcal O_{x'}$ 是环
$\mathcal O_{\alpha(x')}$ 上的平坦模（也就是说，关于
$\mathcal O_{\alpha(x')}$-模 $M$ 的函子
$\mathcal O_{x'}\otimes_{\mathcal O_{\alpha(x')}}M$ 正合）。一个态射若
平坦且满射，就称为\emph{忠实平坦}。

例如，若 $S=\operatorname{Spec}(A)$、$S'=\operatorname{Spec}(A')$，则
$S'$ 在 $S$ 上平坦，当且仅当 $A'$ 是平坦 $A$-模；$S'$ 在 $S$ 上忠实
平坦，当且仅当 $A'$ 是忠实平坦 $A$-模（也就是说，关于 $A$-模 $M$ 的
函子 $A'\otimes_A M$ 正合且忠实）。用 Serre [5] 的术语，这也就是说环对
$(A,A')$ 平坦。若 $S'$ 在 $S$ 上忠实平坦，则 $S$ 上拟凝聚层的逆像函子
正合且忠实；换言之，$S$ 上拟凝聚层同态所成的一个序列

\src{315}
正合，当且仅当它在 $S'$ 上的逆像序列正合（特别地，$S$ 上拟凝聚层同态
是单态射、满态射或同构，当且仅当其在 $S'$ 上的逆像分别是单态射、满态射
或同构）。把讨论限制在 $S$ 的任意开集上，这个性质仍然成立；在这种形式
下，它刻画忠实平坦态射。

\result{定义}{1.2} 态射 $\alpha:S'\longrightarrow S$ 称为
\emph{拟紧}，若 $S$ 的每个拟紧开子集 $U$ 的逆像都是拟紧的（也就是说，
是有限多个仿射开集的并）。

显然，只须对 $S$ 的仿射开集验证这一性质。例如，仿射态射（即仿射开集的
逆像仍是仿射开集的态射）是拟紧的。

平坦态射、忠实平坦态射、拟紧态射各自所成的类都对复合及\fquote{基变换}
封闭，并且当然都包含同构。

\result{定理}{1} 设 $\alpha:S'\longrightarrow S$ 是预概形间的忠实平坦
拟紧态射。则对拟凝聚层所成的纤维范畴 $\mathcal F$（A，第1节，例2），
$\alpha$ 是\emph{严格下降态射}（A，定义1.7）。

这个陈述包含两件事：
\begin{enumerate}
\item[(i)] 若 $F,G$ 是 $S$ 上两个拟凝聚层，$F',G'$ 是它们在 $S'$ 上的
逆像，则自然同态
\[
  \operatorname{Hom}(F,G)\longrightarrow\operatorname{Hom}(F',G')
\]
把左端双射到右端的如下子群：其中的同态 $F'\longrightarrow G'$ 与这些层
上的典范下降数据相容；也就是说，经 $S''=S'\times_SS'$ 到 $S'$ 的两个投影
取逆像后，它们给出同一个同态 $F''\longrightarrow G''$。
\item[(ii)] $S'$ 上每个配备了相对于态射 $\alpha:S'\longrightarrow S$ 的
下降数据（A，定义1.6）的拟凝聚层 $F'$，都同构于 $S$ 上某个拟凝聚层 $F$
的逆像（并且这个同构连同所给下降数据一起成立）。
\end{enumerate}
在（i）中令 $F=\mathcal O_S$，得到：

\result{推论}{1} 设 $G$ 是 $S$ 上的拟凝聚层，$G',G''$ 分别是它在 $S'$
及 $S''=S'\times_SS'$ 上的逆像，$p_1,p_2$ 是 $S''$ 到 $S'$ 的两个投影。
则下列集合映射图

\src{316}
\[
  \Gamma(G)\xrightarrow{\alpha^*}\Gamma(G')
  \underset{p_2^*}{\overset{p_1^*}{\rightrightarrows}}\Gamma(G'')
\]
正合（A，定义1.(a)）。

另一方面，把（i）、（ii）与定义1.1结合起来，得到：

\result{推论}{2} 设 $G$ 如推论1。则 $G$ 的拟凝聚子层与 $G'$ 的如下拟凝聚
子层之间存在一一对应：经两个投影 $p_1,p_2$ 取到 $S''$ 上的逆像给出
$G''$ 的同一个子层。

当然，也有用商层表述的等价陈述。如 A 第4节（e）所述，定理1应看作 Hilbert
\fquote{定理90}的推广，并且把用一次上同调表述的若干结果包含为特殊情形。
为作证明，很容易归结到 $S=\operatorname{Spec}(A)$、
$S'=\operatorname{Spec}(A')$ 的情形；而对（i），又很容易归结为证明推论1，
也就是证明：对每个 $A$-模 $M$，图
\[
  M=A\otimes_A M\longrightarrow A'\otimes_A M
  \rightrightarrows A'\otimes_AA'\otimes_A M
\]
正合。这由下述更一般的引理推出：

\result{引理}{1.1} 设 $A'$ 是忠实平坦 $A$-代数。则对每个 $A$-模 $M$，
以 $M$ 为增广的复形 $C^*(A'/A,G_a)\otimes_A M$（参见 A 第4节（e））
是 $M$ 的一个\emph{消解}。

只须证明：把上述复形从 $A$ 基变换到 $A'$ 后得到的增广复形满足同样结论。
这就要在以 $A'$ 代替 $A$、以 $A'\otimes_AA'$ 代替 $A'$ 后验证该陈述，
所以归结到存在 $A$-代数同态 $A'\longrightarrow A$ 的情形（用几何术语，
就是 $S'$ 在 $S$ 上有一个截面的情形）。在这种情形下，结论来自 A 第4节
（a）的一般事实。顺便指出如下推论，它推广了 Galois 上同调中一个熟知的
陈述（比较 A 第4节（e））：

\result{推论}{} 若 $A'$ 在 $A$ 上忠实平坦，则 $H^0(A'/A,G_a)=A$，且
\[
  H^i(A'/A,G_a)=0\qquad\text{对 $i\geq1$}.
\]

为证明定理1的（ii），同（i）一样归结到 $S'$ 在 $S$ 上有截面的情形；
此时结论由（i）得出（参见 A 第1节（c））。

显然，可以任意改变定理1及其推论，在所考察的拟凝聚层（或层系）上加入
各种附加结构。例如，在 $S$ 上给定一个拟凝聚交换代数层，\fquote{等价于}
在 $S'$ 上给定一个这样的层以及相对于 $\alpha:S'\longrightarrow S$ 的下降
数据。考虑到 $S$ 上这类拟凝聚层与 $S$ 上仿射预概形之间的函子性对应，便
得到下一个定理的第二个断言：

\src{317}
\result{定理}{2} 设 $\alpha:S'\longrightarrow S$ 如定理1。则 $\alpha$
是\emph{下降态射}（一般不严格）（A，定义2.4）；而对预概形上的仿射概形
所成纤维范畴，它是\emph{严格下降态射}（A，定义1.7）。

定理的第一个断言是说：设 $X,Y$ 是 $S$ 上两个预概形，$X',Y'$ 是它们在
$S'$ 上的逆像，$X'',Y''$ 是它们在 $S''=S'\times_SS'$ 上的逆像；则下列
自然映射图
\[
  \operatorname{Hom}_S(X,Y)\xrightarrow{\alpha^*}
  \operatorname{Hom}_{S'}(X',Y')
  \underset{p_2^*}{\overset{p_1^*}{\rightrightarrows}}
  \operatorname{Hom}_{S''}(X'',Y'')
\]
正合。也就是说，$\alpha^*$ 把 $\operatorname{Hom}_S(X,Y)$ 双射到
$\operatorname{Hom}_{S'}(X',Y')$ 中如下态射所成的部分：它们与 $X',Y'$
上的典范下降数据相容（即经 $S''$ 到 $S'$ 的两个投影取逆像后彼此相等）。
当把 $Y$ 限制为 $S$ 上的仿射概形时，这很容易由定理1、推论1得到；在一般
情形下，还必须把定理1与下述结果结合起来：

\result{引理}{1.2} 设 $\alpha:S'\longrightarrow S$ 是忠实平坦拟紧态射。
则 $S$ 可与 $S'$ 的一个\emph{商拓扑空间}等同；也就是说，$S$ 的每个子集
$U$，若 $\alpha^{-1}(U)$ 是开集，则 $U$ 也是开集。

为完成定理2，还须给出 $S'$-预概形 $X'$ 上下降数据的有效性判据（此处不
假设 $X'$ 在 $S'$ 上仿射）。首先指出，\emph{这样的下降数据未必有效}；
即使 $S$ 是域 $k$ 的谱，$S'$ 是 $k$ 的一个二次扩张 $k'$ 的谱，而 $X'$
是 $S'$ 上的一个2维固有代数概形，情形仍然如此（按 Serre 的方法，利用
Nagata 的非射影曲面即可看出）。\emph{$X'/S'$ 上相对于忠实平坦拟紧态射
$\alpha:S'\longrightarrow S$ 的一个下降数据有效，当且仅当 $X'$ 可写成
在 $S'$ 上仿射、并在 $X'$ 的该下降数据下\fquote{稳定}的开集 $X'_i$ 的并。}
若态射 $\alpha:S'\longrightarrow S$ 是\emph{纯不可分态射}（即单射，并且其
剩余域扩张

\src{318}
都是纯不可分的）。还可以证明，若 $\alpha:S'\longrightarrow S$
\emph{有限}，并且 $X'$ 中含于 $X'$ 在 $S$ 上某个纤维的每个有限子集，都
含于 $X'$ 的一个在 $S$ 上仿射的开集，那么结论仍然成立（这就是
\emph{Weil 判据}）。特别地，若 $X'/S'$ 拟射影，便是这种情形；这时还可
证明，\fquote{下降}所得的预概形 $X/S$ 也拟射影（若 $X'/S'$ 射影，则它也
射影）。综上：

\result{定理}{3} 设 $\alpha:S'\longrightarrow S$ 是预概形间的忠实平坦
拟紧态射。若 $\alpha$ 是\emph{纯不可分态射}，则它是\emph{严格下降态射}。若
$\alpha$ 有限，则相对于预概形上的拟射影（或射影）预概形所成纤维范畴，
它是严格下降态射。

\result{注}{} 报告人不知道，上述第二个断言中 $\alpha$ 为\emph{有限}态射
的假设是否确实必要；无论如何，可以\fquote{形式地}验证，它可由下面这个
表面上更弱的假设代替：\emph{$S$ 的每一点都有一个开邻域 $U$，存在一个在
$U$ 上有限且忠实平坦的 $U'$，并且存在从 $U'$ 到 $S'$ 的 $S$-态射。}
前一情形未涵盖的一个典型例子是
$S=\operatorname{Spec}(A)$、$S'=\operatorname{Spec}(\widehat A)$，其中
$A$ 是诺特局部环，$\widehat A$ 是其完备化；另一个例子是 $S'$ 在 $S$
上拟有限（即局部同构于一个有限 $S$-概形的开子集），但并不有限。在这两种
情形下，报告人也不知道下述问题的答案：设 $X$ 是 $S$-概形，并且
$X'=X\times_SS'$ 在 $S'$ 上射影，$X$ 是否在 $S$ 上射影？\footnote{1962年
勘误给出了否定答案：拟有限、平展且满射的态射
$S'\longrightarrow S$，或态射
$\operatorname{Spec}(\widehat A)\longrightarrow\operatorname{Spec}(A)$，
未必是严格下降态射；即使 $A$ 是代数闭域 $k$ 上一条代数曲线的局部环，且
$S=\operatorname{Spec}(A)$，情形仍然如此。确实可以找到固有光滑态射
$f:X\longrightarrow S$，使 $X$ 成为以椭圆曲线 $E$ 为结构群、以 $S$ 为基
的主丛，并且 $f':X'\longrightarrow S'$ 射影而 $f$ 不射影。因此，这也给出
阿贝尔概形下非等平凡的主齐性丛的一个例子。}

\subsection*{2. 某些态射性质的下降}
\addcontentsline{toc}{subsection}{2. 态射性质的下降}

设 $P$ 是预概形态射的一个类。设 $\alpha:S'\longrightarrow S$ 是预概形
态射，$f:X\longrightarrow Y$ 是 $S$-预概形间的态射，而
$f':X'\longrightarrow Y'$ 是 $f$ 经 $\alpha$ 的逆像。可以问：关系
\fquote{$f'\in P$} 是否蕴含\fquote{$f\in P$}。可以看出，在许多重要
情形下，只要假设 $\alpha$ \emph{忠实平坦且拟紧}，答案就是肯定的（其中
后一假设有时是多余的）。当 $P$ 分别是满射的态射、纯不可分态射（这两种情形来自
$\alpha$ 的满性）、平坦态射、忠实平坦态射、光滑态射（这几种情形来自
$\alpha$ 的忠实平坦性），或有限型态射所成的类时，这一点直接可见而毫无
困难。利用定理1、定理2与引理1.2，还可看出，当 $P$ 是以下各类之一时也是
如此：同构、开浸入、闭浸入、浸入（若 $f$ 有限型且 $Y$ 局部诺特）、
仿射态射、有限态射、拟有限态射、开态射、闭态射、同胚、分离态射，

\src{319}
以及固有态射。唯一尚未弄清的重要情形，就是第1节注记已经指出的射影与
拟射影态射。

\subsection*{3. 沿有限忠实平坦态射的下降}
\addcontentsline{toc}{subsection}{3. 沿有限忠实平坦态射的下降}

设 $\alpha:S'\longrightarrow S$ 是有限态射，对应于 $S$ 上的代数层
$\mathcal A'$；作为模层，$\mathcal A'$ 局部自由、有限型并且处处 $\ne0$。
于是 $\alpha$ 是忠实平坦拟紧态射，可以应用上述结果。在 $S'$ 上给定拟凝聚
层 $F'$，等价于在 $S$ 上给定拟凝聚层 $\alpha_*(F')$，连同它的
$\mathcal A'$-模结构（注意 $\mathcal A'=\alpha_*(\mathcal O_{S'})$）。
为简便起见，仍以 $F'$ 表示 $S$ 上这个层。$F'$ 在
$S'\times_SS'$ 上的两个逆像 $p_i^*(F')$，同样对应于下列
$(\mathcal A'\otimes_{\mathcal O_S}\mathcal A')$-模的拟凝聚层：
\[
  F'\otimes_{\mathcal O_S}\mathcal A'
  \quad\text{和}\quad
  \mathcal A'\otimes_{\mathcal O_S}F'.
\]
从前者到后者给定一个 $(S'\times_SS')$-同态，等价于给定一个
$(\mathcal A'\otimes_{\mathcal O_S}\mathcal A')$-模同态；又因为
$\mathcal A'$ 局部自由，这等价于给定一个
$(\mathcal A'\otimes\mathcal A')$-模同态
\[
  \mathcal U=
  \operatorname{Hom}_{\mathcal O_S}(\mathcal A',\mathcal A')
  =\mathcal A'\otimes\check{\mathcal A}'
  \longrightarrow
  \operatorname{Hom}_{\mathcal O_S}(F',F'),
\]
也就是，对 $\mathcal U$ 在开集 $V$ 上的每个截面 $\xi$，给定一个满足下列
条件的 $\mathcal O_S$-模同态
$T_\xi:F'|_V\longrightarrow F'|_V$：
\begin{equation}
  T_{f\xi}(x)=fT_\xi(x),
  \qquad T_{\xi f}(x)=T_\xi(fx),
  \tag{3.1}
\end{equation}
其中 $f,x$ 分别是 $\mathcal A',F'$ 在 $S$ 的某个含于 $V$ 的开集上的
截面。于是下降数据的条件（i）、（ii）（A，第1节，（c））分别写成
\begin{equation}
  T_{1_{\mathcal U}}x=x,
  \qquad\text{即}\qquad
  T_{1_{\mathcal U}}=\operatorname{id}_{F'},
  \tag{3.2}
\end{equation}
\begin{equation}
  T_{\xi\eta}=T_\xi T_\eta.
  \tag{3.3}
\end{equation}
换言之，\emph{$F'$ 上的一个下降数据，等价于给出 $\mathcal O_S$-代数层
$\mathcal U=\operatorname{Hom}_{\mathcal O_S}(\mathcal A',\mathcal A')$
在 $F'$ 上的一个表示，即一个 $\mathcal O_S$-代数同态
$\mathcal U\longrightarrow\operatorname{Hom}_{\mathcal O_S}(F',F')$，并满足（3.1）的两个线性
条件。}若在 $S'$ 上有拟凝聚层的一个配对
\[
  F'_1\times F'_2\longrightarrow F'_3
\]
（也可以把它解释为 $S$ 上 $\mathcal A'$-模层的配对），

\src{320}
并且 $F'_i$ 上有由同态
\[
  T_i:\mathcal U\longrightarrow
  \operatorname{Hom}_{\mathcal O_S}(F'_i,F'_i)
  \qquad (i=1,2,3),
\]
定义的粘合数据，那么这些数据与所给配对按显然意义相容，当且仅当满足下列
条件。

对 $\mathcal U$ 在某个开集上的每个截面 $\xi$，以
\[
  \Delta\xi=\sum_i\xi'_i\otimes_{\mathcal A'}\xi''_i
\]
表示 $\mathcal U\otimes_{\mathcal A'}\mathcal U$ 的截面（这里把
$\mathcal U$ 按其左模结构看作 $\mathcal A'$-模）；它由公式
\[
  \xi\mathbin{\cdot}(fg)=\sum_i\xi'_i(f)\xi''_i(g)
\]
定义（$f,g$ 是 $\mathcal A'$ 在某个更小开集上的两个截面）。那么，对
$F'_1,F'_2$ 分别在某个更小开集上的每一对截面 $x,y$，都有公式
\begin{equation}
  T^{(3)}_\xi(x\mathbin{\cdot}y)
  =\sum_iT^{(1)}_{\xi'_i}x\mathbin{\cdot}T^{(2)}_{\xi''_i}y
  \tag{3.4}
\end{equation}
成立。\footnote{印本作\fquote{$\mathcal A'$ 的截面}。$T^{(1)}$ 与
$T^{(2)}$ 的定义域以及配对 $F'_1\times F'_2\longrightarrow F'_3$ 都要求
$x,y$ 分别是 $F'_1,F'_2$ 的截面。}（也可以说，相对于所给配对，诸同态
$T^{(i)}$ 与 $\mathcal U$ 的对角映射相容。）特别地，公式（3.1）至（3.4）
使我们能够用带对角映射的代数表示来解释 $S'$ 上拟凝聚代数层的下降数据；
因而，限制到交换代数后，也能解释 $S'$ 上仿射概形的下降数据。

由此可过渡到 $S'$ 上任意预概形 $X'$ 的下降数据的类似解释：给定这样的
$X'$，等价于给定 $S$ 上一个预概形 $X'$，连同一个 $\mathcal O_S$-代数
同态
\[
  \mathcal A'\longrightarrow\mathcal O_{X'},
\]
而给定 $X'$ 上的下降数据，等价于给定一个与态射
$h:X'\longrightarrow S'$ 相容、并满足上述（3.1）至（3.4）之类似条件的
层同态
\[
  \mathcal U\longrightarrow
  \operatorname{Hom}_{h^{-1}(\mathcal O_S)}
  (\mathcal O_{X'},\mathcal O_{X'}),
\]

\textsc{例1（\fquote{Weil}）。}——假设 $S'/S$ 是 Galois 群为
$\Gamma$ 的 Galois 平展覆盖（参见 A 第3节和第4.d节）。那么，$S'$ 上拟凝聚
层 $F'$（相应地，$S'$-预概形 $X'$）的下降数据，等价于用
$(S',F')$（相应地，$(S',X')$）的自同构给出 $\Gamma$ 的表示，并且该表示
与 $\Gamma$ 在 $S'$ 上的作用相容。这个结论

\src{321}
是\fquote{形式的}，也就是说，可以用范畴语言证明；但从本节的观点来看，
它也来自 $\mathcal U$ 的显式结构（包括其环结构、环同态
$\mathcal A'\longrightarrow\mathcal U$ 以及对角映射），而下述结果完全
确定了这一结构：作为左 $\mathcal A'$-模，$\mathcal U$ 有一个基，基中
各截面恰好对应于 $\Gamma$ 的元素。

\textsc{例2（\fquote{Cartier}）。}——设 $p$ 是素数，并假设
$p\mathcal O_S=0$（即 $\mathcal O_S$ 的特征为 $p$）、
$(\mathcal A')^p\subset\mathcal O_S=\mathcal A$（即 $S'/S$ 是高为1的
纯不可分态射），且 $\mathcal A$-代数层 $\mathcal A'$ 局部有一个 $p$-基，
也就是说，有一族截面 $(x_i)$，使 $\mathcal A'$ 由 $x_i$ 作为代数生成，
它们只服从关系 $x_i^p=0$。假设指标 $i$ 所成集合有限，其基数为 $n$。令
$\mathfrak g$ 为 $\mathcal A'$ 在 $\mathcal A$ 上的导子层；它是秩为 $n$
的局部自由 $\mathcal A'$-模层，并且还是 $\mathcal A$ 上的限制李代数层（$p$-李代数层）
（但不是 $\mathcal A'$ 上的），满足条件
\begin{equation}
  [X,fY]=X(f)Y+f[X,Y].
  \tag{3.5}
\end{equation}

\textsc{引理。}——$\mathcal U=
\operatorname{Hom}_{\mathcal O_S}(\mathcal A',\mathcal A')$ 作为一个带有
代数同态 $\mathcal A'\longrightarrow\mathcal U$ 的 $\mathcal O_S$-代数，
由左 $\mathcal A'$-子模 $\mathfrak g$ 生成，并附加关系
\begin{equation}
  \left\{
  \begin{aligned}
    Xf-fX&=X(f),\\
    XY-YX&=[X,Y],\\
    X^p&=X^{(p)}.
  \end{aligned}
  \right.
  \tag{3.6}
\end{equation}

由上述引理，$S'$ 上拟凝聚层 $F'$ 的下降数据等价于：对每个
$X\in\mathfrak g$ 给定 $F'$ 的一个 $\mathcal O_S$-自同态
$\overline X$，满足条件
\begin{align}
  \overline{fX}&=f\overline X, \tag{3.7}\\
  \overline X(fx)&=X(f)x+f\overline X(x), \tag{3.8}\\
  \overline{[X,Y]}&=[\overline X,\overline Y], \tag{3.9}\\
  \overline{X^{(p)}}&=\overline X^p. \tag{3.10}
\end{align}
（这可以称为 $F'$ 上一个无曲率、并与 $p$ 次幂相容的线性联络。）同样可以
具体写出 $S'$-预概形 $X'$ 上下降数据的概念；这里关系（3.4）由如下条件
代替：诸 $\overline X$ 都是 $\mathcal O_{X'}$ 的导子。由于态射
$S'\longrightarrow S$ 是纯不可分态射，定理3保证每个这样的下降数据都

\src{322}
有效，因而定义一个 $S$-预概形 $X$。

应当注意，我们没有对 $F'$（相应地 $X'$）作任何平坦性、非奇异性或有限性
假设。

\subsection*{4. 有理性判据的应用}
\addcontentsline{toc}{subsection}{4. 有理性判据的应用}

设 $X$ 是 $S$-预概形，并且 $\mathcal O_X$ 在 $S$ 上的正像是
$\mathcal O_S$；经过基 $S$ 的任何平坦扩张 $S'\longrightarrow S$，这个
性质仍然成立。若 $F$ 是 $X$ 上的可逆层（即秩为1的局部自由层），则 $F$
的自同构可与 $\mathcal O_X$ 的可逆截面等同，从而与 $\mathcal O_S$ 的
可逆截面一一对应。现在设 $s$ 是 $X$ 在 $S$ 上的一个截面；我们形象地把
$S$ 上可逆层 $s^*(F)$ 的一个截面称为 $F$ 在 $s$ 上方的一个截面。由前述
事实可知，若 $F_i$（$i=1,2$）是 $X$ 上两个可逆层，各自配有一个在 $s$
上方的截面，并且 $F_1,F_2$ 同构，那么恰有一个从 $F_1$ 到 $F_2$、与所述
截面相容的同构（即把第一个截面变成第二个截面的同构）。此外，与截面 $s$
无关地，约定把 $X$ 上两个可逆层 $F_1,F_2$ 看作等价，若 $S$ 的每一点都有
一个开邻域 $U$，使 $F_1,F_2$ 在 $X|U$ 上的限制同构。那么，$X$ 上每个
可逆层 $F$ 都等价于一个配备了 $s$ 上方标记截面的可逆层 $F_1$（取
$F_1=F\,s^*(F)^{-1}$），而 $F_1$ 在同构意义下唯一确定。换言之，$X$
上可逆层在等价意义下的分类，与配备标记截面的可逆层在同构意义下的分类
相同。

因为这些性质经过基的平坦扩张 $\alpha:S'\longrightarrow S$ 后仍成立
（以 $s$ 经 $\alpha$ 的逆像 $s'$ 代替 $s$），结合定理1便得到：

\emph{设预概形 $X/S$ 如上并有截面 $s$，设
$\alpha:S'\longrightarrow S$ 是忠实平坦拟紧态射，并设 $F'$ 是
$X'=X\times_SS'$ 上的可逆层。则 $F'$ 等价于 $X$ 上某个可逆层 $F$ 在
$X'$ 上的逆像，当且仅当它在
$X'\times_XX'=X\times_S(S'\times_SS')$ 上的两个逆像
$p_1^*(F')$ 与 $p_2^*(F')$ 等价。若此条件成立，$F$ 在等价意义下唯一
确定。}（这时称 $F'$ 在 $S$ 上有理。）

受这一原理启发，当 $\alpha:S'\longrightarrow S$ 如上一节例1或例2时，
便得到 Weil 或 Cartier 的有理性判据。（应当注意，这两位作者只考察
$S,S'$

\src{323}
都是域的谱的情形；这时 $S$ 更是局部环的谱，而上面引入的等价关系无非就是
同构关系。）在第一种情形下，$F'$ 在 $S$ 上有理，当且仅当它经 $\Gamma$
的各个变换都与 $F'$ 等价。为表述第二种情形下的有理性判据，一般地考察
$X'/X$ 的对角态射
\[
  X'\longrightarrow X''=X'\times_XX'
\]
及 $X'\times_XX'$ 上相应的理想层 $\mathcal I$，再考察层
$\mathcal I/\mathcal I^2$；它可与其在 $X'$ 上的逆像
$\Omega^1_{X'/X}$ 等同。因为 $F''_i=p_i^*(F')$（$i=1,2$）在对角线上的
限制彼此同构（都同构于 $F'$），也就是说
$F''_1(F''_2)^{-1}=F''$ 在对角线上的限制平凡，所以 $F''$ 在
$(X'',\mathcal O_{X''}/\mathcal I^2)$ 上的限制，在同构意义下由下列群中
唯一确定的一个元素 $\xi$ 给出：
\[
  H^1(X'',\mathcal I/\mathcal I^2)
  =H^1(X',\Omega^1_{X'/X}).
\]
此外，在当前情形下
\[
  \Omega^1_{X'/X}
  =\Omega^1_{S'/S}\otimes_{\mathcal O_S}\mathcal O_X,
\]
因而若 $\Omega^1_{S'/S}$ 在 $S$ 上局部自由（如 Cartier 情形），则 $\xi$
定义 $S'$ 上层
\[
  R^1f'_*(\mathcal O_{X'})\otimes\Omega^1_{S'/S}
\]
的一个截面，称为 $X'/S$ 上可逆层 $F'$ 的 Atiyah--Cartier 类。这个类为零，
当且仅当 $F'$ 经
\[
  (X'',\mathcal O_{X''}/\mathcal I^2)
  =X\times_S(S'',\mathcal O_{S''}/\mathcal J^2)
\]
到 $X'$ 的两个投影所得的逆像彼此等价（这里 $\mathcal J$ 是
$S''=S'\times_SS'$ 上由对角态射
$S'\longrightarrow S'\times_SS'$ 定义的理想层）。因此，要使 $F'$ 在
$X''=X\times_SS''$ 本身上的两个逆像等价，这个类为零显然是必要的；从而，
要使 $F'$ 等价于 $X$ 上某个可逆层 $F$ 的逆像，这个条件也是必要的。
Atiyah--Cartier 类还可解释为如下障碍：在 $S'$ 上局部地，$F'$ 上是否存在
一个相对于 $X'/X$ 的导子的联络；而当知道 $S'/S$ 的各导子自然延拓到
$X'$ 后所对应的 $F'$ 的导子时，这个联络也随之确定。由此以及上一节的展开，
容易推出：在上一节例2的情形下，并且当 $X/S$ 有截面时，Atiyah--Cartier
类为零也足以保证 $F'$ 在 $S$ 上有理。

\subsection*{5. 阿贝尔概形中基概形限制的应用}
\addcontentsline{toc}{subsection}{5. 阿贝尔概形中基概形的限制}

设 $S$ 是预概形。所谓 $S$ 上的阿贝尔概形，是指在 $S$ 上光滑且固有的概形
$X$，并且对每个 $x\in S$，其纤维都是 $\kappa(x)$ 上的阿贝尔簇。

\src{324}
假设 $S$ 诺特且正则（即各局部环正则）。利用 Murre [4] 的连通性定理，
可以证明（至少在\fquote{等特征}情形下，而目前所引定理已在该情形证明）
$X$ 在 $S$ 上的每个有理截面都处处有定义（即确为一个截面）；这推广了
Weil 的一个经典定理。更一般地，由此可知，若 $X'$ 是 $S$ 上的光滑概形，
则从 $X'$ 到 $X$ 的每个有理 $S$-映射都处处有定义。于是得到下述结果，它
推广了 Chow--Iang 的一个结果：\footnote{印本作\fquote{Igusa--Lang}；1962年
勘误指定改为\fquote{Chow--Iang}。}设 $S$ 诺特且正则，$K$ 表示其有理函数
环（若干域的直积），并设 $X$ 是 $K$ 上的阿贝尔概形。若 $X$ 同构于形如
$X_0\times_S\operatorname{Spec}(K)$ 的 $K$-概形，其中 $X_0$ 是 $S$
上的阿贝尔概形，则 $X_0$ 在唯一同构意义下确定。

利用上述唯一性结果，可以看出，$X$ 的基概形限制问题在 $S$ 上是局部的
（因而只须知道如何在各 $\operatorname{Spec}(\mathcal O_x)$ 上实施这种限制，
其中 $x\in S$）。同理，若 $\alpha:S'\longrightarrow S$ 是有限型光滑态射，
$K'$ 是 $S'$ 的有理函数环，并且 $X\otimes_KK'$ 形如
$X'_0\times_{S'}\operatorname{Spec}(K')$，则 $X'_0$ 配有相对于
$\alpha$ 的典范下降数据。结合定理3，得到：

\result{命题}{5.1} 设 $S$ 是诺特、正则、不可约的预概形，其有理函数域为
$K$；设 $K'$ 是 $K$ 的一个在 $S$ 上非分歧的有限扩张，$S'$ 是 $S$ 在
$K'$ 中的正规化（所以是 $S$ 的平展覆盖）；并设 $X$ 是 $K$ 上的阿贝尔
概形，使 $X\otimes_KK'$ 形如
$X'_0\times_{S'}\operatorname{Spec}(K')$，其中 $X'_0$ 是 $S'$ 上的
射影阿贝尔概形。则 $X$ 形如 $X_0\times_S\operatorname{Spec}(K)$，
其中 $X_0$ 是 $S$ 上的射影阿贝尔概形。

\result{注}{} 报告人不知道，能否把 $S'\longrightarrow S$ 是满射的平展
覆盖这一假设（它使定理3可用），换成它是满射的有限型光滑态射这一假设（即使
再假设它平展）；也不知道，若不假设 $X'_0$ 在 $S'$ 上射影，命题是否仍
成立（这一条件或许会自动满足）。

\subsection*{6. 局部平凡性与等平凡性的判据}
\addcontentsline{toc}{subsection}{6. 局部平凡性与等平凡性判据}

设 $S$ 是预概形，$G$ 是 $S$ 上的一个\fquote{群预概形}，$P$ 是
$S$ 上的预概形，\fquote{$G$ 在其上作用}（取右作用）。若由 $G$ 在 $P$
上的作用得到的熟知态射
\[
  G\times_SP\longrightarrow P\times_SP
\]
是同构，就称 $P$ 是 $G$ 作用下的形式主齐性空间。

\src{325}
以下假设 $G$ 在 $S$ 上平坦（因而在 $S$ 上忠实平坦），并把 $G$ 作用下
的\emph{主齐性丛}这一名称专用于这样的形式主齐性空间 $P$：它在 $S$ 上
忠实平坦且拟紧。显然，这也等价于能够找到基 $S$ 的一个忠实平坦拟紧扩张
$S'\longrightarrow S$，使 $G'=G\times_SS'$ 作用下的形式主齐性丛
\[
  P'=P\times_SS'
\]
平凡，即同构于 $G'$（也就是有截面）；特别地，可以取 $S'=P$。还要指出，
若 $S$ 局部诺特，则 $P$ 忠实平坦的假设等价于：对每个 $s\in S$，
\[
  \widehat P_s=P\times_S\operatorname{Spec}(\widehat{\mathcal O}_s)
\]
在 $\widehat{\mathcal O}_s$ 上忠实平坦；这里
$\widehat{\mathcal O}_s$ 表示局部环 $\mathcal O_s$ 的完备化。这个等价
来自 $\widehat{\mathcal O}_s$ 在 $\mathcal O_s$ 上忠实平坦这一事实。
此外，若 $P$ 在局部诺特的 $S$ 上有限型，则满足上述条件的点 $s$ 所成集合
是可构造的；所以，若 $S$ 是\fquote{Jacobson 预概形}（例如域上的有限型
概形，或更一般地，Jacobson 环的谱），只须对 $S$ 的闭点验证所述条件。
这把问题归结到基为完备局部环 $A$ 的谱的情形。当
$S=\operatorname{Spec}(A)$（$A$ 为诺特完备局部环），而 $P$ 在 $S$ 上
有限型时，$P/S$ 忠实平坦也等价于存在一个在 $S$ 上有限且平坦的 $S'$，
使 $P'$ 平凡；若再假设 $G$ 在 $S$ 上光滑，则可以令 $S'$ 在 $S$ 上
平展。因此，若 $A$ 的剩余域还代数闭（\fquote{几何情形}），则 $P$ 在
$A$ 上忠实平坦，当且仅当它平凡。所以，若 $S$ 是代数闭域上的代数预概形，
$G$ 是 $S$ 上有限型光滑群预概形，则在 $S$ 上忠实平坦这一条件等价于
Serre 的解析平凡性条件（SLF）（[6]，第1--12页）。

还可以在 $P$ 上引入具有\fquote{局部平凡性}性质的其他更强条件。特别地，
若对每个 $s\in S$，都存在 $s$ 的一个开邻域 $U$ 和有限忠实平坦态射
（相应地，一个满射的平展覆盖）$U'\longrightarrow U$，使
$P'=P\times_SU'$ 平凡，就称 $P$ \emph{等平凡}（相应地，\emph{严格
等平凡}）。（我们的术语不同于 Serre [6]，他把我们称作严格等平凡的情形
称为局部等平凡。）当 $G$ 在 $S$ 上光滑时，严格等平凡性尤其有用；在其他
情形下，这个概念反而并不合适。

若 $G$ 在 $S$ 上仿射，则由第2节可知，$G$ 作用下的每个主齐性丛 $P$ 都
仿射；因此，利用定理2，可以沿忠实平坦拟紧态射\fquote{下降}

\src{326}
这样的丛。特别取 $G=\operatorname{Gl}(n)_S$；它由如下条件定义：
$S$-预概形上的函子
\[
  S'\longmapsto\operatorname{Hom}_S(S',G)
\]
（取值于群范畴）可与 A 第4.e节的函子
\[
  \operatorname{Gl}(n)(S')
  =\operatorname{Gl}\bigl(n,H^0(S',\mathcal O_{S'})\bigr)
\]
等同。利用如下三个事实：（i）$G$ 作用下的每个主齐性丛（相应地，$S$ 上
秩为 $n$ 的每个局部自由层）经过 $S$ 的某个适当忠实平坦拟紧扩张后，都
同构于\fquote{平凡}对象 $G$（相应地，$\mathcal O_S^n$）；（ii）可以经
这类态射下降所考察类型的对象（$G$ 作用下的主齐性丛，相应地，秩为 $n$
的局部自由层）；（iii）$S'/S$ 上平凡丛的自同构群函子性同构于 $S'$ 上
秩为 $n$ 的平凡局部自由层的自同构群。于是\fquote{形式地}得到：在 $S$
上（或在某个 $S'/S$ 上）给定一个群为 $G$ 的主齐性丛，与在那里给定一个
秩为 $n$ 的局部自由层\fquote{是同一回事}。（更确切地说，存在纤维范畴
之间的等价。）特别得到：

\result{命题}{6.1} 群为 $\operatorname{Gl}(n)_S$ 的每个主齐性丛都局部
平凡。

由熟知的论证，对 $\operatorname{Sl}(n)_S$、$\operatorname{Sp}(n)_S$
以及这类群的乘积等结构群，也得到相同结论。还可推出：若 $F$ 是
$G=\operatorname{Gl}(n)_S$ 的闭子群，在 $S$ 上平坦，并且商 $G/F$ 存在；
若 $G$ 是 $G/F$ 上以 $F\times_S(G/F)$ 为群的等平凡（相应地，严格
等平凡）主齐性丛，则群为 $F$ 的每个主齐性丛都是等平凡的（相应地，严格
等平凡的）。这适用于此前所用的所有 $S$ 上的\fquote{线性群}，特别适用于
$G=S\times_k\Gamma$ 的情形，其中 $S$ 是域 $k$ 上的预概形，而 $\Gamma$
是 $k$ 上经典意义下的线性群，特别是光滑群。因此，对这类群，Serre 的一个
问题得到解决（出处同上）。

还应指出，对 $S$ 上已知的大多数光滑群（无论线性与否），并且至少对上述
形如 $S\times_k\Gamma$ 的群，可以证明每个等平凡主齐性丛都严格等平凡。
特别地，考虑到经 Cartier 型下降所得的主齐性丛（参见第3节，例2）显然
等平凡，这就解决了 Serre 的另一个问题（出处同上，第1--14页）。

\result{注}{} 这些问题中的一个本质困难（除去商概形 $G/F$ 是否存在的
问题）是：对于沿忠实平坦但非有限的态射给出的下降数据，缺少有效性判据。

\src{327}
\begin{center}
  \textbf{参考文献}
\end{center}

\begin{description}
  \item[{[1]}] Dieudonné (J.) 与 Grothendieck (Alexander)。---
    \emph{Éléments de géométrie algébrique}，即将发表于 Institut des
    Hautes Études Scientifiques 的 Publications mathématiques。
  \item[{[2]}] Grauert (H.) 与 Remmert (R.)。--- \emph{Komplexe Räume}，
    Math. Annalen，第136卷，1958年，第245--318页。
  \item[{[3]}] Grothendieck (Alexander)。--- \emph{Géométrie formelle et
    géométrie algébrique}。Séminaire Bourbaki，第11卷，1958/59年，第182号。
  \item[{[4]}] Murre (J. P.)。--- \emph{On a connectedness theorem for a
    birational transformation at a simple point}，Amer. J. Math.，第80卷，
    1958年，第3--15页。
  \item[{[5]}] Serre (Jean-Pierre)。--- \emph{Géométrie algébrique et
    géométrie analytique}，Ann. Inst. Fourier Grenoble，第6卷，1955--56年，
    第1--42页。
  \item[{[6]}] Serre (Jean-Pierre)。--- \emph{Espaces fibrés algébriques}，
    Séminaire Chevalley，第2卷，1958年：Anneaux de Chow et applications，
    第1号。
\end{description}

% END 190.tex

% BEGIN 195.tex
\unit{第195次报告——下降技术 II}
\sid{195}
\src{369}
\seminarzh{第12年度，1959/60，第195号\\（1960年2月）}
\paperhead{下降技术与代数几何中的存在性定理\\
II. 形式模理论中的存在性定理}%
  {亚历山大·格罗滕迪克 著}

\fsection{A}{可表函子与 pro-可表函子}

\subsection*{1. 可表函子}
\addcontentsline{toc}{subsection}{1. 可表函子}

设 $\mathcal C$ 是一个范畴。对每个 $X\in\mathcal C$，以 $h_X$ 表示从
$\mathcal C$ 到集合范畴 $(\operatorname{Ens})$ 的反变函子
\[
  h_X:\mathcal C\longrightarrow(\operatorname{Ens}),
\]
其定义公式为
\[
  h_X(Y)=\operatorname{Hom}(Y,X).
\]
若在 $\mathcal C$ 中给定一个态射 $X\longrightarrow X'$，便以显然方式得到
函子同态 $h_X\longrightarrow h_{X'}$：$h_X$ 关于 $X$ 协变。换言之，我们
定义了一个典范协变函子
\[
  h:\mathcal C\longrightarrow
  \operatorname{Hom}\bigl(\mathcal C^o,(\operatorname{Ens})\bigr)
\]
它从 $\mathcal C$ 取值于从 $\mathcal C$ 的对偶范畴 $\mathcal C^o$ 到集合
范畴的协变函子所成的范畴。回顾：

\result{命题}{1.1} 函子 $h$ 是\emph{全忠实}的；也就是说，对 $\mathcal C$
中任意一对对象 $X,X'$，自然映射
\[
  \operatorname{Hom}(X,X')\longrightarrow
  \operatorname{Hom}(h_X,h_{X'})
\]
是双射。

特别地，若函子
$F\in\operatorname{Hom}(\mathcal C^o,(\operatorname{Ens}))$ 同构于某个形如
$h_X$ 的函子，则 $X$ 在唯一同构意义下确定。这时称函子 $F$ 是
\emph{可表的}。上述命题意味着，典范函子 $h$ 给出 $\mathcal C$ 与
$\operatorname{Hom}(\mathcal C^o,(\operatorname{Ens}))$ 中由可表函子组成的
全子范畴之间的等价。这一事实是\fquote{泛问题的解}这一概念的基础：这类问题
总是在考察从 $\mathcal C$ 到 $(\operatorname{Ens})$ 的某个给定函子（如这里
一样为反变函子，或在对偶情形中为协变函子）是否可表。

\src{370}
此外，由范畴中积的概念本身的定义[1]，函子 $h:X\longmapsto h_X$ 在积存在时
与积交换（更一般地，在有限或无限逆极限存在时与之交换；特别包括纤维积、
\fquote{核}[2]的形成等等）：只要 $X\times X'$ 存在，就有函子同构
\[
  h_{X\times X'}\xrightarrow{\sim}h_X\times h_{X'}
\]
也就是说，有关于 $Y$ 具有函子性的双射
\[
  h_{X\times X'}(Y)\xrightarrow{\sim}h_X(Y)\times h_{X'}(Y).
\]
特别地，在 $\mathcal C$ 中给定一个态射
\[
  X\times X'\longrightarrow X''
\]
（也就是在 $\mathcal C$ 中给定 $X,X',X''$ 之间的一个\fquote{合成律}），
等价于给定一个态射
\[
  h_{X\times X'}=h_X\times h_{X'}\longrightarrow h_{X''},
\]
也等价于对每个 $Y\in\mathcal C$ 给定一个集合的合成律
\[
  h_X(Y)\times h_{X'}(Y)\longrightarrow h_{X''}(Y)
\]
并使得对 $\mathcal C$ 中每个态射 $Y\longrightarrow Y'$，集合映射系
\[
  h_{X^{(i)}}(Y)\longleftarrow h_{X^{(i)}}(Y')
  \qquad(\text{对 }i=0,1,2)
\]
是分别相对于 $Y$ 与 $Y'$ 的两个合成律之间的态射。由此可见，在
$\mathcal C$ 的对象 $X$ 上给定\fquote{$\mathcal C$-群}、
\fquote{$\mathcal C$-环}等结构，最便利的表述（无论理论上还是实践中）是：
对每个 $Y\in\mathcal C$，在集合 $h_X(Y)$ 上给定通常意义下的群律（相应地，
环律等等），并要求与态射 $Y\longrightarrow Y'$ 相应的映射
$h_X(Y)\longleftarrow h_X(Y')$ 是群同态（相应地，环同态等等）。例如，在任意
基预概形 $S$ 上定义各种经典群 $G_a$、$G_m$、$\operatorname{Gl}(n)$ 等并写出
这些群之间的经典关系；在由拟凝聚层 $\mathcal F$ 定义的 $S$ 上仿射概形
$V(\mathcal F)$ 上赋予\fquote{向量丛}结构；以及定义并研究相应的各种旗簇
（Grassmann 簇、射影丛）等等，采用这种做法最为直观便利。一般的思路就是借助
典范函子 $h$，直接把 $\mathcal C$ 的对象等同于某些特殊反变函子，也就是
可表函子，

\src{371}
它们从 $\mathcal C$ 取值于集合范畴。

通常的反转箭头手续——例如对仿射概形而言，从几何语言转到交换代数语言时
必须如此——引导我们把上述讨论对偶化；特别地，还要引入从
$\mathcal C\longrightarrow(\operatorname{Ens})$ 的可表协变函子，也就是形如
\[
  Y\longmapsto\operatorname{Hom}(X,Y)=h'_X(Y).
\]
的函子。

\subsection*{2. pro-可表函子与 pro-对象}
\addcontentsline{toc}{subsection}{2. pro-可表函子与 pro-对象}

设 $\widetilde X=(X_i)_{i\in I}$ 是 $\mathcal C$ 中对象的一个逆系统；与之相应
有协变函子
\[
  h'_{\widetilde X}=\varinjlim_i h'_{X_i},
\]
更明确地说，
\[
  h'_{\widetilde X}(Y)=\varinjlim_i h'_{X_i}(Y)
  =\varinjlim_i\operatorname{Hom}(X_i,Y),
\]
它从 $\mathcal C$ 取值于 $(\operatorname{Ens})$。若从 $\mathcal C$ 到
$(\operatorname{Ens})$ 的一个函子同构于这类函子，且 $I$ 是滤过的，就称它是
\emph{pro-可表的}。根据上一小节，这恰好就是同构于可表函子的滤过直极限的
那些函子。设 $\widetilde X'=(X'_j)_{j\in J}$ 是 $\mathcal C$ 中第二个滤过
逆系统（建立在第二个滤过预序指标集 $J$ 上）；容易验证存在典范双射（它推广
命题1.1）
\[
  \operatorname{Hom}(h'_{\widetilde X'},h'_{\widetilde X})
  =\varprojlim_j\varinjlim_i\operatorname{Hom}(X_i,X'_j).
\]
这就引出 $\mathcal C$ 的 pro-对象所成的范畴
$\operatorname{Pro}(\mathcal C)$：它的对象是 $\mathcal C$ 中对象的逆系统
（其指标集可为任意滤过预序集）。若
$\widetilde X=(X_i)_{i\in I}$ 与 $\widetilde X'=(X'_j)_{j\in J}$ 是两个这类
对象，则定义
\[
  \operatorname{ProHom}(\widetilde X,\widetilde X')
  =\varprojlim_j\varinjlim_i\operatorname{Hom}(X_i,X'_j),
\]
pro-同态的合成是显然的。由构造本身，
$\widetilde X\longmapsto h'_{\widetilde X}$ 可视为关于 $\widetilde X$ 的反变
函子；它给出 $\mathcal C$ 的 pro-对象范畴的对偶范畴与从 $\mathcal C$ 到
$(\operatorname{Ens})$ 的 pro-可表协变函子范畴之间的等价。当然，
$\mathcal C$ 的一个对象 $X$ 典范地定义一个 pro-对象，

\src{372}
仍记为 $X$，从而 $\mathcal C$ 等价于
$\operatorname{Pro}(\mathcal C)$ 的一个全子范畴。
若 $\widetilde X=(X_i)_{i\in I}$ 是 $\mathcal C$ 的任意 pro-对象，则（在上述
等同下）
\[
  \widetilde X=\varprojlim_i X_i,
\]
这里的逆极限是在 $\operatorname{Pro}(\mathcal C)$ 中取的（因为
$h'_{\widetilde X}=\varinjlim_i h'_{X_i}$）。须注意，即使 $X_i$ 的逆极限在
$\mathcal C$ 中存在，它通常也不同构于 $\operatorname{Pro}(\mathcal C)$ 中的
逆极限 $\widetilde X$；当 $\mathcal C$ 是集合范畴时，这一点已经十分明显。
还要指出，按定义，$\varprojlim_{\mathcal C}X_i=L$ 的含义是：关于
$Y\in\mathcal C$、取值于 $(\operatorname{Ens})$ 的函子
\[
  \varprojlim_i\operatorname{Hom}_{\mathcal C}(Y,X_i)
  =\operatorname{Hom}_{\operatorname{Pro}(\mathcal C)}(Y,\widetilde X)
\]
可由 $L$ 表示，也就是同构于 $\operatorname{Hom}_{\mathcal C}(Y,L)$。因此，
$\varprojlim_{\mathcal C}X_i$ 已经由 pro-对象 $\widetilde X$ 定义；只要它有
定义，它就精确地、函子性地依赖于 pro-对象 $\widetilde X$，所以不妨把它记为
$\varprojlim_{\mathcal C}(\widetilde X)$。若 $\mathcal C$ 中的逆极限总是存在，
则 $\varprojlim_{\mathcal C}(\widetilde X)$ 是从
$\operatorname{Pro}(\mathcal C)$ 到 $\mathcal C$ 的函子，并且有一个函子间的
典范同态
$\varprojlim_{\mathcal C}(\widetilde X)\longrightarrow\widetilde X$。

任取一个函子（为确定起见，设为协变函子）
\[
  F:\mathcal C\longrightarrow\mathcal C'
\]
它都以显然方式延拓为函子
\[
  \operatorname{Pro}(F):\operatorname{Pro}(\mathcal C)
  \longrightarrow\operatorname{Pro}(\mathcal C'),
\]
由此可见，若 $\mathcal C'$ 中的逆极限总是存在，则 $F$ 还典范地定义复合函子
\[
  \overline F:\operatorname{Pro}(\mathcal C)\longrightarrow\mathcal C',
  \qquad
  \overline F=\varprojlim_{\mathcal C'}\circ\operatorname{Pro}(F),
\]
它把 $\widetilde X=(X_i)_{i\in I}$ 变为
$\varprojlim_{\mathcal C'}F(X_i)$。

若 pro-对象 $\widetilde X$ 同构于某个 pro-对象 $(X_i)_{i\in I}$，且其转移
态射 $X_i\longrightarrow X_j$ 都是满态射，则称 $\widetilde X$ 是一个
\emph{严格 pro-对象}；由这类对象定义的函子称为\emph{严格 pro-可表的}。
这时还可以要求 $I$ 是滤过偏序集，并且每个满态射
$X_i\longrightarrow X'$ 都等价于一个满态射

\src{373}
$X_i\longrightarrow X_j$，其中 $j\in I$ 适当（并由此条件唯一确定）。在这些
条件下，逆系统 $(X_i)_{i\in I}$ 在唯一同构意义下确定（这里同构取通常的
逆系统同构之意）。由此可知，在 $\operatorname{Pro}(\mathcal C)$ 中，严格
pro-对象的逆系统 $\widetilde X^{(\alpha)}$ 的逆极限总是存在；并且沿用上面对
$F$、$\overline F$ 的记号，有
\[
  \overline F\left(\varprojlim_\alpha\widetilde X^{(\alpha)}\right)
  =\varprojlim_{\alpha,\mathcal C'}
  \overline F\left(\widetilde X^{(\alpha)}\right).
\]
特别地，若 $\mathcal C$ 的每个 pro-对象都是严格的（参见下一小节），则延拓
函子 $\overline F$ 与逆极限交换。

\subsection*{3. pro-可表函子的刻画}
\addcontentsline{toc}{subsection}{3. pro-可表函子的刻画}

设 $\mathcal C$、$\mathcal C'$ 是两个具有有限逆极限的范畴（即相对于有限
预序集而取的逆极限，指标集不必滤过）；等价地说，它们具有有限积和有限
纤维积（这特别蕴涵存在一个\fquote{右单位对象} $e$，使得对每个 $X$，
$\operatorname{Hom}(X,e)$ 只有一个元素）。设 $F$ 是从 $\mathcal C$ 到
$\mathcal C'$ 的协变函子。下列条件等价：

\begin{enumerate}
\item[(i)] $F$ 与有限逆极限交换；
\item[(ii)] $F$ 与有限积和有限纤维积交换；
\item[(iii)] $F$ 与有限积交换，并且对 $\mathcal C$ 中每个正合图
\[
  X\longrightarrow X'\rightrightarrows X''
\]
（[3]，A，定义2.1），经 $F$ 变换所得的图
\[
  F(X)\longrightarrow F(X')\rightrightarrows F(X'')
\]
是正合的。
\end{enumerate}

这时称 $F$ 是\emph{左正合}的。

以下假设 $\mathcal C$ 中存在有限逆极限。由定义立即可见，可表函子左正合；
再取极限可知，pro-可表函子也左正合。

为得到逆命题，设
\[
  F:\mathcal C\longrightarrow(\operatorname{Ens})
\]
是一个协变函子，并设 $X\in\mathcal C$、$\xi\in F(X)$。若对每个
$X'\in\mathcal C$、每个 $\xi'\in F(X')$ 以及每个严格单态射（[3]，A，2）
$u:X'\longrightarrow X$，只要 $\xi=F(u)(\xi')$ 就必有 $u$ 为同构，则称
$\xi$（或二元组 \src{374}$(X,\xi)$）是\emph{极小的}。另一方面，若存在态射
$v:X\longrightarrow X''$ 使 $\xi''=F(v)(\xi)$，就称二元组 $(X,\xi)$
\emph{支配} $(X'',\xi'')$（其中 $\xi\in F(X)$、$\xi''\in F(X'')$）。
若 $\xi$ 极小且 $F$ 左正合，则态射 $v$ 是唯一的；若 $\xi''$ 极小，则
$v$ 是满射。由此容易推出下列命题：

\result{命题}{3.1} $F$ 严格 pro-可表，当且仅当它满足下列两个条件：
\begin{enumerate}
\item[(i)] $F$ 左正合；
\item[(ii)] 每个二元组 $(X,\xi)$（其中 $\xi\in F(X)$）都由某个极小二元组
支配。
\end{enumerate}

若 $\mathcal C$ 的每个对象都是阿廷对象，则后一条件为空条件：只须在 $X$
的所有子对象 $X'$ 中，选取一个极小者，使得存在 $\xi'\in F(X')$ 且
$\xi$ 是 $\xi'$ 的像。因此有：

\result{推论}{} 设 $\mathcal C$ 是一个对象均为阿廷对象且具有有限逆极限的
范畴。那么，从 $\mathcal C$ 到 $(\operatorname{Ens})$ 的 pro-可表函子恰好
是左正合函子，并且它们实际上都是严格 pro-可表的。

最后这一事实也就是说，\emph{$\mathcal C$ 的每个 pro-对象都是严格的}。

\subsection*{4. 例：Galois 型群与 pro-代数群}
\addcontentsline{toc}{subsection}{4. Galois 型群与 pro-代数群}

若 $\mathcal C$ 是通常有限群的范畴，则 $\operatorname{Pro}(\mathcal C)$
等价于紧全不连通拓扑群的范畴。这类群以及把通常有限群换成给定基预概形上的
有限群概形（例如域 $k$ 上的有限代数群）所得的推广，将充当预概形的绝对与
相对基本群、同伦群和同调群。在所有这些例子中，命题3.1的推论都适用，而且
$\pi_1$ 确实应当借助相应函子来定义[2]。从域上（或更一般地，从诺特预概形上）
的代数群或拟代数群范畴出发时也是如此：由此得到 Serre [4] 的
\fquote{pro-代数群}。

\subsection*{5. 例：“形式簇”}
\addcontentsline{toc}{subsection}{5. 例：“形式簇”}

设 $\Lambda$ 是诺特环，$\mathcal C$ 是所有 $\Lambda$-代数 $A$ 所成的
范畴，其中 $A$ 作为 $\Lambda$-模是有限型阿廷模（简言之，

\src{375}
是阿廷 $\Lambda$-代数）。此时满足命题3.1之推论的条件。这里的范畴
$\operatorname{Pro}(\mathcal C)$ 等价于下列 $\Lambda$-拓扑代数 $O$ 所成的
范畴：$O$ 同构于拓扑逆极限
\[
  O=\varprojlim_i O_i
\]
，其中 $O_i\in\mathcal C$；也就是说，$O$ 的拓扑是线性的、分离的且完备的，
并且对 $O$ 的每个开理想 $J_i$，商 $O/J_i$ 都是 $\Lambda$ 上的阿廷代数。
与这种代数相应的函子 $\mathcal C\longrightarrow(\operatorname{Ens})$
正是
\begin{align*}
  F(A)=h'_O(A)
  &=\{\text{拓扑 $\Lambda$-代数的连续同态 }
       O\longrightarrow A\}\\
  &=\varinjlim_i
    \operatorname{Hom}_{\Lambda\text{-代数}}(O_i,A).
\end{align*}

还要指出，范畴 $\mathcal C$ 本质上是若干相应范畴的积，这些范畴对应于
$\Lambda$ 的各个极大理想 $\mathfrak m$ 所给出的
$\Lambda_{\mathfrak m}$ 的完备化局部环；因此，如有需要，可以只考察
$\Lambda$ 是这类完备局部环的情形。无论如何，$O$ 典范地分解成若干局部
分量的拓扑积，这些分量对应于由 $O$ 定义的形式概形[2]的\fquote{点}。
这类点由某个 $F(K)$ 中的对象 $\xi$ 定义，其中 $K\in\mathcal C$ 是一个域
（例如所考察的局部分量的剩余域）；两个二元组 $(\xi,K)$ 与 $(\xi',K')$
定义同一点，当且仅当它们由同一个 $(\xi'',K'')$ 支配，或者等价地，它们
支配同一个 $(\xi''',K''')$。（若各 $\Lambda/\mathfrak m$ 都是代数闭的，
只须取各 $F(\Lambda/\mathfrak m)$ 的不交并。）

关键在于给出条件，保证 $O$ 中与 $\xi\in F(K)$ 相应的局部分量 $O_\xi$
是诺特环。当 $\Lambda$ 是完备局部环（这里回顾，它是诺特的）时，这等价于
$O_\xi$ 同构于 $\Lambda$ 上的形式幂级数环
$\Lambda[[t_1,\ldots,t_n]]$ 的一个商环。为给出这种判据，对每个环 $A$
引入 $A$ 的\fquote{对偶数}所成的 $A$-代数 $I_A$：
\[
  I_A=A[t]/t^2A[t].
\]
设 $\varepsilon:I_A\longrightarrow A$ 是增广同态；若
$A\in\mathcal C$，它定义映射
\[
  F(\varepsilon):F(I_A)\longrightarrow F(A),
\]

\src{376}
再利用 $F$ 左正合这一事实，就能在 $F(I_A)$ 的子集
\[
  F(I_A,\xi)=F(\varepsilon)^{-1}(\xi)
\]
上内在地定义 $A$-模结构；这个子集由所有\fquote{约化为 $\xi$}的
$\xi'\in F(I_A)$ 组成。利用 $F$ 关于 $O$ 的明确形式可知，这个 $K$-模
可等同于
\[
  \operatorname{Hom}_{\Lambda}(\mathfrak m_\xi/\mathfrak m_\xi^2,A),
\]
其中 $\mathfrak m_\xi$ 是同态 $\xi:O\longrightarrow A$ 的核；也就是说，
若 $A$ 是域，它就是 $O$ 的局部分量 $O_\xi$ 的极大理想。由此立即得到：

\result{命题}{5.1} 设 $\xi\in F(K)$，其中 $K\in\mathcal C$ 是一个域。
$O$ 中相应的局部分量 $O_\xi$ 是诺特环，当且仅当 $F(I_K)$ 中约化为
$\xi$ 的元素所成的集合 $F(I_K,\xi)$ 是 $K$ 上有限维向量空间。在这些
条件下，有典范同构
\[
  F(I_K,\xi)=
  \operatorname{Hom}\!\left(
    \frac{\mathfrak m_\xi/\mathfrak m_\xi^2}
         {\operatorname{im}(\mathfrak n_\xi/\mathfrak n_\xi^2)},K
  \right),
\]
其中 $\mathfrak n_\xi$ 是 $\Lambda$ 的极大理想，也就是同态
$\Lambda\longrightarrow K$ 的核。
\footnote{印刷文本在这里给出的公式只有当 $\Lambda$ 是域时才正确。
1962年的勘误规定，在一般情形中，须将
$\mathfrak m_\xi/\mathfrak m_\xi^2$ 除以
$\mathfrak n_\xi/\mathfrak n_\xi^2$ 的像。}
特别地，$K$-向量空间 $F(I_K,\xi)$ 的维数等于上述修正后的向量空间在域
$O_\xi/\mathfrak m_\xi=K(\xi)$ 上的维数。

假设 $O_\xi$ 是诺特环；为简化记号，再假设 $\Lambda$ 是完备局部环且
$O=O_\xi$。若 $O$ 是 $\Lambda[t_1,\ldots,t_n]$ 在该环某个诱导出
$\Lambda$ 极大理想的极大理想处局部化之后的完备化代数上的有限平展代数，
则称 $O$\emph{在 $\Lambda$ 上光滑}。若 $O$ 在 $\Lambda$ 上的剩余扩张
平凡（例如 $\Lambda$ 的剩余域代数闭），这等价于说 $O$ 本身同构于这类
形式幂级数代数。
\footnote{1962年的勘误指出，这个\fquote{在 $\Lambda$ 上光滑}的定义只有
在剩余扩张 $k'/k$ 可分时才正确；一般情形参见 SGA III，1.1。}
最后，若不再假设 $O$ 为诺特环，只要 $O$ 同构于若干商
$\Lambda$-代数的拓扑逆极限，而这些商代数是诺特的并且按上述意义在
$\Lambda$ 上光滑，仍称 $O$ 在 $\Lambda$ 上光滑。这立即推广到不再假设
$\Lambda$、$O$ 为局部环的情形。于是有：

\result{命题}{5.2} $O$ 在 $\Lambda$ 上光滑，当且仅当相应函子 $F$
把满态射变为满态射。

这就是说，对 $\mathcal C$ 中每个满同态 $A\longrightarrow A'$，
$F(A)\longrightarrow F(A')$ 都是满射。当然，只须在 $A$ 为局部环时验证
这一条件；并且逐步进行时，只须考虑 $A\longrightarrow A'$ 的核理想被
$A$ 的极大理想零化的情形。在实践中，这归结为验证某个与无穷小不变量有关的
障碍

\src{377}
为零；这里的无穷小不变量属于产生函子 $F$ 的情形。这是一个上同调性质的问题。

最后，在上述背景下简要谈谈\emph{定义环}。仍设 $F$ 是从 $\mathcal C$ 到
$(\operatorname{Ens})$ 的函子，并由拓扑 $\Lambda$-代数 $O$ pro-表示。
那么对每个 $A\in\mathcal C$ 和每个 $\xi\in F(A)$，都存在 $A$ 的一个最小
子环 $A'$，使得 $\xi$ 是 $F(A')$ 中某个元素 $\xi'$ 的像（此时 $\xi'$
唯一确定）。事实上，只须以从 $O$ 到 $A$ 的同态来解释 $\xi$，再取
$A'$ 为 $O$ 在该同态下的像。此时称 $A'$ 是对象 $\xi\in F(A)$ 的定义环。
若 $u:A\longrightarrow B$ 是代数同态且 $\eta=F(u)(\xi)$，则 $\eta$ 的
定义环就是 $\xi$ 的定义环在 $u$ 下的像。从一个给定函子
$F:\mathcal C\longrightarrow(\operatorname{Ens})$ 出发时，定义环的存在性
及其上述性质大致等价于 $F$ 可 pro-表示；因此，这些性质通常远非平凡。

\fsection{B}{两个存在性定理}

沿用 A 第5节的记号，并从一个协变函子
\[
  F:\mathcal C\longrightarrow(\operatorname{Ens}) ;
\]
出发；我们要寻找便于使用的判据，使 $F$ 可 pro-表示，也就是能像上面那样用
一个 $\Lambda$-代数 $O$ 表出。根据 A 中命题3.1的推论，这当且仅当 $F$
左正合。就下降技术目前的状况而言（参见[3]第9页提出的问题），在最重要的
情形中无法直接以这个形式验证该判据，因而需要表面看来要求较弱的判据。

\result{定理}{1} 函子 $F$ 可 pro-表示，当且仅当它满足下列两个条件：
\begin{enumerate}
\item[(i)] $F$ 与有限积交换；
\item[(ii)] 对每个代数 $A\in\mathcal C$ 以及 $\mathcal C$ 中每个同态
$A\longrightarrow A'$，只要图
\[
  A\longrightarrow A'\rightrightarrows A'\otimes_A A'
\]
是正合的（[3]，A，定义1.2），其变换后的图
\[
  F(A)\longrightarrow F(A')\rightrightarrows F(A'\otimes_A A')
\]

\src{378}
就是正合的。
\end{enumerate}
此外，只须在 $A$ 为局部环并且又属于下列两种情形之一时验证（ii）：
\begin{enumerate}
\item[(a)] $A'$ 是 $A$ 上的自由模；
\item[(b)] 商模 $A'/A$ 是长度为1的 $A$-模。
\end{enumerate}

该定理的证明颇为精细，无法在这里勾勒。我们只指出，它主要依赖于对阿廷代数
之谱中的等价关系（取范畴意义）的研究；这种研究仍提出若干问题，而其解决看来
是理论进一步发展所不可缺少的。

在应用中，条件（i）的验证总是平凡的。条件（ii）的验证分成两部分：
当 $A'$ 是自由 $A$-模时，这属于沿平坦态射的下降技术（[3]，定理1、2、3），
并无困难。为处理情形（b），采用下列结果：

\result{定理}{2} 设 $A$ 是极大理想为 $\mathfrak m$ 的阿廷局部环，
$A'$ 是包含 $A$ 的 $A$-代数，满足 $\mathfrak mA'\subset A$，并且
\[
  A\longrightarrow A'\rightrightarrows A'\otimes_A A'
\]
正合（特别地，当 $A'/A$ 是长度为1的 $A$-模时便是如此）。设
$\mathcal F$ 是变动预概形上的平坦拟凝聚层所成的纤维范畴（[3]，A，
定义1.1）。那么态射
$\operatorname{Spec}(A')\longrightarrow\operatorname{Spec}(A)$ 是严格
$\mathcal F$-下降态射（[3]，A，定义1.7）。

换言之，给定一个平坦 $A$-模 $M$，完全等价于给定一个平坦
$A'$-模 $M'$，以及一个从 $M'\otimes_AA'$ 到
$A'\otimes_AM'$ 的 $A'\otimes_AA'$-同构；该同构满足下降数据通常的传递性
条件（见前引处）。

定理2的证明首先要证
\[
  H^i(A'/A,G_a)=0\qquad\text{对 }i\geq 1
\]
（[3]，A，4(e)）；假设 $\mathfrak mA'\subset A$ 使问题很容易化归到
$A$ 是域 $k$ 的情形（即 $A/\mathfrak m$）。然后应用[3]第16页指出的各个
等价。

\fsection{C}{若干特殊情形的应用}

\subsection*{1. 由预概形表示的函子：一般评注}
\addcontentsline{toc}{subsection}{1. 由预概形表示的函子：一般评注}

设 $S$ 是局部诺特预概形。称 $S$ 上的预概形 $X$
\src{379}在 $S$ 上局部有限型，如果对每个投影到 $y\in S$ 的
$x\in X$，\footnote{印刷文本此处作\fquote{$y\in Y$}，并在本段末尾再次
使用 $Y$。由于所考察的态射是 $X\longrightarrow S$，两处均按 $S$ 读取。}
存在 $y$ 的一个仿射邻域（其环为 $A$）以及 $x$ 在该邻域上的一个仿射邻域
（其环为 $B$），使得 $B$ 是有限型 $A$-代数。经典泛问题的解提供了许多重要
例子：有些预概形在 $S$ 上局部有限型，却不在 $S$ 上有限型。因此，必须能够
把曲线的 Picard 概形看作无穷多个连通分支之并；须注意不要把这些分支与单位元
的连通分支，也就是\fquote{Picard 簇}混淆。有时，把自己置于 $S$ 上局部
有限型预概形所成的范畴 $\mathcal C$ 中，并在其中考察反变函子 $F$ 的可表性，
会很方便。这一系列报告的主要目的，是发展一种一般技术，用来识别这类函子
$F$ 是否可表，并借助 $F$ 的性质研究相应 $S$-预概形 $X$ 的性质。顺便指出，
在这一研究中会遇到一些并不病态的、在 $S$ 上非分离的预概形例子，特别是作为
优良 $S$-概形之\fquote{Picard 预概形}的例子。因此，不应把不是概形的预概形
逐出代数几何。

设 $X$ 是 $S$ 上局部有限型的预概形，并令
\[
  F(Y)=\operatorname{Hom}_S(Y,X)
\]
为相应的反变函子。可以考察 $F$ 在 $\mathcal C$ 的子范畴
$\mathcal C_0$ 上的限制 $F_0$；$\mathcal C_0$ 由 $S$ 上阿廷且有限的
预概形 $Y$ 组成。因此，当 $S=\operatorname{Spec}(\Lambda)$ 时，
$\mathcal C_0$ 是 B 中所考察的阿廷 $\Lambda$-代数范畴的对偶范畴。
若 $Y=\operatorname{Spec}(A)$，其中 $A$ 是阿廷局部环，则 $Y$ 只有一个点
$y$，它位于 $S$ 的某个闭点 $s$ 之上；从 $Y$ 到 $X$ 的 $S$-态射（即
$F(Y)$ 的一个元素）由下列数据确定：位于 $s$ 上方的一点 $x\in X$，以及
从 $\mathcal O_{x,X}$ 到 $A$ 的一个 $\mathcal O_{s,S}$-代数同态。若存在这种
同态，则 $x$ 必为 $X$ 的闭点（其剩余域是 $s$ 的剩余域的代数扩张）。
由此可见，$F$ 在\fquote{阿廷 $\Lambda$-代数}上的限制 $F_0$ 是 pro-可表的，
并由如下拓扑 $\Lambda$-代数表示：其局部分量是 $X$ 在各闭点 $x$ 处的局部环
之完备化 $\widehat{\mathcal O}_{x,X}$，这里 $x$ 位于 $S$ 的闭点之上。
这说明，仅知道 $F_0$ 就能得到关于 $X$ 结构的精确信息，即它在上述各点处
局部环之完备化的结构。还应指出，即使 $S$

\src{380}
是一个代数闭域的谱，也只有系统地考察这样的\fquote{簇} $Y$，其中
$\mathcal O_Y$ 可以含有幂零元，特别是使用阿廷局部环的谱，才能得到经典泛问题
的\fquote{正确表述}，并把握其\fquote{无穷小}方面。

给定一个预先提出的函子 $F$，要判断它是否可表时，借助定理1和定理2研究函子
$F_0$，将给出几乎完备的判据。可能出现的第一种情形是，正如经常发生的那样
（例如只需检验集合 $F(I_K,\xi)$ 的性质及其函子性行为，参见 A），便发现
$F_0$ 已经不是 pro-可表的；这解释了迄今那些试图以多少自然的方式定义秩
$>1$ 向量丛之分类模空间的尝试为何失败。第二种情形是，可以验证 $F_0$ 确实
可表，但向量空间 $F(I_K,\xi)$ 不是有限维的，此时只能满足于\fquote{形式}
解。最后一种情形是，$F_0$ 确实由完备诺特局部环的一个乘积表示；这为 $F$
本身可表提供了极强的理由，再结合我们打算随后发展的、性质类似但更为整体的
条件，无疑足以推出它确实可表。最后，还会遇到一些有趣的几何问题（见下文
第4、5节），其中只有函子 $F_0$ 可用（它并不来自任何\fquote{整体}函子
$F$）；若能为它配上一个\fquote{形式模概形}，就已经值得庆幸。

作为这些一般评注的结尾，我们说明概形理论如何解释一些表面上的反常现象，
例如 Igusa 曲面 $V$：它的\fquote{Picard 簇} $P$ 缩为一点，但仍有
\[
  H^1(V,\mathcal O_V)\ne 0 ;
\]
在此情形，$P$ 是一个\fquote{纯无穷小}群，并未约化为单位群；也就是说，
它由基域 $k$ 上有限秩的一个局部代数 $O$ 定义，并带有对应于 $P$ 乘法结构
的对角映射。若 $\mathfrak m$ 是 $O$ 的极大理想，则
$\mathfrak m/\mathfrak m^2$ 的对偶与 $H^1(V,\mathcal O_V)$ 典范同构
（参见下文第3节）。只有当 Picard 群是经典意义下的代数群，也就是在基域
$k$ 上光滑时，$H^1(V,\mathcal O_V)$ 的维数（它始终等于
$\mathfrak m/\mathfrak m^2$ 的维数）才等于 Picard 群的维数。

\subsection*{2. 概形 $\underline{\operatorname{Hom}}_S(X,Y)$、
$\Pi_{X/S}Z$、$\underline{\operatorname{Aut}}_S(X)$ 等}
\addcontentsline{toc}{subsection}{2. Hom、Pi 与 Aut 概形}

设 $X,Y$ 是 $S$ 上的两个预概形；对每个 $S$ 上的预概形 $T$，

\src{381}
令 $X_T=X\times_ST$、$Y_T=Y\times_ST$，并把集合
\[
  F(T)=\operatorname{Hom}_T(X_T,Y_T)
      =\operatorname{Hom}_S(X_T,Y)
      =\operatorname{Hom}_S(X\times_ST,Y)
\]
看作关于 $T$ 的反变函子。若它可表，就用
$\underline{\operatorname{Hom}}_S(X,Y)$ 表示代表它的 $S$-预概形；因此有
函子同构
\[
  \operatorname{Hom}_S
  \bigl(T,\underline{\operatorname{Hom}}_S(X,Y)\bigr)
  \xrightarrow{\sim}
  \operatorname{Hom}_S(T\times_SX,Y).
\]
这一泛问题有许多可以化归于它的变体：$S$-预概形 $X$ 的 $S$-自同构预概形
（它将是一个群预概形）；从一个 $S$-群预概形到另一个群预概形的
$S$-同态预概形（若第二个群概形可交换，它将是一个可交换群预概形）；等等。
还可以推广 $\underline{\operatorname{Hom}}_S(X,Y)$ 的定义：考虑位于
$S$ 上预概形 $X$ 之上的一个预概形 $Z$，以及函子
\[
  F(T)=\operatorname{Hom}_{X_T}(X_T,Z_T)
\]
（即 $Z_T$ 在 $X_T$ 上的\fquote{截面}所成的集合）。若此函子可表，代表它的
$S$-预概形记为 $\Pi_{X/S}Z$；依定义便有函子同构
\[
  \operatorname{Hom}_S(T,\Pi_{X/S}Z)
  =\operatorname{Hom}_{X_T}(X_T,Z_T).
\]
取 $Z=Y\times_SX$，便重新得到
$\underline{\operatorname{Hom}}_S(X,Y)$。由这些定义可为新引入的预概形
导出一套虽平凡但有用的公式；这里不再列出，因为它们在任何存在积和纤维积的
范畴中都成立。更实质的问题是
$\underline{\operatorname{Hom}}_S(X,Y)$ 这一类概形是否存在。首先可以看出，
固定 $X$ 后，$\underline{\operatorname{Hom}}_S(X,Y)$ 要对每个 $S$ 上的
$Y$ 都存在（或者 $\Pi_{X/S}Z$ 要对每个 $X$ 上的 $Z$ 都存在），$X$ 必须
在 $S$ 上平坦。其次可以确信，若希望对相当一般的 $Y$ 都有解，还应要求 $X$
在 $S$ 上固有。另一方面，这些条件似乎足以保证
$\underline{\operatorname{Hom}}_S(X,Y)$ 和 $\Pi_{X/S}Z$ 存在，只须在必要
时分别对 $Y/S$ 或 $Z/X$ 加上某种\fquote{拟射影}假设；至少在相当多的情形
可以验证这一点（例如 $Y$ 在 $S$ 上仿射，或以直接的初等构造处理 $X$ 在
$S$ 上有限的情形）。定理1和定理2给出如下结论：

\src{382}
\result{命题}{2.1} 设 $\Lambda$ 是诺特环，$X$ 和 $Y$ 是
$\Lambda$ 上任意两个预概形。在阿廷 $\Lambda$-代数范畴
$\mathcal C_0$ 上考虑函子
\[
  F(A)=\operatorname{Hom}_A(X_A,Y_A)
\]
若 $X$ 在 $\Lambda$ 上平坦，则此函子 pro-可表。

此外，对每个 $A\in\mathcal C_0$ 和 $\xi\in F(A)$，有典范同构
\[
  F(I_A,\xi)=H^1\!\left(
    X_A,
    \underline{\operatorname{Hom}}_{\mathcal O_{X_A}}
    \bigl(\xi^*(\Omega^1_{Y_A/A}),\mathcal O_{X_A}\bigr)
  \right),
\]
其中 $\Omega^1_{Y_A/A}$ 是 $Y_A$ 相对于 $A$ 的 Kähler $1$-微分层。令
$A$ 为域，并使用 A.5.1 与 [2] 的有限性定理，便得到如下推论：

\result{推论}{} 设 $X$ 在 $S$ 上平坦且固有，$Y$ 在 $S$ 上局部有限型。
则 $F$ pro-可表，并且相应拓扑 $\Lambda$-代数的各个局部分量都是诺特环。

\result{评注}{} 本节所考察的问题以及许多其他问题，在经典代数几何的框架中，
通常借助射影空间内循环的\fquote{Chow 坐标}来研究；这种坐标使得可以把这些
循环看成适当射影簇的点。从概形观点看，这一方法以及一般的 Chow 坐标用法似乎
有不可补救的缺陷，因为它抹去了参数簇中的幂零元，尤其不适于令人满意地研究
循环的无穷小变化（更不用说它依赖于射影空间而不具有内蕴性）。遗憾的是，许多
代数几何学家研究簇族或循环族时只使用 Chow 坐标的语言；尽管这种方法确有
（大概只是暂时的）技术价值，它似乎严重妨碍了这些概念的澄清。若要在概形理论
中得到 Chow 簇的对应物，就会引向如下泛问题：设 $X$ 是 $S$ 上的预概形；对
每个 $S$ 上的预概形 $T$，令 $F(T)$ 为
$X_T=X\times_ST$ 中在 $T$ 上平坦的闭子预概形所成的集合。试图用一个
$S$ 上的预概形表示这个关于 $T$ 的函子。更一般地，可以在 $X$ 上给定一个
拟凝聚层 $G$，并令 $F(T)$ 为

\src{383}
$G_T$ 的所有商层中在 $T$ 上平坦者所成的集合。当 $X$ 在局部诺特预概形
$S$ 上固有，且
$G$ 还是凝聚层时，这一问题似乎有解，并可由一个在 $S$ 上局部有限型的概形
$C$ 表示。无论如何，只假设 $S$ 局部诺特，$F$ 在
\fquote{阿廷 $S$-代数}上的限制就是 pro-可表的；若再假设 $X$ 在 $S$
上固有且 $G$ 凝聚，则相应拓扑环 $O$ 的各个局部分量都是诺特环。当然，即使
已经证明 $X$ 在 $S$ 上的\fquote{Chow 概形}存在，仍须把它分解成互不相交的
开集 $C_i$（它们对应于固定所变动循环的次数、维数等连续不变量），使它们
在 $S$ 上有限型，而不只是局部有限型；还须阐明这一概形与经典 Chow
簇之间的关系，并确定何时某个 $C_i$ 在 $S$ 上甚至是射影的，或者至少是
拟射影的。\footnote{1962 年勘误说明，射影情形已由 Hilbert 概形（TDTE IV）
完全解决；而 Nagata 与 Hironaka 的例子表明，在没有射影性假设时，这些函子
可能不可表，哪怕只考察一个三维非奇异完备簇的零维子簇也是如此；这一现象与
对称平方可能不存在有关。}

\subsection*{3. Picard 概形}
\addcontentsline{toc}{subsection}{3. Picard 概形}

设 $f:X\longrightarrow S$ 是一个 $S$-预概形，\footnote{关于更完整的研究，
1962 年勘误转引 TDTE V。}考虑 $X$ 的结构层中可逆元所成的乘法层
$\mathcal O_X^*$，以及群
\[
  P(X/S)=H^0\bigl(S,R^1f_*(\mathcal O_X^*)\bigr),
\]
称为 $X/S$ 的相对 Picard 群。这个群的一个元素可由如下数据给出：$S$ 的一个
开覆盖 $(U_i)$，以及每个 $f^{-1}(U_i)$ 上的一个可逆层 $L_i$；并且对每一对
$(i,j)$，限制 $L_i|f^{-1}(U_i\cap U_j)$ 与
$L_j|f^{-1}(U_i\cap U_j)$ 至少在 $U_i\cap U_j$ 上局部同构（即这两个层在
[3]，B，4 的意义下\fquote{等价}）。若 $X/S$ 有截面，则 $P(X/S)$ 就是
$X/S$ 上可逆层模\fquote{等价}（同上引文）所得的类的集合。现在对每个
$S$ 上的 $T$，令
\[
  F(T)=P(X_T/T) ;
\]
便得到一个关于 $T$ 的反变函子，可以称为 $X/S$ 的 Picard 函子。若此函子
可表，代表它的 $S$-预概形称为 $X/S$ 的 Picard 预概形，记为
$\underline P(X/S)$。于是有函子同构
\[
  \operatorname{Hom}_S\bigl(T,\underline P(X/S)\bigr)
  \xrightarrow{\sim}P(X_T/T).
\]
Picard 预概形的形成与基扩张相容。特别地，$X$ 在 $S$ 上各纤维的 Picard
预概形（它们是各点 $s\in S$ 的剩余域 $K(s)$ 上的预概形），正是
$\underline P(X/S)$ 的各纤维。当然，由于 $P(X_T/T)=F(T)$ 是可交换群，
Picard 预概形都是群预概形。还应指出，Rosenlicht 的广义 Jacobi 簇正是可以
带有奇点的完备曲线之 Picard 概形中单位元的连通分支；一旦存在性得到证明，
这应使它们的大多数性质变得显然。

\result{评注}{} 这里采用的定义只有在 $S$ 的每个点
\footnote{印刷文本作\fquote{$Y$}；此处没有定义任何基 $Y$，而这一条件针对
的是基 $S$。}都有一个开邻域 $U$，使得 $X$ 在其上有截面时才是合理的。
在一般情形，必须稍微修改 Picard 函子的定义，才能仍然得到存在性定理。

这里 Picard 预概形存在的合理条件如下：$X$ 在 $S$ 上固有且平坦，
$f_*(\mathcal O_X)=\mathcal O_S$，并且 $X$ 在 $S$ 上局部有截面。最后这个
条件在应用下降技术时自然出现：给 $X$ 上的可逆层 $L$ 赋予一个沿截面 $s$
的标记截面，便可消去它的自同构（[3]，B，4）。特别地，得到：

\src{384}
\result{命题}{3.1} 设 $X$ 在
$S=\operatorname{Spec}(\Lambda)$ 上平坦，其中 $\Lambda$ 是诺特环，并且对
每个在 $S$ 上有限型的 $T$，都有
\[
  f_{T*}(\mathcal O_{X_T})=\mathcal O_T
\]
（若 $f$ 固有且可分，\footnote{1962 年勘误明确规定应加上
\fquote{$f$ 固有且可分}这一假设。}或者若 $S$ 是一个域的谱，则 Künneth
定理表明，最后一个条件等价于
$f_*(\mathcal O_X)=\mathcal O_S$）。于是，$X/S$ 的 Picard 函子在阿廷
$\Lambda$-代数范畴上 pro-可表。

此外，此时有
\[
  F(I_A,\xi)=H^1(X_A,\mathcal O_{X_A}),
\]
特别地：

\result{推论}{} 若 $X$ 在 $S$ 上固有，则与 Picard 函子相应的拓扑
$\Lambda$-代数 $O$ 的各个局部分量都是诺特环。

\result{评注}{} 本节的定义和结果可以推广到 $X$ 上主丛的分类，其中结构群
是 $S$ 上一个仿射、在 $S$ 上平坦且可交换的群概形 $G$。若 $G$ 不可交换，
则一个主丛的伴随群丛（其截面就是该主丛的自同构）不再平凡，因而命题3.1
不再以原有形式成立。不过，可以修改这一泛问题，使其仍然有解（至少目前在
形式几何中如此）。在本
\src{385}节及以下各节的语境中，以及每当我们为仅在同构意义下确定的对象类寻找
\fquote{模概形}时，都应牢记如下黄金法则：必要时引入辅助结构（标记点或
标记截面的元素、指定微分形式等），以消去待分类对象可能具有的自同构；所取
辅助结构应足够温和，不致实质性地改变原问题。

\subsection*{4. 簇的形式模}
\addcontentsline{toc}{subsection}{4. 簇的形式模}

设 $\Lambda$ 是剩余域为 $k$ 的诺特局部环（最常见的情形是
$\Lambda$ 等于 $k$，或为 Cohen $p$-环），并设 $X_0$ 是 $k$ 上的
预概形。对每个阿廷局部 $\Lambda$-代数 $A$，令 $F(A)$ 为下列对象在同构
意义下所成的类的集合：在 $A$ 上平坦的 $A$-预概形 $X$，连同一个同构
\[
  (*)\qquad
  X\otimes_A\kappa(A)\xleftarrow{\sim}
  X_0\otimes_k\kappa(A),
\]
其中 $\kappa(A)$ 是 $A$ 的剩余域；当然，这类平坦 $A$-预概形之间的同构
必须保持结构中给定的上述同构。若 $A$ 是不必为局部环的阿廷
$\Lambda$-代数，其局部分量为 $A_i$，就把 $F(A)$ 定义为各 $F(A_i)$ 的
乘积。这样，$F$ 就成为关于 $A$ 的乘法函子，可以称为 $X_0$（及
$\Lambda$）的模函子。若此函子可表，则与之相应的是一个剩余域为 $k$ 的
局部拓扑 $\Lambda$-代数 $O$；$O$ 的形式谱称为 $X_0$（及 $\Lambda$）
的形式模概形。关于这一概念的若干细节，参见 [2]。

这里应用下降技术时，$X_0$ 的\fquote{有限}自同构是无害的，因为它们不影响
$A$-预概形 $X$ 是否具有上述意义下的自同构。当 $A$ 不只是一个域时，$X$
不具有非平凡 $A$-自同构的充要条件是
\[
  H^0(X_0,\Theta_{X_0/k})=0,
\]
其中 $\Theta_{X_0/k}$ 表示 $X_0$ 的 $k$-导子层（即切层）。此外容易求得
（至少当 $X_0$ 在 $k$ 上光滑时）
\[
  F(I_A,\xi)=H^1(X_A,\Theta_{X_A/A}).
\]
于是照例得到：

\src{386}
\result{命题}{4.1} 设
$H^0(X_0,\Theta_{X_0/k})=0$。则 $X_0$ 的形式模概形存在。若进一步
$X_0$ 在 $k$ 上固有，则该形式模概形是诺特的。

\result{评注}{} 1\textsuperscript{o} 当不假设 $X_0$ 在 $k$ 上光滑时，
$F(I_A,\xi)$ 可等同于下式的一个 $A$-子模：
\[
  \operatorname{Ext}^1_{\mathcal O_{P_A}}
  \bigl(P_A;\mathcal I_{X_A},\mathcal O_{X_A}\bigr),
\]
其中令 $P_A=X_A\times_\Lambda X_A$；借助对角态射
$X_A\longrightarrow P_A$，将 $\mathcal O_{X_A}$ 视为 $P_A$ 上的凝聚层；
而 $\mathcal I_{X_A}$ 表示由该对角态射定义的 $P_A$ 上凝聚理想层。更确切地
说，Hochschild 理论的一个直接整体化表明，上述 $\operatorname{Ext}^1$ 可
等同于下列对象在同构意义下所成的类的集合：$X_A$ 上的平坦
$I_A$-代数层 $\mathcal O$，连同一个增广同构
\[
  \mathcal O\otimes_{I_A}A\longrightarrow\mathcal O_{X_A}
\]
（回顾一下，约定 $I_A=A[t]/(t^2)$）。子模 $F(I_A,\xi)$ 对应于可交换
代数层。因此，光滑性假设在模理论中并非本质性的，这一点与 [2] 所暗示的
不同。

2\textsuperscript{o} 回顾前引处的结果：特别地，任一在 $k$ 上光滑且固有
的代数曲线 $X_0$ 都有一个形式模概形；该概形在 $\Lambda$ 上光滑，并且当
亏格 $g$ 为 $\geq2$ 时，其相对维数为 $3g-3$，当 $g=0$ 或 $1$ 时则为 $g$。
后两种情形已不能直接归入命题4.1。不过，在椭圆曲线情形（$g=1$）中，可以
借助以下讨论将其归入该命题。

当然，还可以通过考察带有各种结构的 $k$ 上概形系统，对命题4.1 作任意变形。
例如，设 $X_0$ 是 $k$ 上带有标记原点的阿贝尔概形（即把 $X_0$ 视为 $k$
上的群概形），并令 $F(A)$ 为下列对象在同构意义下所成的类的集合：$A$
上的阿贝尔概形（即 $A$ 上固有且光滑的群概形），连同阿贝尔概形之间的一个
同构 $(*)$。可以验证，赋予乘法结构（甚至只赋予一个\fquote{单位截面}）
便会消去无穷小自同构，因而存在一个形式模概形，它对应于一个完备诺特局部环
$O$。此外，设 $X$ 是局部诺特预概形 $S$ 上的固有光滑概形，并具有
\fquote{绝对连通}的纤维；则 $X$ 上任何具有单位截面的乘法结构必然是结合
且可交换的（只要它至少在某个纤维上如此，并且 $S$ 连通），而且还唯一地
由对
\src{387}单位截面的了解所确定。此外，设 $S$ 是剩余域为 $k$ 的阿廷局部环
$A$ 的谱，$X$ 在 $A$ 上固有并配有截面 $s$，且 $X\otimes_Ak$ 配有 $k$
上的阿贝尔概形结构，其单位元为 $X\otimes_Ak$ 中与 $s$ 对应的点。一个直接
的障碍计算，结合 Serre 的一个论证，可以证明 $X$ 上确实存在以截面 $s$ 为
单位截面的乘法结构。（由此出发，先利用 [2] 的\fquote{存在性定理}过渡到
$A$ 为完备诺特局部环的情形，再利用 [3] 的下降技术处理一般情形，便可对任意
连通的局部诺特 $S$ 证明相应结论。）这就证明了这里所考察的函子 $F(A)$ 与
本节开头定义的相应函子同构；定义后者时忽略了 $X_0$ 上的乘法结构。特别地，
若 $\mathfrak m$ 是 $O$ 的极大理想，则 $\mathfrak m/\mathfrak m^2$ 典范
同构于 $H^1(X_0,\Theta_{X_0/k})$ 的对偶，因而其维数为 $n^2$，其中
$n=\dim X_0$。确定 $O$ 是否确实在 $\Lambda$ 上光滑，也就是是否同构于
$\Lambda$ 上含 $n^2$ 个不定元的形式幂级数代数，将是一个十分有趣的问题。
第1节命题5.2 将这个问题等价地表述为如下存在性问题：构造约化为某个给定
阿贝尔概形的阿贝尔概形。无论如何，超越方法表明，当 $k$ 的特征为零时，
答案是肯定的。在特征 $p\ne0$ 时，显然只需考虑 $\Lambda$ 是由代数闭域 $k$
构造的 Witt 向量环这一情形。也许现在正是\fquote{Greenberg 函子}证明其
威力的时候了\ldots\footnote{1962 年勘误指出，域上阿贝尔簇的形式模概形
在 Witt 环上光滑：任一阿廷局部商环上的阿贝尔概形都可提升到该商环所来自的
环上。证明使用了第182次报告第12页引入的提升障碍类的变换性质。勘误还回顾了
Mumford 对亏格为 $g$ 的曲线以及极化阿贝尔概形所作的模概形构造
（参见 SHMT）。}

\subsection*{5. 覆盖的延拓}
\addcontentsline{toc}{subsection}{5. 覆盖的延拓}

设 $\mathfrak X$ 是一个诺特形式预概形 [2]，$U$ 是 $\mathfrak X$ 的一个
开子集；它局部地由 $\mathcal O_{\mathfrak X}$ 的某个截面的
\fquote{不消失}所定义，该截面不仅在 $\mathfrak X$ 上，而且在每个
$X_n$ 上都是非零因子，\footnote{这一限定由 1962 年勘误加入。}因而这个开集
充分大，使得 $\mathcal O_{\mathfrak X}$ 在任一开集 $V$ 上的截面，只要在
$U\cap V$ 上为零，就必为零。设 $\mathcal J$ 是 $\mathfrak X$ 的一个
\fquote{定义理想}，并令
\[
  X_n=(\mathfrak X,\mathcal O_{\mathfrak X}/\mathcal J^{n+1}),
\]
后者于是是一个普通的诺特预概形。若
$\mathfrak X',\mathfrak X''$ 是 $\mathfrak X$ 的两个平坦覆盖（即
$\mathfrak X$ 上由代数层定义的两个预概形，这些代数层作为模层是凝聚且局部
自由的），并且它们在 $U$ 上非分歧，则自然映射
\[
  \operatorname{Hom}_{\mathfrak X}(\mathfrak X',\mathfrak X'')
  \longrightarrow
  \operatorname{Hom}_{X_0}(X'_0,X''_0)
\]
是单射；特别地，若 $\mathfrak X'$ 的一个自同构在 $X'_0$ 上诱导恒等，则
它本身就是恒等。这使下降技术可以应用于当前情形。具体地，从 $X_0$ 的一个
平坦覆盖 $X'_0$ 出发，设它在

\src{388}
$U_0$ 上非分歧；令 $G(\mathfrak X)$ 为 $\mathfrak X$ 的平坦覆盖
$\mathfrak X'$ 的同构类所组成的集合，其中这些覆盖在 $X_0$ 上诱导出
$X'_0$（因而必然在 $U$ 上非分歧），而所取的同构须在 $X'_0$ 上诱导恒等。
对于 $\mathfrak X$ 的任意开子集 $V$，以同样方式定义 $G(V)$；更一般地，
对于 $\mathfrak X$ 上的任意形式预概形 $\mathfrak Y$，定义
$G(\mathfrak Y)$。于是 [2] 和 [3] 的结果首先推出以下结论：

\noindent a. 当 $V$ 遍历 $\mathfrak X$ 的开子集时，$G(V)$ 在
$\mathfrak X$ 上组成一个层，记作 $G_{\mathfrak X}=G$。该层在 $U$ 上的
限制是各茎均仅含一个元素的常值层。

更一般地，确定 $G_{\mathfrak X}$ 的各茎归结为完备局部环的问题。确切地说：

\noindent b. 对任意 $x\in\mathfrak X$，有
\[
  G_x=G\bigl(\operatorname{Spec}(\mathcal O_{\mathfrak X,x})\bigr)
  \subset
  G\bigl(\operatorname{Spec}(\widehat{\mathcal O}_{\mathfrak X,x})\bigr),
\]
（右端集合由 $\widehat{\mathcal O}_{\mathfrak X,x}$ 上有限自由代数 $B$
的同构类组成，并配有
$B\otimes_{\widehat{\mathcal O}_{\mathfrak X,x}}\mathcal O_{X_0,x}$
与 $\widehat{\mathcal O}'_{0x}$ 之间的同构，其中 $\mathcal O'_0$ 是
$X_0$ 上定义 $X'_0$ 的代数层）。\footnote{印刷公式把 $G_x$ 直接等同于
完备化后的项，并漏掉了 $\mathcal O'_{0x}$ 上的帽号。这里的包含关系和
完备化均依照 1962 年勘误。}

\noindent c. 存在典范同构
\[
  G_{\mathfrak X}=\varprojlim_nG_{X_n} ;
\]
换言之，对 $\mathfrak X$ 的任意开子集 $V$，有
\[
  G(V)=\varprojlim_nG(V_n).
\]

\noindent d. 设 $X$ 是一个普通固有概形，定义在带有定义理想 $\mathfrak m$
的完备诺特局部环 $\Lambda$ 上，并且 $\mathfrak X$ 由它得到：
对 $\mathcal O_X$ 作 $\mathcal J$-进完备化，其中
$\mathcal J=\mathfrak m\mathcal O_X$。那么 $G(\mathfrak X)$ 典范同构于
普通概形 $X$ 的那些\fquote{约化为 $X'_0$}的平坦覆盖的同构类所组成的集合。

形象地说，(a) 和 (b) 建立了问题的局部方面与整体方面之间的基本联系，
(c) 给出了\fquote{有限}方面与\fquote{无穷小}方面之间的联系，而 (d) 则在
上述条件下重申了\fquote{形式}方面与\fquote{代数}方面的一致性。

现在设 $\mathfrak X$ 定义在完备诺特局部环 $\Lambda$ 上，并且
$\mathcal J=\mathfrak m\mathcal O_{\mathfrak X}$（因而 $X_0$ 是
$\Lambda/\mathfrak m$ 上的预概形）。对任意有限 $\Lambda$-代数 $A$，令
\[
  F(A)=G(\mathfrak X\times_\Lambda A).
\]
这是一个以 $A$ 为自变量、取值于集合范畴的协变函子；由 (c)，只要知道它
在阿廷代数 $A$ 上的取值，就能完全确定该函子。说这个函子 pro-可表，等价
于说它具有如下形式：
\[
  F(A)=\operatorname{Hom}_{\Lambda\text{-拓扑代数}}(O,A),
\]

\src{389}
其中 $O$ 是 A 第5节所考察类型的拓扑 $\Lambda$-代数；而这又等价于只对
阿廷代数 $A$ 要求同一表示。因此，定理1与定理2合在一起确实给出：

\result{命题}{5.1} 前述函子 pro-可表。

当然，由 (a) 可知，若 $U=\mathfrak X$，则对 $\mathfrak X$ 上的任意
$\mathfrak Y$，$G(\mathfrak Y)$ 都只有一个元素；此时函子 $F$ 并无多大
意义（将有 $O=\Lambda$）。看来，在几乎所有其他情形中，拓扑局部环 $O$
都不是诺特环。然而，它的存在以格外鲜明的方式揭示了由解构成的集合
$G(\mathfrak X)$ 的\fquote{连续}性质（直观上，这对应于：对 $X'_0$ 作
扩张时，分歧如何传播存在一种\fquote{连续}的选择）。可将这一结果与
J.-P. Serre 在局部类域论中的观点 [4] 相比较；后者同样揭示了一个
\fquote{几何}局部域的最大阿贝尔扩张之拓扑 Galois 群所具有的连续性质。
在那里，其对偶群（在 Pontryagin 意义下）表现为代数群（或至少是拟代数群）
的归纳极限；在这里，扩张的分类同样由无限维的\fquote{簇}给出。此外，在
上述讨论中，可以取 $\mathfrak X$ 为某个完备局部环的形式谱（例如，
$\Lambda$ 可取为它的一个 Cohen 子环）；可以期望，本节的结果能够用于研究
维数 $>1$ 的完备局部环的扩张。无论在局部情形还是整体情形中，这些结果或许
都能使人精确表述高阶分歧现象与零特征现象（后者可用超越方法处理）之间的
关系。无论如何，正是命题5.1之前的分析，使 [2] 中研究基本群的方法能够推广
到\fquote{驯分歧}的情形，并使\fquote{三点问题}能够用超越方法解决。

最后指出，当 $X_0$ 的维数为1时，情形会得到简化：由 (a) 与 (b)，
\[
  G(\mathfrak X)\simeq
  \prod_iG\bigl(\operatorname{Spec}(\widehat{\mathcal O}_{\mathfrak X,x_i})\bigr),
\]
其中 $x_i$ 是 $X_0-U$ 的各点：可以任意给定各分歧点处的\fquote{局部}
扩张。此外，当 $X_0$ 正规时，由5.1所保证存在的形式模概形在
$\operatorname{Spec}(\Lambda)$ 上光滑。

\src{390}
\begin{center}
  \textbf{参考文献}
\end{center}

\begin{description}
  \item[{[1]}] Grothendieck (Alexander). --- \emph{关于同调代数的若干问题},
    Tohoku Math. J., 第 9 卷, 1957, 第 119--221 页。
  \item[{[2]}] Grothendieck (Alexander). --- \emph{形式几何与代数几何},
    Bourbaki 讨论班, 第 11 卷, 1958/59, 分册 3,
    n\textsuperscript{o} 182。
  \item[{[3]}] Grothendieck (Alexander). --- \emph{下降技术与代数几何中的
    存在性定理，I：一般论；沿忠实平坦态射的下降}, Bourbaki 讨论班,
    第 12 卷, 1959/60, 分册 1, n\textsuperscript{o} 190。
  \item[{[4]}] Serre (Jean-Pierre). --- \emph{局部域与同源},
    Bourbaki 讨论班, 第 11 卷, 1958/59, 分册 3,
    n\textsuperscript{o} 185。
\end{description}

% END 195.tex

% BEGIN 212.tex
\unit{第212次报告——商预概形}
\sid{212}
\src{99}
\seminarzh{第13年度，1960/61，第212号\\（1961年2月）}
\paperhead{构造技术与代数几何中的存在性定理\\
III. 商预概形}%
  {亚历山大·格罗滕迪克 著}

\subsection*{引言}
\addcontentsline{toc}{subsection}{引言}

本讲所处理的问题与前两讲所考察的问题不同，因为这里试图表示的是以变动概形为
变量的某些协变函子，而不再是反变函子。不过，取商过程在代数几何的许多构造问题中
都是本质性的，其中也包括讲稿 I 和 II 中的问题（[1]、[2]）。例如，在 $T$-预概形
$X$ 上，相对于忠实平坦且拟紧的态射 $T\longrightarrow S$，一项\emph{下降数据的有效性}
问题，等价于如下问题：由该下降数据在 $X$ 中定义的平坦等价关系是否存在一个
$X$ 的商（并满足下文考察的合理性质）；[1]，A，\S~2 c 中提出的问题，大概会与
本讲第 2 节提出的问题同时得到解决。同样，$S$-概形 $X$ 的
\emph{Picard 概形}（定义见 [2]，C，\S~3）可以用多种方式定义为某些其他概形
（正因子概形，或到某个射影概形中的浸入概形）关于平坦等价关系的商；而这些辅助
概形的定义和构造在技术上更为简单：事实上，它们是 [2]，C，\S~2 中定义的
$\underline{\operatorname{Hom}}_S(X,Y)$ 型概形及其变体；在适当的射影性假设下，
它们的构造将是下一讲的主题。因此，结合本讲与下一讲的结果，便可在适当假设下构造
Picard 概形。

在预概形范畴中取商的问题仍包含若干尚未解决的疑问。其中最重要的一个见第 8 节。
目前，它仍是构造\emph{整数环上任意亏格曲线、极化阿贝尔簇等的模概形}的唯一障碍。
由此可见，这个问题值得代数群专家为之努力。

\src{100}
\subsection*{1. 等价关系与有效等价关系}
\addcontentsline{toc}{subsection}{1. 等价关系与有效等价关系}

设 $\mathcal C$ 是一个范畴，$X$ 是 $\mathcal C$ 的一个对象。若一对态射
\[
  p_1,p_2:R\rightrightarrows X
\]
满足如下条件，就称它为 $\mathcal C$ 中一个\emph{等价对}，其\emph{靶}为 $X$，
\emph{源}为 $R$：对 $\mathcal C$ 的每个对象 $T$，相应的两个映射
\[
  p_1(T),p_2(T):R(T)\rightrightarrows X(T)
\]
（这里对 $\mathcal C$ 的每个对象 $Y$，约定
$Y(T)=\operatorname{Hom}(T,Y)$）定义一个映射
\[
  R(T)\longrightarrow X(T)\times X(T)
\]
它把 $R(T)$ 双射地映到集合 $X(T)$ 上某个等价关系的图。靶为 $X$ 的各个等价对
之间有一个显然的等价关系；关于后一等价关系的一个等价类称为 $X$ 中的一个
\emph{$\mathcal C$-等价关系}；不会引起混淆时，简称等价关系。

若 $X\times X$ 存在，则在 $X$ 中给定一个等价关系，等价于给定
$X\times X$ 的一个子对象 $R$，使得对 $\mathcal C$ 的每个对象 $T$，在
\[
  (X\times X)(T)=X(T)\times X(T)
\]
中与 $R(T)$ 相应的子集，是 $X(T)$ 上某个等价关系的图。以 $p_1$ 和 $p_2$
表示由投影 $\operatorname{pr}_1$ 和 $\operatorname{pr}_2$ 诱导的从 $R$
到 $X$ 的态射，则上述条件正表示 $(p_1,p_2)$ 是一个等价对。在
$X\times X$ 以及纤维积 $(R,p_2)\times_X(R,p_1)$ 存在的前提下，也可以按照
[2]，A，\S~1 中说明的一般原则，在 $\mathcal C$ 中用交换图表示 $X(T)$ 中
$R(T)$ 所满足的集合论等价关系公理。我们不需要采用这种表述。

每当给定一对具有共同源 $R$ 和共同靶 $X$ 的态射 $(p_1,p_2)$ 时，可以把这一对的
\emph{余核}定义为 $\mathcal C$ 中表示如下关于 $Z$ 协变的函子
\[
  \operatorname{Hom}_{p_1,p_2}(X,Z),
\]
的对象 $Y$；该函子的值是所有态射 $u$（从 $X$ 到 $Z$）中满足
$up_1=up_2$ 者所成的集合。若 $Y$ 存在，则它在唯一同构意义下确定。将它记为
$X/(p_1,p_2)$，或滥用记号记为 $X/R$；后一记号主要用于
\src{101}$(p_1,p_2)$ 是等价对的情形：此时习惯上在记号中把该对定义的等价关系
与 $R$ 视为同一对象。注意，若把 $Y$ 看作 $X$ 的商，则它事实上只依赖于
$(p_1,p_2)$ 所定义的等价关系。

现在从一个态射
\[
  f:X\longrightarrow Y
\]
出发；它使我们可以把 $X$ 看作一个“$Y$ 上的对象”。假设纤维积
\[
  R(f)=X\times_Y X
\]
存在，并以 $p_1$、$p_2$ 表示它的两个投影。于是 $(p_1,p_2)$ 是一个等价对，
称为与态射 $f$ \emph{相伴的}等价对。因此，它定义一个等价关系，称为与 $f$
相伴的等价关系。

设 $(p_1,p_2)$ 是一对靶为 $X$、共同源为 $R$ 的态射。若
\begin{enumerate}
  \item[(i)] 余核 $X/(p_1,p_2)=Y$ 存在；
  \item[(ii)] 纤维积 $X\times_Y X$ 存在；
  \item[(iii)] 以 $p_1$、$p_2$ 为分量的态射
  $R\longrightarrow X\times_Y X$ 是同构，
\end{enumerate}
就称 $(p_1,p_2)$ 为一个\emph{有效等价对}。此时 $(p_1,p_2)$ 确实是一个
等价对。它所定义的等价关系也称为\emph{有效等价关系}。

若态射 $f:X\longrightarrow Y$ 满足
\begin{enumerate}
  \item[(i)] 纤维积 $X\times_Y X=R$ 存在；
  \item[(ii)] 商 $X/(p_1,p_2)$ 存在，其中 $p_1$、$p_2$ 是从 $R$ 到 $X$
  的两个投影；
  \item[(iii)] 由 $f$ 诱导的态射
  $X/(p_1,p_2)\longrightarrow Y$ 是同构，
\end{enumerate}
就称 $f$ 为\emph{有效满态射}。此时 $f$ 确实是满态射，甚至是严格满态射
（见 [1]，A，\S~2.3）；若纤维积 $X\times_Y X$ 存在，则逆命题也成立。
由满态射 $f$ 定义的 $X$ 的商对象也称为 $X$ 的一个\emph{有效商}。

上述定义蕴含如下“Galois 对应”：

\src{102}
\result{命题}{1.1} $X$ 中的有效等价关系 $R$ 所成的集合，与 $X$ 的有效商 $Y$
所成的集合之间，存在一个保持自然偏序的双射对应：$R$ 对应于有效商 $X/R$；
$Y$ 对应于由典范投影 $X\longrightarrow Y$ 定义的有效等价关系
（该关系由带有两个投影的纤维积 $X\times_Y X$ 定义）。

在性质极好的范畴中（集合、集合层等），每个商都是有效的，每个等价关系也都是
有效的。但在给定预概形 $S$ 上的预概形范畴等范畴中，情况不再如此；即使 $S$
是一个域的谱，甚至即使只考虑在 $S$ 上有限的概形，也是这样。有效性问题，
乃至商的存在性问题（对于在 $S$ 上不有限的预概形），通常都相当棘手。

\subsection*{2. 例：$S$ 上有限的预概形}
\addcontentsline{toc}{subsection}{2. 例：S 上有限的预概形}

设 $\mathcal C$ 是在 $S$ 上有限的预概形所成的范畴，并假设 $S$ 局部
诺特。它等价于 $S$ 上凝聚交换代数层所成范畴的反范畴；或者，当 $S$ 是环
$A$ 的仿射概形时，等价于在 $A$ 上有限的 $A$-代数所成范畴的反范畴
（即这些代数作为 $A$-模是有限生成的）。立即可知，$\mathcal C$ 中有限逆极限
与有限直极限都存在。对前者而言，这是熟知的事实（而且不需要任何有限性假设
\ldots）。例如，$S$ 上预概形 $X,Y$ 的纤维积，对应于相应代数的张量积
$B\otimes_A C$；由两个 $A$-代数同态
$u,v:C\rightrightarrows B$ 所定义的两个态射 $X\rightrightarrows Y$ 的核，
对应于 $B$ 关于由所有 $u(c)-v(c)$ 生成的理想的商，等等。对于有限直极限，
只须一方面考察有限和——它们对应于 $A$-代数的通常直积——另一方面考察态射对
$X\rightrightarrows Y$ 的余核；后者事实上对应于 $C$ 中由所有使同态
$u,v:C\rightrightarrows B$ 取相同值的元素组成的子环（这一点立即可见，而由
诺特假设，该子环在 $A$ 上有限）。还可以利用诺特假设证明：在
$S$ 上有限的预概形范畴 $\mathcal C$ 中如此构造的有限直极限，特别是这些商，
事实上也是\emph{所有}预概形所成范畴中的商。

\src{103}
如同我们在 [1] 中所说，$\mathcal C$ 中存在非有效满态射
（或非严格满态射\footnote{印刷文本在“严格”之前漏掉了“非”。在这里考察的范畴中，
纤维积存在，而上一段恰好证明了有效性与严格满态射性等价。}；因为纤维积存在，
这两种说法等价）。在不作平坦性假设时，\emph{我不知道等价关系是否总是有效的}。
在这个方向上，我只得到了一些很不完全的正面结果；这些结果对于证明形式模理论中的
基本定理是不可缺少的（见 [2]，B，定理 1）。应当指出，对所提出的问题，很容易
约化到如下情形：$S$ 是一个阿廷局部环 $A$ 的谱，而 $A$ 的剩余域代数闭。
但即使 $A$ 是一个域，答案仍然未知。

也可以考察 $S$ 上一个不再假设在 $S$ 上有限的预概形 $X$，同时考虑 $X$ 中一个
等价关系 $R$，使得 $p_1:R\longrightarrow X$ 是有限态射。这时称 $R$ 为
有限等价关系。为简单起见，假设 $S$ 和 $X$ 都是仿射的（于是 $R$ 也是仿射的，
从而情形归结为纯粹交换代数中的情形），则\emph{即使在这种情况下，我们也不知道
商 $X/R=Y$ 是否存在，以及典范态射 $X\longrightarrow Y$ 是否有限}。
（最简单的情形是：$S$ 是一个域 $k$ 的谱，而 $X$ 是 $k[t]$ 的谱，即仿射直线。）
当然，若前述两个问题都有肯定答案，则在当前情形中可断定 $R$ 是有效的。注意，
若把 $X$ 中等价关系的图换成第 4 节意义下 $X$ 中的一个等价预图，则商 $Y$ 的
\emph{存在性}以及 $f:X\longrightarrow Y$ 的有限性问题，仍以完全相同的形式
出现。

通过某个给定的闭子预概形将若干预概形 $X_i$ “粘合”起来以构造预概形时，会出现
关于任意程度的有限等价关系取商的问题；粘合律恰好由各 $X_i$ 的和预概形 $X$ 中
一个有限等价关系表示。还应预期，要沿 [2] 所开创的方向建立一般的非射影构造技术，
首先必须解决这里提出的问题及其各种变体。

撰稿人所知唯一的一般性正面结果如下：

\result{命题}{2.1} 设 $S$ 是局部诺特预概形，$s$ 是 $S$ 的一个点，
$\Omega$ 是 $k(s)$ 的一个代数闭扩域。考虑
\src{104}相应的“纤维函子” $F$：它把每个在 $S$ 上有限的 $S$-概形 $X$
映到 $X_s$ 的所有 $\Omega$-值点所成的集合。这个函子（显然左正合）
\emph{右正合}，即它与有限直极限交换，特别是与态射对的余核交换。

把这一结果用于 $S$ 的所有“几何点”，可推出 $\mathcal C$ 的“商”范畴
$\mathcal C'$——它通过“模去满射纯不可分态射”进行推理而得到
（即形式地为所讨论的态射添加逆）——是一个“几何”范畴，也就是说，它满足与
集合范畴相同的“有限性质”。特别地，其中每个等价关系都是有效的。这意味着，
若 $R$ 是在 $S$ 上有限的 $X$ 中的一个等价关系，则典范态射
\[
  R\longrightarrow X\times_Y X\qquad (\text{其中 }Y=X/R)
\]
是\emph{纯不可分且满射的}（事实上，因为它是单态射，所以它是一个满射闭浸入）。

\subsection*{3. 作用群的情形}
\addcontentsline{toc}{subsection}{3. 作用群的情形}

再次假设 $\mathcal C$ 是任意范畴。设 $G,X$ 是 $\mathcal C$ 的对象，并假设
$G$ 是作用于 $\mathcal C$ 的对象 $X$ 的一个 $\mathcal C$-群。这意味着
（见 [2]，A，\S~1），对 $\mathcal C$ 的每个对象 $T$，在 $G(T)$ 上给定了
群结构，并在 $X(T)$ 上给定了带 $G(T)$ 作用的集合结构，而且当 $T$ 变动时，
这些结构随 $T$ “函子性地变动”。若乘积 $G\times G$ 和 $G\times X$ 在
$\mathcal C$ 中存在，则这种结构也可以定义为一对态射
\[
  G\times G\longrightarrow G,
  \qquad \pi:G\times X\longrightarrow X,
\]
并要求对 $\mathcal C$ 的每个对象 $T$，集合 $G(T)$ 与 $X(T)$ 上相应的合成律
使 $G(T)$ 成为一个作用于 $X(T)$ 的群。把这条公理翻译成 $\mathcal C$ 中某些
图表的交换性并不困难，却很繁琐；事实上，在我所知的一切情形中，这样做都完全
没有必要。

假设 $G\times X$ 存在，并考虑两个态射
\[
  p_1,p_2:G\times X\rightrightarrows X
\]
其中
\src{105}
\[
  p_1=\operatorname{pr}_2,\qquad p_2=\pi.
\]
\footnote{原印本及法、英可编辑见证均作
$p_1=\operatorname{pr}_1$；但
$\operatorname{pr}_1:G\times X\to G$，而这里两个态射的靶都是 $X$，
且下一个公式要求 $(g,x)\mapsto(x,gx)$，故此处采用类型正确的第二投影。}
立即可见，$(p_1,p_2)$ 是等价对，当且仅当对 $\mathcal C$ 的每个对象 $T$，
由这一对定义的映射
\[
  G(T)\times X(T)\simeq(G\times X)(T)
  \longrightarrow X(T)\times X(T),
\]
是单射；也就是说，当且仅当群 $G(T)$ \emph{自由地}作用于集合 $X(T)$；
换言之，若 $g\in G(T)$、$x\in X(T)$ 且 $g\cdot x=x$，则 $g$ 必为群
$G(T)$ 的单位元。这时称 $G$ \emph{自由作用于 $X$}
（或者称 $X$ 是一个\emph{$G$ 下的 $\mathcal C$-主空间}）。这时，与
$(p_1,p_2)$ 相伴的等价关系称为\emph{由自由作用于 $X$ 的群 $G$ 所定义的
等价关系}。若 $X\times X$ 也存在，并考虑由对 $(p_1,p_2)$ 定义的态射
\[
  p:G\times X\longrightarrow X\times X
\]
则 $G$ 自由作用这一条件表示 $p$ 是\emph{单态射}。

当然，即使 $G$ 在 $X$ 上的作用并不自由，我们仍希望得到 $X$ 关于 $G$ 的商
存在的判据，也就是上述对 $(p_1,p_2)$ 的余核存在的判据。

所讨论的余核通常记为 $X/G$；若 $G$ 从左作用，则最好记为 $G\backslash X$
（前一记号留给 $G$ 从右作用的情形）。注意，即使 $G\times X$ 在 $p$ 下的
“像”存在——例如，把这个像定义为 $X\times X$ 中使 $p$ 可以通过它分解的最小
子对象——并将其记为 $R$，它通常也不是 $X$ 中的等价关系。如果这时试图直接
关于 $R$ 取商（更确切地说，关于由两个投影 $\operatorname{pr}_i$ 诱导的从
$R$ 到 $X$ 的态射对取商），就会不可挽回地丢失原始对 $(p_1,p_2)$ 的特殊性质。
因此，必须找到等价关系概念的一种推广，使它能够直接应用于由一个
$\mathcal C$-作用群定义的态射对。

\subsection*{4. 预等价关系}
\addcontentsline{toc}{subsection}{4. 预等价关系}

回顾：若一个范畴的所有态射都是同构，就称它为\emph{群胚}。另一方面，一个范畴
应当定义为由两个基础集合 $(X,R)$ 组成：\emph{对象}集合与\emph{箭头}集合；
并赋予如下结构：
\src{106}
\begin{enumerate}
\item[(i)] 一对映射
\[
  p_1,p_2:R\rightrightarrows X
\]
分别称为\emph{源映射}和\emph{靶映射}，

\item[(ii)] 一个映射
\[
  \pi:(R,p_2)\times_X(R,p_1)\longrightarrow R
\]
称为\emph{复合映射}。
\end{enumerate}
这些数据必须满足一些众所周知的公理，这里不再复述；这些公理可以明确表述为
若干图的交换性，以及存在一个映射 $D:X\to R$（它必然唯一），使另外两个图
交换。$D$ 对应于把一个对象送到相应恒等箭头的运算，并满足
\[
  p_1D=p_2D=\operatorname{id}_X.
\]
说这个范畴是群胚，是指存在一个映射
\[
  s:R\longrightarrow R
\]
（它必然唯一），称为 $R$ 的\emph{对称映射}，它把每支箭头变为一支逆箭头。
这一条件可以明确表述为由 $s$、$\Delta$ 和前述数据构成的另外四个图的
交换性，其中前两个写作
\[
  p_1s=p_2,\qquad p_2s=p_1.
\]

回顾以上各点后，([2]，A，\S~1) 中的一般定义特别说明了，应当如何理解任意
范畴 $\mathcal C$ 的一对对象 $X,R$ 上的
$\mathcal C$-\emph{范畴}结构，或相应的 $\mathcal C$-\emph{群胚}结构：
按照定义，这是对 $\mathcal C$ 的每个对象 $T$ 给出一个集合论意义下的范畴
结构或群胚结构，其对象集为 $X(T)$、箭头集为 $R(T)$，而且这些结构随变化的
$T$ “以函子方式变化”。因此，这些数据给出
两个态射
\[
  p_1,p_2:R\rightrightarrows X
\]
分别称为\emph{源态射}和\emph{靶态射}；当所需的纤维积存在时，还给出一个态射
\src{107}
\[
  \pi:(R,p_2)\times_X(R,p_1)\longrightarrow R
\]
称为\emph{复合态射}。这三项数据足以确定 $X,R$ 上的范畴结构；须对它们施加
的公理如下：对每个 $T$，在 $X(T),R(T)$ 上所得的相应三项数据，都在这对
集合上定义一个范畴结构（或群胚结构）。适当时，这可以表述为若干图的交换性；
这些图涉及一个唯一确定的态射
\[
  D:X\longrightarrow R
\]
以及在群胚情形下一个唯一确定的态射
\[
  s:R\longrightarrow R,
\]
而这些图可完全按照“集合论”情形写出。所幸在实际使用中，无须对这些公理作
如此繁琐的解释。能够借助某些纤维积之间的态射及其等式来表述上述数据与公理，
其唯一的理论意义如下：若函子 $F:\mathcal C\to\mathcal C'$ 左正合（即与
有限积和纤维积交换），它就把一个 $\mathcal C$-范畴（或
$\mathcal C$-群胚）变为一个 $\mathcal C'$-范畴（或
$\mathcal C'$-群胚）（这里要求 $\mathcal C$ 中存在有限积和纤维积）。

在实际使用中，重要的是把态射 $p_1,p_2,\pi,D,s$ 解释为 $\mathcal C$ 中某个
适当半单纯对象或单纯对象的\emph{单纯运算}，至少在 $\mathcal C$ 中存在
纤维积时如此。为固定术语，引入\emph{标准单形范畴 $\mathcal S$}：其对象是
形如
\[
  \Delta_n=[0,n],\qquad\text{其中 }n\geq 0
\]
的有限集合（即从 $0$ 到 $n$ 的整数区间），其态射或箭头则是这些有限集合
之间的\emph{任意映射}。注意，范畴 $\mathcal S$ 等价于以有限集合之间的
映射为态射的\emph{非空}有限集合范畴。在 $\mathcal S$ 中，任意有限
\emph{非空}对象族的余积显然存在，两个对象在第三个对象之下的推出也存在
（这是纤维积的对偶运算）。以 $\mathcal S'$ 表示 $\mathcal S$ 的如下
子范畴：它具有相同的对象，但其态射仅为各 $\Delta_n$ 之间的
\emph{保序映射}。这个范畴等价于非空有限全序集的范畴。在这个范畴中，
\src{108}
两个对象的余积从不存在，而两个对象 $A,B$ 在第三个对象 $C$ 之下的推出
一般也不存在（例如取 $C=\Delta_0$、$A=B=\Delta_1$，并令两个结构映射
$u:C\to A$ 和 $v:C\to B$ 相同）。不过，它在某些情形下确实存在，例如
\[
  C=\Delta_0,\qquad A=\Delta_m,\qquad B=\Delta_n,\qquad
  u(0)=m,\qquad v(0)=0,
\]
此时有
\[
  \Delta_m\amalg_{\Delta_0}\Delta_n=\Delta_{m+n}.
\]

范畴 $\mathcal C$ 中的一个\emph{单纯对象}（相应地，一个
\emph{半单纯对象}），按照定义，是从 $\mathcal S$（相应地，从
$\mathcal S'$）到 $\mathcal C$ 的反变函子 $K$。因此，单纯对象经限制给出
半单纯对象；但前一个概念与后一个概念的本质区别，在于
$K_n=K(\Delta_n)$ 上存在\emph{对称运算}，它们对应于 $n+1$ 个字母上的
对称群各元素在函子 $K$ 下的像（该对称群视为 $\Delta_n$ 在
$\mathcal S$ 中的自同构群）。

现在，对每个 $n$，令 $\Delta'_n$（相应地，$\Delta''_n$）为如下有限范畴：
其对象集为 $\Delta_n$，并由 $\Delta_n$ 上的混沌序关系（相应地，自然全序
关系）定义（箭头集即所述序关系的图）。显然，$\Delta'_n$（相应地，
$\Delta''_n$）以函子方式依赖于 $\mathcal S$（相应地，$\mathcal S'$）的
对象 $\Delta_n$。若 $Z$ 是一个范畴，则当 $\Delta_n$ 变化时，
$\operatorname{Hom}(\Delta'_n,Z)$（相应地，
$\operatorname{Hom}(\Delta''_n,Z)$）给出从范畴 $\mathcal S$
（相应地，$\mathcal S'$）到集合范畴的一个函子，即一个单纯集（相应地，
半单纯集），称为\emph{与范畴 $Z$ 相伴的}单纯集（相应地，半单纯集）；
分别记作 $Z'$ 和 $Z''$。此外，还存在一个显然的自然映射，从 $Z'$ 所给出的
半单纯集到 $Z''$；这个映射当且仅当 $Z$ 是群胚时为同构。由此得到：

\result{命题}{4.1} 从范畴所成的范畴到半单纯集范畴的函子
$Z\longmapsto Z''$ 是全忠实的，并且在\emph{范畴}所成的范畴与下述
半单纯集所成的范畴之间定义一个等价：这些半单纯集即从 $\mathcal S'$ 到
$(\mathrm{Ens})$ 的反变函子 $K$，并且把上述类型的推出
$A\amalg_C B$ 变为集合的纤维积。同样，从群胚范畴到单纯集范畴的函子
$Z\longmapsto Z'$ 是全忠实的，并且在\emph{群胚}范畴与
\src{109}
下述单纯集所成的范畴之间定义一个等价：这些单纯集即从 $\mathcal S$ 到
$(\mathrm{Ens})$ 的反变函子 $K$，并且把推出变为纤维积。

因此，可以把范畴视为特殊的半单纯集，把群胚视为特殊的单纯集；当然，在用
一种结构解释另一种结构时，必须按照惯例“在同构意义下”进行论证。通常的
化归到集合论情形的方法于是推出：

\result{推论}{4.2} 若把前一陈述中的范畴、群胚、单纯集分别换成
$\mathcal C$-范畴、$\mathcal C$-群胚、$\mathcal C$-单纯对象，并假设
$\mathcal C$ 中存在纤维积，则前一陈述仍然成立。

范畴 $\mathcal C$ 中与一个 $\mathcal C$-范畴 $(X,R,\ldots)$ 相伴的
半单纯对象 $K$ 可以明确写出：
可把 $K$ 的分量 $K_n=K(\Delta_n)$ 解释为 $(R,p_1)$ 在 $X$ 上的第
$(n+1)$ 次纤维幂；更好地，也可使用递推公式
\[
  K_0=R,\qquad
  K_n=(K_{n-1},p_n^{(n-1)})\times_X(R,p_1),
\]
\footnote{原印本及法、英可编辑见证均作 $K_0=R$，并在递推式中使用
$p_n^{(n-1)}$；但前文 $K_n=\operatorname{Hom}(\Delta''_n,Z)$ 的定义给出
$K_0=X$、$K_1=R$，且所声明的投影下标范围为
$0\leq i\leq n-1$。这里保留原式；该下标冲突已作为源文疑点记录。}
其中 $p_i^{(n-1)}$（$0\leq i\leq n-1$）是从 $K_{n-1}$ 到 $X$ 的自然
投影（它们本身也由递推定义）。这样，$p_1,p_2,\pi,D,s$ 就分别解释为对应于
$\mathcal S$ 中下列态射的单纯运算：$\Delta_1$ 的面 $0$、
$\Delta_1$ 的面 $1$、$\Delta_2$ 的面 $(0,2)$、退化
$\Delta_1\to\Delta_0$，以及 $\Delta_1$ 的对称。其他所有半单纯运算
（相应地，单纯运算）都可通过复合和纤维积，由前述四个（相应地，五个）
运算形式地导出。

现在，把一个范畴对象 $X$ 上的\emph{预等价关系}定义为如下数据：一个群胚，
其对象之对象（不妨这样说）为 $X$。因此，这种数据尤其包括一个对象 $R$
以及两个态射
\[
  p_1,p_2:R\rightrightarrows X.
\]
不过要注意，与等价对的情形不同，仅凭这些数据并不能确定所考虑的结构。
本报告研究这个概念，是为了得到能够取商的判据，即能够形成态射对
$(p_1,p_2)$ 的余核的判据。因此，问题的陈述并不涉及群胚固有的附加数据。
但在随后结果的证明中，我们仍将
\src{110}
使用这些附加数据，特别是直至维数 $3$ 的单纯运算（其中包括对称运算，并且
涉及 $R$ 在 $X$ 上的四重纤维幂）。

$\mathcal C$ 的对象 $X$ 上的一个等价关系定义一个预等价关系：只需在集合论
情形下验证这一点；对集合 $X$ 上的一个等价关系，取对象集为 $X$、箭头集为
该等价关系之图的群胚即可。

作用在 $\mathcal C$ 的对象 $X$ 上的一个 $\mathcal C$-幺半群 $G$ 定义一个
$\mathcal C$-范畴，其底层对象为 $R=G\times X$ 和 $X$（这里要求
$G\times X$ 存在）；它当且仅当 $G$ 是群时为 $\mathcal C$-群胚。仍只须在
集合论情形下验证这一点。此时，箭头 $(g,a)$ 与 $(g',g\cdot a)$ 的复合定义为
\[
  (g',g\cdot a)(g,a)=(g'g,a),
\]
也就是说，若 $a,b\in X$，则 $\operatorname{Hom}(a,b)$ 按定义是把 $a$
送到 $b$ 的输运子，而态射的复合由 $G$ 中元素的复合给出。

\medskip
\noindent\textbf{注. ---} 对于 (4.1) 这类陈述，只要约定所考虑的全部对象
都位于某个固定的“宇宙”中（该宇宙本身是一个集合），就可以避开它所引起的
逻辑困难。

\fsection{5}{按有限平坦等价关系取商}

\result{定理}{5.1} 设 $X=\operatorname{Spec}(B)$ 是仿射概形，$R$ 是 $X$
中的一个预等价关系，其分量 $R_1$ 是仿射的，记
$R_1=\operatorname{Spec}(C)$。假设第一投影 $p_1:R_1\to X$ 是有限局部自由
态射；换言之，相应的环同态 $p_1^*:B\to C$ 使 $C$ 成为有限生成投射
$B$-模。令 $A$ 为 $B$ 的如下子环：它是同态对
\[
  p_1^*,p_2^*:B\rightrightarrows C
\]
的核（即所有满足 $p_1^*(b)=p_2^*(b)$ 的元素 $b$ 所成的集合）。令
$Y=\operatorname{Spec}(A)$，并令 $f:X\to Y$ 为由 $A$ 到 $B$ 的嵌入所定义
的态射。在这些条件下：
\begin{enumerate}
\item[(i)] $B$ 在 $A$ 上整，即 $f$ 是整态射。

\item[(ii)] 态射 $f$ 是满射，其纤维是下述等价类
\src{111}
$p_2\allowbreak(p_1^{-1}(x))$，即 $X$ 模 $R$ 的集合论等价类；而且 $Y$
的拓扑是 $X$ 的拓扑的商。

\item[(iii)] $Y$ 是预概形范畴中 $X$ 按 $R$ 所取的商。

\item[(iv)] 若 $R$ 来自一个\emph{等价关系}，则态射 $f:X\to Y$ 是有限
局部自由的（即 $B$ 是有限生成投射 $A$-模），且该等价关系是有效的，即
$R_1\to X\times_Y X$ 是同构。
\end{enumerate}

这个定理推广了如下熟知定理：有限群 $G$ 以自同构作用于环 $B$，而 $A$ 是
不变量环；其证明也与熟知的证明类似。可将 (iii) 具体化如下：

\result{推论}{5.2} 典范态射
$R_1\to X\times_Y X$ 是\emph{满射}。

仍设 $R$ 是 $X$ 中一个“有限且局部自由”的预等价关系，但现在允许 $X$ 为
任意预概形。假设能够找到一个预概形 $Y$ 和一个态射 $f:X\to Y$，使得
$fp_1=fp_2$；并且还假设 $Y$ 上环层同态序列
\[
  \mathcal O_Y\longrightarrow f_*(\mathcal O_X)
  \rightrightarrows g_*(\mathcal O_R)
\]
正合（其中 $g=fp_i$）。由上述定理可推出与定理中 (i) 至 (iv) 类似的结论；
特别地，由 (iii)，$Y$ 是 $X$ 按 $R$ 所取的商，因而在唯一同构意义下确定。
在这些条件下，称 $X$ 中的预等价关系 $R$ 是\emph{可容许的}。采用这一定义：

\result{定理}{5.3} 设 $X$ 是预概形，$R$ 是 $X$ 中的预等价关系，并且
$p_1:R_1\to X$ 是有限局部自由态射。则 $R$ 可容许，当且仅当 $X$ 模 $R$
的每个集合论等价类 $p_2(p_1^{-1}(x))$ 都包含在某个仿射开集内（若 $X$
的每个有限子集都包含在某个仿射开集内，则这一条件总是成立；例如，$X$ 在
某个仿射概形上拟射影时即如此）。

事实上，不难证明 $X$ 模 $R$ 的每个等价类都包含在某个\emph{在 $R$ 下稳定}
的仿射开集内；再把利用定理 5.1 得到的各个部分粘合起来，便构造出商 $Y$。
\src{112}

\result{推论}{5.4} 假设上述条件成立，并进一步假设 $R$ 来自一个等价关系。
那么，该等价关系是有效的，即 $R_1\to X\times_Y X$ 是同构，而且
$f:X\to Y$ 是有限局部自由态射。

由“下降”立即推出：

\result{推论}{5.5} 在 (5.4) 的条件下，$X$ 在 $Y$ 上处处秩为 $n$，当且
仅当 $(R_1,p_1)$ 在 $X$ 上处处秩为 $n$。若 $X,R_1$ 是 $Z$-预概形，
$p_1,p_2$ 是 $Z$-态射，因而 $Y$ 是 $Z$-预概形，则 $X$ 在 $Z$ 上平坦
当且仅当 $Y$ 在 $Z$ 上平坦。

综上：

\result{附论}{} 给出预概形之间的一个有限、局部自由且满的态射
$f:X\to Y$，等价于给出一个配有等价关系 $R$ 的预概形 $X$，使得
$p_1:R\to X$ 有限且局部自由，并且每个等价类
$p_2(p_1^{-1}(x))$ 都包含在某个仿射开集内。

\medskip
\noindent\textbf{注 5.6.}

\smallskip
\noindent $1^\circ$ 以上没有使用任何诺特性假设。

\smallskip
\noindent $2^\circ$ 这种取商概念把 CARTIER 的“不可分下降”包括为特殊情形；
后者对应于确定如下有限局部自由态射 $f:X\to Y$：$f_*(\mathcal O_X)$
具有相对于 $\mathcal O_Y$ 的 $p$-基（其中 $X$ 是给定的预概形，且其层
$\mathcal O_X$ 被素数 $p>0$ 零化）。注意，这一结果很容易在不对局部环施加
正则性假设、也不假设 $X$ 是域上的代数概形的情况下表述。取 $X$ 为特征
$p$ 的域的谱，便得到 JACOBSON--BOURBAKI 理论。

\smallskip
\noindent $3^\circ$ GABRIEL 先前已在域 $k$ 上有限交换群的取商理论中得到了
定理 (5.3) 的一个特殊情形。[参见 (7.4)]。

\fsection{6}{按固有平坦等价关系取商}

\result{定理}{6.1} 设 $S$ 是局部诺特预概形，$X$ 是拟射影 $S$-概形，
$R$ 是 $X$ 中的一个预等价关系，并且满足：
\src{113}
\begin{enumerate}
\item[(a)] $p_1:R_1\to X$ 是固有且平坦的；
\item[(b)] $R_1\to X\times_S X$ 是有限态射；或者，由 (a)，等价地说，
它的纤维是有限的（若 $R$ 来自一个等价关系，则此条件自动满足）。
\end{enumerate}
在这些条件下：
\begin{enumerate}
\item[(i)] $X/R=Y$ 存在，并且（若 $S$ 是诺特的且 $R$ 来自一个等价关系）
在 $S$ 上拟射影。\footnote{印刷文本没有第二个假设。与报告 212 相联系的勘误
指出，此时似乎只有在 $R$ 来自一个等价关系时，拟射影性才得到了证明。}

\item[(ii)] 典范态射 $f:X\to Y$ 是满射、固有且开的；它的纤维是 $X$
模 $R$ 的等价类 $p_2(p_1^{-1}(x))$。因此，$Y$ 可与 $X$ 关于 $R$ 所定义
之集合论等价关系的商拓扑空间等同。最后，
$R_1\to X\times_Y X$ 是满射。

\item[(iii)] 若 $R$ 来自一个等价关系，则该等价关系是有效的，即
$R_1\to X\times_Y X$ 是同构；而且 $f:X\to Y$ 是平坦的
（从而是忠实平坦的）。
\end{enumerate}

在证明中，借助 $X$ 关于 $R$ 的适当准截面可将问题化归到 (5.1)；
其证明类似于 Chevalley 讨论班中商代数群的构造。

概括如下：

\result{附论}{} 设 $X$ 在局部诺特的 $S$ 上拟射影。给定一个从 $X$ 到
某个 $S$-预概形 $Y$ 的固有、忠实平坦且满的态射 $f:X\to Y$，等价于
给定 $X$ 上的一个等价关系 $R$，使得 $p_1:R\to X$ 固有且平坦。

同一方法给出下列结果：

\result{定理}{6.2} 设 $S$ 是诺特预概形，$X$ 是 $S$ 上有限型的预概形，
$R$ 是 $S$-预概形 $X$ 中的一个预等价关系。假设：
\begin{enumerate}
\item[(a)] $p_1:R_1\to X$ 平坦且为有限型；
\item[(b)] 态射 $R_1\to X\times_S X$ 是拟有限的（即纤维有限）。
\end{enumerate}
则 $X$ 中存在一个关于 $R$ \emph{饱和}的稠密开子集 $U$，使得：
\begin{enumerate}
\item[(i)] 若 $R_U$ 是 $R$ 在 $U$ 中诱导的预等价关系，则 $U/R_U$
存在，并且在 $S$ 上为有限型。

\item[(ii)] 典范态射 $U\to U/R_U$ 是满射且为开态射；它的纤维是
$R_U$ 的集合论等价类（因而 $U/R_U$ 是 $U$ 关于 $R_U$ 所定义之
集合论等价关系的商拓扑空间）。最后，态射
\[
  (R_U)_1\longrightarrow U\times_{U/R_U}U
\]
是满射。
\src{114}

\item[(iii)] 若 $R$ 来自一个等价关系，则可以使
$U\to U/R_U$ 忠实平坦，并使 $R_U$ 有效。
\end{enumerate}

这是一个本质上具有“双有理”性质的结果。

\medskip
\noindent\textbf{评注 6.3。}

\smallskip
\noindent $1^\circ$ 我不知道在 (6.1) 和 (6.2) 中假设 (b) 是否多余。
在取群的商时，它实际上迫使我们限于所有稳定子都是有限群的情形。

\smallskip
\noindent $2^\circ$ 可以问，在没有平坦性假设时，是否仍有与 (6.1)
和 (6.2) 类似的结果。我不知道这方面的任何反例。另一方面，即使保留
平坦性假设，并且只考虑满足 $p_1:R\to X$ 平坦且拟有限（但不有限）
的等价关系，当 $X$ 仿射时，$R$ 仍可能不是有效的：取覆盖 Nagata 簇的
仿射开集上所诱导的等价关系即可（在该簇上，一个二阶群以“非容许”的方式
作用）。

\fsection{7}{应用}

正如引言中所说，(6.1) 最重要的应用是构造 Picard 概形，以及求解其他
各种“模问题”；我们以后还会回到这些问题。

由此可对 SHIMURA 的下列结果给出一个简单证明：

\result{命题}{7.1} 设 $A$ 是定义在离散赋值环 $V$ 上的阿贝尔概形，
$K$ 是 $V$ 的分式域。那么，$K$ 上每个与 $A\otimes_V K$ 的某个商同源的
阿贝尔概形 $B'$ 都“在 $V$ 上具有良好约化”，即同构于某个
$B\otimes_V K$，\footnote{印刷文本作 $B\otimes_V A$。由于 $B$ 是
$V$-概形，而命题所讨论的是它的泛纤维，基因子必为 $K$。}其中 $B$ 是
$V$ 上的阿贝尔概形（回顾一下，它在本质上是唯一的）。

可以设 $B'$ 是 $A_K$ 关于某个子群概形 $C'$ 的商。（注意：$C'$ 一般
并不是“既约”的，即它的局部环中会有幂零元。）考虑 $A$ 的闭子概形 $C$，
它是 $C'$ 的“闭包”，即满足 $C_K$ 包含 $C'$ 的最小闭子概形。将有
$C_K=C'$；由 $V$ 是离散赋值环，很容易推出 $C$ 是 $V$ 上 $A$ 的
子群概形。由于 $A$ 在 $\operatorname{Spec}(V)=S$ 上固有，$C$ 也如此；
另一方面，$A$ 在 $S$ 上甚至是射影的。因此可以应用 (6.1) 构造
$A/C=B$，
\src{115}
这正是所求的 $B$。

最后，利用本质上已知的论证，可以从 (6.2) 推出下列结果：

\result{定理}{7.2} 设 $S$ 是一个阿廷环的谱，$F$ 和 $G$ 是 $S$ 上
有限型的群概形，$u:F\to G$ 是 $S$ 上群概形的同态。假设：
\begin{enumerate}
\item[(i)] $F$ 在 $S$ 上平坦；
\item[(ii)] $u$ 的核是有限的。
\end{enumerate}
在这些条件下，商概形 $G/F$ 存在；典范态射 $G\to G/F$ 是满射且为
开态射，其纤维是由 $F$ 在 $G$ 上的右作用所定义的集合论等价类。最后，
若 $u$ 是单态射，则态射 $G\to G/F$ 平坦，并且态射
\[
  G\times F\longrightarrow G\times_{G/F}G
\]
是同构。换言之，$G$ 是 $G/F$ 上的主齐性空间，其结构群（从右作用）为
$F$；更确切地说，是把 $F\times_S(G/F)$ 看作 $G/F$ 上的群概形
（参见 [1]，B，\S~6）。

\result{推论}{7.3} 在这些条件下，$G$ 在 $S$ 上平坦，当且仅当 $G/F$
在 $S$ 上平坦。若此条件成立，则关于 $F$ 取商与 $S$ 的任意基扩张交换；
若 $F$ 是 $G$ 的正规子群，则可在 $G/F$ 上赋予 $G$ 关于 $F$ 的
\emph{商群}结构。

当 $S$ 是一个域的谱时，情形尤其简单，因为每个 $S$-预概形都自动在
$S$ 上平坦。由此得到：

\result{推论}{7.4} 设 $F,G$ 是域 $k$ 上有限型的群概形，
$u:F\to G$ 是 $k$-群同态。则 $u$ 可分解为 $F\to F'\to G$，其中
$F\to F'$ 是关于 $F$ 的闭子群 $\operatorname{Ker}u$ 取商的同态，
而 $F'\to G$ 是一个闭浸入的群同态。商 $G/F=G/F'$ 存在。在 $k$ 上的
代数群中，通常的形式体系（如诺特型定理）成立。

这一结果使我们能够统一处理经典意义下代数群的取商（即在 $k$ 上不可约且
在 $k$ 上光滑的代数群），以及关于 CARTIER 所考察的“无穷小”子群
\src{116}
取商。把 CARTIER 在 DIEUDONNÉ 关于形式群的工作之后引入的“超代数”
看作 $k$ 上形式概形范畴中的群，是有利的；在适用时（即它们对应于
$k$ 上有限秩的超代数时），还可把它们看作 $k$ 上的有限代数群。

\fsection{8}{一个猜想}

这一猜想源于如下需要：对作用在“Hilbert 概形”的某些子概形上的射影群
取商（在概形理论中，Hilbert 概形取代了 Chow 簇）。

设 $S$ 是预概形，$n$ 是整数。对 $S$ 上的每个预概形 $S'$，令其对应于
群 $\operatorname{Gl}(n,\Gamma(S',\mathcal O_{S'}))$，即元素取自
$\mathcal O_{S'}$ 的截面环的 $n$ 行 $n$ 列可逆矩阵所成的群。由此得到
关于 $S'$ 的反变函子；容易证明它是可表的，因而对应于 $S$ 上的一个群概形，
而且该群概形在 $S$ 上仿射，记作 $\operatorname{Gl}(n)_S$。它的构造与
基变换相容，因此实际上这一切都来自 $\mathbf Z$ 上的一个群概形，记作
$\operatorname{Gl}(n)$。群 $\operatorname{Gl}(1)$ 称为乘法群，常记作
$\mathbf G_m$；它对应于函子
$S\longmapsto\Gamma(S,\mathcal O_S)^*$，即 $S$ 上“单位”所成的群。
有一个显然的同态 $\operatorname{Gl}(1)\to\operatorname{Gl}(n)$；
容易构造其商群，记作 $\operatorname{GP}(n-1)$，称为
\emph{$\mathbf Z$ 上次数为 $n-1$ 的射影群}。它表示如下函子：把 $S$
映到层
$\widetilde{\operatorname{Gl}(n)}_S/
\widetilde{\operatorname{Gl}(1)}_S$ 的截面群，其中
$\widetilde{\operatorname{Gl}(n)}_S$ 表示
$\operatorname{Gl}(n)_S$ 在 $S$ 上的截面芽所成的层。（须注意：
$\operatorname{GP}(n-1)_S$ 在 $S$ 上的截面一般并非来自
$\operatorname{Gl}(n)_S$ 在 $S$ 上的截面！）还应指出，至少在 $S$
诺特时，可以证明 $\operatorname{GP}(n-1)$ 也表示函子
$S\longmapsto\operatorname{Aut}_S(\mathbf P_S^{n-1})$，其中
$\mathbf P_S^{n-1}$ 是 $S$ 上相对维数为 $n-1$ 的标准射影概形。
它正是以这种方式进入模理论的。

设 $S$ 是诺特概形；如果愿意，可以假设它是仿射的。设 $X$ 是拟射影
$S$-预概形，并配备一个相对于 $S$ 甚丰沛的可逆层 $\mathcal L$。
假设群 $G=\operatorname{GP}(n)_S$ 作用在 $X$ 上，同时也作用在
$\mathcal L$ 上（并与它在 $X$ 上的作用相容），而且它在 $X$ 上
\emph{自由地}作用。\footnote{印刷文本作“在 $S$ 上”。这里所考察的
作用和取商条件都是关于 $X$ 的。}
\src{117}

\result{猜想}{8.1}\footnote{与报告 212 相联系的勘误指出，这一猜想
“终究是错误的”。它还指出，评注 8.1 中的肯定性结果后来已经得到证明，
并提及 Mumford 的工作及其改进。} 在上述条件下：

\smallskip
\noindent $1^\circ$ \emph{$G$ 所定义的等价关系是有效的，商
$Y=X/G$ 在 $S$ 上为有限型，并且典范态射 $f:X\to Y$ 平坦且为满射}
（因而 $X$ 成为 $Y$ 上的主齐性丛，其群为
$G\times_S Y=\operatorname{GP}(n)_Y$）。

\smallskip
\noindent $2^\circ$ 设 $\mathcal L'$ 是由 $\mathcal L$ 沿 $f$ 作
“忠实平坦下降”而得到的 $Y$ 上可逆层（参见 [1]，B，\S~定理 1）。
则 $\mathcal L'$ \emph{相对于 $S$ 在 $Y$ 上预丰沛}，即存在一个整数
$m$ 和一个从 $Y$ 到某个适当的标准射影概形 $\mathbf P_S^N$ 的拟有限态射，
使 $(\mathcal L')^{\otimes m}$ 同构于
$\mathcal O_{\mathbf P_S^N}(1)$ 的逆像。

应当指出，即使 $X$ 在 $S$ 上是分离的，$Y$ 仍可能不在 $S$ 上分离
（这种情形会出现在完全不病态的“模问题”中）。若 $1^\circ$ 成立，则
$Y$ 分离，当且仅当 $G$ 所定义的等价关系具有闭图，即
$G\times X\to X\times X$ 的像是闭的（此时该态射是闭浸入）。
若 $Y$ 分离，则 $\mathcal L'$ 相对于 $S$ 在 $Y$ 上预丰沛，当且仅当
它是丰沛的，即某个适当的张量幂给出一个射影浸入。在引言所述的模问题中，
可以证明最终得到的等价关系确实具有闭图。

\medskip
\noindent\textbf{评注 8.1。---} 为了使讨论具体，并且因为目前它在实践中
是最重要的情形，我们假设了 $G=\operatorname{GP}(n)_S$。对 $G$ 更合理的
假设似乎应当是：$G$ 是某个 Tohokû 群在 $S$ 上的一种“形式”
（CHEVALLEY 也已在整数上构造了这些群）。就上述猜想而言，我所知道的
唯一肯定性结果如下：设 $X$ 是\emph{特征为 $0$ 的域 $k$ 上的仿射概形，
群 $\operatorname{Gl}(n)_k$ 或 $\operatorname{GP}(n-1)_k$ 在其上
自由地作用。则 $G$ 所定义的等价关系是有效的，商 $X/G$ 也是仿射的，
并且态射 $X\to X/G$ 平坦且为满射}。证明使用下述事实（目前仅在特征
$0$ 时得到证明）：让 $G$ 作用在 $G$ 的仿射坐标环 $A$ 上，并把 $A$
看作 $k$ 上的向量空间，则 $G$ 的平凡表示只出现一次（在某个 $k$ 上
有限维且在 $G$ 下稳定的向量子空间的合成列中）。看来，系统地运用
\src{118}
$G$ 的线性表示理论，最终可能给出这一猜想的证明，至少当基底是一个域时
如此。当基底不再是域时，报告人对此一无所知。

\section*{参考文献}
\addcontentsline{toc}{section}{参考文献}

\begin{description}
\item[{[1]}] GROTHENDIECK (Alexander). --- 下降技术与代数几何中的
存在性定理，I：一般论，Séminaire Bourbaki，1959/60，
n\textsuperscript{o} 190，29 页。

\item[{[2]}] GROTHENDIECK (Alexander). --- 下降技术与代数几何中的
存在性定理，II：形式模理论中的存在性定理，Séminaire Bourbaki，
1959/60，n\textsuperscript{o} 195，22 页。

\item[{[3]}] GROTHENDIECK (A.) 与 DIEUDONNÉ (J.). --- 代数几何基础，
I：概形的语言。--- 巴黎，Presses universitaires de France，1960
（Institut des hautes Études scientifiques, Publications
mathématiques, 4）。
\end{description}

% END 212.tex

% BEGIN 221.tex
\unit{第221次报告——Hilbert 概形}
\sid{221}
\src{249}
\seminarzh{第13年度，1960/61，第221号\\（1961年5月）}
\paperhead{代数几何中的构造技术与存在性定理\\
IV. Hilbert 概形}%
  {亚历山大·格罗滕迪克 著}

\subsection*{引言}
\addcontentsline{toc}{subsection}{引言}

[2] 的 I、II 两篇报告中阐述的方法，就其实质而言，并不依赖于对所考察概形所作的任何射影性假设。遗憾的是，目前这些方法尚不能解决 [2] 的报告 II 中提出的存在性问题。在本篇报告及下一篇报告中，我们将在射影性假设下解决这些问题。所用方法具有典型的射影性质，并且实际上几乎不使用 [2] 的 I、II 两篇报告中的结果。这里我们将构造用来取代 Chow 坐标的“Hilbert 概形”，正如 [2]，II，n\textsuperscript{o}~2 中所说。在下一篇报告中，III 中所发展的概形取商理论与 Hilbert 概形理论结合起来，将使我们例如能够在相当一般的条件下构造 Picard 概形（定义见 [2]，II，n\textsuperscript{o}~3）。

总而言之，可以说我们现在已有一种大致令人满意的射影构造方法；所欠缺的仍然是\footnote{见本篇报告末尾的补遗。}一个对射影群这类“无不动点”作用的群取商的定理（参见 [2]，III，n\textsuperscript{o}~8）。在解析几何中，情形似乎甚至稍好一些（这里限于研究给定解析空间上的“射影”解析空间），因为对解析空间而言，由一个作用良好的群取商所造成的困难消失了。另一方面，无论在代数几何还是解析几何中，仍有待建立一种无须射影性假设也适用的构造方法。

\fsection{1}{层的有界族：稳定性性质}

设 $k$ 为域，$X$ 为一个 $k$-预概形；为简单起见，设其为有限型。于是，对每个扩张 $K/k$，得到一个 $K$-预概形
$X_K=X\otimes_k K$。若 $\mathcal F$ 是 $X_K$ 上的凝聚层，且 $K'$ 是 $K$ 的一个扩张，则
$\mathcal F\otimes_K K'=\mathcal F_{K'}$ 是
$X_K\otimes_K K'=X_{K'}$ 上的拟凝聚层。现在，若 $K$ 与 $K'$ 是 $k$ 的任意两个扩张，$\mathcal F$ 是 $X_K$ 上的拟凝聚层，而
$\mathcal F'$
\src{250}
是 $X_{K'}$ 上的拟凝聚层，则称 $\mathcal F$ 与 $\mathcal F'$ \emph{等价}，如果存在从 $K$、$K'$ 到 $k$ 的某个共同扩张 $K''$ 的 $k$-同态，使得
$\mathcal F_{K''}$ 与 $\mathcal F'_{K''}$ 在 $X_{K''}$ 上同构。这确实是一个等价关系；我们将研究层关于这一关系的等价类，以及等价类的集合。注意，若 $X$ 在 $k$ 上为有限型，\footnote{印刷文本此处作 $X_0$；本段中只定义了 $X$。}则凝聚层的每个等价类都可由某个 $X_K$ 上的凝聚层定义，其中 $K$ 是 $k$ 的有限型扩张。因此，在凝聚层类的定义中，可以只考虑 $k$ 的代数闭扩张；也可以只考虑一个固定的、超越次数为无穷的 $k$ 的代数闭扩张 $\Omega$。此时，$X_\Omega$ 上的两个凝聚层 $\mathcal F$ 与 $\mathcal F'$ 等价，当且仅当存在 $\Omega$ 的一个 $k$-自同构 $\sigma$，使
$\mathcal F\otimes_\Omega(\Omega,\sigma)$ 与 $\mathcal F'$ 同构。容易看出，按上述两种定义得到的凝聚层类之间存在一一对应。

设 $E$、$E'$ 为 $X$ 上凝聚层类的两个集合。考虑所有形如
$\mathcal F\otimes\mathcal F'$ 的层的类，其中 $\mathcal F$ 与
$\mathcal F'$ 是同一个 $X_K$ 上的凝聚层，$\mathcal F$ 的类属于
$E$，$\mathcal F'$ 的类属于 $E'$。这样得到一个凝聚层类的集合，记作
$E\otimes E'$。类似地定义
$\operatorname{Tor}_i(E,E')$ 等。更一般地，设函数 $U$ 对同一个
$X_K$ 上的任意 $n$ 个凝聚层序列
$\mathcal F_1,\ldots,\mathcal F_n$ 指定 $X_K$ 上凝聚层的一个集合
$U(\mathcal F_1,\ldots,\mathcal F_n)$，并且显然地与层同构以及基变换下的逆像相容；由此可定义一个仍以 $U$ 表示的函数，它对凝聚层类的任意 $n$ 个集合序列
$E_1,\ldots,E_n$ 指定一个凝聚层类的集合
$U(E_1,\ldots,E_n)$。

本节的目标，是定义层类集合中的某些所谓\emph{有界}集合，并证明最常见的运算 $U$ 作用于有界集合后，所得集合仍然有界。

设 $X$ 是诺特预概形 $S$ 上的有限型预概形。对每个 $s\in S$，纤维
$X_s$ 是 $k(s)$ 上的有限型预概形；我们按前述意义考虑 $X_s$ 上的凝聚层类。于是，“$X/S$ 的某条纤维上的凝聚层类”这一说法以及类似说法都有了意义。同样地，在每条纤维上分别进行构造，仍可考虑
$E\otimes E'$ 等运算，它们由 $X/S$ 各纤维上凝聚层类的若干集合体系
\src{251}
得到 $X/S$ 各纤维上凝聚层类的另一个集合。

\result{定义}{1.1} 设 $E$ 为 $X/S$ 各纤维上凝聚层类的一个集合。若存在 $S$ 上的一个有限型预概形 $S'$，以及
$X'=X\times_S S'$ 上的一个凝聚层 $\mathcal F'$，使 $E$ 包含于由
$\mathcal F'$ 所定义的 $X/S$ 各纤维上的层类集合中，则称 $E$ \emph{有界}。

按照定义，后一集合对每个 $s\in S$ 给出诸层
$\mathcal F'\otimes_{S'}k(s')$ 的类，其中 $s'$ 遍历 $S'$ 中位于
$s$ 上方的点（因而 $k(s')$ 是 $k(s)$ 的扩张，并且
$X_s\otimes_{k(s)}k(s')$ 可与 $X'$ 在 $s'$ 处的纤维
$X'\otimes_{S'}k(s')=X'_{s'}$ 等同）。因此可以说，有界族正是那些包含于一个由 $S$ 上有限型的 $S'$ 参数化的凝聚层\emph{代数族}中的族。

有限多个有界族之并仍然有界（取定义各个包含族的参数预概形
$S_i$ 的不交并预概形）。基变换 $T\longrightarrow S$ 将相对于
$X/S$ 的有界族变为相对于 $X_T/T$ 的有界族；若
$T\longrightarrow S$ 满射，则其逆命题也成立（更一般地，只须它的像包含给定的 $X/S$ 上的族 $E$ 中实际出现的那些 $s$）。从理论上说，这把有界族的确定归结为 $S$ 是整数环 $\mathbf Z$ 上有限型代数之谱的情形。

若 $E$ 与 $E'$ 是相对于 $X/S$ 的层类有界族，则
$E\otimes E'$ 也有界。事实上，设 $E$ 与 $E'$ 分别包含在由
$T\longrightarrow S$ 及 $X_T$ 上的 $\mathcal F$、以及
$T'\longrightarrow S$ 及 $X_{T'}$ 上的 $\mathcal F'$ 所定义的代数族中，则 $E\otimes E'$ 包含在由
$T''=T\times_S T'\longrightarrow S$ 以及 $X_{T''}$ 上的层
$\mathcal F''$ 所定义的代数族中，其中 $\mathcal F''$ 是
$\mathcal F$ 与 $\mathcal F'$ 在 $X_{T''}$ 上的逆像之张量积。这一论证之所以成立，仅仅是因为关于 $\mathcal F$、$\mathcal F'$ 的函子
$\mathcal F\otimes\mathcal F'$ 是右正合的，因而与基扩张可交换（特别是与取纤维可交换）。这一论证不能原样用于
$\operatorname{Tor}_i(E,E')$、
$\operatorname{Hom}(E,E')$、
$\operatorname{Ext}^i(E,E')$ 这类局部运算。不过仍可证明，这些运算也把有界集合变为有界集合：其方法与处理
$E\otimes E'$ 时相同，但还要使用下述类型的结果（均见 [3]，IV，6.11）：有界族 $E$ 总可以包含在一个由 $X_T$ 上的凝聚层
$\mathcal F$ 所定义的代数族中（其中 $T$ 在 $S$ 上为有限型），并且
$\mathcal F$ 相对于 $T$ 平坦。（即把原参数空间“分割成若干片”。）这类
\src{252}
适当层的平坦性性质，确实保证
$\operatorname{Tor}_i(\mathcal F,\mathcal F')$ 这类运算与任意基变换可交换。同一方法也适用于全局性质的运算：凝聚层在固有态射下的正像和导出正像、相对于固有态射的全局
$\operatorname{Ext}$（参见 [5]，III，\S~6）等；所有这些运算都将层的有界族变为层的有界族。（注意——这里，在所考虑的这类运算下，承载各个层的预概形可以发生变化。）

以下两个命题本质上也由同一平坦性方法证明；关于有限型态射纤维上的准素分解，尤其可参见 [5]，IV。

\result{命题}{1.2} 设 $E$、$E'$ 为 $X/S$ 各纤维上的层类有界集合，并假设 $X$ 在 $S$ 上固有。则：

\medskip
\noindent (i) 所有同态
$\mathcal F\longrightarrow\mathcal F'$ 的核、余核与像所成的族有界，其中 $\mathcal F$ 的类属于 $E$，$\mathcal F'$ 的类属于 $E'$。

\smallskip
\noindent (ii) 所有可置于短正合列
$0\longrightarrow\mathcal F\longrightarrow\mathcal F''
\longrightarrow\mathcal F'\longrightarrow0$
中的层 $\mathcal F''$ 所成的族有界，其中 $\mathcal F$ 的类属于
$E$，而 $\mathcal F'$ 的类属于 $E'$。

经过适当的基变换，可以假设 $E$ 与 $E'$ 分别由某个
$X_T/T$ 上的凝聚层 $\mathcal G$ 与 $\mathcal G'$ 定义，其中 $T$ 在
$S$ 上为有限型。还可假设满足某些平坦性条件，这些条件保证层
$f_{T*}(\operatorname{Hom}_{X_T}(\mathcal G,\mathcal G'))$ 与
$\operatorname{Ext}^{1}_{f_T}(\mathcal G,\mathcal G')$ 的形成同任意态射
$T'\longrightarrow T$ 的基变换可交换。此外，可以假设上述 $T$ 上的凝聚层是局部自由的。此时，设 $T_0$ 与 $T_1$ 为 $T$ 上的向量丛，其截面芽层分别为上述两个层。于是可以典范地定义
$X_{T_0}$ 上凝聚层的同态
$\mathcal G_{T_0}\longrightarrow\mathcal G'_{T_0}$，以及
$X_{T_1}$ 上凝聚层的扩张
\[
0\longrightarrow\mathcal G_{T_1}\longrightarrow\mathcal G''
\longrightarrow\mathcal G'_{T_1}\longrightarrow0,
\]
它具有显然的泛性质。第二个层定义了一个包含 (ii) 中所考察之族的代数族。上述同态的核、余核与像也同样如此；
\src{253}
考察它们即得 (i)（只须假设余核相对于 $T_0$ 平坦，而这一情形仍可通过把 $T_0$ 分割成若干片来归结到\ldots）。

\result{命题}{1.3} 设 $E$ 为 $X/S$ 各纤维上的层类有界族。则
$(\operatorname{supp}\mathcal F)_{\mathrm{réd}}$ 的结构层的类构成一个有界族，其中 $\mathcal F$ 是某个 $X_K$ 上的凝聚层，$K$ 代数闭，且 $\mathcal F$ 的类属于 $E$。

这里，$(\operatorname{supp}\mathcal F)_{\mathrm{réd}}$ 表示
$\mathcal F$ 的支集赋予其诱导约化结构；也就是说，它的结构层是
$\mathcal O_{X_K}$ 对定义 $\operatorname{supp}\mathcal F$ 的最大理想层之商。对于通过准素分解理论由 $\mathcal F$ 典范导出的层，也可以证明与 (1.3) 类似的结果。例如可取
$\mathcal F/\mathcal F_i$，其中各 $\mathcal F_i$ 分别是关于
$\mathcal F$ 支集某一分支的准素子层，并在具有该性质的子层中极小；也可取
$\mathcal O_{X_K}/\mathfrak p$，其中 $\mathfrak p$ 是与
$\mathcal F$ 相伴的素理想层；还可取
$\mathcal O_{X_K}/\mathfrak q$，其中 $\mathfrak q$ 是与
$\mathcal F$ 支集的某一分支相伴的准素理想层（这里基域为代数闭域）。

\fsection{2}{有界族与 Hilbert 多项式}

以下假设 $X$ 在 $S$ 上射影，并配备一个甚丰沛层，记作
$\mathcal O_X(1)$。对 $S$ 中一点 $s$ 的剩余域 $k(s)$ 的任一扩张
$K$，在 $X_K$ 上考虑相应的层 $\mathcal O_{X_K}(1)$；它仍然甚丰沛。

对 $X_K$ 上的每个凝聚层 $\mathcal F$，赋予函数
\[
P_{\mathcal F}(n)=\text{$X_K$ 上 $\mathcal F(n)$ 的 Euler--Poincaré 示性数},
\]
它是整数 $n$ 的一个多项式，称为 $\mathcal F$ 的\emph{Hilbert 多项式}。当 $n$ 充分大时，$P(n)$ 就是
$H^0(X_K,\mathcal F(n))$ 在 $K$ 上的维数，因为当 $i>0$ 且 $n$ 充分大时，
$H^i(X_K,\mathcal F(n))$ 均为零。

现在若 $\mathcal F$ 是 $X$ 上相对于 $S$ 平坦的凝聚层，则在属于
$S$ 的同一个连通分支的各纤维 $X_s$ 上诱导出的层
$\mathcal F_s$，其 Hilbert 多项式全都相等 [5]，III，\S~7。由此可知（无须平坦性假设），对 $X$ 上的任一凝聚层 $\mathcal F$，诸层
$\mathcal F_s$（$s\in S$）的 Hilbert 多项式所成的集合是有限的。
\src{254}
另一方面，回忆若 $\mathcal F$ 是 $X$ 上的凝聚层，则当 $n$、$N$ 充分大时，它同构于某个层 $\mathcal O_X(-n)^N$ 的商层。因此，纤维上诱导出的层 $\mathcal F_s$ 也都是该纤维上的层
$\mathcal O(-n)^N$ 的商层。\footnote{印刷文本及法、英可编辑见证在此
第二次出现时都漏去了指数 $N$；上一句明确给出
$\mathcal O_X(-n)^N$，取纤维仍保留同一重数。}

由这两点可得下述定理的“必要性”部分：

\result{定理}{2.1} 设 $X$ 在诺特预概形 $S$ 上射影，
$\mathcal O_X(1)$ 在 $X$ 上相对于 $S$ 甚丰沛。设 $E$ 为
$X/S$ 各纤维上的层类集合。$E$ 有界的充分必要条件\footnote{印刷文本仅作“必要”；1962 年关于报告 221 的勘误规定改为“充分必要”。}是满足以下条件：

\medskip
\noindent (a) 存在 $X$ 上的一个凝聚层 $\mathcal L$（可假设其形如
$\mathcal O_X(-n)^N$），使得 $E$ 包含于所有形如 $\mathcal L_K$ 的层之凝聚商层的类所成的族中。

\smallskip
\noindent (b) 类属于 $E$ 的诸层 $\mathcal F$ 的 Hilbert 多项式
$P_{\mathcal F}$ 都属于同一个有限多项式集合。

尚须证明“充分性”部分；它将是一个更精确结果的特例。对域 $K$ 上有限型预概形上的任一凝聚模 $\mathcal F$ 以及任一整数 $r$，令
$N_r$ 为 $\mathcal F$ 的子模，其在每个开集上的截面，正是
$\mathcal F$ 在该开集上支集维数 $<r$ 的那些截面。于是，当
$r>\dim\operatorname{supp}\mathcal F$ 时
$N_r=\mathcal F$；当 $r\leq0$ 时 $N_r=0$。由此得到
$\mathcal F$ 的一个有限递增滤过，其因子
$N_{r+1}/N_r$\footnote{印刷文本把此商颠倒写成
$N_r/N_{r+1}$，这与刚定义的递增滤过不相容。}的相伴素循环，恰为与
$\mathcal F$ 相伴且维数为 $r$ 的素循环。置
\[
\mathcal F_{(r)}=\mathcal F/N_r.
\]
于是 $\mathcal F_{(r)}$ 的相伴素循环，恰为与 $\mathcal F$ 相伴且维数
$\geq r$ 的素循环；特别地，$\mathcal F_{(r)}=\mathcal F$ 当且仅当
$\mathcal F$ 的全部相伴素循环的维数均 $\geq r$。在此记号下：

\result{定理}{2.2} 在 (2.1) 的条件下，设 $s$ 为整数，并假设 $E$ 满足条件 (a)，以及下列条件 (b) 的弱化形式：

\medskip
\noindent $(b_s)$ 类属于 $E$ 的诸层 $\mathcal F$ 的 Poincaré 多项式
$P_{\mathcal F}$，其次数 $\geq s-1$ 的系数保持有界。

在这些条件下，层 $\mathcal F_{(s)}$（其中 $\mathcal F$ 的类遍历
$E$）构成有界族。此外，$P_{\mathcal F}$ 的 $s-2$ 次系数保持有下界。
\src{255}
因此：

\result{推论}{2.3} 假设类属于 $E$ 的诸层 $\mathcal F$ 的全部相伴素循环，其维数 $d$ 均满足 $s\leq d\leq r$。则在 (2.1) 的条件
(b) 中，只须考虑 $P_{\mathcal F}$ 的 $s-1$ 次至 $r$ 次系数。

本节余下部分用于概述 (2.2) 的证明。关键是以下两个引理，其中第一个是熟知的（并概括了 Chow 坐标中有用的数学内容）：

\result{引理 2.4（CHOW）}{} 考虑各纤维 $X_K$ 的子概形 $Y$ 的结构层，其中 $K$ 遍历 $S$ 的各剩余域的代数闭扩张；假设 $Y$ 约化，且它的所有分支都有同一维数 $r$（将 $\mathcal O_Y$ 视为
$\mathcal O_X$ 的商层）。若 $Y$ 的次数保持有界，则这些 $Y$ 构成一个有界族。

这里，最方便的做法是把 $Y$ 的次数 $a$ 定义为下式首项中的系数：
\[
P_{\mathcal O_Y}=a\frac{n^r}{r!}+\cdots.
\]

\result{引理}{2.5} 设 $\mathcal L$ 为 $X$ 上的凝聚层，$E$ 为形如
$\mathcal L_K$ 的层之商层 $\mathcal F$ 的类所成的集合，其中 $K$ 是
$S$ 的某个剩余域的扩张。假设 $X$ 在 $S$ 上的纤维维数
$\leq r$，并写成
\[
P_{\mathcal F}(n)=a_{\mathcal F}\frac{n^r}{r!}
+b_{\mathcal F}\frac{n^{r-1}}{(r-1)!}
+\text{次数 $<r-1$ 的项}.
\]
则系数 $a_{\mathcal F}$ 保持有界，而 $b_{\mathcal F}$ 保持有下界。若
$b_{\mathcal F}$ 保持有界，则诸 $\mathcal F_{(r)}$ 构成有界族。

\emph{证明。}——以一族覆盖 $S$ 的子概形之并替代 $S$ 后，可以假设存在有限态射
$f:X\longrightarrow\mathbf P^r_S$，使 $\mathcal O_X(1)$ 同构于
$\mathcal O_{\mathbf P^r_S}(1)$ 的逆像。因此，对 $X$ 上的每个凝聚层
$\mathcal F$，有
$P_{\mathcal F}=P_{f_*(\mathcal F)}$。另一方面，容易验证（使用上一节的方法），$X$ 上诸层 $\mathcal F$ 所成的集合有界，当且仅当诸
$f_*(\mathcal F)$ 所成的集合有界。最后有
\[
f_*(\mathcal F)_{(r)}=f_*(\mathcal F_{(r)}).
\]
\src{256}
因此可归结到 $X=\mathbf P^r_S$ 的情形。此外，可以假设
$\mathcal L=\mathcal O_{\mathbf P^r_S}(k)^s$，其中 $k$、$s$ 适当。系数
$a_{\mathcal F}$ 满足
\[
0\leq a_{\mathcal F}\leq s
\]
因而保持有界。现在，断言 $P_{\mathcal F}(n)$ 的 $n^{r-1}$ 系数保持有下界（相应地，保持有界），等价于对
$P_{\mathcal F}(n-k)=P_{\mathcal F(-k)}(n)$ 作同一断言。因此可归结到
\[
\mathcal L=\mathcal O_{\mathbf P^r_S}^{,s}.
\]

考虑正合列
\[
0\longrightarrow N_r\longrightarrow\mathcal F
\longrightarrow\mathcal F_{(r)}\longrightarrow0,
\]
从而
\[
P_{\mathcal F}=P_{\mathcal F_{(r)}}+P_{N_r}.
\]
由于 $P_{N_r}$ 中 $n^{r-1}$ 的系数为正（因为
$\dim\operatorname{supp}N_r\leq r-1$），故
\[
b_{\mathcal F_{(r)}}\leq b_{\mathcal F}.
\]
因此在证明本引理的各断言时，可以用 $\mathcal F_{(r)}$ 替代
$\mathcal F$；亦即，可以假设所考察的 $\mathcal L$ 的商层
$\mathcal F$ 均无挠。

由于 $\mathbf P^r_K$ 正规，可知 $\mathcal F$ 在开集
$U=\mathbf P^r_K-Y$ 上是秩为 $a=a_{\mathcal F}$ 的局部自由层，其中
$Y$ 的余维 $\geq2$。于是 $\bigwedge^a\mathcal F$ 是
$\mathbf P^r_K$ 上的层，其在 $U$ 上的限制可逆；因而（由于
$\mathbf P^r_K$ 正则且 $Y$ 的余维 $\geq2$）它同构于
$\mathbf P^r_K$ 上某个可逆层在 $U$ 上的限制，而后者在同构意义下唯一确定。该可逆层形如
$\mathcal O_{\mathbf P^r_K}(d)$，其中整数 $d$ 唯一确定。由于
$\bigwedge^a\mathcal F$ 是
$\bigwedge^a\mathcal O_{\mathbf P^r_K}^{,s}\simeq
\mathcal O_{\mathbf P^r_K}^{,N}$\footnote{印刷文本此处以一个未定义的
$n$ 代替指数 $s$；$s$ 是满射到 $\mathcal F$ 的自由层之秩。}的商，其中
$N=\binom{s}{a}$，它有 $N$ 个典范截面，因而在 $U$ 上定义了
\src{257}
$\mathcal O_{\mathbf P^r_K}(d)$ 的截面；这些截面是
$\mathcal O_{\mathbf P^r_K}(d)$ 的截面 $s_i$（$1\leq i\leq N$）在
$U$ 上的限制（因为 $\mathbf P^r_K$ 正规且 $Y$ 的余维 $\geq2$）。这些截面在 $U$ 的各点生成
$\mathcal O_{\mathbf P^r_K}(d)$，所以它们不全为零；由此推出
$d\geq0$。另一方面，简单计算给出
\[
b_{\mathcal F}=a_{\mathcal F}(r+1)/2+d.
\]
特别地，这表明 $b_{\mathcal F}>0$，故 $b_{\mathcal F}$ 保持有下界。
$b_{\mathcal F}$ 保持有界，当且仅当 $d$ 保持有界；下面证明在这种情形下 $\mathcal F$ 属于一个有界族。可以假设
$a_{\mathcal F}$ 与 $b_{\mathcal F}$ 固定，分别为 $a$ 与 $b$，因而
$d$ 固定。知道 $\mathcal O_{\mathbf P^r_K}(d)$ 的这 $N$ 个截面
$s_i$，亦即知道一个同态
\[
s:\bigwedge^a\mathcal L_K\longrightarrow\mathcal O_{\mathbf P^r_K}(d),
\]
便可把 $\mathcal F$ 恢复为下述相应复合同态的余像：
\[
\mathcal L_K\longrightarrow
\operatorname{Hom}\!\left(\bigwedge^{a-1}\mathcal L_K,
\bigwedge^a\mathcal L_K\right)\longrightarrow
\operatorname{Hom}\!\left(\bigwedge^{a-1}\mathcal L_K,
\mathcal O_{\mathbf P^r_K}(d)\right),
\]
其中第一个同态是由外积给出的典范同态，第二个同态由 $s$ 导出。于是由
(1.2) (i) 得到结论。

结合以上两个引理即可证明：

\result{引理}{2.6} 在 (2.1) 的预备假设下，设对每个 $\mathcal F$ 均有
\[
P_{\mathcal F}(n)=a_{\mathcal F}\frac{n^r}{r!}
+b_{\mathcal F}\frac{n^{r-1}}{(r-1)!}
+\text{次数 $<r-1$ 的项},
\]
并且系数 $a_{\mathcal F}$ 保持有界。则系数 $b_{\mathcal F}$ 保持有下界。若 $b_{\mathcal F}$ 保持有界，则诸 $\mathcal F_{(r)}$ 构成有界族。

可以假设诸层 $\mathcal F$ 的基域 $K$ 都是代数闭域。对每个
$\operatorname{supp}\mathcal F_{(r)}$ 赋予诱导约化结构；它是
$\operatorname{supp}\mathcal F$ 的维数为 $r$ 的各分支之并。\footnote{印刷文本作“次数 $r$”；滤过 $\mathcal F_{(r)}$ 是按支集维数而非次数定义的。}于是
$\operatorname{supp}\mathcal F_{(r)}$ 的次数以 $a$ 为上界，所以由
(2.4)，诸 $\operatorname{supp}\mathcal F_{(r)}$ 构成有界集合。此外，对
$\operatorname{supp}\mathcal F_{(r)}$ 的每个分支，$\mathcal F_{(r)}$ 沿该分支的长度都 $\leq a$；因此，若 $\mathcal J_{\mathcal F}$ 是定义
$\operatorname{supp}\mathcal F_{(r)}$ 的理想，则可把
$\mathcal F_{(r)}$
\src{258}
视为 $X$ 的子概形 $Y_{\mathcal F}$ 上的模，其中
$Y_{\mathcal F}$ 由 $\mathcal J_{\mathcal F}^{,a}$ 定义。与上一引理相同，还可归结到 $\mathcal F=\mathcal F_{(r)}$ 的情形，因此
$\mathcal F$ 来自 $Y_{\mathcal F}$ 上的一个模。诸 $Y_{\mathcal F}$ 对应于
$\mathcal O_{X_K}$ 的商模所成的一个有界族，因而来自概形
$X\times_S T$ 的某个闭子概形 $Y$。于是可以把 (2.5) 应用于
$Y/T$ 及 $\mathcal L\otimes_XY$，从而得到结论。

现在可以对 $\dim\operatorname{supp}\mathcal F$ 的上界 $r$ 作归纳来证明
(2.2)。当 $r<0$ 时结论显然。现设 $r\geq0$，并假设结论已对所有
$r'<r$ 得证。由 (2.6)，诸 $\mathcal F_{(r)}$ 构成有界集合；因此再由
(1.2) (i)，诸同态
$\mathcal L_K\longrightarrow\mathcal F_{(r)}$ 的核也构成有界集合。所以存在 $X$ 上的凝聚模 $\mathcal L'$，使得上述各核以及
\[
N_r(\mathcal F)=\operatorname{Ker}(\mathcal F\longrightarrow\mathcal F_{(r)})
\]
也都是诸模 $\mathcal L'_K$ 的商。由于诸 $\mathcal F_{(r)}$ 有界，
$P_{\mathcal F_{(r)}}$ 保持有界；公式
\[
P_{\mathcal F}=P_{\mathcal F_{(r)}}+P_{N_r}
\]
于是表明，$P_{N_r}$ 满足与 $P_{\mathcal F}$ 相同的条件 $(b_s)$。因此
$N_r$ 满足条件 (a) 与 $(b_s)$；由归纳假设，诸
$(N_r)_{(s)}$ 保持有界。另一方面，有短正合列
$0\longrightarrow(N_r)_{(s)}\longrightarrow\mathcal F_{(s)}
\longrightarrow\mathcal F_{(r)}\longrightarrow0$，
所以由 (1.2) (ii)，诸 $\mathcal F_{(s)}$ 保持有界。为证明 (2.2)
的最后一个断言，注意同态
$\mathcal F\longrightarrow\mathcal F_{(s)}$ 的核 $N_s$ 由 (1.2) (i) 保持有界，而且 $P_{N_s}$ 中 $n^{s-1}$ 项的系数保持有界；引理 2.6 于是证明下一项的系数保持有下界。证明完毕。

\fsection{3}{Hilbert 概形：定义与存在性定理}

设 $X$ 是另一个预概形 $S$ 上的预概形，$\mathcal F$ 是 $X$ 上的拟凝聚模。记
\[
\operatorname{Quot}(\mathcal F/X/S)
\]
为 $\mathcal F$ 的、在 $S$ 上平坦的拟凝聚商模所成的集合。现在设
$S'\longrightarrow S$ 为一个基变换态射，置
$X'=X\times_S S'$ 以及
$\mathcal F'=\mathcal F\otimes_{\mathcal O_S}\mathcal O_{S'}$，于是
$X'$ 是 $S'$ 上的预概形，并配备了一个

\src{259}
拟凝聚模 $\mathcal F'$，并可考虑
$\operatorname{Quot}(\mathcal F'/X'/S')$。置
\[
\underline{\operatorname{Quot}}_{\mathcal F/X/S}(S')
=\operatorname{Quot}(\mathcal F'/X'/S')
\qquad\text{（其中 $X'=X\times_SS'$）}.
\]
现在若 $S''\longrightarrow S'$ 是一个 $S$-态射，则
$X''=X\times_SS''$ 同构于 $X'\times_{S'}S''$，而 $\mathcal F''$ 同构于
$\mathcal F'\otimes_{\mathcal O_{S'}}\mathcal O_{S''}$。由于从 $X'$ 上
拟凝聚模范畴到 $X''$ 上拟凝聚模范畴的逆像函子
\[
\mathcal G'\longmapsto
\mathcal G'\otimes_{\mathcal O_{S'}}\mathcal O_{S''}
\]
是右正合的，并把 $S'$-平坦模变为 $S''$-平坦模，因而得到一个自然映射
\[
\operatorname{Quot}(\mathcal F'/X'/S')\longrightarrow
\operatorname{Quot}(\mathcal F''/X''/S''),
\]
所以，以 $S'$（$S$ 上的预概形）为变量，
$\underline{\operatorname{Quot}}_{\mathcal F/X/S}(S')$ 是一个取值于
集合范畴的反变函子。下文中，我们假设 $X$ 在诺特预概形 $S$ 上射影，
$\mathcal F$ 凝聚，并且为简便起见，只考虑 $S$ 上局部诺特的 $S'$。

\result{定理}{3.1} 在这些条件下，局部诺特 $S$-预概形范畴上的反变函子
$\underline{\operatorname{Quot}}_{\mathcal F/X/S}$ 可由一个 $S$-预概形
$\operatorname{Quot}_{\mathcal F/X/S}$ 表示；该预概形是一列射影
$S$-概形的不交并（特别地，
$\operatorname{Quot}_{\mathcal F/X/S}$ 在 $S$ 上局部有限型）。

我们将如下得到这种分解。设 $\mathcal O_X(1)$ 是 $X$ 上相对于 $S$
甚丰沛的可逆层。对于任意有理系数多项式 $P(n)$，令
$\operatorname{Quot}^P(\mathcal F/X/S)$ 为
$\operatorname{Quot}(\mathcal F/X/S)$ 中如下部分：它由 $\mathcal F$
的凝聚商 $\mathcal G$ 构成，这些商在 $S$ 上平坦，并且在每个
$s\in S$ 处的 Hilbert 多项式均等于 $P$。于是置
\[
\underline{\operatorname{Quot}}^P_{\mathcal F/X/S}(S')
=\operatorname{Quot}^P(\mathcal F'/X'/S'),
\]
从而得到 $\underline{\operatorname{Quot}}_{\mathcal F/X/S}$ 的一个
子函子。第2节中回顾的 Hilbert 多项式不变性意味着如下结论：
$\underline{\operatorname{Quot}}_{\mathcal F/X/S}$ 可表，当且仅当各个
$\underline{\operatorname{Quot}}^P_{\mathcal F/X/S}$ 均可表；此时，
表示它的 $S$-预概形 $\operatorname{Quot}_{\mathcal F/X/S}$ 同构于
由各个 $\operatorname{Quot}^P_{\mathcal F/X/S}$ 构成的不交并预概形，
其中 $\operatorname{Quot}^P_{\mathcal F/X/S}$ 表示函子
$\underline{\operatorname{Quot}}^P_{\mathcal F/X/S}$。因此，(3.1)
是下述定理的推论：
\src{260}
\result{定理}{3.2} 沿用上述记号，函子
$\underline{\operatorname{Quot}}^P_{\mathcal F/X/S}$ 可由射影
$S$-预概形 $\operatorname{Quot}^P_{\mathcal F/X/S}$ 表示。

本节余下部分用于证明 (3.2)。

设 $\nu$ 为一个整数。对于 $S$ 上的任意 $S'$，以 $A_\nu(S')$
表示 $\mathcal F'=\mathcal F\otimes_{\mathcal O_S}\mathcal O_{S'}$
的如下商
$\mathcal G=\mathcal F'/\mathcal K$ 的集合：这些商是凝聚的、在
$S'$ 上平坦，并满足下列条件：
\begin{align*}
\text{a.}\quad &R^if'_*(\mathcal G(n))=0 &&\text{当 $i>0$ 且 $n\geq\nu$ 时},\\
\text{b.}\quad &R^if'_*(\mathcal K(n))=0 &&\text{当 $i>0$ 且 $n\geq\nu$ 时},\\
\text{c.}\quad &f'_*(\mathcal K(\nu+k))
=\mathcal S'_k f'_*(\mathcal K(\nu)) &&\text{当 $k\geq0$ 时}.
\end{align*}
在最后一个关系中，假设 $X$ 已写成 $S$ 上一个拟凝聚正分次代数
$\mathcal S_*$ 的齐次素谱；该代数由 $\mathcal S_1$ 生成。置
$\mathcal S'_* =\mathcal S_*\otimes_{\mathcal O_S}\mathcal O_{S'}$，
于是 $X'$ 是 $\mathcal S'_*$ 的齐次素谱。为证明 (3.2)，容易归约到
$X=\mathbf P^r_S$ 的情形（因为 $S$ 是一些开集 $U$ 的并，在每个
$U$ 上，$X|U$ 是 $\mathbf P^r_U$ 的闭子概形，且
$\mathcal O_X(1)$ 由 $\mathcal O_{\mathbf P^r_U}(1)$ 诱导），并归约到
$\mathcal F$ 形如 $\mathcal O_{\mathbf P^r_S}^{,N}$ 的情形；此时
$\mathcal F$ 在 $S$ 上平坦。于是，在上述讨论中，各层 $\mathcal K$
也在 $S'$ 上平坦。由 Künneth 关系（[5]，III，\S~7）可知，条件
(a) 和 (b) 在基变换下稳定，并且蕴涵：当 $n\geq\nu$ 时，
$f'_*(\mathcal G(n))$ 和 $f'_*(\mathcal K(n))$ 的形成与基扩张可交换
（同上）；因此 (c) 在基扩张下也稳定。

换言之，以 $S'$ 为变量的 $A_\nu(S')$ 是一个反变函子；更确切地说，
它是
$A(S')=\underline{\operatorname{Quot}}^P_{\mathcal F/X/S}(S')$
的一个子函子。当 $\nu$ 变化时，由此得到 $A(S')$ 的递增子集列
$A_\nu(S')$；根据 Serre 的一个著名定理 [5]，III，\S~2，其并集为
$A(S')$。现在设 $\mathcal G$ 是 $\mathcal F'$ 的凝聚商，并在 $S'$
上平坦；设 $s$ 是 $S'$ 的一个点，使得基变换
$\operatorname{Spec}(k(s))\longrightarrow S'$ 给出 $\mathcal F_s$
的一个商 $\mathcal G_s$，且该商满足条件 (a)、(b)、(c)，亦即属于
$A_\nu(\operatorname{Spec}(k(s)))$。那么存在 $s$ 的一个开邻域 $U$，
使同样的条件对 $\mathcal G|f'^{-1}(U)$ 成立，亦即这个商属于
$A_\nu(U)$。对于 (a) 和 (b)，这确由“全纯函数定理”[5]，III，\S~7
推出；而 (c) 则由 Nakayama 引理以及下述总是成立的事实推出：
$f'_*(\mathcal K(n+k))=\mathcal S'_k f'_*(\mathcal K(n))$
\src{261}
其中 $n$ 足够大且 $k\geq0$（参见 [5]，III，\S~2）。由这些说明可得
如下结论（比较 [4]，IV）：函子 $A$ 可表，当且仅当各函子 $A_\nu$
均可表；此时，表示 $A$ 的 $S$-预概形 $Q$ 是表示各 $A_\nu$ 的开
子预概形 $Q_\nu$ 的递增并。

令
\[
M_* =\bigoplus_{n\geq0}f_*(\mathcal F(n))=\mathcal S_*^{,N},
\]
于是
\[
M'_* =M_*\otimes_{\mathcal O_S}\mathcal O_{S'}
=\bigoplus_{n\geq0}f'_*(\mathcal F'(n))=\mathcal S'_*{}^{N}.
\]
由 (a) 可得：

\noindent (a') 当 $n\geq\nu$ 时，$f'_*(\mathcal G(n))$ 局部自由，秩为
$P(n)$，[5]，III，\S~7；而由 (b) 取 $i=1$ 可得：

\noindent (a'') $f'_*(\mathcal G(n))$ 是 $M'_n$ 的一个商模。

此外，由 (c)，只要知道 $n=\nu$ 时的这个商模，就能知道所有
$n\geq\nu$ 时 $M'_n$ 的子模 $f'_*(\mathcal K(n))$，从而知道
$\mathcal K$，进而知道 $\mathcal G$。由此得到一个单射
\[
A_\nu(S')\longrightarrow\operatorname{Grass}_{P(\nu)}(M'_\nu),
\]
它把 $A_\nu(S')$ 映入 $M'_\nu$ 的秩为 $P(\nu)$ 的局部自由商模集合；
因而得到函子性同态
\[
i_\nu:A_\nu(S')\longrightarrow
\underline{\operatorname{Grass}}_{P(\nu)}(M_\nu)(S'),
\]
其中右端函子由 Grassmann 概形
$\operatorname{Grass}_{P(\nu)}(M_\nu)$ 表示（比较 [4]，V），该概形
在 $S$ 上射影。我们断言：

\result{引理}{3.3} 函子 $A_\nu(S')$ 可表，并且表示同态 $i_\nu$ 的态射
$Q_\nu\longrightarrow\operatorname{Grass}_{P(\nu)}(M_\nu)$ 是一个
浸入（这蕴涵 $Q_\nu$ 在 $S$ 上拟射影）。

该断言等价于如下断言（比较 [4]，IV）：设已给定 $M'_\nu$ 的一个商模
$N$，它局部自由且秩为 $P(\nu)$。那么存在 $S'$ 的一个子预概形 $Z$，
使得对于 $S'$ 上的任意局部诺特预概形 $T'$
\src{262}
，$N$ 在 $T'$ 上的拉回属于 $\operatorname{Im}A_\nu(T')$，当且仅当
$T'\longrightarrow S'$ 通过子预概形 $Z$ 分解。更换记号后，可以
假设 $S'=S$；亦即，我们给定了 $M_\nu$ 关于一个子模 $R_\nu$ 的商
$N_\nu$。它来自 $A_\nu(S)$ 的一个元素，当且仅当满足下列两个条件：

\medskip
\noindent (i) 对于 $k\geq0$，
$M_{\nu+k}/\mathcal S_kR_\nu$ 局部自由且秩为 $P(\nu+k)$。

\smallskip
\noindent (ii) 由 $M_*$ 的分次子模
\[
R_* =\sum_{k\geq0}\mathcal S_kR_\nu
\]
（[5]，II，\S~3）所定义的 $\mathcal F$ 的子层 $\mathcal K$，以及商
$\mathcal G=\mathcal F/\mathcal K$，满足上面的条件 (a) 和 (b)。

这些条件显然是必要的。反之，若它们成立，则 (ii) 中定义的层
$\mathcal G$ 同构于分次 $\mathcal S_*$-模 $N_*$ 的相伴层，其中
$N_*$ 是各个 $M_{\nu+k}/\mathcal S_kR_\nu$ 的直和；因此
$\mathcal G$ 在 $S$ 上平坦（因为它的各纤维都是 $N_*$ 在
$\mathcal S_*$ 的齐次素理想处局部化所得模的直和项），并且由 (i)
可知其 Hilbert 多项式为 $P$。再考虑 (ii)，可知 (a') 和 (a'')
均成立；所以当 $n\geq\nu$ 时，自然同态
$N_n\longrightarrow f_*(\mathcal G(n))$ 是相同秩的局部自由模之间
的满同态，因而是同构。因此
\[
f_*(\mathcal K(\nu+k))=\mathcal S_kR_\nu
\quad\text{对于每个 $k\geq0$},
\]
这就证明了 $\mathcal G\in A_\nu(S)$，并且 $R_\nu$ 正是由
$\mathcal G$ 所定义的
$\operatorname{Grass}_{P(\nu)}(M_\nu)$ 的元素。

条件 (i)、(ii) 所构成的这个判据同样适用于基变换
$S'\longrightarrow S$ 之后得到的情形。我们先来证明：基变换
$S'\longrightarrow S$ 后条件 (i) 成立，可以表述为
$S'\longrightarrow S$ 通过 $S$ 的某个子预概形 $Z$ 分解。得到这一
结果后，以 $Z$ 代替 $S$，便归约到条件 (i) 已在 $S$ 上成立的情形；
由于该条件在基变换下稳定，余下只须表述条件 (ii)。此时，令 $U$
表示所有如下 $s\in S$ 的集合：$\mathcal G(n)$ 和 $\mathcal K(n)$
在纤维 $X_s$ 上所诱导的层，其所有次数 $>0$ 的上同调在
$n\geq\nu$ 时均为零。我们已经指出 $U$ 是开集，而基变换
$S'\longrightarrow S$ 后条件 (ii) 成立，当且仅当
$S'\longrightarrow S$ 通过 $U$ 分解。这就证明了引理 3.3。现在只需
证明下述引理：

\result{引理}{3.4} 设 $S$ 是一个局部诺特预概形，配备一个
由 $\mathcal S_1$ 生成的拟凝聚正分次代数 $\mathcal S_*$；设 $M_*$
是一个有限型拟凝聚分次 $\mathcal S_*$-模，$P$ 是一个有理系数
多项式，$\nu$ 是一个整数。那么存在一个（显然唯一的）$S$ 的
子预概形 $Z$
\src{263}
，具有如下性质：对于 $S$ 上的任意预概形 $S'$，
$M_n\otimes_{\mathcal O_S}\mathcal O_{S'}$ 对每个 $n\geq\nu$
均局部自由且秩为 $P(n)$，当且仅当 $S'\longrightarrow S$ 通过 $Z$
分解。

显然可以假设 $S$ 是仿射的，因而是诺特的。此时：

\result{引理}{3.5} 对每个整数 $N\geq\nu$，令 $U_N$ 为 $S$ 中由如下
$s\in S$ 构成的开集：
\[
\operatorname{rang}_{k(s)}
M_{ns}\otimes_{\mathcal O_{S,s}}k(s)\leq P(n)
\quad\text{对于每个 $\nu\leq n\leq N$}.
\]
那么开集的递减列 $U_N$ 最终稳定。

已知 [3]，IV，$S$ 有一个由约化子预概形 $S_i$ 构成的有限分割，使得
每个 $M\otimes_{\mathcal O_S}\mathcal O_{S_i}$ 在 $S_i$ 上平坦。
因此可以假设 $M$ 平坦，从而各 $M_n$ 也平坦。最后，显然还可以假设
$S$ 连通。于是存在一个整数 $n_0$ 和一个多项式 $Q$，使得
\[
\operatorname{rang}_{k(s)}
M_{ns}\otimes_{\mathcal O_{S,s}}k(s)=Q(n)
\quad\text{当 $n\geq n_0$ 时},
\]
见 [5]，III，\S~7。先假设 $P<Q$\footnote{印刷文本为
$P\neq Q$；1962年的勘误规定应为 $P<Q$，随后把“$P=Q$”的情形替换为
相反情形，即当 $n$ 足够大时 $P(n)\geq Q(n)$，并把 (*) 中的严格
不等式替换为非严格不等式。}，因而当 $n$ 足够大时
$P(n)<Q(n)$；于是当 $N$ 足够大时，显然有 $U_N=\varnothing$，
所以 $U_N$ 列最终稳定。在相反情形下，当 $n$ 足够大时有
$P(n)\geq Q(n)$，因而当 $N\geq n_0$ 时 $U_N=U_{n_0}$，所以
$U_N$ 列同样最终稳定。

特别地，由如下 $s\in S$ 构成的集合 $U_\infty$
\[
\tag{*}
\operatorname{rang}_{k(s)}
M_{ns}\otimes_{\mathcal O_{S,s}}k(s)\leq P(n)
\quad\text{对于每个 $n\geq\nu$}
\]
是各 $U_N$ 的交，因而是开集。为证明 (3.4)，于是可以用开集
$U_\infty$\footnote{印刷文本及法、英可编辑见证在这里作未定义的
$U$；本段唯一刚刚定义且由不等式 (*) 给出的开集是 $U_\infty$。}
代替 $S$，这就归约到不等式 (*) 对每个 $s\in S$ 均成立的情形。
此外：

\result{引理}{3.6} 设 $M$ 是局部诺特预概形 $S$ 上的一个模，$r$ 是
一个整数。那么存在 $S$ 的一个（显然唯一的）子预概形 $Z$，具有
如下性质：对于 $S$ 上的任意 $S'$，
$M\otimes_{\mathcal O_S}\mathcal O_{S'}$ 局部自由且秩为 $r$，当且
仅当 $S'\longrightarrow S$ 通过 $Z$ 分解。若
\[
\operatorname{rang}_{k(s)}
M_s\otimes_{\mathcal O_{S,s}}k(s)\leq r
\quad\text{对于每个 $s$},
\]
则 $Z$ 是 $S$ 的闭子预概形（这里假设 $M$ 凝聚）。
\src{264}
事实上，前面的论证把我们归约到不等式 (*) 对每个 $s\in S$ 均成立
的情形（必要时，可用由满足该不等式的各点 $s$ 构成的开部分代替
$S$）。于是可以假设 $M$ 出现在一个正合列中：
\[
\mathcal O_S^q\longrightarrow\mathcal O_S^r\longrightarrow M\longrightarrow0.
\]
对于 $S$ 上的 $S'$，所考察的条件也意味着：在相应的正合列
$\mathcal O_{S'}^q\longrightarrow\mathcal O_{S'}^r
\longrightarrow M'\longrightarrow0$ 中，第二个箭头是同构，亦即
第一个箭头为零。于是可见，由定义同态
$\mathcal O_S^q\longrightarrow\mathcal O_S^r$ 的矩阵各项所生成的
理想定义的 $S$ 的闭子预概形 $Z$，满足所要求的条件。

现在回到先前停下的 (3.4) 的证明。对于每个 $n$，令 $Z_n$ 为由
(3.6) 应用于模 $M_n$ 和整数 $r=P(n)$ 所得到的 $S$ 的闭子概形；
令 $Z'_N$ 为 $\nu\leq n\leq N$ 时各 $Z_n$ 的
$\operatorname{Inf}$。于是，各 $Z'_N$ 构成 $S$ 的闭子概形的递减列，
\footnote{印刷文本及法、英可编辑见证在此作“$Z$ 的闭子概形”；但
$Z$ 直到下一句才被定义为该列的稳定值，而各 $Z_n$ 都是 $S$ 的闭子概形。}
因而必然最终稳定。令 $Z$ 为 $N$ 足够大时各 $Z'_N$ 的恒定值。这就是
(3.4) 中所求的 $Z$。由此完成了 (3.4) 的证明，也完成了 (3.3) 的
证明。

因此，我们已经证明
$\underline{\operatorname{Quot}}^P_{\mathcal F/X/S}$ 可由一个
$S$-预概形 $Q$ 表示，而 $Q$ 是一列在 $S$ 上拟射影的开子预概形
$Q_\nu$ 的递增并。为继续推进，必须援用定理 2.1；由此容易推出 $Q$
是拟紧的（因为它是一个在 $S$ 上有限型、参数化具有 Hilbert 多项式
$P$ 的各 $\mathcal F_K$ 的商层族的预概形 $S'$ 的像）。因此，$Q$
等于某个 $Q_\nu$，所以在 $S$ 上拟射影。为了证明它在 $S$ 上射影，
只需再证明它在 $S$ 上固有；为此，只须援用 [5]，II，7.3.8
所给形式的固有性赋值判据。只需验证如下结论：

\result{引理}{3.7} 设 $S$ 是一个离散赋值环的谱，$s$ 是其一般点，
$X$ 是 $S$ 上的一个预概形，$\mathcal F$ 是 $X$ 上的一个拟凝聚模，
$\mathcal G_s$ 是 $X_s$ 上
$\mathcal F_s=\mathcal F\otimes_{\mathcal O_S}k(s)$ 的一个拟凝聚商模。
那么存在唯一的 $\mathcal F$ 的拟凝聚商模 $\mathcal G$，它在 $S$
上平坦，并且在 $X_s$ 上的限制是 $\mathcal G_s$。
\src{265}
事实上，若 $\mathcal G_s=\mathcal F_s/\mathcal K_s$，只须取
$\mathcal F$ 中诱导出 $\mathcal K_s$ 的最大子层 $\mathcal K$，
[5]，I，9.4.2，并置 $\mathcal G=\mathcal F/\mathcal K$。容易验证
该层满足要求。

至此，定理 3.2 以及随之而来的 (3.1) 均已完全证明。

该证明同时证明了如下结论：

\result{命题}{3.8} 在 (3.2) 的条件下，令
$Q=\operatorname{Quot}^P_{\mathcal F/X/S}$，$X_Q=X\times_SQ$，
$\mathcal F_Q=\mathcal F\otimes_{\mathcal O_S}\mathcal O_Q$；并令
$\mathcal G$ 为 $\mathcal F_Q$ 的凝聚商，它在 $Q$ 上平坦，具有相对
Hilbert 多项式 $P$，且 $(Q,\mathcal G)$ 表示函子
$\underline{\operatorname{Quot}}^P_{\mathcal F/X/S}$。那么存在一个
整数 $\nu$，使得当 $n\geq\nu$ 时，
$(f_Q)_*(\mathcal G(n))$ 是 $Q$ 上秩为 $P(n)$ 的局部自由模，并且
相对于 $S$ 甚丰沛；亦即，它定义了 $Q$ 到 $S$ 上 Grassmann 概形
$\operatorname{Grass}_{P(n)}(M)$ 的一个浸入。特别地，当
$n\geq\nu$ 时，$Q$ 上的层
\[
\bigwedge^{P(n)}(f_Q)_*(\mathcal G(n))
\]
是相对于 $S$ 甚丰沛的可逆层。

事实上，与 (3.2) 中一样，可以归约到 $\mathcal F$ 在 $S$ 上平坦的
情形；此时只须取一个使 $A_\nu=A$ 的整数 $\nu$（沿用上述记号）。

(3.2) 最重要的应用情形是 $\mathcal F=\mathcal O_X$。此时记
\[
\operatorname{Quot}_{\mathcal O_X/X/S}=\operatorname{Hilb}_{X/S},
\qquad
\operatorname{Quot}^P_{\mathcal O_X/X/S}=\operatorname{Hilb}^P_{X/S},
\]
从而得到分解
\[
\operatorname{Hilb}_{X/S}=\coprod_P\operatorname{Hilb}^P_{X/S}.
\]
依定义，$\operatorname{Hilb}_{X/S}$ 表示函子
$\underline{\operatorname{Hilb}}_{X/S}(S')$，后者是
$X'=X\times_SS'$ 中在 $S'$ 上平坦的闭子预概形之集合；而
$\operatorname{Hilb}^P_{X/S}$ 表示由具有给定 Hilbert 多项式 $P$
的闭子预概形构成的相应子函子。这些预概形也分别称为 $X$ 在 $S$
上的 \emph{Hilbert 预概形} 和 \emph{指标为 $P$ 的 Hilbert
预概形}。这一术语因 Hilbert 多项式在理论中所起的作用而得名。
它们与经典 Chow 簇（后者用于参数化循环，而不是簇）在性质上的
差别，正如一个簇的循环类 Chow 环与该簇的层类环（即
Riemann--Roch 定理 [1] 中出现的环）之间的差别；事实上，当
\src{266}
$X=\mathbf P^r_S$ 且 $S$ 是一个域的谱时，知道 $X$ 上凝聚模
$\mathcal F$ 的 Hilbert 多项式，也等价于知道 $\mathcal F$ 的
Chern 类，或知道 $\mathcal F$ 在 $X$ 上凝聚层类环中的类。

\result{注记}{3.9} 还应注意，
$\operatorname{Quot}_{\mathcal F/X/S}$ 和
$\operatorname{Quot}^P_{\mathcal F/X/S}$ 的构造已经归约到
$X=\mathbf P^r_S$、$\mathcal F=\mathcal O_X^N$ 的情形，其中
$\mathcal O_X(1)$ 是通常的甚丰沛层；更确切地说，一般的
$\operatorname{Quot}_{\mathcal F/X/S}$ 可实现为上述预概形的
闭子预概形。由于 $\operatorname{Quot}^P_{\mathcal F/X/S}$ 的形成
显然与基变换 $S'\longrightarrow S$ 相容，可见还可归约到
$S=\operatorname{Spec}(\mathbf Z)$ 的情形。归根结底，我们只须研究
下列 $\mathbf Z$ 上的射影概形：
\[
\mathfrak Q^P_{r,N}=
\operatorname{Quot}^{P}_{(\mathcal O_{\mathbf P^r_{\mathbf Z}})^N/
\mathbf P^r_{\mathbf Z}/\operatorname{Spec}(\mathbf Z)},
\]
尤其是绝对 Hilbert 概形：
\[
\operatorname{Hilb}^P_r=\mathfrak Q^P_{r,1}.
\]
对这些概形作更细致的研究会非常有意义；首先应确定其连通分支
（它们是否连通？）\footnote{1962年的勘误指出，Hartshorne 已证明
代数闭域上的 $\operatorname{Hilb}^P_r$ 是连通的，并确定了满足
$\operatorname{Hilb}^P_r\neq\varnothing$ 的各对 $(r,P)$。}及其
不可约分支（根据 Serre [7]，可能存在完全位于某个素数 $p\neq0$
之上的不可约分支）。回顾 Weil 的问题：在
$s\in\operatorname{Spec}(\mathbf Z)$ 上，$\operatorname{Hilb}^P_r$
各纤维的不可约分支是否对应于素域的“正则”扩张，亦即它们是否
“相对连通”。对于 Hilbert 概形，这些问题或许比对于“Chow 簇”
更容易处理。

\fsection{4}{变体}

\noindent a. 在 (3.1) 的条件下，设 $U$ 是 $X$ 中的开集，并令 $A'$
为 $A=\underline{\operatorname{Quot}}_{\mathcal F/X/S}$ 的如下
子函子：$A'(S')$ 是 $\mathcal F'$ 的、在 $S'$ 上平坦且支集包含于
$U'$ 的商模 $\mathcal G$ 之集合。立即可见，$A'$ 可由表示 $A$ 的
预概形 $\operatorname{Quot}_{\mathcal F/X/S}$ 的一个开部分表示。
由此可知，定理 (3.1) 和 (3.2)
\src{267}
在假设 $X$ 在 $S$ 上拟射影而非射影时仍然成立，只须在结论中也把
“射影”替换为“拟射影”，并且现在以
$\underline{\operatorname{Quot}}_{\mathcal F/X/S}(S')$ 表示
$\mathcal F'$ 的如下凝聚商 $\mathcal G$ 之集合：这些商在 $S'$ 上
平坦，且其支集在 $S'$ 上固有。

\medskip
\noindent b. 更一般地，可以对
$\mathcal F'=\mathcal F\otimes_{\mathcal O_S}\mathcal O_{S'}$
的、在 $S'$ 上平坦的商 $\mathcal G$ 施加各种在基变换下稳定的
附加自然条件，由此得到
$\underline{\operatorname{Quot}}_{\mathcal F/X/S}$ 的相应子函子，
并考察这些子函子的可表性。在许多情形下，通常的判据可以证明，
所得函子仍可由 $\operatorname{Quot}_{\mathcal F/X/S}$ 的开部分
表示。特别地，当施加下列附加性质之一时，情形正是如此：

\smallskip
\noindent $1^\circ$ 在各纤维 $X'_{s'}$ 上诱导的模
$\mathcal G_{s'}$（$s'\in S'$）之相伴素循环的维数，属于某个给定的
整数集合。

\smallskip
\noindent $2^\circ$（当 $\mathcal F=\mathcal O_X$ 时，$\mathcal G$
对应于 $X'$ 的一个闭子预概形 $Y$）：$Y$ 是 $S'$ 上的光滑预概形
（[3]，IV），或在 $S'$ 上正规（亦即各纤维 $Y_{s'}$ 在
“$k(s')$ 上”正规，也就是说在基域的任意扩张后均正规），或
（当 $X$ 在 $S$ 上平坦时）相对于 $S'$ 是 $X'$ 中的局部完全交
（亦即各纤维 $Y_{s'}$ 是 $X'_{s'}$ 中的局部完全交）。

其他条件可以涉及在 $X'_{s'}$ 上诱导的各模 $\mathcal G_{s'}$ 的
上同调性质，等等。当然，若若干条件分别由
$\operatorname{Quot}_{\mathcal F/X/S}$ 的开集 $U_i$ 表示，则它们的
合取由这些开集的交表示。例如，对 $S$ 上的每个 $S'$，考虑
$X'=X\times_SS'$ 的如下闭子预概形 $Y$ 的集合：$Y$ 是 $S'$ 的给定
秩 $r$ 的平展覆盖（[3]，I）。由此得到一个关于 $S'$ 的可表反变函子。

\medskip
\noindent c. [2]，II，C，n\textsuperscript{o} 2 中定义的预概形
$\operatorname{Hom}_S(X,Y)$、$\Pi_{X/S}(Z/X)$、
$\operatorname{Isom}_S(X,Y)$ 在适当的射影性假设下存在，并可实现为
适当 Hilbert 预概形中的开部分。由于
\[
\operatorname{Hom}_S(X,Y)=\Pi_{X/S}((X\times_SY)/X),
\]
$\operatorname{Hom}_S(X,Y)$ 的情形可归约到
$\Pi_{X/S}(Z/X)$ 的情形。于是注意到，对于 $S$ 上的每个 $S'$，
$Z'=Z\times_SS'$ 在 $X'=X\times_SS'$ 上的截面集合，与 $Z'$ 的如下
子预概形 $\Gamma$ 的集合一一对应：由 $Z'\longrightarrow X'$ 诱导的
态射 $\Gamma\longrightarrow X'$ 是同构（若 $Z$ 在 $X$ 上分离，
则这些 $\Gamma$ 必然是闭的）。由此，

\src{268}
当 $X$ 在 $S$ 上平坦且固有，而 $Z$ 在 $S$ 上拟射影时，
$\Pi_{X/S}(Z/X)$ 存在，并实现为
$\operatorname{Hilb}_{Z/S}$ 的一个开子预概形。因此，当 $X$ 在 $S$ 上
射影且平坦，而 $Y$ 在 $S$ 上拟射影时，$\operatorname{Hom}_S(X,Y)$
存在，并实现为
$\operatorname{Hilb}_{(X\times_SY)/S}$ 的一个开子预概形。当 $X$ 和
$Y$ 都在 $S$ 上射影时，立即可知 $\operatorname{Isom}_S(X,Y)$ 也存在，
并由 $\operatorname{Hom}_S(X,Y)$ 的一个开部分表示。同样，当 $X$ 在
$S$ 上平坦且射影，而 $Y$ 在 $S$ 上拟射影时，与
$\operatorname{Hom}_S(X,Y)$ 所表示函子的如下子函子相对应的
$S$-预概形 $\operatorname{Imm}_S(X/Y)$ 也由
$\operatorname{Hom}_S(X,Y)$ 的一个开部分表示：该子函子的元素是作为
浸入的 $S'$-同态 $X'\longrightarrow Y'$。

设 $\mathcal L$（相应地，$\mathcal M$）是 $X$（相应地，$Y$）上相对于
$S$ 甚丰沛的可逆层，于是
$\mathcal L\otimes_{\mathcal O_S}\mathcal M$ 是 $X\times_SY$ 上相对于
$S$ 甚丰沛的层。因此，对每个有理系数多项式 $P$，
$\operatorname{Hilb}^P_{(X\times_SY)/S}$ 均有定义，并且是 $S$ 上的
拟射影预概形。它因而在 $\operatorname{Hom}_S(X,Y)$ 上诱导出一个
在 $S$ 上拟射影的既开又闭部分，记作
$\operatorname{Hom}_S(X,Y)^P$。所以
$\operatorname{Hom}_S(X,Y)^P$ 在 $S$ 上的截面，正是满足下述条件的
$S$-态射 $g:X\longrightarrow Y$：对每个整数 $n$，都有
\[
\chi\!\left(\left(\mathcal L\otimes_{\mathcal O_X}g^*(\mathcal M)\right)^{
\otimes n}\right)=P(n).
\]

由此可得到 Matsusaka 定理的若干推广；该定理断言，一个“极化”射影簇
的自同构组成一个代数群。这里这一断言显然具有更精确的含义，因为我们
已将这个群定义为某个泛问题的解。还应注意，在代数闭域上，以往所考察
的自同构群，是由这里定义的“真正的”自同构群模去幂零元而得到的。这也
说明，较早的构造很难在非完美基域上成立：基域扩张后出现的幂零元理想
未必“在 $k$ 上有定义”。同一说明也适用于大多数旧式构造。

\fsection{5}{Hilbert 概形的微分研究}

它由下述结果推出：

\src{269}
\result{命题}{5.1} 设 $S$ 是预概形，$S_0$ 是由平方为零的拟凝聚理想
$\mathcal J$ 定义的子预概形，$X$ 是 $S$-预概形，$\mathcal F$ 是
$X$ 上的拟凝聚模，并令
$X_0=X\times_SS_0$ 及
$\mathcal F_0=\mathcal F\otimes_{\mathcal O_S}\mathcal O_{S_0}$；
最后，设
$\mathcal G_0=\mathcal F_0/\mathcal K_0$ 是 $\mathcal F_0$ 的一个
在 $S_0$ 上平坦的拟凝聚商模。对 $X$ 的每个开集 $U$，令
$\mathcal E(U)$ 为 $\mathcal F|U$ 的所有拟凝聚商模 $\mathcal G$
所成的集合，其中 $\mathcal G$ 在 $S$ 上平坦，并且满足
$\mathcal G\otimes_{\mathcal O_S}\mathcal O_{S_0}=\mathcal G_0$；
于是，当 $U$ 变化时，各 $\mathcal E(U)$ 是 $X$ 上某个层
$\mathcal E$ 的截面。\footnote{印刷文本及法、英可编辑见证写成层
“在 $U$ 上”；由于 $U$ 变化而 $\mathcal E(U)$ 构成开集上的截面函子，
该层的底空间是固定的 $X$。}在此设定下，群层
\[
\mathcal A=\operatorname{Hom}_{\mathcal O_{X_0}}
(\mathcal K_0,\mathcal G_0\otimes_{\mathcal O_{S_0}}\mathcal J)
\]
自然作用于 $\mathcal E$，从而使 $\mathcal E$ 成为一个“在
$\mathcal A$ 作用下形式主齐性的层”（即对 $X$ 中的每个开集 $U$，
$\mathcal E(U)$ 或为空，或为在 $\mathcal A(U)$ 作用下的主齐性集）。

由此得到：

\result{推论}{5.2} 假设在 $X$ 上局部存在 $\mathcal G_0$ 的一个
延拓 $\mathcal G$，它是 $\mathcal F$ 的一个在 $S$ 上平坦的商
（即层 $\mathcal E$ 的各纤维均非空）。那么存在一个典范障碍类
\[
c(\mathcal G_0)\in H^1(X,\mathcal A),
\]
它的消失是存在 $\mathcal G_0$ 的全局延拓 $\mathcal G$、使之成为 $\mathcal F$ 的一个
在 $S$ 上平坦的商的充分必要条件。若此类为零，则所有可能延拓组成的
集合 $\mathcal E(X)$ 是在
\[
\mathcal A(X)=\operatorname{Hom}_{\mathcal O_{X_0}}
(\mathcal K_0,\mathcal G_0\otimes_{\mathcal O_{S_0}}\mathcal J)
\]
作用下的主齐性集。
\footnote{在最后这个公式中，印刷文本写的是 $\mathcal O_X$，而不是
紧接在前面的 $\mathcal A$ 定义中已经确定的 $\mathcal O_{X_0}$。}
特别地，若 $H^1(X,\mathcal A)=0$，则全局延拓必定存在。

\result{推论}{5.3}\footnote{印刷文本将此推论编号为“3.3”；但它在
第 5 节中接在 (5.1) 和 (5.2) 之后。} 假设
$Q=\operatorname{Quot}_{\mathcal F/X/S}$ 存在（参见 4 (a)）——例如，
$X$ 在局部诺特的 $S$ 上拟射影，且 $\mathcal F$ 凝聚。设
$x\in Q$ 对应于 $k(s)$（$s\in S$）的剩余域扩张 $K=k(x)$：
\src{270}
于是，$x$ 由 $K$-预概形 $X_K$ 上的模
$\mathcal F_0=\mathcal F_K$ 的一个凝聚商模
$\mathcal G_0=\mathcal F_0/\mathcal K_0$ 定义。设 $\mathcal A$ 是
$X_K$ 上由
\[
\mathcal A=\operatorname{Hom}_{\mathcal O_{X_0}}(\mathcal K_0,\mathcal G_0)
\]
定义的凝聚层。那么，纤维 $Q_s$ 在点 $x$ 处的 Zariski 切空间
（即 $\mathfrak m/\mathfrak m^2$ 关于 $K$ 的对偶，其中
$\mathfrak m$ 是 $\mathcal O_{Q_K,x}$ 的极大理想）典范同构于
$H^0(X_K,\mathcal A)$。

给出 Zariski 切空间的这一结果还可推广。对于给定的 $S$-态射
$g:S'\longrightarrow Q$，即 $Q'=Q\times_SS'$ 在 $S'$ 上的一个截面
$g'$，模
\[
\Omega=g^*(\Omega^1_{Q/S})=g'^*(\mathcal J/\mathcal J^2)
\]
（其中 $\mathcal J$ 是由 $Q'$ 在 $S'$ 上的截面 $g'$ 所定义的
$Q'$ 上的理想）可由下述公式刻画；该公式关于 $S'$ 上的凝聚模
$\mathcal M$ 具有函子性：
\[
\operatorname{Hom}_{\mathcal O_{S'}}(\Omega,\mathcal M)
\simeq H^0(X',\mathcal A),
\]
其中 $\mathcal A$ 仍是 $X'=X\times_SS'$ 上由
\[
\mathcal A=\operatorname{Hom}_{\mathcal O_{X'}}
(\mathcal K,\mathcal G\otimes_{\mathcal O_{S'}}\mathcal M)
\]
定义的模。
\footnote{印刷文本在这个张量积中写的是 $\mathcal O_S$；但这里
$\mathcal G$ 和 $\mathcal M$ 都是 $S'$ 上的模，而平方零形变的构造
要求在 $\mathcal O_{S'}$ 上取张量积。}
这里，$\mathcal G=\mathcal F'/\mathcal K$ 是与 $g$ 对应的
$\mathcal F'=\mathcal F\otimes_{\mathcal O_S}\mathcal O_{S'}$ 的商模。
事实上，只须在 (5.1) 中以 $S'$ 代替 $S_0$，并以预概形
$D(\mathcal M)=(S',\mathcal O_{S'}+\mathcal M)$ 代替 $S$，其中
$\mathcal M$ 被视为平方为零的理想。

在 (5.1) 中取 $\mathcal F=\mathcal O_X$ 时，给定 $\mathcal G_0$
就等价于给定 $X_0$ 的一个在 $S_0$ 上平坦的闭子预概形 $Y_0$，
其定义理想为 $\mathcal J_0=\mathcal K_0$。此时 (*) 给出
\[
\mathcal A=\operatorname{Hom}_{\mathcal O_{X_0}}
(\mathcal J_0/\mathcal J_0^2,
\mathcal O_{Y_0}\otimes_{\mathcal O_{S_0}}\mathcal J),
\]
其中 $\mathcal J_0/\mathcal J_0^2$ 被解释为 $Y_0$ 在 $X_0$ 中的余法层，
\footnote{印刷文本及法、英可编辑见证在说明句中漏去了下标 $0$；
紧接在前面的公式及 $Y_0\subset X_0$ 的定义表明余法层是
$\mathcal J_0/\mathcal J_0^2$，而 $\mathcal J$ 是底的平方零理想。}
也记作
\src{271}
$\check{\mathcal N}_{Y_0/X_0}$；此时，将 $\mathcal A$ 看作 $Y_0$
上的模，并在 $Y_0$ 上计算 $H^0$ 和 $H^1$，会更为方便。此外，若
$Y_0$ 在 $X_0$ 中局部为完全交，且 $X$ 在 $S$ 上平坦，则 (5.1)
中的局部延拓必定存在；另一方面，$\mathcal J_0/\mathcal J_0^2$ 在
$Y_0$ 上局部自由，并可写成
\[
\mathcal A=\mathcal N_{Y_0/X_0}\otimes_{\mathcal O_{S_0}}\mathcal J,
\]
其中右端第一个因子是 $Y_0$ 在 $X_0$ 中的法层。利用光滑性的基本
判别准则（[3]，III，3.1），例如可得：

\result{推论}{5.4} 在 (5.3) 的条件下，假设
$\mathcal F=\mathcal O_X$，$X$ 在 $S$ 上平坦，并且与
$\mathcal G_0$ 对应的 $X_0$ 的闭子预概形 $Y_0$ 在局部为完全交。
那么，$Q_s$ 在 $x$ 处的 Zariski 切空间典范同构于
$H^0(Y_0,\mathcal N_{Y_0/X_0})$。若
$H^1(Y_0,\mathcal N_{Y_0/X_0})=0$，则 Hilbert 预概形
$\operatorname{Hilb}_{X/S}$ 在点 $x$ 处相对于 $S$ 光滑。
（$\mathcal N_{Y_0/X_0}$ 是 $Y_0$ 在 $X_0$ 中的法层。）

\result{注}{5.5} 这一陈述特别适用于如下情形：$Y_0$ 是 $X_0$ 中由
一个方程定义的完全交，即一个有效“Cartier 除子”。这时，
$\mathcal N_{Y_0/X_0}$ 同构于 $Y_0$ 上由 $X_0$ 上的可逆层
$\mathcal J^{-1}$ 所诱导的层，其中该可逆层由除子 $Y_0$ 定义。
这种情形尤其出现在对非奇异射影簇 $X_0$ 上有效除子族的研究中。
$Q$ 在点 $x$ 处的 Zariski 切空间（或者说，$Q$ 中与除子对应的开集
$D$ 在 $x$ 处的 Zariski 切空间）与
$H^0(Y_0,\mathcal N_{Y_0/X_0})$ 之间的同构，在经典代数几何中称为
由前者到后者的“特征同态”。它仅在 $x$ 是某个除子“完全连续族”的
参数簇 $T$ 的光滑点时才有定义；从我们的观点看，$T$ 就是概形
$D$ 的一个不可约分支，并赋予诱导的约化结构。$T$ 在 $x$ 处的切空间
于是是 $D$ 在 $x$ 处切空间的子空间，所以古典理论中的特征同态确实
是单射，但只有在附加条件下才是满射，例如 $D$ 在 $x$ 处为整时。
事实上，Zappa [8] 构造了一个例子
\src{272}
（其中 $X$ 是复数域上的非奇异射影曲面），在这个例子中，即使在
$T$ 的一般点处，特征同态也不是满射。因此，$D$ 即使在所考察的
不可约分支的一般点处也不是整的。这格外鲜明地表明，要理解最经典的
曲面理论中的现象，带幂零元的簇是不可或缺的。
\footnote{1962 年勘误补充说，Mumford 当时刚刚构造出
$\mathbf P^3_{\mathbf C}$ 的 Hilbert 概形的一个不可约分支，其一般点
表示次数为 14、亏格为 24 的非奇异曲线，而该分支在其一般点处
非约化；沿其中一条曲线作爆破，还可得到一个三维正则射影概形，其
形式模概形在一般点处非约化。}

\subsection*{5.6}

我们在 (5.4) 中给出了一个光滑性判别准则，它特别适用于除子概形。
Kodaira 在 [6] 中给出了另一个判别准则，即
$H^1(X_0,\mathcal L)$ 的消失，其中
$\mathcal L=\mathcal J_0^{-1}$ 是 $X_0$ 上由除子 $Y_0$ 定义的可逆层；
该判别准则在 $S$ 是特征 0 域的谱时成立，而在基域为 $\mathbf C$
的情形下，[6] 通过超越方法证明了它。这里还应指出，更一般地，再次令
$S$ 任意，则 Kodaira 条件是使典范态射
\[
D\longrightarrow\operatorname{Pic}_{X/S}
\]
从除子预概形到 $X/S$ 的 Picard 预概形在所考察点 $x$ 处光滑的
充分条件（一旦知道 $\operatorname{Pic}_{X/S}$ 存在，这一点很容易
由通常的光滑性判别准则验证）。因此，若此外
$\operatorname{Pic}_{X/S}$ 在 $x$ 的像点处相对于 $S$ 光滑
（例如，若 $\operatorname{Pic}_{X/S}$ 在 $S$ 上光滑），则 $D$
在 $x$ 处相对于 $S$ 光滑。另一方面，Cartier 已证明，特征 0 域
$k$ 上每个局部有限型的群预概形都在 $k$ 上光滑。结合这两个结果，
即可重新得到 Kodaira 的结果。还应注意，这些说明表明，在特征
$p>0$ 的域 $K$ 上，当 $\operatorname{Pic}_{X/S}$ 在 $K$ 上不光滑
\footnote{印刷文本及法、英可编辑见证在这里由已经指定的域 $K$
突然改作 $k$；本句余下部分仍以 $K$ 为基域，故恢复为 $K$。}
时（$X$ 为 Igusa 曲面时即属此情形），条件
$H^1(X_0,\mathcal L)=0$ 反而蕴含 $D$ 在 $x$ 处不光滑；若 $K$
代数闭，则它在 $x$ 处甚至不是约化的。

最后再指出下述结果；它在纤维空间的微分研究中起着重要作用：

\result{命题}{5.7} 设 $X$ 是局部诺特预概形 $S$ 上有限平坦的
预概形，$Z$ 是 $S$ 上的预概形，并且 $\Pi_{X/S}(Z/X)$ 存在
（当 $Z$ 在 $X$ 上拟射影时即如此）。若 $Z$ 在 $X$ 上光滑，则
$\Pi_{X/S}(Z/X)$ 在 $S$ 上光滑。

这是定义及通常的光滑性判别准则（[3]，III，3.1）的直接推论。注意，
当 $X$ 在 $S$ 上有限平坦时，$\Pi_{X/S}(Z/X)$ 的存在性问题
\footnote{印刷文本在这里写的是 $Z/K$；但文中没有定义 $K$，而且从
(5.7) 开始所考察的截面函子是 $\Pi_{X/S}(Z/X)$。}
\src{273}
可用非常初等的方法处理，无须使用 Hilbert 概形理论。例如，若 $X$
在 $S$ 上纯不可分，则 $\Pi_{X/S}(Z/X)$ 无论对怎样的 $Z$ 都存在。
再如，设 $T$ 是一个 $S$-预概形，并令 $T_n$ 为 $T\times_ST$ 的对角
在 $T\times_ST$ 中的“$n$ 阶无穷小邻域”，配备由两个投影诱导的态射
$p_1,p_2:T_n\longrightarrow T$。通过 $p_1$ 将 $T_n$ 看作 $T$
上的有限预概形，并进一步假设它在 $T$ 上平坦（当 $T$ 在 $S$ 上
光滑时即如此）。对每个 $T$ 上的预概形 $X$，令
\[
(X/T/S)^{(n)}=
\Pi_{T_n/T}\bigl(p_2^*(X/T)/T_n\bigr),
\]
\footnote{印刷文本在符号 $\Pi$ 下写的是 $T_n/S$。这里将命题 5.7
应用于有限平坦态射 $p_1:T_n\to T$，并且所得对象明确是 $T$ 上的
预概形；故分母应为 $T_n/T$。}
这是 $T$ 上的一个预概形，称为 \emph{$X$ 在 $T$ 上的 $n$ 阶截面芽丛}
（相对于 $S$）。它函子性地依赖于 $X$；若 $X$ 在 $T$ 上光滑，
它也在 $T$ 上光滑。

\fsection{6}{与范数概念及对称积的关系}

设 $S$ 是预概形，$X$ 和 $Y$ 是 $S$-预概形，并设
\[
u:(X/S)^n\longrightarrow Y
\]
是从 $X/S$ 的 $n$ 次笛卡儿幂到 $Y$ 的对称 $S$-态射。为简单起见，
假设 $S$ 局部诺特，并且 $X$ 和 $Y$ 在 $S$ 上有限型。于是，对于
$X$ 上每个满足下列条件的凝聚模 $\mathcal F$：其支集在 $S$ 上有限，
它在 $S$ 上平坦，并且在 $S$ 上的相对秩等于 $n$（即
$f_*(\mathcal F)$ 是 $S$ 上秩为 $n$ 的局部自由模），可以自然地对应
$Y$ 在 $S$ 上的一个截面：
\[
\pi^u_{X/S}(\mathcal F)\in\Gamma(Y/S).
\]
这里不写出正式定义，只指出由此得到的形式体系是通常范数与迹的形式
体系的一个自然推广。当 $X$ 在 $S$ 上的 $n$ 次对称幂存在时
（例如，当对称群 $\mathfrak S_n$ 在 $(X/S)^n$ 上作用的各轨道都包含
于仿射开集中时），可取这个对称幂
$\operatorname{Symm}^n_S(X)$ 作为 $Y$，并得到一个典范元素
\[
\pi_{X/S}(\mathcal F)\in\Gamma(\operatorname{Symm}^n_S(X)/S),
\]
\src{274}
由它可以恢复各 $\pi^u_{X/S}(\mathcal F)$。另一重要情形是：$X$
是 $S$ 上的交换幺半群，$X=Y$，而态射 $u$ 来自 $X$ 的合成律。
这时，将与 $X$ 上的模 $\mathcal F$ 相对应的 $X$ 在 $S$ 上的截面
简记为 $\pi(\mathcal F)$。

现在假设 $X$ 上有一个凝聚模 $\mathcal F$，使得
$\operatorname{Quot}_{\mathcal F/X/S}$ 存在；或者至少，函子
$\underline{\operatorname{Quot}}^n_{\mathcal F/X/S}$ 可由某个
$S$-预概形 $\operatorname{Quot}^n_{\mathcal F/X/S}$ 表示。这里该
函子把 $S$ 上的每个 $S'$ 映为下述凝聚商层 $\mathcal M$ 所成的集合：
$\mathcal M$ 是
$\mathcal F'=\mathcal F\otimes_{\mathcal O_S}\mathcal O_{S'}$ 的商，
在 $S'$ 上平坦，并具有相对秩 $n$。（当 $X$ 在 $S$ 上拟射影时，
$\operatorname{Quot}^n_{\mathcal F/X/S}$ 的确存在；采用第 3 节的
记号，它就是 $\operatorname{Quot}^P_{\mathcal F/X/S}$，其中 $P$
是恒等于 $n$ 的常数多项式。）由于
$\pi^u_{X/S}(\mathcal M)$ 的构造与基变换相容，因而得到一个典范态射
\[
\pi^u_{X/S}:\operatorname{Quot}^n_{\mathcal F/X/S}\longrightarrow Y,
\]
特别地，若 $X$ 在 $S$ 上的 $n$ 次对称幂存在，则得到
\[
\pi_{X/S}:\operatorname{Quot}^n_{\mathcal F/X/S}
\longrightarrow\operatorname{Symm}^n_S(X).
\]
最重要的情形是 $\mathcal F=\mathcal O_X$，它给出态射
\[
\pi_{X/S}:\operatorname{Hilb}^n_{X/S}
\longrightarrow\operatorname{Symm}^n_S(X).
\]
当 $n=0$ 或 $n=1$ 时，这显然是同构。但当
$n>1$\footnote{印刷文本再次写成 $n\geq1$，这与刚刚断言的
$n=1$ 时为同构相矛盾。} 时，即使 $S$ 是域 $k$ 的谱，并且 $X$
在 $S$ 上光滑，它一般也不是同构，甚至不是单射态射；这是因为，
$X$ 的一个零维子概形（例如，它对应于局部环 $\mathcal O_{X,x}$
中的一个以该极大理想为根的准素理想 $\mathfrak i$，其中 $x$ 是 $X$
的闭点）并不能由它所定义的循环确定（在这里，即不能由
$\mathfrak i$ 在 $\mathcal O_{X,x}$ 中关于 $k$ 的余维确定）。
当 $S$ 再次任意时，只能作如下断言：

\medskip
\noindent a. 若 $X$ 在 $S$ 上光滑，则范数态射在
$\operatorname{Hilb}^n_{X/S}$ 的如下开部分与
$\operatorname{Symm}^n_S(X)$ 的如下开部分之间诱导一个同构：前者
分类包含于 $X$ 中的秩为 $n$ 的平展覆盖（参见
n\textsuperscript{o} 4，(b)），后者对应于次数为 $n$ 且无重分量的循环。
\src{275}

\medskip
\noindent b. 若此外 $X$ 在 $S$ 上的相对维数为 1，则范数态射甚至
给出
$\operatorname{Hilb}^n_{X/S}$ 与 $\operatorname{Symm}^n_{X/S}$
之间的同构。

这是因为，非奇异代数曲线的零维子概形可由相应的循环确定。同一说明
更一般地适用于非奇异代数概形上的有效 Cartier 除子（而且不能排除：
在这一极特殊的情形下，Chow 簇给出的对象与 Hilbert 概形相同）。

\fsection{7}{补充与问题}

正如 J.-P. Serre 所指出的，Nagata 的一个著名例子表明，可以找到
如下对象：一个概形 $S$，它是域 $k$ 的谱；一个 $S$-概形 $S'$，它是
$k$ 的二次扩张 $k'$ 的谱；以及一个维数为 3 的固有光滑但非射影的
$S'$-概形 $X$，使得 $\Pi_{S'/S}(X/S')$ 不存在。由此更可知，Hilbert
概形 $\operatorname{Hilb}^2_{X/S}$ 不存在（甚至连表示包含于 $X$
中的 $S$ 的秩 2 平展覆盖的 $k$-概形也不存在，从而 $X$ 的对称平方
当然更不存在；参见上一节）。这给代数几何中进行非射影构造的可能性
施加了严重限制。（不过，可以设想在解析几何中不会出现这种限制，
正如在形式几何中不会出现一样（参见 [2]，II）。）另一方面，若 $X$
是域 $k$ 的谱 $S$ 上的固有概形，而 $Z$ 在 $X$ 上拟射影，则
$\Pi_{X/S}(Z/X)$ 存在，并且是一个由一列 $S$ 上拟射影概形的不交并
组成的概形（如射影情形 (3.1)）。事实上，可以用一个射影
$S$-概形 $X'$ 支配 $X$，从而归约到 $X$ 本身射影的情形。这里不写出
证明细节；证明还使用 [2]，I，A，n\textsuperscript{o} 2 (b) 中所述
有限态射的分解结果。该方法之所以奏效，是因为当 $S$ 是域的谱时，
Chow 引理中出现的 $X'$ 自动在 $S$ 上平坦。

我不知道在不对 $S$ 作任何假设、而只假设 $X$ 在 $S$ 上固有且平坦、
$Z$ 在 $X$ 上拟射影时，这一结果是否仍然成立。应用中的一个重要情形
是 $Z$ 为 $X$ 的闭子概形；这时，若 $\Pi_{X/S}(Z/X)$ 存在，它必为
$S$ 的闭子概形。当 $X$ 在 $S$ 上射影时，可以
\src{276}
直接而且相当简单地构造它，无须使用 Hilbert 概形理论。所用方法还更
一般地表明：若 $Z$ 在 $X$ 上仿射，则 $\Pi_{X/S}(Z/X)$ 存在，并在
$S$ 上仿射。它还表明：若 $X$ 在 $S$ 上固有且平坦（不必在 $S$ 上
射影），则对 $X$ 上每个局部平凡向量丛 $Z$，
$\Pi_{X/S}(Z/X)$ 都存在，并且是 $S$ 上的向量丛。最好能重新研究并
统一这些结果。

\section*{参考文献}
\addcontentsline{toc}{section}{参考文献}

\begin{description}
\item[{[1]}] BOREL (A.) 和 SERRE (J.-P.). -- Le théorème de Riemann--Roch,
\emph{Bull. Soc. math. France}, 第 86 卷，1958 年，第 97--136 页。

\item[{[2]}] GROTHENDIECK (Alexander). -- Technique de construction et
théorèmes d'existence en géométrie algébrique, I, II, III, \emph{Séminaire
Bourbaki}, 第 13 卷，1959/60 年，n\textsuperscript{o} 190，29 页，
n\textsuperscript{o} 195，22 页；以及第 14 卷，1960/61 年，
n\textsuperscript{o} 212，20 页。

\item[{[3]}] GROTHENDIECK (Alexander). -- Séminaire de Géométrie algébrique,
I, II, III, IV. -- 巴黎，Institut des hautes Études scientifiques，1960/61
年（油印本）。

\item[{[4]}] GROTHENDIECK (Alexander). -- Technique de construction en
géométrie analytique, IV, V, \emph{Séminaire Cartan}, 第 13 卷，1960/61 年，
n\textsuperscript{o} 11 和 12。

\item[{[5]}] GROTHENDIECK (A.) 和 DIEUDONNÉ (J.). -- \emph{Éléments de
géométrie algébrique}, I : Le langage des schémas. -- 巴黎，Presses
universitaires de France，1960 年（Institut des hautes Études scientifiques,
Publications mathématiques, 4）；II，III，IV（即将出版）。

\item[{[6]}] KODAIRA (K.). -- Characteristic linear systems of complete
continuous systems, \emph{Amer. J. of Math.}, 第 78 卷，1956 年，
第 716--744 页。

\item[{[7]}] SERRE (Jean-Pierre). -- Exemples de variétés projectives en
caractéristique $p$ non relevables en caractéristique 0, \emph{Proc. Nat. Acad.
Sc. U.S.A.}, 第 47 卷，1961 年，第 108--109 页。

\item[{[8]}] ZAPPA (G.). -- Sull'esistenza, sopra le superficie algebriche, di
sistemi continui completi infiniti, la cui curva è a serie caratteristica
incompleta, \emph{Pont. Acad. Sc. Acta}, 第 9 卷，1945 年，第 91--93 页。
\end{description}

\section*{补遗}
\addcontentsline{toc}{section}{补遗}
\begin{center}
{\small [校样修订时加入]}
\end{center}

现在看来，[2]，III，n\textsuperscript{o} 8 中的猜想是错误的；即使
对于特征 0 域上的非奇异簇，无论就商的存在性还是拟射影性而言，这些
猜想都不成立，甚至在 $G$ 的作用具有闭图像时也是如此。

% END 221.tex

% BEGIN 232.tex
\unit{第232次报告——Picard 概形}
\sid{232}
\src{143}
\seminarzh{第14年度，1961/62，第232号\\（1962年2月）}
\paperhead{下降技术与代数几何中的存在性定理\\
V. Picard 概形：存在性定理}%
  {亚历山大·格罗滕迪克 著}

\fsection{1}{相对 Picard 群与函子}

对任意预概形（更一般地，任意赋环空间）$X$，我们把 $X$ 上可逆模
（即局部同构于 $\mathcal O_X$ 的模）的同构类所成的群称为 $X$ 的
（绝对）\emph{Picard 群}，并记作 $\operatorname{Pic}(X)$。因此有典范
同构
\begin{equation}
\operatorname{Pic}(X)\xrightarrow{\sim}H^1(X,\mathcal O_X^*),
\tag{1.1}
\end{equation}
其中 $\mathcal O_X^*$ 表示 $\mathcal O_X$ 的单位元层（它确实可与标准
可逆模 $\mathcal O_X$ 的自同构层等同）。注意，
$X\longmapsto\operatorname{Pic}(X)$ 显然是关于 $X$ 的反变函子，并且
同构 (1.1) 是函子性的。

若 $X$ 是另一个预概形 $S$ 上的预概形，则当 $S'$ 在 $S$ 上预概形范畴
$(\mathrm{Sch})/S$ 中变化时，由上可得反变函子
$S'\longmapsto\operatorname{Pic}(X\times_S S')$。这个函子不可能“可表”
（[4]，II，A），因为所要分类的可逆模具有非平凡自同构，所以该函子
不具有“局部性质”（[5]，IV，5.4）。因此应将它“局部化”；一般地，对
任意相对预概形 $X/S$，引入一个相对性质的群
\begin{equation}
\operatorname{Pic}'(X/S)=H^0\bigl(S,R^1f_*(\mathcal O_X^*)\bigr),
\tag{1.2}
\end{equation}
其中 $f:X\longrightarrow S$ 是结构态射（比较 [4]，II，C 3）。在前引文献中，
这个群称为相对 Picard 群；由于即将说明的理由，这里最好把它称为
$X/S$ 的\emph{限制相对 Picard 群}。当 $S'$ 在 $(\mathrm{Sch})/S$
中变化时，
$S'\longmapsto\operatorname{Pic}'(X\times_S S'/S')$ 是关于 $S'$ 的
反变函子，也记作 $\operatorname{Pic}'_{X/S}$，因而基本上由下式定义：
\begin{equation}
\operatorname{Pic}'_{X/S}(S')=
\operatorname{Pic}'(X\times_S S'/S').
\tag{1.3}
\end{equation}

\src{144}
由于我们已经作了所需的处理，这个函子现在具有“局部性质”。直观地说，
(1.3) 右端可解释为 $X/S$（的各纤维）上可逆层类的“代数族”所成的
集合，这些代数族由参数预概形 $S'/S$ 标号。当函子
$\operatorname{Pic}'$ 可表时，表示它的 $S$ 上预概形记作
$\operatorname{Pic}'_{X/S}$，称为 $X$ 在 $S$ 上的
\emph{Picard 预概形}；于是有
\begin{equation}
\operatorname{Hom}_S(S',\operatorname{Pic}'_{X/S})
\simeq\operatorname{Pic}'_{X/S}(S')
=\operatorname{Pic}'(X\times_S S'/S').
\tag{1.4}
\end{equation}

然而，在一些重要情形中，$\operatorname{Pic}'_{X/S}$ 不可表（例如域
$k$ 上没有 $k$-有理点的“Brauer--Severi”簇），但相对 Picard 预概形
仍有一个自然定义。原因在于，从绝对 Picard 群
$\operatorname{Pic}(X\times_S S'/S')$ 出发定义函子
$\operatorname{Pic}'$ 时，局部化尚不充分；更确切地说，
$\operatorname{Pic}'$ 一般并不“与忠实平坦下降相容”。下面明确说明。

令 $\mathfrak M$ 为所有\emph{忠实平坦且拟紧}的预概形态射所成的类；
这个类在基变换和复合下稳定。设 $P$ 是从 $(\mathrm{Sch})/S$ 到集合范畴
的反变函子；对每个满足 $u\in\mathfrak M$ 的 $S$-态射
$u:T'\longrightarrow T$，考虑图式
\begin{equation}
P(T)\longrightarrow P(T')\rightrightarrows P(T'\times_TT')
\tag{1.5}
\end{equation}
它是由 $P$ 作用于图式
\[
T\longleftarrow T'\mathrel{\substack{\xleftarrow{\,\mathrm{pr}_1\,}\\[-0.35ex]
\xleftarrow[\,\mathrm{pr}_2\,]{}}}T'\times_TT'.
\]
所得。若 $P$ 可表，则由下降理论（[4]，I，B，定理 2），对每个
$u\in\mathfrak M$，图式 (1.5) 都正合。我们把这一事实表述为 $P$ 与
$\mathfrak M$ 相容；在这里，即称 $P$“与忠实平坦下降相容”，也可说
$(\mathrm{Sch})/S$ 上的“预层”$P$ 对由 $\mathfrak M$ 给出的局部化概念
而言是一个“层”。当 $P$ 任意时，一个在通常拓扑局部化情形中熟知的
标准过程，可以给它配上一个“层”$\mathcal P$ 和一个函子同态
$P\longrightarrow\mathcal P$，并使该同态在显然意义下具有泛性。
$\mathcal P$ 的计算可明确写成如下形式：为定义 $\mathcal P(T)$，对每个
$T$ 上满足 $u:T'\longrightarrow T$ 属于 $\mathfrak M$ 的 $T'$，记

\src{145}
$\overline H^0(T'/T,P)$ 为 $P(T')$ 的如下子集：其元素 $\xi$ 在
$P(T'\times_TT')$ 中的两个像 $\xi_1,\xi_2$ 满足，存在一个
$v\in\mathfrak M$ 的态射 $v:T''\longrightarrow T'\times_TT'$，使
$\xi_1$ 与 $\xi_2$ 在 $P(T'')$ 中的像相同。[注意——如此定义的集合
$\overline H^0$ 因而大于 [4]，I，A，4(a) 中引入的集合
$H^0(T'/T,P)$。] 当 $T'$ 在固定的 $T$ 上变化时（始终要求
$u\in\mathfrak M$），若以支配关系定义的预序赋予所有 $T'$ 所成的
集合，则诸 $\overline H^0(T'/T,P)$ 构成一个归纳系；置
\begin{equation}
\mathcal P(T)=\varinjlim_{T'}\overline H^0(T'/T,P).
\tag{1.6}
\end{equation}
这一表达式关于 $T$ 的函子律可显然地写出。

当
\[
P(T)=\operatorname{Pic}(X\times_ST),
\]
时，由 (1.6) 定义的 $(\mathrm{Sch})/S$ 上的反变函子称为 $X$ 在 $S$
上的\emph{相对 Picard 函子}，记作 $\operatorname{Pic}_{X/S}$；群
$\operatorname{Pic}_{X/S}(S)$ 称为 $X$ 在 $S$ 上的
\emph{相对 Picard 群}，记作 $\operatorname{Pic}(X/S)$。于是有显然的
双射
\begin{equation}
\operatorname{Pic}_{X/S}(T)\simeq\operatorname{Pic}(X\times_ST/T).
\tag{1.7}
\end{equation}

因此，$\operatorname{Pic}(X/S)$ 的一个元素由某个群
$\operatorname{Pic}(X\times_SS')$ 的元素 $\xi'$ 给出，其中
$S'\longrightarrow S$ 忠实平坦且拟紧，并且存在一个忠实平坦且拟紧的
态射 $S''\longrightarrow S'\times_SS'$，使 $\xi'$ 在
$\operatorname{Pic}(X\times_SS'')$ 中的两个逆像相同。元素 $\xi'$ 属于 $\operatorname{Pic}(X\times_SS')$，而元素
$\xi_1$ 属于 $\operatorname{Pic}(X\times_SS_1)$（二者满足刚才明确说明的
条件）定义 $\operatorname{Pic}(X/S)$ 中的同一个元素，当且仅当存在
忠实平坦且拟紧的态射
$S'_1\longrightarrow S'\times_SS_1$，使上述两个元素在
$\operatorname{Pic}(X\times_SS'_1)$ 中的像相等。继续使用上面引入的
函子 $P'=\operatorname{Pic}'_{X/S}$ 往往很方便；立即可见，典范态射
$P\longrightarrow P'$ 定义一个同构
\begin{equation}
\mathcal P\xrightarrow{\sim}\mathcal P',
\tag{1.8}
\end{equation}
由此可用 $\operatorname{Pic}'_{X/S}=P'$ 来描述
$\operatorname{Pic}_{X/S}$，这往往更方便。根据下文 2.3，在刚刚明确
说明的 $\operatorname{Pic}(X/S)$ 的描述中以 $P'$ 代替 $P$ 时，至少
在前引文献中所说明的条件下，确实可以
\src{146}
取 $S''=S'\times_SS'$、$S'_1=S'\times_SS_1$。

当函子 $\operatorname{Pic}_{X/S}$ 可表时，称 $X/S$ 有 Picard
预概形；表示该函子的 $S$ 上预概形称为 $X$ 在 $S$ 上的
\emph{Picard 预概形}，仍记作 $\operatorname{Pic}_{X/S}$。为此，
$P'=\operatorname{Pic}'_{X/S}$ 可表显然已经足够，因为此时 $P'$ 已经是
一个“层”，而公式 (1.8) 表明态射 $P'\longrightarrow\mathcal P'$ 可与
典范态射
\begin{equation}
\operatorname{Pic}'_{X/S}\longrightarrow\operatorname{Pic}_{X/S},
\tag{1.9}
\end{equation}
等同；该态射此时是同构。因此，我们的术语与上面随 (1.4) 引入的
术语相容。一般地，当 $\operatorname{Pic}_{X/S}$ 存在时，它由函子性
同构
\begin{equation}
\operatorname{Hom}_S(S',\operatorname{Pic}_{X/S})\xrightarrow{\sim}
\operatorname{Pic}(X\times_SS'/S').
\tag{1.10}
\end{equation}
定义。

\fsection{2}{各类相对与绝对 Picard 群之间的关系}

\result{命题}{2.1} 设 $f:X\longrightarrow S$ 是满足
$\mathcal O_S\xrightarrow{\sim}f_*(\mathcal O_X)$ 的态射。则有正合列
\[
0\longrightarrow\operatorname{Pic}(S)\longrightarrow\operatorname{Pic}(X)
\longrightarrow\operatorname{Pic}'(X/S).
\]
当 $X$ 在 $S$ 上有截面时，最后一个态射为满射，亦即有同构
\[
\operatorname{Pic}'(X/S)\xrightarrow{\sim}
\operatorname{Pic}(X)/\operatorname{Pic}(S).
\]

该正合列可视为关于 $f$ 与 $\mathcal O_X^*$ 的 Leray 谱序列所对应的
低次数正合列。第二个断言也同样是形式的。

\result{命题}{2.2} 设 $f:X\longrightarrow S$ 是拟紧分离态射，满足
$\mathcal O_S\xrightarrow{\sim}f_*(\mathcal O_X)$；设
$S'\longrightarrow S$ 是忠实平坦且拟紧的态射。则：

\noindent (i) $\operatorname{Pic}'(X/S)\longrightarrow
\operatorname{Pic}'(X\times_SS'/S')$ 是单射；

\noindent (ii) 若 $X$ 在 $S$ 上局部有截面（即每个 $s\in S$ 都有一个
开邻域 $U$，使 $X|U$ 在 $U$ 上有截面），则图式
\[
\operatorname{Pic}'(X/S)\longrightarrow
\operatorname{Pic}'(X\times_SS'/S')\rightrightarrows
\operatorname{Pic}'(X\times_SS''/S'')
\]

\src{147}
是正合的，其中 $S''=S'\times_SS'$。

借助忠实平坦下降的初等性质，第一个断言可由下述一般事实推出。若
$f:X\longrightarrow S$ 满足
$\mathcal O_S\xrightarrow{\sim}f_*(\mathcal O_X)$，则从 $S$ 上有限型
局部自由模范畴到 $X$ 上有限型局部自由模范畴的函子
$\mathcal F\longmapsto f^*(\mathcal F)$ 是充分忠实的；它的本质像由
$X$ 上满足如下条件的模 $\mathcal G$ 构成：$f_*(\mathcal G)$ 局部
自由，且典范同态
\[
f^*\bigl(f_*(\mathcal G)\bigr)\longrightarrow\mathcal G
\]
是同构。第二个断言已在 [4]，I，B，4 中用下降理论证明。

2.2 的结果也可表述如下：

\result{推论}{2.3} 在 2.2 的条件下，典范同态 (1.9)
$\operatorname{Pic}'_{X/S}\longrightarrow\operatorname{Pic}_{X/S}$ 是
单射；若 $X$ 在 $S$ 上局部有截面，则它还是双射。（因此，在后一
情形中，相对 Picard 群 $\operatorname{Pic}(X/S)$ 与限制相对
Picard 群 $\operatorname{Pic}'(X/S)$ 相同。）

与 2.1 结合可得：

\result{推论}{2.4} 在 2.2 的条件下，有正合列
\[
0\longrightarrow\operatorname{Pic}(S)\longrightarrow\operatorname{Pic}(X)
\longrightarrow\operatorname{Pic}(X/S).
\]
当 $X$ 在 $S$ 上有截面时，最后一个同态为满射，亦即此时有同构
\[
\operatorname{Pic}(X/S)\xrightarrow{\sim}
\operatorname{Pic}(X)/\operatorname{Pic}(S).
\]

\result{注记}{2.5} 设 $f:X\longrightarrow S$ 是满足
$\mathcal O_S\xrightarrow{\sim}f_*(\mathcal O_X)$ 的态射，$g$ 是 $X$
在 $S$ 上的一个截面。设 $\mathcal L$ 是 $X$ 上的可逆模；把同构
$\mathcal O_S\xrightarrow{\sim}g^*(\mathcal L)$ 称为 $\mathcal L$ 的
$g$-刚化，并把配有一个 $g$-刚化的 $X$ 上可逆模 $\mathcal L$ 称为
\emph{$g$-刚化可逆模}。这种结构的每个自同构都是平凡的，并且
$\operatorname{Pic}'(X/S)$ 可与 $X$ 上 $g$-刚化可逆模的同构类集合
等同\footnote{印刷文本作“在 $S$ 上”；这里，由截面 $g:S\to X$
进行刚化所涉及的是 $X$ 上的模 $\mathcal L$。}。（正是这一事实使我们
能够使用下降理论来证明 2.2(ii)。）至少当 $f$ 还拟紧且分离时，这给出
$\operatorname{Pic}(X/S)$ 的一种新解释；此时由 2.3，
$\operatorname{Pic}(X/S)\simeq\operatorname{Pic}'(X/S)$。

\src{148}
\result{注记}{2.6} 设 $f:X\longrightarrow S$ 为 2.2 中那样的态射，
$S'\longrightarrow S$ 为忠实平坦且拟紧的态射，并且存在一个
$S$-态射 $S'\longrightarrow X$，亦即 $X'=X\times_SS'$ 在 $S'$ 上有
截面。置 $S''=S'\times_SS'$、$X''=X\times_SS''$，并考虑正合列
\[
\operatorname{Pic}(X/S)\longrightarrow\operatorname{Pic}(X'/S')
\rightrightarrows\operatorname{Pic}(X''/S'').
\]
应用 2.4，得到正合列
\[
\operatorname{Pic}(X/S)\longrightarrow
\operatorname{Pic}(X')/\operatorname{Pic}(S')
\rightrightarrows
\operatorname{Pic}(X'')/\operatorname{Pic}(S''),
\]
特别地，相对 Picard 群的每个元素都已经“来自”
$\operatorname{Pic}(X')$ 的某个元素。因此，这大大简化了上一节给出的
相对 Picard 群的描述，也同样简化了 $X$ 在 $S$ 上的 Picard 函子的
描述，因为对 $S$ 上的每个 $T$，可将上述结果应用于
$X\times_ST/T$ 以及态射 $T'=S'\times_ST\longrightarrow T$。例如，
若 $f$ 本身忠实平坦，则可取 $S'=X$；当 $f$ 还为有限型（相应地，
光滑，等等）时，在相对 Picard 函子
$T\longmapsto\operatorname{Pic}(X\times_ST/T)$ 的描述中，就可以只考虑
有限型（相应地，光滑，等等）的基变换 $T'\longrightarrow T$。当 $f$
射影且平坦、$S$ 局部诺特时，可以证明上面的
$S'\longrightarrow S$ 可取成使 $S'$ 为 $S$ 的诸开集 $S_i$（它们覆盖
$S$）的平坦覆盖 $S'_i$ 的不交并；若 $f$ 甚至可分，则可取 $S'_i$ 在
$S_i$ 上平展。

\fsection{3}{主要存在性定理：陈述}

即使作为猜想，目前也没有一个涵盖所有已知情形的 Picard 预概形存在性
命题。可以说，一个“几乎必要”的条件是 $f:X\longrightarrow S$ 固有
（这保证若干本质的有限性性质）且平坦。这些条件并不充分，即使 $S$
是域 $k$ 上对偶数代数 $k[t]/(t^2)$ 的谱（例如取复数域
$\mathbf C$），且 $X$ 的维数为 1，也仍不充分。在撰写本篇报告时，
关于 Picard 预概形的最重要存在性定理均由下述定理推出：

\result{定理}{3.1} 设 $f:X\longrightarrow S$ 是局部诺特预概形之间的
态射。假设：

\noindent (i) $f$ 射影；

\src{149}
\noindent (ii) $f$ 平坦；

\noindent (iii) $f$ 的几何纤维均为整概形。

在这些条件下，$\operatorname{Pic}_{X/S}$ 存在。

将在以下两节概述的证明，同时还会表明如下事实：设 $\xi$ 是
$\operatorname{Pic}_{X/S}$ 的一个截面，它对应于 $X/S$ 上的甚丰沛层
$\mathcal O_X(1)$（即由射影浸入
$X\longrightarrow\mathbf P(\mathcal E)$ 所诱导）；则存在
$\operatorname{Pic}_{X/S}$ 的一个开部分 $U$，它是若干在 $S$ 上
拟射影的开部分之不交并，使得 $U$ 在由 $\xi$ 作平移下稳定，并且
$\operatorname{Pic}_{X/S}$ 是诸开集 $U-n\xi$（它们全都同构于 $U$）
的递增并。特别地，在 3.1 的条件下，
$\operatorname{Pic}_{X/S}$ 在 $S$ 上分离。

\result{注记}{3.2} 从一些例子可以看出（例如 $S$ 是离散赋值环的谱，
$X$ 在 $S$ 上相对维数为 1），若在 3.1 中去掉假设 (iii)，并以如下
较弱假设代替：对每个 $s\in S$，同态
\[
k(s)\longrightarrow H^0(X_s,\mathcal O_{X_s})
\]
都是同构，则 $\operatorname{Pic}_{X/S}$ 未必在 $S$ 上分离。这既会
发生在 $f$ 的几何纤维均约化、但一个整的几何一般纤维在特殊化时
“裂开”为两个不可约分支的情形，也会发生在 $f$ 的几何纤维均不可约、
但一个整的几何一般纤维特殊化为一个“重纤维”的情形。第一种情形例如
出现在一条二次曲线退化为两条相交直线时；第二种情形的一个例子由
D. Mumford 告诉我，其中一条椭圆曲线退化为一条二重椭圆曲线。这些
例子在任意特征下都成立。

\result{注记}{3.3} 在 3.1 的条件下，我不知道
$\operatorname{Pic}_{X/S}$ 是否是不交开部分之并，而这些开部分在 $S$
上为有限型、因而拟射影。\footnote{1962 年的勘误指出，D. Mumford
利用这里发展的 Picard 概形理论肯定地解决了这个问题。}注意，像
Hilbert 概形的情形（[4]，IV）一样，考察 Hilbert 多项式
$Q\in\mathbf Q[t]$ 可把 $\operatorname{Pic}_{X/S}$ 分解为诸开部分
$\operatorname{Pic}^{Q}_{X/S}$ 的不交并；这些开部分似乎应当在 $S$
上拟射影。至少在下一篇报告中，当 $f$ 是光滑态射时，我们将会看到
这一点。应注意，如果把假设 (i) 替换为“$X$ 在 $S$ 上局部射影”这一
假设（这足以保证 3.1 成立，因为 $\operatorname{Pic}_{X/S}$ 的存在性
显然是 $S$ 上的局部问题），那么另一方面却容易

\src{150}
给出一些例子，其中 $\operatorname{Pic}_{X/S}$ 含有在 $S$ 上并非有限型
的连通分支。例如，设 $X_0$ 是代数闭域 $k$ 上的非奇异射影代数簇，
配备一个自同构 $u$ 和 $X_0$ 的 Néron--Severi 群中的一个元素 $\xi$，
使各 $u^n(\xi)$ 两两不同。例如，可以取 $X_0$ 为一条椭圆曲线 $E$
与自身的积，并取 $u$ 为 $E\times E$ 的如下自同构：
\[
(x,y)\longmapsto(x,y+x)
\]
令 $S$ 为两条在 $a$、$b$ 两点相交的非奇异不可约曲线之并。$S$ 上存在
一个群为 $\mathbf Z$ 的连通主覆盖 $P$；利用由 $u$ 定义的
$\mathbf Z$ 在 $X_0$ 上的作用，可由此得到 $S$ 上一个以 $X_0$ 为纤维
的伴随丛（它在 $S-a$ 和 $S-b$ 上平凡）；在所考察的特殊情形中，它
实际上是 $S$ 上的一个阿贝尔概形。容易看出，
$\operatorname{Pic}_{X/S}$ 也是由 $P$ 以及 $\mathbf Z$ 通过 $u$ 在
$\operatorname{Pic}_{X_0/k}$ 上的作用所得到的伴随丛，并含有一个
同构于
\[
P\times\operatorname{Pic}^{0}_{X_0/k},
\]
的连通分支（其中 $\operatorname{Pic}^{0}$ 表示
$\operatorname{Pic}$ 中单位元所在的连通分支）；该分支在 $S$ 上并非
有限型。[在 3.3 所考察的、$S$ 上非分离的 Picard 预概形的若干情形中，
也会出现类似现象。]

\fsection{4}{相对 Cartier 除子与射影丛}

我们只需使用有效除子，故在本节下文中省略这一附加限定词。

设 $X$ 是一个预概形。$X$ 上的一个 \emph{Cartier 除子}，或为简便起见
简称 $X$ 上的除子，是 $X$ 的一个闭子预概形 $D$；它由一个理想
$\mathcal J$ 定义，而该理想是可逆模，亦即局部由
$\mathcal O_X$ 的一个非零因子截面生成。我们给 $D$ 配置一个可逆模，
即
\[
\mathcal L(D)=\mathcal J^{-1},
\]
而典范单射 $\mathcal J\longrightarrow\mathcal O_X$ 给出一个典范同态
\[
s_D:\mathcal O_X\longrightarrow\mathcal J^{-1}=\mathcal L(D),
\qquad\text{亦即}\quad s_D\in\Gamma(X,\mathcal L(D)).
\]
此外，给定一个除子，本质上等价于给定 $X$ 上的一个可逆模
$\mathcal L$ 以及它的一个处处为非零因子的截面 $s$；反过来，对这样的
一对 $(\mathcal L,s)$，配置 $s$ 的“除子”，记作
$\operatorname{div}(s)$。对于 $X$ 上一个给定的可逆模 $\mathcal L$，
所有

\src{151}
定义 $\mathcal L$ 的除子 $D$ 之集合，与商集
\[
\Gamma(X,\mathcal L)^*/\Gamma(X,\mathcal O_X^*),
\]
一一对应，其中 $\Gamma(X,\mathcal L)^*$ 表示
$\Gamma(X,\mathcal L)$ 中由处处为非零因子的截面构成的部分。

现在假设给定一个局部有限型态射 $f:X\longrightarrow S$，并为简便起见
假设 $S$ 局部诺特。设 $\mathcal J$ 是 $X$ 上的凝聚理想，$D$ 是它
定义的 $X$ 的子概形，$x\in X$，且 $s=f(x)$。可以证明，下列条件
等价：

\noindent (i) $\mathcal J$ 在 $x$ 处可逆（亦即 $\mathcal J_x$ 由
$\mathcal O_{X,x}$ 的一个正则元生成），并且 $D$ 在 $x$ 处于 $S$
上平坦；

\noindent (ii) $X$ 和 $D$ 在 $x$ 处均于 $S$ 上平坦，并且 $D_s$ 在点
$x$ 处是纤维 $X_s$ 上的 Cartier 除子；

\noindent (iii) $X$ 在 $x$ 处于 $S$ 上平坦，并且 $\mathcal J_x$
由一个元素 $f_x$ 生成，而该元素在 $X_s$ 上诱导出一个非零因子芽。

此时称 $D$ 在所考察点处是 $X/S$ 上的一个
\emph{相对 Cartier 除子}（或简称 $X/S$ 上的相对除子）。由 (i)
可见，此时 $D$ 在 $x$ 的邻近各点处也是相对除子。因此，若 $X$ 和
$D$ 均在 $S$ 上平坦，而 $D$ 在 $S$ 上固有，则所有如下 $s\in S$
构成 $S$ 的一个开部分：$D_s$ 是 $X_s$ 中的 Cartier 除子（亦即
$D$ 在 $X_s$ 的各点处是相对 Cartier 除子）。另一方面，上述定义
正是为使相对 Cartier 除子的概念在任意基变换
$S'\longrightarrow S$ 下稳定而作出的。现在考虑 $X/S$ 上相对除子
的集合 $\operatorname{Div}(X/S)$，并定义以 $S$ 上的 $S'$ 为变量的
反变函子
\[
\operatorname{Div}_{X/S}(S')=
\operatorname{Div}(X\times_SS'/S').
\]

假设 $X$ 在 $S$ 上平坦且固有。由相对 Cartier 除子的刻画 (ii)，
$\operatorname{Div}_{X/S}$ 可视为 [4]，IV 中定义的函子
$\operatorname{Hilb}_{X/S}$ 的子函子；由上述说明，包含态射
\[
\operatorname{Div}_{X/S}\longrightarrow\operatorname{Hilb}_{X/S}
\]
“可由开浸入表示”（参见 [5]，IV，3.13）。运用 [4]，IV 的主存在
定理，得到：

\src{152}
\result{命题}{4.1} 假设 $f:X\longrightarrow S$ 射影且平坦。那么函子
$\operatorname{Div}_{X/S}$ 可表；更确切地说，它可由
$\operatorname{Hilb}_{X/S}$ 的一个开部分表示。

对于 $X/S$ 上一个给定的甚丰沛层 $\mathcal O_X(1)$，利用
$\operatorname{Hilb}_{X/S}$ 按照 Hilbert 多项式
$Q\in\mathbf Q[t]$ 典范分解为各开部分
$\operatorname{Hilb}^{Q}_{X/S}$ 的分解，还可推出类似的分解
\[
\operatorname{Div}_{X/S}=
\coprod_{Q\in\mathbf Q[t]}\operatorname{Div}^{Q}_{X/S}
\]
即分解为在 $S$ 上拟射影的两两不交开部分之不交并。

另一方面，映射 $D\longmapsto\mathcal L(D)$ 给出函子性同态
\begin{equation}
\operatorname{Div}_{X/S}\longrightarrow\operatorname{Pic}_{X/S}
\tag{+}
\end{equation}
，我们将研究这个同态；在相当一般的条件下，它将是相对可表的
（[5]，IV，3）。因此，从
$\operatorname{Pic}_{X/S}(S')$ 的一个元素 $\xi$ 出发；为简化记号，
假设 $S'=S$，并证明 $\operatorname{Div}_{X/S}$ 的相应子函子可表。
先考虑 $\xi$ 由 $X$ 上的一个可逆模 $\mathcal L$ 定义的情形。我们
假设 $X$ 在 $S$ 上固有且平坦，并且 $X$ 在 $S$ 上的各几何纤维均为
整的；这还蕴涵（[2]，III，第7段）
$\mathcal O_S\xrightarrow{\sim}f_*(\mathcal O_X)$，并且这一关系在
任意基变换 $S'\longrightarrow S$ 后仍成立。于是，$X/S$ 上满足
下列条件的相对 Cartier 除子 $D$：
$\mathcal L(D)$ 与 $\mathcal L$ 定义
$\operatorname{Pic}(X/S)=\operatorname{Pic}_{X/S}(S)$ 的同一元素，
亦即根据 2.4，$\mathcal L(D)$ 与 $\mathcal L$ 在 $S$ 上局部同构；
它们与商层 $f_*(\mathcal L)^*/\mathcal O_S^*$ 的截面一一对应。
这一对应与基变换相容。另一方面，所引文献中“Künneth”型的一般论证
（另见 [6]）表明，$X/S$ 的固有性以及 $\mathcal L$ 在 $S$ 上的
平坦性蕴涵：存在 $S$ 上一个凝聚模 $\mathcal Q$，它在唯一同构的意义
下确定，并且存在层同构
\[
f_*(\mathcal L)\xrightarrow{\sim}
\operatorname{Hom}_{\mathcal O_S}(\mathcal Q,\mathcal O_S),
\]
\footnote{印刷文本把 $\mathcal O_X$ 写在
$\operatorname{Hom}$ 的下标中；由于 $\mathcal Q$ 是 $S$ 上的模，
其对偶应在 $\mathcal O_S$ 上取。}
而 $\mathcal Q$ 的形成与基变换相容。这里
$f_*(\mathcal L)^*$ 表示 $f_*(\mathcal L)$ 的集合子层；它在 $U$
上的截面，是 $\mathcal L$ 在 $f^{-1}(U)$ 上定义

\src{153}
$f^{-1}(U)/U$ 上相对 Cartier 除子的截面，亦即在各
$X_s$（$s\in U$）上诱导出非零因子截面的截面。利用纤维 $X_s$ 为
整的这一假设，这只意味着在各纤维 $X_s$ 上诱导的截面并非恒等于零；
用局部同态 $\mathcal Q\longrightarrow\mathcal O_S$ 来说，也就是
这些同态为满射（Nakayama 引理）。这表明，
$f_*(\mathcal L)^*/\mathcal O_S^*$ 的截面集合与 $\mathcal Q$ 的
可逆商模集合一一对应；或者，根据与凝聚模 $\mathcal Q$ 相伴的射影丛
$\mathbf P(\mathcal Q)$ 的定义（参见 [5]，V，2），它与
$\mathbf P(\mathcal Q)$ 在 $S$ 上的截面集合一一对应。这一描述与
拉回的形成相容，因此得到下述定理。

\result{定理}{4.3} 设 $f:X\longrightarrow S$ 是一个固有平坦态射，
其几何纤维均为整的，其中 $S$ 局部诺特；设 $\mathcal L$ 是 $X$
上的可逆模。对于 $S$ 上的每个 $S'$，令 $T(S')$ 为
$X\times_SS'/S'$ 上所有满足下列条件的相对除子 $D$ 之集合：
$\mathcal L(D)$ 在 $S'$ 上局部同构于
$\mathcal L\otimes_{\mathcal O_S}\mathcal O_{S'}$（亦即
$\mathcal L(D)$ 与
$\mathcal L\otimes_{\mathcal O_S}\mathcal O_{S'}$ 定义
$\operatorname{Pic}(X\times_SS'/S')$ 的同一元素）。那么存在 $S$
上的一个凝聚模 $\mathcal Q$，它在唯一同构的意义下确定，使函子
$T$ 可由射影丛 $\mathbf P(\mathcal Q)$ 表示。

\result{推论}{4.4} 若进一步假设 $f$ 射影，则函子性同态
$\operatorname{Div}_{X/S}\longrightarrow\operatorname{Pic}_{X/S}$
可由射影态射表示。

事实上，当 $X$ 在 $S$ 上有截面（相应地，局部有截面）时，由 4.3
和 2.1，上述同态可由射影丛（相应地，由局部射影丛）表示。在 $f$
拟射影的情形下，使用 $X$ 在 $S$ 上局部的有限平坦准截面，通过
下降法容易归约到前一种情形。

\result{注}{4.5} 在 4.3 的条件下，$S$ 上的模 $\mathcal Q$ 一般并非
局部自由；这表现为 $\mathcal Q$ 的约化纤维的维数，亦即当
$s\in S$ 变化时 $H^0(X_s,\mathcal L_s)$ 的维数，可能发生跳跃。
给定局部诺特预概形 $S$ 上的一个凝聚模 $\mathcal Q$，容易验证：
对于给定的 $s\in S$，$\mathcal Q$ 在 $s$ 处自由，当且仅当
$\mathbf P(\mathcal Q)$ 在 $s$ 上方各点处于 $S$ 上平坦（在这种
情形下，它甚至在 $s$ 上方各点处于 $S$ 上光滑）。在当前情形中，
$\mathcal Q$ 按上述方式由 $\mathcal L$ 定义；这也意味着正像
$f_*(\mathcal L)$ 的形成“与基变换在 $s$ 的

\src{154}
邻域内可交换”，或者等价地说，
$f_*(\mathcal L)_s\longrightarrow H^0(X_s,\mathcal L_s)$ 为满射。
例如，当 $H^1(X_s,\mathcal L_s)=0$ 时，情况就是如此。在所涉及的
各预概形存在的前提下，这些判据尤其适用于泛情形
$\operatorname{Div}_{X/S}\longrightarrow\operatorname{Pic}_{X/S}$，
并分别给出这个态射在 $\operatorname{Div}_{X/S}$ 的一个给定点处
光滑的充分必要条件和充分条件。

\fsection{5}{主存在定理的证明}

在 3.1 的条件下，选择 $X/S$ 上的一个甚丰沛模
$\mathcal O_X(1)$，并令 $\xi$ 为
$\operatorname{Pic}(X/S)$ 中的相应元素。简记
\[
P(S')=\operatorname{Pic}(X\times_SS'/S'),
\]
并为简便起见，先假设 $X/S$ 有一个截面。令 $P^+(S')$ 为
$P(S')$ 中由 $X\times_SS'$ 上满足下列条件的可逆模 $\mathcal L$
的类所构成的部分：
\[
R^if'_*(\mathcal L(n))=0\quad\text{对 $i>0$ 及所有 $n\geq0$},
\qquad
f'_*(\mathcal L(n))\neq0\quad\text{对所有 $n\geq0$}.
\]
这些条件在基变换下稳定，因而定义了 $P$ 的一个子函子 $P^+$；
它显然在由 $\xi$ 给出的平移下稳定。利用 Serre 的“定理 A 和 B”
（[2]，III，2）以及 [5]，IV，5 中的一般理论，容易看出：$P$
可表，当且仅当 $P^+$ 可表；此时，$P^+$ 可由表示 $P$ 的 Picard
预概形 $\operatorname{Pic}_{X/S}$ 的一个开部分 $U$ 表示，而后者
是各开部分 $U-n\xi$ 的递增并。

简记 $D=\operatorname{Div}_{X/S}$，并令 $D^+$ 为典范态射
$D\longrightarrow P$ 下 $P^+$ 的逆像。因此有态射
\begin{equation}
D^+\longrightarrow P^+,
\tag{+}
\end{equation}
并且已经知道，$D^+$ 可由预概形
$D=\operatorname{Div}_{X/S}$ 的一个开部分 $D^+$ 表示（其存在性由
4.1 保证）。由 4.3 容易推出，态射 (+) 可由光滑、射影且满射的态射
表示（更确切地说，由与处处非零的局部自由模相伴的射影丛表示）。
其原因是：当 $X\times_SS'$ 上的 $\mathcal L$ 如本节开头所述时，
$f'_*(\mathcal L)$ 是一个局部自由模 $\neq0$，其形成与基变换可交换；
采用 4.3 的记号，$\mathcal Q$ 此时同构于
$f'_*(\mathcal L)$ 的对偶。利用一般判据（[5]，IV，4.7），我们将由此
推出 $P^+$ 的可表性。在所引文献中，取 $\mathfrak C$ 为所有忠实

\src{155}
平坦且拟紧的预概形态射所成的类（根据 [4]，I B，它们确实是有效
满态射）。所引文献的条件 (a)，即 (+) 可由属于 $\mathfrak C$ 的
态射表示，已经由上述论证验证；条件 (b) 也同样得到验证，它表示
函子 $P^+$ 与忠实平坦拟紧下降相容，而这一点是立即可见的。只剩下
证明所引文献的条件 (c)：由 $\mathfrak C$-可表态射 (+) 在预概形
$D^+$ 中导出的等价关系 $R$ 是 $\mathfrak C$-有效的，亦即它是有效
的，并且 $D^+\longrightarrow D^+/R$ 属于 $\mathfrak C$。为此，
先注意，与各个虚 Hilbert 多项式 $Q\in\mathbf Q[t]$ 对应的 $D^+$
的开部分 $D^{+Q}$ 在 $R$ 下稳定（因为 $R$ 的纤维连通）；因此，
问题归约为证明：对于每个 $Q$，$D^{+Q}$ 中诱导的等价关系 $R^Q$
是 $\mathfrak C$-有效的。现在，$D^{+Q}$ 是拟射影的，而等价关系
$R^Q$ 是射影且平坦的。因此，[4]，III，6.1 的适用条件均已满足，
并蕴涵所需结果。

在 $X/S$ 不必有截面的一般情形下，可通过下降法容易归约到前一种
情形；也可以在 $P^+$ 的定义中作必要修改后，重复上述证明。

\result{注记}{5.1} 所采用的方法本质上是 Matsusaka 用于射影构造
Picard 簇的方法。为取有效商而援用的 [4]，III 中的结果，也容易由
Hilbert 概形存在定理推出（例如参见 [6]）（经典方法则利用 Chow
坐标来构造这些商）。还应注意，
$\operatorname{Pic}_{X/S}$ 的开部分
$\operatorname{Pic}^{+}_{X/S}$ 的形成，以及它按照定义所考察可逆模
的各除子的 Hilbert 多项式而分解为在 $S$ 上拟射影的开部分
$\operatorname{Pic}^{+Q}_{X/S}$，均与基变换相容（这使下降法得以
应用）。

\result{注记}{5.2} 不排除在下述所有情形中
$\operatorname{Pic}_{X/S}$ 都存在：$f:X\longrightarrow S$ 固有且
平坦，并且各同态
\[
k(s)\longrightarrow H^0(X_s,\mathcal O_{X_s})\qquad(s\in S)
\]
均为同构（后一条件此时还表示
$\mathcal O_S\xrightarrow{\sim}f_*(\mathcal O_X)$，且这一关系在
任意基变换 $S'\longrightarrow S$ 后仍成立）。至少，在 $f$ 进一步
射影时，解析空间的框架中可以用 [5]，IX，3.1 所说明的另一种方法
（若我没有弄错，即 Chow 的方法）证明这一点。
\footnote{1962年的勘误指出，本注记开头提出的存在性猜想是错误的，
但 Mumford 的方法给出一个稍弱的定理。}
在这种方法中，取

\src{156}
在除子概形的某个开集中按固有平坦等价关系取商，被替换为在 $X$
浸入 $\mathbf P^r_S$ 的浸入概形中按射影群取商。利用 Mumford 关于按
射影群取商的结果 [8]，这一方法很可能可以推广到概形的情形；目前，
只有当 $X$ 在 $S$ 上有“许多局部截面”时才有成文的证明，例如，当
$X$ 在一个剩余域代数闭的完备局部环上可分时。
\footnote{印刷文本漏去了“剩余域代数闭”；勘误指出，Mumford 的反例
表明这一限制不可或缺。}
原则上，所述方法的适用范围会更广，因为它还能在 Picard 预概形不分离
的情形下给出其存在性，而第一种方法在这种情形下必然失效。（从技术上
说，困难来自如下事实：当 $f$ 的几何纤维不再为整时，4.3 中所考察的
函子不再由射影丛 $\mathbf P(\mathcal Q)$ 本身表示，而是由它的一个
开部分表示；这就引出了按平坦但非固有的等价关系取商这一棘手问题。）

\result{注}{5.3} 应当注意，这里给出的证明既不借助预先构造曲线或
曲线族的 Jacobian，也不借助阿贝尔簇或阿贝尔概形理论；因此，它与
Lang 的著作 [7] 或 Chevalley 的报告 [1] 这类沿用 A. Weil 所勾勒
途径的传统论述有本质区别。即使对于代数闭域（比如复数域）上非奇异
曲线的 Jacobian，这里给出的构造也是目前所知唯一能赋予它第 1 节中
被我们用作定义的那些很强性质的构造；这些性质本质上就是 Chevalley
的性质，但还把带幂零元的“参数簇”考虑在内。Picard 概形的构造应当
先于而不是后于阿贝尔簇理论，这一点先验地就很清楚，因为 Picard 概形
一般既不是阿贝尔概形，也不能归结为阿贝尔概形。代数闭域上的奇异曲线
已经说明了这一点：此时会出现 Rosenlicht 的“广义 Jacobian”，而它们
并不是阿贝尔簇。此外，一旦有了一般的 Picard 概形理论，阿贝尔簇理论
以及更一般的阿贝尔概形理论都会大为简化。特别地，阿贝尔概形的对偶
理论，尤其是 Cartier 型结果，将由此变得几乎是形式的（例如参见
[10]）。

\src{157}
\result{注}{5.4} Igusa 关于退化为奇异曲线的曲线之 Jacobian 的
“相容性原理”，只有把它理解为不必在 $S$ 上光滑的相对曲线概形
$X/S$ 的 Picard 概形存在性定理，才能得到恰当的说明。因此，
当特化曲线为整时（即用经典术语说，它不可约且重数为 1），这就是主要
存在性定理 3.1 的一个特例。应当注意，目前当特殊曲线可约时，即使各
分支的重数均为 1，也就是说特殊曲线在剩余域上可分，已知存在性定理
仍不能处理这一情形；唯一的例外是基底为一个剩余域代数闭的完备离散
赋值环，参见 5.2。在用概形理论对亏格为 $g$ 的曲线模概形作
Baily--Satake“紧化”的几何构造中，必然会遇到这一存在性问题。
（目前仅当 $g=1$ 时，由于 Igusa 的工作，才知道这种紧化。）

\fsection{6}{相对存在性定理}

这里将概述若干有用情形：某些 Picard 概形的存在性蕴含另一些 Picard
概形的存在性；由此可从主要定理 3.1 推出其他若干存在性定理。

\result{命题}{6.1} 设 $f:X\longrightarrow S$ 是平坦射影态射，并且在
Stein 分解 $f=f''f'$ 中，态射 $f':X\longrightarrow S'$ 平坦且具有
整的几何纤维（因而满足 3.1 的假设），而有限态射
$f'':S'\longrightarrow S$ 平坦。那么 $\operatorname{Pic}_{X/S}$
存在，并且（采用 [4]，II，C 2 中引入的记号）有典范同构
\[
\operatorname{Pic}_{X/S}\xrightarrow{\sim}
\Pi_{S'/S}\operatorname{Pic}_{X/S'}.
\]

为证明这一点，先建立函子的同构
\[
\operatorname{Pic}_{X/S}\xrightarrow{\sim}
\Pi_{S'/S}\operatorname{Pic}_{X/S'},
\]
再使用 3.1；该定理蕴含 $\operatorname{Pic}_{X/S'}$ 可表。最后，
利用第 n\textsuperscript{o} 3 节中指出的
$\operatorname{Pic}_{X/S'}$ 的结构，即
$\operatorname{Pic}_{X/S'}$ 在 $S$ 上每条纤维中的任一有限子集都
包含于某个仿射开集中，得到
$\Pi_{S'/S}\operatorname{Pic}_{X/S'}$ 的存在性。

\src{158}
例如，当 $X$ 是 $S$ 上满足 3.1 条件的诸概形 $X_i$ 的不交并概形时，
陈述 6.1 归结为平凡的陈述
\[
\operatorname{Pic}_{X/S}\xrightarrow{\sim}
\prod_i\operatorname{Pic}_{X_i/S}.
\]

\result{推论}{6.2} 设 $f:X\longrightarrow S$ 是射影平坦态射，其几何
纤维局部为整（例如，该态射射影且正规），则
$\operatorname{Pic}_{X/S}$ 存在。

此外，在这一情形下，$S'$ 是 $S$ 的一个平展覆盖（每当 $f$ 可分时
均如此；所谓可分，是指该态射平坦且几何纤维约化）。利用
$\operatorname{Pic}_{X/S'}$ 的类似结构，可以验证第
n\textsuperscript{o} 3 节对 $\operatorname{Pic}_{X/S}$ 陈述的
结构定理仍然成立。

另一方面，下降法给出一个相对存在性定理；其适用范围取决于 [4]，I，
A (c) 中提出的非平坦下降问题的解决。这里仅明确陈述其中一个特例：

\result{命题}{6.3} 设 $f:X\longrightarrow S$ 是固有态射，$X_1$ 和
$X_2$ 是 $X$ 的两个在 $S$ 上平坦的子预概形，分别由两个凝聚理想
$\mathcal I_1$ 和 $\mathcal I_2$ 定义，并满足
$\mathcal I_1\cap\mathcal I_2=(0)$，且
$\mathcal O_X/(\mathcal I_1+\mathcal I_2)$ 在 $S$ 上平坦（即 $X_1$
与 $X_2$ 在 $X$ 中的上确界为 $X$，而它们的下确界 $Z$ 在 $S$ 上
平坦）。
\footnote{最后一个字母在印刷文本中为“$X$”，但紧接在前面的条件是
$\mathcal O_X/(\mathcal I_1+\mathcal I_2)$ 在 $S$ 上平坦，因此定义
的是 $Z$ 在 $S$ 上平坦。}
再假设对每个 $s\in S$，同态
\[
k(s)\longrightarrow H^0(X_{i,s},\mathcal O_{X_{i,s}}),\qquad i=1,2,
\]
均为双射。那么函子的自然同态
\[
\operatorname{Pic}_{X/S}\longrightarrow
\operatorname{Pic}_{X_1/S}\times\operatorname{Pic}_{X_2/S}
\]
由仿射态射表示。因此，如果 $\operatorname{Pic}_{X_1/S}$ 和
$\operatorname{Pic}_{X_2/S}$ 存在，则 $\operatorname{Pic}_{X/S}$
也存在，并且典范态射
\[
\operatorname{Pic}_{X/S}\longrightarrow
\operatorname{Pic}_{X_1/S}\times\operatorname{Pic}_{X_2/S}
\]
是仿射的。

通过忠实平坦下降，可归结到 $Z$ 在 $S$ 上有截面的情形；该截面因而
定义了 $X$、$X_1$、$X_2$ 在 $S$ 上的截面，并使我们能够像 2.5 中
所说明的那样消去所考察结构中的自同构。此时，证明只须注意，一个
可逆模

\src{159}
$\mathcal L$，它在 $X$ 上带有“刚化”；这等价于给定三元组
$(\mathcal L_1,\mathcal L_2,u)$，其中 $\mathcal L_i$ 是
$X_i$ 上带有“刚化”的可逆模，而 $u$ 是从 $\mathcal L_1|Z$ 到
$\mathcal L_2|Z$ 且与刚化相容的同构。只须验证，在
$\mathcal L_1$、$\mathcal L_2$ 固定时，$u$ 的数据可以表示成 $S$
上某个在 $S$ 上仿射的适当概形的一个截面；这是容易的。由 6.3
容易推出：

\result{推论}{6.4} 设 $X$ 是域 $k$ 上固有且可分的概形，$X_i$ 是
$X$ 的各不可约分支。如果各 $\operatorname{Pic}_{X_i/k}$ 存在，
则 $\operatorname{Pic}_{X/k}$ 也存在，并且典范态射
\[
\operatorname{Pic}_{X/k}\longrightarrow\prod_i\operatorname{Pic}_{X_i/k}
\]
是仿射的。

结合 6.2，例如可知，每当 $X$ 是域 $k$ 上射影且可分的概形时，
$\operatorname{Pic}_{X/k}$ 都存在。当 $X$ 在 $k$ 上不再可分时，
利用 Oort [11] 的一个论证，也可得到一个归约结果。该方法同样适用于
任意基概形（例如，在下一篇报告中证明 3.3 所预告的有限性结果时，
这一情形很有用）。为避免陈述过于技术化，这里仅限于基底为域的情形：

\result{命题}{6.5} 设 $X$ 是域 $k$ 上的固有概形，$X_0$ 是一个
具有相同底集合的子概形（因而由 $X$ 上的一个幂零理想定义）。
那么函子态射
\[
\operatorname{Pic}_{X/k}\longrightarrow\operatorname{Pic}_{X_0/k}
\]
由仿射态射表示。特别地，若 $\operatorname{Pic}_{X_0/k}$ 存在，
则 $\operatorname{Pic}_{X/k}$ 也存在，而且态射
$\operatorname{Pic}_{X/k}\longrightarrow\operatorname{Pic}_{X_0/k}$
是仿射的。

结合 6.4，容易得到：

\result{推论}{6.6} 设 $X$ 是域 $k$ 上的射影预概形，则
$\operatorname{Pic}_{X/k}$ 存在。

\result{注}{6.6} 极有可能，对域 $k$ 上的每个固有概形 $X$，
$\operatorname{Pic}_{X/k}$ 都存在。
\footnote{1962 年勘误指出，Murre 当时刚刚肯定地解决了这一问题。}
前述结果把这一问题归结到 $k$ 代数闭且 $X$ 为整的情形。这时，由
Chow 引理，存在一个在 $k$ 上射影的整概形 $X'$ 以及支配态射
$g:X'\longrightarrow X$。因此，只须证明相应的

\src{160}
函子态射
\[
\operatorname{Pic}_{X/k}\longrightarrow\operatorname{Pic}_{X'/k}
\]
可表（而且无疑由仿射态射表示），因为另一方面已经知道
$\operatorname{Pic}_{X'/k}$ 可表。这引出了目前尚未解决的非平坦下降
问题。应当注意，若只考虑函子 $\operatorname{Pic}_{X/k}$ 在约化
预概形上的限制（其中 $X$ 是代数闭域 $k$ 上的固有整概形），则确实
得到可表函子；Chevalley [1] 在 $X$ 正规时证明了这一点，而 Seshadri
[12] 用下降法证明了一般情形。但采用我们的记号，这些作者构造的概形
并不是 $\operatorname{Pic}_{X/k}$，而是
$(\operatorname{Pic}_{X/k})_{\mathrm{réd}}$，即与
$\operatorname{Pic}_{X/k}$ 对应的约化概形。

\result{注}{6.7} 更一般地，设 $f:X'\longrightarrow X$ 是域 $k$ 上
固有预概形之间的满射态射。非平坦下降方面的考虑使人猜想
\[
\operatorname{Pic}_{X/k}\longrightarrow\operatorname{Pic}_{X'/k}
\]
是仿射态射；
\footnote{1962 年勘误指出，这一问题有肯定答案，并参见下一篇报告
《评论》的最后一段。}
这尤其会蕴含：模去单位元所在的连通分支后，Néron--Severi 群上的
相应同态模挠单射。当 $X$ 和 $X'$ 非奇异时，可由相交理论验证这一点。
在其他任何情形下，答案似乎都还是未知的。

\result{注}{6.8} 与人们可能以为的相反，考察带幂零元的代数概形的
Picard 概形，在各种问题中是有用的，甚至是不可或缺的。例如，设 $X$
是一个射影概形，比如在 $k$ 上光滑，
\footnote{印刷文本作“正则”；勘误规定改为“在 $k$ 上光滑”。}
而 $Y$ 是一个超平面截面；这时必须考察 $Y$ 的所有阶“无穷小邻域”
$X_n$ 以及 Picard 概形 $\operatorname{Pic}_{X_n/k}$。当 $X$
不可约且维数 $\geq4$（相应地，$\geq3$）时，典范态射
\[
\operatorname{Pic}_{X/k}\longrightarrow\operatorname{Pic}_{X_n/k}
\qquad(n\text{ 充分大})
\]
是同构（相应地，在 Néron--Severi 挠子群的逆像上诱导同构）；这一
结果将在下一篇报告中对 Picard 概形的定性研究中用到。同样，考察某些
带幂零元曲线的 Picard 概形，并结合形式几何的基本定理 [3]，可以在
等特征情形下验证 Mumford 的一个猜想，即：对每个维数为 2 的完备
诺特正规局部环 $A$，$A$ 的除子类群都可视为剩余域 $k$ 上某个

\src{161}
代数群 $G$ 的 $k$-有理点集合（而且一旦在 $A$ 中给定一个系数域，
$G$ 就被典范地确定）。当 $A$ 具有任意维数时，很可能存在一个 $k$
上的 pro-代数群，扮演
上述 $G$ 的同样角色；在能够对 $\operatorname{Spec}(A)$ 作“解奇异”
的情形下，它可构造为 $k$ 上适当的带幂零元射影概形之 Picard 概形的
逆极限。

\section*{参考文献}
\addcontentsline{toc}{section}{参考文献}

\begin{description}\small
\item[{[1]}] CHEVALLEY (Claude). -- Sur la théorie de Picard,
\emph{Amer. J. of Math.}, 第 82 卷，1960 年，第 435--490 页。

\item[{[2]}] GROTHENDIECK (A.) 和 DIEUDONNÉ (J.). -- Éléments de géométrie
algébrique, Chapitre I et suivants. -- 巴黎，Presses universitaires de France,
1960--1961 年（Institut des Hautes Études Scientifiques, Publications
mathématiques 4, 8, \ldots）。

\item[{[3]}] GROTHENDIECK (Alexander). -- Géométrie formelle et géométrie
algébrique, \emph{Séminaire Bourbaki}, 第 11 卷，1958/59 年，
n\textsuperscript{o} 182，28 页。

\item[{[4]}] GROTHENDIECK (Alexander). -- Techniques de descente et théorèmes
d'existence en géométrie algébrique, I--IV, \emph{Séminaire Bourbaki}, 第 12
卷，1959/60 年，n\textsuperscript{o} 190，29 页及
n\textsuperscript{o} 195，22 页；第 13 卷，1960/61 年，
n\textsuperscript{o} 212，20 页及 n\textsuperscript{o} 221，28 页。

\item[{[5]}] GROTHENDIECK (Alexander). -- Techniques de construction en
géométrie analytique, I--X, \emph{Séminaire Cartan}, 第 13 卷，1960/61 年：
Déformation des structures analytiques complexes,
n\textsuperscript{o} 7--17。

\item[{[6]}] GROTHENDIECK (Alexander). -- Séminaire de géométrie algébrique,
Lichtenbaum 编纂，Harvard University，1961 年（待刊）。

\item[{[7]}] LANG (Serge). -- \emph{Abelian varieties}. -- New York,
Interscience Publishers, 1959（Interscience Tracts in pure and applied
Mathematics, 7）。

\item[{[8]}] MUMFORD (D.). -- An elementary theorem in geometric invariant
theory, \emph{Bull. Amer. math. Soc.}, 第 67 卷，1961 年，第 483--486 页。

\item[{[9]}] MUMFORD (D.). -- The topology of normal singularities of an
algebraic surface and a criterion for simplicity. -- 巴黎，Presses
universitaires de France，1961 年（Institut des Hautes Études Scientifiques,
Publications mathématiques, 9）。

\item[{[10]}] MUMFORD (D.) 和 TATE (J.). -- Séminaire de géométrie
algébrique, Harvard University, Spring Term 1962（待刊）。

\item[{[11]}] OORT (F.). -- Sur le schéma de Picard,
\emph{Bull. Soc. math. France}, 第 90 卷，1962 年（待刊）。

\item[{[12]}] SESHADRI (C. S.). -- 待刊论文，拟发表于
\emph{Annali di Matematica}。
\end{description}

% END 232.tex

% BEGIN 236.tex
\unit{第236次报告——Picard 概形：一般性质}
\sid{236}
\src{221}
\seminarzh{第14年度，1961/62，第236号\\（1962年5月）}
\paperhead{下降技术与代数几何中的存在性定理\\
VI. Picard 概形：一般性质}%
  {亚历山大·格罗滕迪克 著}

\fsection{0}{上一篇报告（[3]，V）的补充\texorpdfstring{\textsuperscript{(1)}}{（1）}}

关于 [3]，V 中提出的 Picard 预概形存在性问题，已有了一些进展：

\medskip
\noindent a.（Mumford）一般而言，并非只要
$f:X\longrightarrow S$ 是射影可分态射（= 平坦且纤维可分），Picard
预概形 $\operatorname{Pic}_{X/S}$ 就存在；即使 $f$ 的纤维维数为 1，
且 $S$ 是完备离散赋值环的谱，也不成立。可取
$S=\operatorname{Spec}\mathbf R[[t]]$，并取 $X$ 为
$\mathbf P^2_S$ 中由下式定义的子概形（其中 $x,y,z$ 为齐次变量）：
\[
x^2+y^2=tz^2,
\]
这给出一条在几何上退化为两条相交直线的二次曲线，但在域
$\mathbf R$ 上的特殊纤维仍不可约（它由 $\mathbf R$ 上的方程
$x^2+y^2=0$ 给出）。容易看出，经过平展扩张
$S'\longrightarrow S$，其中
$S'=\operatorname{Spec}\mathbf C[[t]]$，$X'/S'$ 的 Picard 预概形存在，
并可明确描述为若干份 $\widetilde S$ 的不交并，其中
$\widetilde S$ 由 $S$ 将原点无限次加倍而得。关于
$S'\longrightarrow S$ 的 $\operatorname{Pic}_{X'/S'}$ 上的下降数据
（这里由 $S'$ 在 $S$ 上的 Galois 群 $\mathbf Z/2\mathbf Z$ 的作用给出）
不是有效的；这一点容易验证，因为该群交换某些加倍后的点（从而有些
轨道不包含在任何仿射开集中）。不过，Mumford 能够证明：若
$f:X\longrightarrow S$ 是射影可分态射，并且对每个 $s\in S$，
$X_s$ 的各不可约分支相对于 $k(s)$ 几何不可约，则
$\operatorname{Pic}_{X/S}$ 存在；其证明依赖于他的取商定理的一个
加强形式，参见 [7]。另一方面应注意，即使不对纤维 $X_s$ 的不可约
分支作任何假设，下文将引入的概形
$\operatorname{Pic}^{\tau}_{X/S}$ 仍有可能存在。

\src{222}
\noindent b.（Murre）设 $X$ 是域上的固有概形，则
$\operatorname{Pic}_{X/k}$ 存在。证明部分沿用 Chevalley 的证明 [1]，
并充分利用 Picard 函子的群结构。

关于 Picard 概形理论的一些补充，特别是其与阿贝尔概形的关系，宜参见
[7]。最后，本篇报告的一个显著缺陷，是没有给出把一个射影概形的
Picard 概形同其超平面截面的 Picard 概形进行比较的“等价判据”；
发展这类判据所需的关键定理见 [5]，再加上 Picard 概形的存在性定理。

\fsection{1}{交换群预概形的拓扑性质}

先设 $k$ 为域，$G$ 为 $k$ 上的群预概形。由于单位元 $e$ 在 $k$ 上
有理，所以它必为闭点；立即可知 $G\times_kG$ 的对角线是闭的，因而
$G$ 分离：域上的每个群预概形都是分离的。以 $G^\circ$ 表示单位元
$e$ 所在的连通分支。由于 $e$ 在 $k$ 上有理，$G^\circ$ 实际上是
几何连通的，亦即 $G^\circ$ 的形成与基域变换相容。还可推出
$G^\circ$ 作为集合在乘法下稳定；当 $G$ 局部诺特、因而
$G^\circ$ 为开集时，可把 $G^\circ$ 视为 $G$ 的开子群。以下假设
$G$ 在 $k$ 上局部有限型；则 $G^\circ$ 几何不可约，并在 $k$ 上
有限型。事实上，可以假设 $k$ 代数闭，进而还可假设 $G$ 约化
（因为此时 $G_{\mathrm{réd}}$ 是 $G$ 的子群，这是由于
$G_{\mathrm{réd}}\times_kG_{\mathrm{réd}}$ 仍然约化）；于是在 $k$
上的某个非空开集上光滑，而通过平移该开集可见它处处光滑。
这样，$G$ 局部不可约，所以其不可约分支就是其连通分支，故
$G^\circ$ 不可约。再令 $U$ 为 $G^\circ$ 中 $e$ 的一个仿射邻域；
利用 $G^\circ$ 的不可约性，立即可见
$U\cdot U=G^\circ$，这证明 $G^\circ$ 拟紧，因而在 $k$ 上有限型。

为简单起见，假设 $G$ 交换。对每个整数 $n>0$，令 $G^{(n)}$ 为
$G^\circ$ 在 $G$ 的 $n$ 次幂同态 $\varphi_n$ 下的逆像，因此
$G^{(n)}$ 是 $G$ 的开子群。置
\[
G^\tau=\bigcup_nG^{(n)}.
\]

\src{223}
\[
G^\sigma=\bigcup_{(n,p)=1}G^{(n)},
\qquad
G^\rho=\bigcup_hG^{(p^h)},
\]
其中 $p$ 是域 $k$ 的特征指数。这样得到 $G$ 的若干开子群，并满足
\[
G^\sigma\cap G^\rho=G^\circ,
\qquad
G^\sigma G^\rho=G^\tau.
\]

\result{注记}{} 可以构造商群概形 $G/G^\circ=N$（参见 [3]，IV），
进而把 $G^\tau$、$G^\sigma$、$G^\rho$ 定义为 $N$ 的挠子群
（相应地，后述准素分量的自然补，即各 $q$-准素分量之和，其中 $q$
遍历满足 $\neq p$ 的素数；再相应地，其 $p$-准素分量）在 $G$ 中的逆像。\footnote{印刷
文本将后两项的说明次序颠倒；依本页紧前给出的定义，译文按与特征指数互素部分、
特征指数准素部分的次序复原。}就此还应注意，$N$
是 $k$ 上的离散可分群概形，所以一旦选定 $k$ 的一个代数闭包
$\overline k$（从而得到一个 Galois 群 $\pi$），它便可与一个普通离散
群等同，而 $\pi$ 通过自同构作用于该群。由此可以显然地解释挠子群的
构造，以及它分解为各 $q$-准素分量的方式。当 $G$ 是 $k$ 上固有概形
$X$ 的 Picard 概形时，$N$ 可称为 $X$ 在 $k$ 上的（约化）
\emph{Néron--Severi 概形}。当
$G^\circ_{\mathrm{réd}}$ 是 $G^\circ$ 的群子概形时——特别地，只要
$k$ 完美，或者 $G^\circ$ 在 $k$ 上固有（例如 $X$ 几何正规），就会
如此——还应引入商
$N'=G/G^\circ_{\mathrm{réd}}$；从特化的观点看，亦即当 $X$ 在代数
概形族中变化时，它往往比 $N$ 表现得更好。

现在设 $S$ 为局部诺特预概形，$G$ 为 $S$ 上的群预概形，并且在 $S$ 上局部有限型。
我们不假设 $G$ 在 $S$ 上有限型，也不假设它在 $S$ 上分离。置
\[
G^\circ=\bigcup_{s\in S}(G_s)^\circ,
\]
并且当 $G$ 交换时，置
\[
G^\tau=\bigcup_{s\in S}(G_s)^\tau,
\qquad
G^\sigma=\bigcup_{s\in S}(G_s)^\sigma,
\qquad
G^\rho=\bigcup_{s\in S}(G_s)^\rho.
\]

\src{224}
这些是 $G$ 的一些在 $G$ 的乘法下稳定的子集；这显然并不意味着它们
能用 $G$ 的群子预概形来定义。特别地，一般似乎不存在底集恰为
$G^\circ$ 的 $G$ 的群子预概形。当然，如果其中某个集合是开的，那么
赋予诱导结构后，它就是 $G$ 的开群子预概形。我们将看到，
$G^\tau$ 总是如此。因此，从可表函子的观点，特别是从“特化”的观点
看，数值等价的表现比代数等价更为令人满意。下面列出刚才定义的这些
集合的主要一般性质：

\result{定理}{1.1} $G^\circ$、$G^\tau$、$G^\sigma$、$G^\rho$ 均局部
可构造。此外：

\noindent (i) $G^\circ$ 在 $S$ 上拟紧。若各 $G_s^\circ$ 均固有，且
$G$ 在 $S$ 上分离，则 $G^\circ$ 在 $S$ 上固有，因而在 $G$ 中为闭。

\noindent (ii) $G^\tau$ 是开的。若 $G^\circ$ 为闭，则 $G^\tau$ 也为
闭。

\noindent (iii) 若 $G^\circ$ 为闭，则在等特征情形，即 $S$ 的所有
剩余域具有同一特征时，$G^\sigma$ 也为闭。若 $G^\circ$ 为闭，并且
$G\longrightarrow S$ 在 $G^\sigma$ 的各点普遍开（参见下文推论
1.5），则 $G^\sigma\longrightarrow S$ 普遍开。

\noindent (iv) 若 $G^\circ$ 为闭，则 $G^\rho$ 也为闭。假设处于等特征
情形，并且对每个满足 $(n,p)=1$ 的整数 $n>0$，$G$ 中的 $n$ 次幂同态
都是开态射，则 $G^\rho$ 为开。

下面给出证明的要点。$G^\circ$ 的局部可构造性以及它在 $S$ 上的
拟紧性，包含在下述引理中：

\result{引理}{1.2} 假设 $S\neq\varnothing$，则 $S$ 中存在一个非空
开集 $U$，$U$ 上存在一个纤维连通的有限型群概形 $H$，并且存在群
概形同态 $H\longrightarrow G|U$，它是以 $G^\circ|U$ 为像的开浸入。

为证明该引理，可以假设 $S$ 不可约，并令 $\eta$ 为其一般点；作基变换
$S'=\operatorname{Spec}(\mathcal O_{S,\eta})\longrightarrow S$，得到
Artin 局部环 $\mathcal O_{S,\eta}=A$ 上的群概形 $G'$，其中有上文所说
的一个 $A$ 上有限型开群子概形 $G'{}^\circ$。因此，它来自 $\eta$ 的
某个开邻域 $U$ 上的有限型群概形 $H$，而典范浸入
$G'{}^\circ\longrightarrow G'$ 来自一个

\src{225}
开浸入 $H\longrightarrow G|U$；当 $U$ 取得足够小时，它是群概形同态。
若把 $U$ 取得足够小，则 $H$ 的纤维连通，且全都具有同一维数，等于
$G$ 的纤维维数（$U$ 足够小时）。因此，对每个 $s\in U$，$H_s$ 在 $G_s$ 中的像恰为
$G_s^\circ$（仍把 $U$ 取得足够小），这证明了 1.2。

(i) 的第二个断言包含在下述引理中（将它应用于 $G$ 中 $G^\circ$ 的
一个拟紧开邻域）：

\result{引理}{1.3} 设 $X$ 是局部诺特预概形 $S$ 上的有限型分离预概形，
$g$ 是 $X$ 在 $S$ 上的一个截面，$X^\circ$ 是各 $X_s$ 中包含
$g(s)$ 的连通分支之并。设 $s\in S$，且 $X_s^\circ$ 在 $k(s)$ 上
固有，则存在 $s$ 的一个开邻域 $U$，使 $X^\circ|U$ 在 $U$ 上固有，
进而在 $X|U$ 中为闭。

通过基的忠实平坦下降，可归结到 $S$ 是完备局部环的谱、$s$ 是其闭点
的情形。应用 [2]，III，5.5.1，可见 $X$ 分解为两个不交开集 $X'$ 与
$X''$ 之和；其中前者在 $S$ 上固有，且 $X'_s=X_s^\circ$。这把问题
归结到 $X=X'$ 的情形，亦即 $X$ 在 $S$ 上固有。此时可用一个标准证明
解决，即使用某个部分的固有性赋值判据（[2] 第二章中遗漏了这一判据）。

下面证明 $G^\tau$ 为开。考虑到 $G^\tau$ 的形成（正如 $G^\circ$、
$G^\sigma$、$G^\rho$ 的形成）与基扩张可交换，这等价于证明：对 $G$
在 $S$ 上的任意截面 $g$，$g^{-1}(G^\tau)$ 为开。这包含两件事：

\medskip
\noindent a. 设 $y\in S$ 且 $g(y)\in G_y^\tau$，则对每个满足
$y'\in\overline{\{y\}}$ 且充分接近 $y$ 的点，有
$g(y')\in G_{y'}^\tau$。

\medskip
\noindent b. 设 $y'\in S$ 且 $g(y')\in G_{y'}^\tau$，则对 $y'$ 的
每个泛化 $y$，有 $g(y)\in G_y^\tau$。

对于 (a)，注意存在整数 $n>0$，使 $g^n(y)\in G_y^\circ$。由于
$G^\circ$ 可构造，$(g^n)^{-1}(G^\circ)$ 也可构造，所以有
$g^n(y')\in G_{y'}^\circ$，其中 $y'\in\overline{\{y\}}$ 且充分接近
$y$。对于 (b)，注意诸
$g(y')^n=g^n(y')$ 都留在 $G_{y'}$ 的某个拟紧开集中（因为它们包含在
模 $G_{y'}^\circ$ 的有限多个类中），故 $G$ 中存在一个拟紧开集 $U$
包含所有 $g^n(y')$，从而也包含它们的各泛化 $g^n(y)$；因此
$g(y)$ 的各次幂都留在 $G_y$ 的一个拟紧开集中，这容易推出
$g(y)\in G_y^\tau$。

\src{226}
假设 $G^\circ$ 为闭；下面证明 $G^\tau$ 也为闭。由于已经知道
$G^\tau$ 为开，因而局部可构造，只须再证明它在特化下稳定。这源于
它是若干闭集之并，即闭集 $G^\circ$ 在各 $n$ 次幂同态
$\varphi_n$ 下的逆像之并。

一旦证明 $G^\sigma$ 与 $G^\rho$ 局部可构造，同一论证便证明：若
$G^\circ$ 为闭，则 $G^\sigma$ 与 $G^\rho$ 也为闭（前一种情形还须
处于等特征）。对 $x\in G^\tau$，令 $\nu(x)$ 为满足如下条件的最小
整数 $n>0$：$n$ 次幂同态 $\varphi_n$ 把 $x$ 映入 $G^\circ$。因此，
$G^\sigma$（相应地，$G^\rho$）由所有满足如下条件的 $x\in G$
\footnote{印刷文本此处作“$X$”，但本段中没有定义这个符号；所考察的四个
部分都是 $G$ 的子集。}构成：$\nu(x)$ 与 $p$ 互素（相应地，是 $p$
的幂）。于是，我们的可构造性断言由下述更精确的结论推出：

\result{引理}{1.4} $G^\tau$ 上的函数 $\nu$ 局部可构造。

这确实意味着：对每个整数 $n>0$，满足 $\nu(x)=n$ 的
$x\in G^\tau$ 所成的集合局部可构造；而该集合是
$\varphi_n^{-1}(G^\circ)$ 与诸
$\varphi_d^{-1}(G^\circ)$ 之并的差，其中 $d$ 遍历 $n$ 的真因子。
由于 $G^\circ$ 局部可构造，各 $\varphi_d^{-1}(G^\circ)$ 也局部
可构造，因而上述差集同样如此。

假设 $G^\circ$ 为闭，且 $G\longrightarrow S$ 在 $G^\sigma$ 的各点
普遍开；下面证明 $G^\sigma\longrightarrow S$ 为开，亦即它把
$G^\sigma$ 中 $x\in G^\sigma$ 的一个邻域映为 $y=f(x)$ 的一个邻域。
由于 $G^\sigma$ 局部可构造，只须证明：对 $y$ 的每个泛化 $y'$，都
存在 $G^\sigma$ 中 $x$ 的一个泛化 $x'$，位于 $y'$ 上方。经基变换，
可归结到 $S$ 是离散赋值环的谱，而 $y,y'$ 分别为闭点与一般点的
情形。利用 $G\longrightarrow S$ 在 $x$ 处为开这一事实（因而存在
$G$ 中 $x$ 的一个泛化 $x_1$，位于 $y'$ 上方），再作一次基变换后，
可以进一步假设存在 $G$ 在 $S$ 上的截面 $g$，使 $x=g(y)$。若
$k(y')$ 的特征为 0，只须取 $G$ 中位于 $y'$ 上方的 $x$ 的任意泛化
$x'$；由于 $G^\tau$ 为开，它属于 $G^\tau$，又由于
$G_{y'}^\sigma=G_{y'}^\tau$，它属于 $G^\sigma$。若 $k(y')$ 的特征为
$p>0$，置
\[
\nu(g(y'))=p^hm\qquad\text{其中 }(m,p)=1 ;
\]

\src{227}
取两个整数 $a$、$b$，使 $ap^h+bm=1$，并置
\[
g_1=g^{ap^h},\qquad g_2=g^{bm},\qquad\text{从而}\qquad g=g_1g_2 ;
\]
于是依构造有 $g_1(y')\in G^\sigma$、$g_2(y')\in G^\rho$。由于
$G^\circ$ 为闭，可推出 $g_1(y)\in G^\sigma$、$g_2(y)\in G^\rho$；
再由
\[
g(y)=g_1(y)g_2(y)\in G^\sigma,
\]
还可推出 $g_2(y)\in G^\sigma$，因而
\[
g_2(y)\in G_y^\sigma\cap G_y^\rho=G_y^\circ.
\]
由假设以及 $S$ 是离散赋值环的谱这一事实，
$G-(G_y-G_y^\circ)$ 是 $G$ 的一个开集，$G\longrightarrow S$ 在其上
诱导开态射；因此，存在一条“准截面”通过 $G_y^\circ$ 的每个点。于是，
再对基 $S$ 作一次扩张后，可以假设存在 $G^\circ$ 在 $S$ 上的截面
$g'_2$，使 $g'_2(y)=g_2(y)$。置 $g'=g_1g'_2$；则依构造，
$g'(y)=g(y)=x$，且
\[
g'(y')=g_1(y')g'_2(y')\in G^\sigma,
\]
所以 $g'(y')=x'$\footnote{印刷文本此处作“$g(y')=x'$”；上一句构造且
此处计算其一般纤维的截面是 $g'$。}是 $G^\sigma$ 中 $x$ 的一个
泛化，位于 $y'$ 上方。这证明了 (iii)。

最后证明 (iv) 的最后一个断言。可归结为证明：若 $x\in G^\rho$，则
$x$ 的每个泛化 $x'$ 都属于 $G^\rho$。可以假设（对适当的 $h$ 取
$\varphi_{p^h}$ 下的像后）\footnote{印刷文本作
$\varphi_{n^h}$；对 $x\in G^\rho$，第 223 页给出的定义提供了一个
特征指数 $p$ 的幂，把 $x$ 映入 $G^\circ$。}甚至有
$x\in G^\circ$。于是，对每个与该特征互素的整数 $n$，$x$ 都属于
$\varphi_n$ 的像（因为在特征与 $n$ 互素的域上，有限型连通群中的
$n$ 次幂映射是满射）。由于 $\varphi_n$ 为开，$x'$ 也属于
$\varphi_n$ 的像。更确切地说，设 $U$ 是 $G^\tau$ 中包含
$G_y^\circ$ 的一个拟紧开集，则对每个与 $p$ 互素的 $n$，都有
$x'\in\varphi_n(U)$。取 $n$ 为诸 $\nu(z)$ 中与 $p$ 互素的因子的
一个公倍数，其中 $z\in U$，可得 $x'\in G^\rho$。

(1.1) 的断言 (iii) 可补充如下：

\result{推论}{1.5} 设 $n>1$ 为整数，使 $n$ 次幂同态
$\varphi_n:G\longrightarrow G$ 普遍开（例如平展），置
\[
G^{(n)}=\varphi_n^{-1}(G^\circ),
\]
并假设各连通纤维 $G_s^\circ$ 不“包含任何加性成分”（亦即由

\src{228}
把域 $k(s)$ 扩张到代数闭包所得到的群不含同构于 $\mathbf G_a$ 的
子群）。则 $G\longrightarrow S$ 在各 $G^{(n^h)}$ 的点普遍开。特别
地，若各 $G_s^\circ$ 不含加性成分，并且对每个整数 $n>1$，同态
$\varphi_n$ 在剩余特征与 $n$ 互素的各点普遍开
\footnote{印刷文本作“与 $p$ 互素”。在这个由 $n$ 量化的命题中，
平展性判据以及使用核 ${}_nG$ 的论证要求剩余特征与 $n$ 互素。}，则
$G\longrightarrow S$ 在 $G^\sigma$ 的各点普遍开；因而由 (1.1)，
(iii)，若 $G^\circ$ 为闭，则
$G^\sigma\longrightarrow S$ 普遍开。在这些情形中，
$s\longmapsto\dim G_s$ 是 $S$ 上的局部常值函数。

事实上，由假设可知 $\varphi_n$ 的核 ${}_nG$ 在 $S$ 上普遍开，所以
$G\longrightarrow S$ 在 ${}_nG$ 的各点普遍开；因而以 $n^h$ 代替
$n$ 后，它在 ${}_{n^h}G$ 的各点普遍开。关于纤维 $G_s^\circ$ 的
假设恰好保证，$G_s$ 中阶为 $n$ 的幂的点在诸
$G^{(n^h)}$ 之并中稠密；由此容易推出，$G\longrightarrow S$ 在该并
的每个点普遍开，特别是沿单位截面普遍开，进而容易推出
$s\longmapsto\dim G_s$ 是局部常值函数。

\result{注记}{1.6} 回忆，若态射 $X\longrightarrow S$ 把每个开集映成
开集，并且在每次基变换 $S'\longrightarrow S$ 后仍保持这一可贵性质，
则称它普遍开。这也意味着（当 $S$ 局部诺特且
$X\longrightarrow S$ 局部有限型时），$X$ 的每个不可约分支都支配
$S$，并且这一性质在每次基变换 $S'\longrightarrow S$ 后都保持。在
这两个断言中（在上述有限性假设下），其实只须用如下基变换
$S'\longrightarrow S$ 作检验：$S'$ 是离散赋值环的谱（如果愿意，
还可要求完备且剩余域代数闭~……）。该定义显然可推广到 $X$ 的一个
子集 $Z$（例如 $G$ 的子集 $G^\sigma$）。即使从具有几何连通纤维的
射影光滑态射 $X\longrightarrow S$ 出发（例如
$S=\operatorname{Spec}\mathbf C[j]$ 上椭圆曲线模族的纤维平方），
$\operatorname{Pic}_{X/S}$ 在 $S$ 上并非普遍开，也是完全正常的现象；
亦即，$\operatorname{Pic}_{X/S}$ 可能有一些不可约分支完全位于 $S$
的某一个点之上。这与 $X/S$ 各纤维的 Néron--Severi 群的秩可能向上
跳跃有关（“复乘”现象）。相反，(1.5) 保证，在良好情形中，
$\operatorname{Pic}_{X/S}^\sigma$（而且看来往往甚至
$\operatorname{Pic}_{X/S}^\tau$）在 $S$ 上普遍开。


\src{229}
最后，给出一个有用的情形：在这一例外情形下，$G^\circ$ 确为开集：

\result{推论}{1.7} 设 $G\longrightarrow S$ 在 $G^\circ$ 的各点处普遍开
（参见 (1.5)），并且各纤维 $G_s$ 可分，因而在 $k(s)$ 上光滑（依 Cartier
的一个结果，后一条件在剩余特征为零时自动满足）。于是 $G^\circ$ 在 $G$
中是开的。若进一步 $S$ 是约化的，则 $G^\circ$ 在 $S$ 上光滑（特别地，
是平坦的）。

还可将第一个断言细化如下：若对某个给定的 $s\in S$，
$G\longrightarrow S$ 在 $G_s^\circ$ 的各点处普遍开，并且 $G_s^\circ$ 在
$k(s)$ 上可分，则 $G^\circ$ 是 $G_s^\circ$ 的一个邻域（而且，在 $s$
处所作的假设在邻近点处仍将成立）。这一陈述实际上与任何群结构无关，将在
[2], IV, 第 7 节中证明。(1.7) 的最后一个断言同样与任何群结构以及任何连通
分支问题无关，它是 [2], IV, 第 5 节中一个平坦性判据的特例；该判据更一般地
推出：

\result{推论}{1.8} 设 $U$ 是某一纤维 $G_s$ 的一个开部分，并且
$G\longrightarrow S$ 在 $U$ 的各点处普遍开（参见 (1.5)）。若 $G_s$ 在
$k(s)$ 上可分且 $S$ 在 $s$ 处约化，则 $G$ 在 $U$ 的各点处平坦于 $S$
（因而在这里，在 $U$ 的各点处光滑于 $S$）。

\result{注记}{1.9} 我不知道，当 $G$ 在 $S$ 上分离时，$G^\circ$ 或
$G^\tau$ 是否总在 $G$ 中闭；这似乎不大可能。无论如何，若去掉分离性假设，
很容易得到反例。例如，取一条仿射直线，并将其原点复制可数无穷多次；由此得
到仿射直线上的一个群概形 $G$，除一条纤维为 $\mathbf Z$ 外，其余各纤维均
约化为单位群。这个群概形是相应于一个退化为两条相交直线的二次曲线族的
$S$-概形 $X$ 的 Picard 预概形中的一个开群子概形，即单位截面的闭包。

此外，即使从离散赋值环 $V$ 的谱 $S$ 上一个有限平坦群概形 $G$ 出发，因而
$G^\circ$ 约化为单位截面并且是闭的，只要舍弃某些假设，各种结论仍会失效。
设 $V$ 具有等特征 $p>0$，并令 $G$ 为同态
$\mathbf G_a\longrightarrow\mathbf G_a$ 的核；该同态由函子同态
\[
f\longmapsto f^p-tf,
\]
定义，其中 $t$ 是 $V$ 的一个一致化元。

\src{230}
（回忆：依定义，$S$ 上的“加法群”$\mathbf G_a$ 表示函子
$S'\longmapsto\Gamma(S',\mathcal O_{S'})$。）于是，作为 $S$-概形，
\[
G=\operatorname{Spec}V[x]/(x^p-tx)
\]
是 $p$ 条相交“直线”的并，它刻画一个退化为加性型无穷小群的群
$\mathbf Z/p\mathbf Z$。在这个例子中，$G^\circ=G^\sigma$（并且它是这
$p$ 条直线之一），故在 $G$ 中不是开的。在剩余特征为 $p>0$ 的混合特征情形
下，从群概形 $H=\boldsymbol\mu_p$ 出发；它是 $\mathbf G_m$ 中 $p$ 次幂同态
的核，因而
\[
H=\operatorname{Spec}V[x]/(x^p-1) ;
\]
这仍是 $p$ 条相交直线的并，它刻画一个（在特征零时的）群
$\mathbf Z/p\mathbf Z$，后者在特征 $p$ 时退化为乘性型无穷小群。令 $H'$
为由通常的有限群 $\mathbf Z/p\mathbf Z$ 定义的“常值群概形”，即
$\mathbf Z/p\mathbf Z$ 份 $S$ 的不交并，也就是
$\operatorname{Spec}V^{\mathbf Z/p\mathbf Z}$。于是 $G=H\times_S H'$ 在特征
零时描述一个群 $(\mathbf Z/p\mathbf Z)^2$，它在特征 $p$ 时退化为一个无穷小
群与 $\mathbf Z/p\mathbf Z$ 的乘积。这里 $G^\rho$ 是单位截面与特殊纤维的
并，故 $G^\rho$ 不是开的；而 $G^\sigma$ 是单位截面与泛纤维的并，故不是闭
的，这与 (1.1) 的 (iii) 和 (iv) 中针对等特征情形所断言的结论相反。当然，
这些都是与特征 $p>0$ 有关的现象。前面的结果事实上给出：

\result{推论}{1.10} 设 $S$ 的所有剩余特征均为零，因而
$G^\tau=G^\sigma$ 且 $G^\circ=G^\rho$。于是 $G^\tau=G^\sigma$ 是开的；若
$G$ 在 $S$ 上分离且各 $G_s^\circ$ 均固有，则它甚至既开又闭。在同一假设
下，$G^\circ=G^\rho$ 在 $S$ 上固有，因而在 $G$ 中闭。最后，设对每个整数
$n>1$，$G$ 中的 $n$ 次幂同态均普遍开，则 $G^\circ$ 是开的；若进一步各
$G_s^\circ$ 没有加性分量（例如各 $G_s^\circ$ 如上均为固有的），则
$G^\tau\longrightarrow S$ 普遍开；若 $S$ 约化，它甚至是光滑的。

最后指出如下容易的结果：

\result{命题}{1.11} 存在 $S$ 的一个开部分 $U$，使得 $G$ 中那些 $G$ 在
$S$ 上光滑（相应地，平坦）的点所成的集合 $G'$，是 $G|U$ 的一个开群子概形
的底层集合。此外，$G$ 在 $U$ 上的每个截面都是 $G'$ 在 $U$ 上的截面。

\src{231}
\result{推论}{1.12} 若 $G$ 在单位截面的各点处光滑（相应地，平坦）于
$S$，则它在 $G$ 于 $S$ 上的每个截面的各点处以及 $G^\circ$ 的所有点处也
具有此性质。若进一步，对每个整数 $n>0$，$n$ 次幂同态
$\varphi_n:G\longrightarrow G$ 在特征与 $n$ 互素的各点处平展，则 $G$ 在
$G^\sigma$ 的所有点处光滑（相应地，平坦）于 $S$。

\fsection{2}{关于 Picard 概形局部性质的应用}

\result{定理}{2.1}

\noindent(i) 设 $f:X\longrightarrow S$ 是固有光滑态射，并设
$\operatorname{Pic}_{X/S}$ 存在（例如，$f$ 是射影态射）。则
$\operatorname{Pic}_{X/S}$ 在 $S$ 上分离，并且对于
$\operatorname{Pic}_{X/S}$ 的每个在 $S$ 上有限型的闭子集 $Z$，$Z$ 在
$S$ 上固有。

\noindent(ii) 设 $X$ 是域 $k$ 上固有且几何正规的预概形。则
$\operatorname{Pic}_{X/k}^\circ$\footnote{印刷文本作
$\operatorname{Pic}_{X/S}^\circ$，但在这一关于域 $k$ 的陈述中并未出现
$S$。} 在 $k$ 上固有。

\begin{proof}
(i) 赋值判据（[2], II, 第 7 节）将问题归结为证明如下事实：若 $S$ 是一个
完备离散赋值环的谱，$U$ 是只含 $S$ 的泛点的开集，则
$\operatorname{Pic}_{X/S}$ 在 $S$ 上的每个有理截面，即 $U$ 上的每个截面，
均能唯一延拓为 $S$ 上的截面。考虑到 $\operatorname{Pic}_{X/S}$ 的定义，
这等价于下述陈述：对于 $V=f^{-1}(U)$ 上的每个可逆模 $\mathcal L$，存在
$X$ 上的一个可逆模延拓 $\mathcal L$，且它在同构意义下唯一。由 $V$ 上以及
$X$ 上的可逆模分别以 Cartier 除子类描述这一事实，并注意到 $X$ 的各局部环
均正则（因为 $X$ 在正则概形 $S$ 上光滑），故依 Auslander 的结果均为唯一
分解环，这一点很容易推出；于是 $X$\footnote{印刷文本此处作“在 $S$ 上”；
这里所用的唯一分解性是 $X$ 的局部环的唯一分解性。} 上的每个除子都是
Cartier 除子。事实上，只需取闭包，$V$ 上的每个除子都可延拓为 $X$ 上的
除子。
\end{proof}

\result{注记}{2.2} 只假设 $f$ 平坦，并且其纤维 $X_s$ 局部为完全交且在
$k(s)$ 上于余维 $\leq2$ 处光滑时，上述证明仍然有效。这是因为 [5] 中证明
了如下结果：在余维 $\leq3$ 处正则的 Noether 局部完全交环是唯一分解环
（“Samuel 猜想”）。注意，这一

\src{232}
结果若将“余维 $\leq2$”替换为“余维 $\leq1$”，即替换为“正规”这一假设，
便不再成立；一个退化为二次锥的非奇异二次超曲面族即足以说明这一点。

(ii) 利用 Chow 引理，可归结到 $X$ 为射影概形的情形，因而嵌入
$\mathbf P_k^n$；可设 $X$ 连通。若 $\dim X=1$，则 $X$ 在 $k$ 上光滑，
从而应用 (i)。若 $\dim X\geq2$，可使用已知“等价判据”的一个变体；它推出
存在有限条 $X$ 上的光滑曲线 $Y_i$（它们由 $X$ 与 $\mathbf P_k^n$ 的适当
线性子空间相交而得），使得
\[
\operatorname{Pic}_{X/k}^\tau\longrightarrow
\prod_i\operatorname{Pic}_{Y_i/k}^\tau
\]
为单态射，因而在连通分支上当然也诱导单态射。由前述结果，右端的连通分支在
$k$ 上固有；而这里涉及的是群概形的同态，它必然是闭浸入。因此
$\operatorname{Pic}_{X/k}^\circ$ 在 $k$ 上也固有。也可以利用代数闭域上交换
代数群的结构（归功于 Chevalley--Borel）来避免诉诸精细的等价判据；问题归结
为证明从去掉原点的仿射直线\footnote{印刷文本作“仿射直线”。相应勘误
第~303 页规定为“去掉原点的仿射直线”。} 到
$\operatorname{Pic}_{X/k}^\tau$ 的每个态射均为常值。这等价于说，
$X[t,t^{-1}]$\footnote{印刷文本作 $X[t]$；相应勘误第~303 页规定为
$X[t,t^{-1}]$。} 上的每个可逆模均来自 $X$ 上的一个可逆模。这一结果众所周知
且十分初等，甚至不需要 $X$ 在 $k$ 上固有（$X$ 正规这一假设使问题立刻归结
到 $X$ 正则的情形）。

\result{推论}{2.3} 设 $f:X\longrightarrow S$ 是固有正规态射（即平坦且具有
几何正规纤维），并设 $\operatorname{Pic}_{X/S}$ 存在。则
$\operatorname{Pic}_{X/S}^\circ$ 在 $S$ 上固有，因而在
$\operatorname{Pic}_{X/S}$ 中闭；此外，$\operatorname{Pic}_{X/S}^\tau$ 和
$\operatorname{Pic}_{X/S}^\rho$ 也均为闭的；在“等特征”情形下，
$\operatorname{Pic}_{X/S}^\sigma$ 也为闭的。

应用 (1.1) 和 (2.1), (ii)。

\result{推论}{2.4} 设 $f:X\longrightarrow S$ 是固有光滑态射，使得
$\operatorname{Pic}_{X/S}$ 存在并且是 $S$ 上有限型概形 $P^{(i)}$ 的不交并
（参见 [3], V, 4.1）。则每个 $P^{(i)}$ 都在 $S$ 上固有。

这由 (2.1), (i) 得出。如 (2.2) 中所指出的，可通过对 $f$ 的纤维作较弱的
假设而推广这一结果；但若只假设 $f$ 正规，结论便不成立。在这种情形下，即使
假定 $\operatorname{Pic}_{X/S}^\tau$ 在 $S$ 上有限型，我也不知道它是否仍在
$S$ 上固有。

\src{233}
\result{定理}{2.5} 设 $f:X\longrightarrow S$ 是固有平坦态射，使得
$\operatorname{Pic}_{X/S}$ 存在；对每个整数 $n$，以 $\varphi_n$ 表示这一
群预概形中的 $n$ 次幂同态。则 $\varphi_n$ 在每个剩余特征与 $n$ 互素的点
$x\in\operatorname{Pic}_{X/S}$\footnote{印刷文本作“$x\in X$”；
$\varphi_n$ 是 Picard 预概形的幂自同态。} 处均为平展的。

由平展态射的无穷小刻画，这一断言等价于下述断言：

\result{引理}{2.6} 设 $S$ 是局部 Artin 环 $A$ 的谱，其极大理想
$\mathfrak m$ 的第 $(\nu+1)$ 次幂为零，并置
\[
A_{\nu-1}=A/\mathfrak m^\nu,
\qquad X_{\nu-1}=X\otimes_A A_{\nu-1},
\]
设 $\mathcal L$ 是 $X$ 上的可逆模，而 $\mathcal L'_{\nu-1}$ 是
$X_{\nu-1}$ 上的可逆模，其第 $n$ 次张量幂同构于
\[
\mathcal L_{\nu-1}=\mathcal L\otimes_A A_{\nu-1}.
\]
那么，存在 $X$ 上的可逆模 $\mathcal L'$，其第 $n$ 次张量幂同构于
$\mathcal L$（这里假定 $n$ 与 $k=k(A)$ 的剩余特征互素）。

证明该引理。置
\[
V=\mathfrak m^\nu=\mathfrak m^\nu/\mathfrak m^{\nu+1} ;
\]
这是 $k=k(A)$ 上的向量空间。先考察将 $\mathcal L'_{\nu-1}$ 延拓为 $X$
上的某个可逆模 $\mathcal L'$ 的问题。这样做的障碍位于
\[
H^2(X_0,\mathcal O_{X_0})\otimes_k V,
\]
但由关于 $\mathcal L'_{\nu-1}$ 的假设以及 $\mathcal L_{\nu-1}$ 可以延拓
这一事实可见，此障碍的 $n$ 倍为零；由于 $n$ 与特征互素，障碍本身也为零。
此外，所选延拓的任意性位于
\[
H^1(X_0,\mathcal O_{X_0})\otimes_k V,
\]
而 $\mathcal L'^{\otimes n}$ 与 $\mathcal L$ 之间的偏差 $\xi$ 也位于同一个
模中；若试图修正 $\mathcal L'$ 以使这一偏差为零，就归结为在该模中找到一个
$\eta$，使得 $n\eta=\xi$。由于 $n$ 与特征互素，这仍然可以做到。

\result{推论}{} 在 (2.5) 的条件下，进一步设各纤维 $X_s$ 的 Picard 概形
均不含加性分量（例如各 $X_s$ 几何正规，参见 (2.1), (ii)）。则
$\operatorname{Pic}_{X/S}\longrightarrow S$ 在
$\operatorname{Pic}_{X/S}^\circ$ 的各点处普遍开。若
$\operatorname{Pic}_{X/S}^\circ$ 是闭的（例如各 $X_s$ 几何正规，参见
(2.3)），则 $\operatorname{Pic}_{X/S}^\sigma$ 本身在 $S$ 上普遍开。最后，
在等特征情形下，$\operatorname{Pic}_{X/S}^\rho\longrightarrow S$ 普遍开。

应用 (1.5) 和 (1.1)。

\src{234}
\result{推论}{2.7} 设 $f:X\longrightarrow S$\footnote{印刷文本作
“$f:X\longrightarrow Y$”，但下面的 Picard 预概形和函数的基都是 $S$。}
是固有平坦态射，并设 $\operatorname{Pic}_{X/S}$ 存在。则 $S$ 上的函数
\[
s\longmapsto\dim\operatorname{Pic}_{X_s/k(s)}
\]
是上半连续的（即它可以向上跳跃，但不能向下跳跃）；若各
$\operatorname{Pic}_{X_s/k(s)}$ 均不含加性分量，则它甚至连续（即局部常值）。

第一个断言对于局部 Noether 基上局部有限型的任一群预概形而言是显然的，
或近乎显然的，因为只需沿单位截面考察即可。第二个断言由 (2.5) 得出。

\result{注记}{2.8} 设 $s,s'\in S$，且 $s$ 是 $s'$ 的一个特化。于是 (2.7)
等价于 $X_{s'}$ 的 Picard 概形与其“特化”$X_s$ 的 Picard 概形的维数之间的
一个不等式（相应地，一个等式）。事实上，早在 Picard 概形构造出来之前，
Serre 就已指出：在光滑态射 $f:X\longrightarrow S$ 的情形下，各 $X_s$ 的
Picard 概形维数不变性，是基本群的特化理论（[4], X）以及 Picard 簇上的有限
阶点与阿贝尔化基本群之间的经典 Kummer 关系（[4], XI）的形式推论。一般地，
以 $\alpha,\mu,\lambda$ 分别表示
$\operatorname{Pic}_{X_s/k(s)}^\circ$ 的阿贝尔部分、乘性部分和加性部分的
维数，并类似地定义 $\alpha',\mu',\lambda'$。已知关系可表为如下不等式：
\[
\tag{*}\alpha+\mu+\lambda\geq\alpha'+\mu'+\lambda'
\]
（只要 $\operatorname{Pic}_{X/S}$ 存在，该式就成立，因而它很可能总是成立）；
当 $\lambda=\lambda'=0$ 时，这个不等式化为一个等式，并在同样的存在性假设下
成立：
\[
\alpha+\mu=\alpha'+\mu' ;
\]
另一方面，有
\[
\tag{**}2\alpha+\mu\leq2\alpha'+\mu'
\]

\src{235}
根据 Serre 的论证，若各 $X_s$ 可分（甚至无须假设
$\operatorname{Pic}_{X/S}^\tau$ 存在），或者若
${}_n\operatorname{Pic}_{X/S}$（Picard 函子中 $\varphi_n$ 的核）在 $S$ 上
分离（注意到依 (2.5)，它们在 $S$ 上平展），则 (**) 成立。人们倾向于猜想
(*) 总是等式，或至少在各 $X_s$ 可分时如此，并且有不等式
\[
\tag{***}\alpha\leq\alpha',\qquad\lambda\geq\lambda',
\]
只要存在一个局部 Noether 的 $S$ 上局部有限型、且纤维维数恒定的群预概形，
这些不等式就应成立（这一方向上的一个肯定结果见 (1.3)）。

\result{注记}{2.9} 在所有已知情形中，$\operatorname{Pic}_{X/S}^\tau$ 在
$S$ 上普遍开；但即使 $f:X\longrightarrow S$ 光滑，大概也不应据此抱有过度
幻想。无论如何，Mumford 构造了一个
$\operatorname{Pic}_{X/S}^\tau$ 不在 $S$ 上平坦的例子（诚然其中 $S$ 非约化，
事实上 $S$ 是一个 Artin 环的谱），其方法是对 Igusa 曲面作无穷小变形。
Mumford 所考察的点实际上位于 $\operatorname{Pic}_{X/S}^\rho$ 中；并且仍有可能
（当 $f:X\longrightarrow S$ 光滑时）$\operatorname{Pic}_{X/S}^\tau$ 在
$\operatorname{Pic}_{X/S}^\sigma$ 的各点处平坦于 $S$。不过，报告人怀疑这是否
总是成立，即便只限于 $\operatorname{Pic}_{X/S}^\circ$ 的各点，并假定 $S$ 是
离散赋值环的谱。这一问题还与有限自同构群作用下阿贝尔概形的不动点研究有关；
在这种情形下似乎缺少例子。

看来，即使只考虑光滑射影态射 $f:X\longrightarrow S$，本节关于
$\operatorname{Pic}_{X/S}$ 的局部正则性结果以及 (2.8) 中指出的猜想，也几乎
穷尽了在不对 $f$ 的纤维性质作更特殊假设时所能说的一切。不过仍需回忆：若
$\operatorname{Pic}_{X/S}$ 的几何纤维均约化且不含加性分量，则由 (1.8) 和
(2.5) 可知，当 $S$ 约化时，$\operatorname{Pic}_{X/S}$ 在
$\operatorname{Pic}_{X/S}^\sigma$ 的各点处光滑于 $S$；若
$f:X\longrightarrow S$ 正规，则即使不对 $S$ 作任何假设，这一结果仍成立，
如我们将在 (3.5) 中看到的。为此指出：

\result{命题}{2.10}

\noindent(i) 若 $\operatorname{Pic}_{X/S}$ 在单位截面的各点处光滑（相应地，
平坦）于 $S$，则它在 $\operatorname{Pic}_{X/S}^\sigma$ 的所有点处以及
$\operatorname{Pic}_{X/S}$ 于 $S$ 上任一截面的各点处也具有此性质。

\src{236}
\noindent(ii) 设 $s\in S$ 满足
\[
H^2(X_s,\mathcal O_{X_s})=0,
\]
则存在 $s$ 的一个开邻域 $U$，使得 $\operatorname{Pic}_{X/S}|U$ 在 $U$ 上
光滑。

\noindent(iii) 设 $X$ 是域上的固有概形，则有
\[
\dim\operatorname{Pic}_{X/k}\leq\dim H^1(X,\mathcal O_X),
\]
等号成立当且仅当 $\operatorname{Pic}_{X/k}$ 在 $k$ 上光滑；若 $k$ 的特征为
零，等号总是成立。

(i) 由 (2.5) 和 (1.12) 得出；(ii) 由光滑态射的无穷小判据和一个熟知的障碍
计算得出，同时要注意，（由“半连续性定理”）在 $s$ 处所作的假设在邻近点处
仍然成立。最后，(iii) 由 $H^1(X,\mathcal O_X)$ 同构于
$\operatorname{Pic}_{X/k}$ 在单位元处的 Zariski 切空间这一事实得出；最后一个
断言是 Cartier 的一个定理的特例，该定理称特征零的“形式群”在 $k$ 上形式
光滑。

\fsection{3}{\texorpdfstring{$\operatorname{Pic}_{X/S}$}{Picard 概形} 的典范阿贝尔子概形与 Albanese 概形}

\result{命题}{3.1} 设 $k$ 是域，$G$ 是 $k$ 上有限型的交换群概形，并且
“不含加性分量”。则 $G_{\mathrm{réd}}^\circ$ 在 $k$ 上可分，因而是 $k$ 上的
光滑群子概形。

当 $k$ 完美时，这是显然的；特别地，对 $G_{\bar k}$ 亦然，其中 $\bar k$
是 $k$ 的代数闭包。因此只需说明 $(G_{\bar k}^\circ)_{\mathrm{réd}}$ 下降自
$G$ 的一个子概形。由关于 $G_{\bar k}$ 的假设（不含加性分量）\footnote{印刷
文本作“Or l'hypothèse ...”；相应勘误第~303 页规定为“Or, de l'hypothèse
...”。} 很容易推出，存在一个整数 $m$，使得
$(G_{\bar k}^\circ)_{\mathrm{réd}}$ 是 $G_{\bar k}$ 中 $m$ 次幂同态的概形论像。
由于后一个同态来自 $G$ 上的相应同态，后者的概形论像就是所求子概形。

\result{推论}{3.2} 设 $X$ 是 $k$ 上的几何正规固有概形，并且
$\operatorname{Pic}_{X/k}$ 存在。则存在 $\operatorname{Pic}_{X/k}$ 的一个
阿贝尔子概形 $A$，其底层集合是 $\operatorname{Pic}_{X/k}^\circ$。

事实上，由 (2.1), (ii)，$\operatorname{Pic}_{X/k}^\circ$ 在 $k$ 上固有，
因而满足 (3.1) 的条件。


\src{237}
因此，前述结果表明，在某些情形下，经典的“Picard 簇”（在现代理论中应记为
$(\operatorname{Pic}_{X/k})_{\mathrm{réd}}^\circ$）“定义在 $k$ 上”，而不必假设域
$k$ 是完美域。

现设 $f:X\longrightarrow S$ 为一个固有且平坦的态射。为简单起见，设有
\[
\mathcal O_S\xrightarrow{\sim}f_*(\mathcal O_X)
\]
并假设 $\operatorname{Pic}_{X/S}$ 存在，且
$\operatorname{Pic}_{X/S}^\circ$ 在 $S$ 上固有。对 (3.3) 的 (ii)，我们还假设
$\operatorname{Pic}_{X/S}$ 中存在一个包含 $\operatorname{Pic}_{X/S}^\circ$、
且在 $S$ 上拟射影的开集；前面已经看到，当 $f$ 是射影态射，且其几何纤维可分、
不可约时，这一条件成立。回忆：$S$ 上的阿贝尔概形，是指一个在 $S$ 上固有、
光滑并具有几何连通纤维的 $S$-群概形。我们要考察，$\operatorname{Pic}_{X/S}$
中是否存在一个阿贝尔子概形 $A$，其底层集合为
$\operatorname{Pic}_{X/S}^\circ$。刚才已经看到，当 $S$ 是某个域的谱时，
这样的子概形总是存在。在这里所考虑的一般情形下，已知如下结论：

\result{定理}{3.3} 在上述条件下：

\noindent(i) 如果 $\operatorname{Pic}_{X/S}$ 中存在一个底层集合为
$\operatorname{Pic}_{X/S}^\circ$ 的阿贝尔子概形，那么它是唯一的。因此，
它的形成与基变换相容。

\noindent(ii) 为使这样的阿贝尔子概形存在，只需它在每个基变换
$S'\longrightarrow S$ 后都存在，其中 $S'$ 为局部 Artin 概形；如果 $S$ 是一个
局部环的谱，那么甚至只需对 $S'=\operatorname{Spec}(A_n)$ 检验，其中
$A_n=A/\mathfrak m^{n+1}$。如果 $S$ 是约化的，那么同样只需对基变换
$S'\longrightarrow S$ 检验，其中 $S'$ 是离散赋值环的谱（如果愿意，可要求
该环完备且剩余域代数闭）。

\noindent(iii) 假设 $A$ 存在，并令 $B=\operatorname{Alb}^0(X/S)$ 为其对偶
阿贝尔概形（即 $B=\operatorname{Pic}_{A/S}^\circ$ [7]）。那么可以典范地构造
$B$ 下的主齐性空间 $P=\operatorname{Alb}^1(X/S)$，以及一个 $S$-态射
$X\longrightarrow P$；它对从 $X$ 到 准阿贝尔概形（即阿贝尔概形下的
主齐性空间）的 $S$-态射是泛的。$\operatorname{Alb}^0(X/S)$、
$\operatorname{Alb}^1(X/S)$ 以及 $X\longrightarrow\operatorname{Alb}^1(X/S)$
的形成都与基变换相容。

我们概述其证明：

(i) 是交换群概形的阿贝尔子概形所具有的一般刚性性质：如果两个这样的子概形
在一点 $s\in S$ 上集合论地重合，那么它们在 $s$ 所在的整个

\src{238}
连通分支上都重合 [7]。（这一结果推广了 Chow 的一个经典定理。）

(ii) 利用 Hilbert 概形可知，下述函子可由一个 $S$ 上有限型概形 $T$ 表示：
它把每个 $S$ 上的 $S'$，映到 $(\operatorname{Pic}_{X/S})\times_S S'$ 的
典范阿贝尔子概形所成的集合（该集合有一个或零个元素）。由 (i)，
$T\longrightarrow S$ 是单态射；由 (3.2)，它是满射。说
$\operatorname{Pic}_{X/S}$ 存在典范阿贝尔子概形，就是说 $T$ 在 $S$ 上有
一个截面，也就是说 $T\longrightarrow S$ 是同构。这也等价于说
$T\longrightarrow S$ 是平展态射；而当 $S$ 约化时，还等价于说
$T\longrightarrow S$ 是固有态射。由此立即得到 (ii)。

(iii) 仅利用 $\operatorname{Pic}_{X/S}$ 的定义即可看出：对于 $S$ 上的任意
阿贝尔概形 $C$，给定一个从 $X$ 到 $C$ 下主齐性空间的 $S$-态射，等价于给定
一个群同态 $C'\longrightarrow\operatorname{Pic}_{X/S}$，其中 $C'$ 是 $C$
的对偶阿贝尔概形。现在，如果 $\operatorname{Pic}_{X/S}$ 的典范阿贝尔
子概形 $A$ 存在，这些同态必然经由 $A$ 分解（也可再次利用有限阶点看出）。
由此立即得到 (iii)。

\result{注}{3.4} 如果 $\operatorname{Pic}_{X/S}$ 的典范阿贝尔子概形存在，
我们把它记为 $\operatorname{Pic}_{X/S}^{\circ\circ}$。遗憾的是，它并不总是
存在；把 Igusa 曲面作无穷小变动（通过一阶模变形）便可看出这一点。不过，
$\operatorname{Pic}_{X/S}^{\circ\circ}$ 至少在 $S$ 约化时仍可能存在；由
(ii)，这等价于 $S$ 是一个离散赋值环的谱时它存在。现设 $X_0,X_1$ 分别为
$X$ 的特殊纤维和泛纤维，并令
\[
A_1=\operatorname{Pic}_{X_1/K}^\circ,
\]
其中 $K$ 是赋值环 $V$ 的分式域。由 Koizumi 的定理，$S$ 上存在一个本质唯一的
阿贝尔概形 $A$，其泛纤维为 $A_1$；与 (2.1) 的 (i) 中一样，很容易看出
（从现在起假设 $X$ 在 $S$ 上光滑），$A_1$ 的恒等态射可延拓为一个态射
\[
A\longrightarrow\operatorname{Pic}_{X/S}.
\]
从而得到同态
\[
\tag{*}A_0\longrightarrow\operatorname{Pic}_{X_0/k}^{\circ\circ},
\]
容易证明它是满同态，其核为有限 $p$-群，其中 $p$ 是剩余域 $k$ 的特征指数

\src{239}
（仍然利用有限阶点）。在此约定下，关于 $X/S$ 的下列条件等价：

\medskip
\noindent a. 上述同态 (*) 是同构（Shimura 会把这一点表述为“Picard 簇”的
形成“与特殊化相容”）。

\noindent b. $\operatorname{Pic}_{X/S}^{\circ\circ}$ 存在（此时它正是 $A$）。

\noindent c. （备考）各 $\operatorname{Pic}_{X_n/S}^{\circ\circ}$ 均存在。

根据前面对 (*) 的核所作的说明，当剩余特征为零时，条件 (a) 成立；但这一结果
将在 (3.5) 中得到显著推广。

当然，如果 $\operatorname{Pic}_{X/S}$ 在 $\operatorname{Pic}_{X/S}^\circ$
的各点处相对于 $S$ 光滑，那么后者在 $\operatorname{Pic}_{X/S}$ 中是开的
（参见 (1.7)）；因此，赋予诱导结构后，它就是 $\operatorname{Pic}_{X/S}$
的一个阿贝尔子概形，因而等于此时确实存在的
$\operatorname{Pic}_{X/S}^{\circ\circ}$。但还有强得多的结论：

\result{定理}{3.5} 在 (3.3) 的条件下，设 $s\in S$，并且
$\operatorname{Pic}_{X_s/k(s)}$ 在 $k(s)$ 上光滑（等价地，
\[
\dim\operatorname{Pic}_{X_s/k(s)}=\dim H^1(X_s,\mathcal O_{X_s})).
\]
那么存在 $s$ 的一个开邻域 $U$，使得 $\operatorname{Pic}_{X/S}$ 在
$\operatorname{Pic}_{X/S}^\circ|U$ 的各点处相对于 $S$ 光滑；因而后者是
$\operatorname{Pic}_{X/S}|U$ 的一个开阿贝尔子概形。更进一步，
$\operatorname{Pic}_{X|U/U}^{\circ\circ}$ 存在。

下面说明证明的原理。前面的讨论把问题归约到 $S$ 是一个局部 Artin 环 $A$
的谱的情形；我们对 $A$ 的无穷小阶数作归纳。因此可以假设
$\operatorname{Pic}_{X_n/A_n}^\circ$ 在 $A_n$ 上光滑，而只需证明
$\operatorname{Pic}_{X_{n+1}/A_{n+1}}^\circ$ 在 $A_{n+1}$ 上光滑。为此，
只要构造一个 $A_{n+1}$ 上的阿贝尔概形 $P_{n+1}$，使它延拓
\[
P_n=\operatorname{Pic}_{X_n/A_n}^\circ,
\]
同时再构造 $X_{n+1}\times_{A_{n+1}}P_{n+1}$ 上的可逆模
$\mathcal L_{n+1}$，使它延拓 $X_n\times_{A_n}P_n$ 上的可逆模
$\mathcal L_n$；后者作为一个泛问题的解，出现在 Picard 概形
$\operatorname{Pic}_{X_n/A_n}$ 的定义中。（注意——我们可以假设 $X$ 在
$S$ 上带有一个截面……）为完成这一构造，必须使用如下关键结果：定义在局部
Artin 环的商上的每个阿贝尔概形都可以延拓（换言之，代数闭域上的阿贝尔
概形的绝对“形式模概形”([3], II) 在

\src{240}
$k$ 的 Witt 向量环上光滑）；这一结果本身可由提升障碍的一般形式性质和群运算
得到。
承认这一结果后，先任意把 $P_n$ 延拓为 $P_{n+1}$；于是会遇到提升
$\mathcal L_n$ 的一个障碍，该障碍位于
\[
H^2(X_0\times P_0,\mathcal O_{X_0\times P_0})\otimes_k V
\qquad(\text{其中 }V=\mathfrak m^{n+1}/\mathfrak m^{n+2}),
\]
而且实际上位于子空间
\[
H^1(X_0,\mathcal O_{X_0})\otimes
H^1(P_0,\mathcal O_{P_0})\otimes_k V
\]
中（这是因为 $\mathcal L_n$ 在两个因子 $X_n$ 和 $P_n$ 上的限制均平凡）。
而这后一空间正是
\[
H^1(P_0,\mathcal T_{P_0/k})\otimes V,
\]
其中 $\mathcal T_{P_0/k}$ 是 $P_0/k$ 的切层；因此，它也正是刻画把 $P_n$
提升为 $P_{n+1}$ 时所具有的不定性的空间（[3], II）。所以，可以修正这一
提升（而且按理应当恰有一种修正方式），从而消去提升 $\mathcal L_n$ 的障碍。

\result{推论}{3.6} 在 (3.5) 的条件下，$R^1f_*(\mathcal O_X)$ 在 $s$ 附近
是 $S$ 上的局部自由模，且它的形成与基变换相容。

事实上，这个模就是 $\operatorname{Pic}_{X/S}$ 沿单位截面的切模。

\result{注}{3.7} 利用与 (3.5) 相同的论证，可以证明：如果 $S'$ 是由一个
凝聚幂零理想定义的 $S$ 的子概形，并且仅假设
$\operatorname{Pic}_{X'/S'}^\circ$ 存在且具有光滑纤维，那么
$\operatorname{Pic}_{X/S}^\circ$ 必然存在，并且是 $S$ 上的阿贝尔概形。
这使我们在缺少射影性假设的某些情形下仍能构造 Picard 概形；例如，Artin 环
上的阿贝尔概形的对偶阿贝尔概形总是存在。另一方面，当 $X$ 是 $S$ 上的
阿贝尔概形时应用 (3.6)，并利用 $H^*(X_s,\mathcal O_{X_s})$ 是由
$H^1(X_s,\mathcal O_{X_s})$ 生成的外代数这一已知结构（Rosenlicht--Serre），
可得对每个 $i$，$R^if_*(\mathcal O_X)$ 都是局部自由的；更确切地说，
它同构于 $R^1f_*(\mathcal O_X)$ 的第 $i$ 次外幂。

\result{注}{3.8} 当 $f:X\longrightarrow S$ 是光滑射影态射，$S$ 约化且所有
剩余特征均为零时，结果 (3.6) 作为 Hodge 理论的推论，早已由超越方法得到。
事实上，此时所有 $R^pf_*(\Omega^q_{X/S})$ 都是局部自由的。另一方面，存在反例

\src{241}
表明在混合特征情形下 $R^1f_*(\mathcal O_X)$ 未必局部自由；这些反例由 Serre
簇得到（[8]，第 50 页）。在等特征情形下似乎尚无反例。

\result{推论}{3.8}\footnote{印刷文本在紧前一条注记之后仍将本推论编号为
3.8；译文保留原编号。} 在 (3.5) 的条件下，$\operatorname{Pic}_{X/S}|U$ 在
$\operatorname{Pic}_{X/S}^\sigma|U$ 的每一点处相对于 $U$ 光滑。

应用 (2.10) 的 (i)。特别地，再计及 (2.10) 的 (iii)，得到：

\result{推论}{3.9} 假设 $S$ 的所有剩余特征均为零。那么
$\operatorname{Pic}_{X/S}^\tau$ 在 $S$ 上光滑。

\result{注}{3.10} 例如，如果 $\operatorname{Pic}_{X/S}^\tau$ 在 $S$ 上固有，
则由此可知，$f$ 的各几何纤维的 Néron--Severi 群的挠子群，在 $S$ 的每个
连通分支上都保持不变（当 $f$ 光滑且射影时，这一点也可明显地由超越方法看出）。
还应指出，直接应用 (2.5) 可以更一般地证明：如果
$\operatorname{Pic}_{X/S}^\tau$ 在 $S$ 上固有（例如
$f:X\longrightarrow S$ 光滑且射影），而 $q$ 是不同于 $S$ 的所有剩余特征
的素数，那么 $X/S$ 的各几何纤维的 Néron--Severi 群的挠子群之 $q$-准素
分量，在 $S$ 的每个连通分支上都保持不变。不过，在特征 $p>0$ 时，对挠子群
的 $p$-准素分量，结论不再成立。尽管如此，
\[
\operatorname{Pic}_{X_s/k(s)}^\tau/
\operatorname{Pic}_{X_s/k(s)}^{\circ\circ}=T_{X_s/k(s)}
\]
在 $k(s)$ 上的秩仍可能局部恒定；当 $S$ 约化时，可以证明，这等价于
$\operatorname{Pic}_{X/S}^{\circ\circ}$ 存在且
$\operatorname{Pic}_{X/S}^\tau$ 在 $S$ 上平坦；并且只需在 $S$ 是离散赋值环
的谱时检验这一问题。这正是我在所考察的少数例子中验证过的情形；不过，因为
$S$ 为 Artin 概形时相应的陈述为假（参见注 (2.9) 和 (3.4)），大概不应抱有
过度乐观的幻想。

\fsection{4}{Picard 概形的有限性定理}

设 $f:X\longrightarrow S$ 是一个射影平坦态射，并且
$\operatorname{Pic}_{X/S}$ 存在。考察“Hilbert 多项式”
$Q\in\mathbf Q[t]$，可以把 $\operatorname{Pic}_{X/S}$ 分解为开集
$\operatorname{Pic}_{X/S}^Q$ 的不交并。如果不作附加假设，例如不假设
$\operatorname{Pic}_{X/S}$ 在 $S$ 上分离，那么这些开集一般未必在 $S$ 上
为有限型；当 $X$ 是“一条退化为两条直线的二次曲线”时即得到一个反例。不过，
仍有

\src{242}
可能在 $f$ 可分且具有几何不可约纤维时，它们确实为有限型。这个问题与
$\operatorname{Pic}_{X/S}^\tau$ 是否在 $S$ 上为有限型有关；后者可能无需对
$X/S$ 的纤维作任何假设便成立。当 $f:X\longrightarrow S$ 光滑时，注意到
$\operatorname{Pic}_{X/S}^\tau$ 包含在某个 $\operatorname{Pic}_{X/S}^Q$
中（因为在非奇异射影簇上，除子的“挠等价”比数值等价更细——事实上，由
Matsusaka 的结果，两者“相同”）；所以，如果各
$\operatorname{Pic}_{X/S}^Q$ 在 $S$ 上为有限型，那么前者亦然。还应注意，
即使不假设
$\operatorname{Pic}_{X/S}$ 存在，刚提出的有限性问题仍有显然的意义，因为
它们等价于说，某些可逆模族在 [3], IV 的意义下是“有界”的：对每个代数闭域
$k$，考虑 $\mathbf P_k^r$ 中维数为 $n$、次数为 $d$ 的整子概形，以及这种
预概形上 Hilbert 多项式为 $Q$ 的可逆模，其中 $r,n,d,Q$ 固定；问题是证明
这些模所成的族（把它们视为各 $\mathbf P_k^r$ 上的凝聚模）有界，即可由一个
$\mathbf Z$ 上的有限型概形参数化。

利用 Matsusaka 的方法 [6]，通过一个颇具技术性的证明（以适当形式使用
“等价性判据”），可以在只考虑 $\mathbf P_k^r$ 的非奇异子簇时给出肯定答案。
更确切地说，有如下结果：

\result{定理}{4.1} 设 $f:X\longrightarrow S$\footnote{印刷文本作“设：
$X\longrightarrow S$”；下文所用的态射是 $f$。} 是一个具有几何连通纤维的
光滑射影态射，$S$ 是 Noether 概形；设 $\mathcal O_X(1)$ 是 $X$ 上相对于
$S$ 的甚丰沛模，$E$ 是 $\operatorname{Pic}_{X/S}$ 的一个子集，对应于
$X/S$ 的几何纤维 $\overline X_{s_i}$ 上的一族可逆模 $(\mathcal L_i)$；
设 $D_i$ 是 $\overline X_{s_i}$ 上定义 $\mathcal L_i$ 的一个除子（不必为
有效除子），
\[
a_n^{(i)}t^n+\cdots+a_0^{(i)}
\]
是 $\mathcal L_i$ 的 Hilbert 多项式；设 $\xi=\xi_i$ 是定义
$\mathcal O_X(1)$ 的一个除子，即一个超平面截面。下列条件等价：

\noindent a. $E$\footnote{印刷文本此处作“$Q$”；陈述中引入的子集是 $E$，
其拟紧性正表示该族有界。} 是拟紧的，即包含在 $S$ 上的某个有限型开集中；
也就是说，族 $(\mathcal L_i)$ 有界。

\noindent b. $E$ 包含在有限多个集合 $\operatorname{Pic}_{X/S}^Q$ 的并中，
$Q\in\mathbf Q[t]$；也就是说，$\mathcal L_i$ 的 Hilbert 多项式只取有限多个值。

\noindent b'. （如果 $X/S$ 的所有纤维都具有相同的维数 $n$。）
$\mathcal L_i$ 的 Hilbert 多项式的两个系数 $a_{n-1}^{(i)}$ 和
$a_{n-2}^{(i)}$ 都只取有限多个值。

\src{243}
\noindent b''. （如果 $X/S$ 的所有纤维都具有相同的维数 $n$。）整数
\[
\xi^{n-1}D_i\qquad\text{和}\qquad\xi^{n-2}D_i^2
\]
（$D_i$ 和 $D_i^2$ 的射影次数）有上界。

\result{推论}{4.2} 设 $f:X\longrightarrow S$ 是一个具有几何连通纤维的
光滑射影态射，那么概形 $\operatorname{Pic}_{X/S}^Q$ 和
$\operatorname{Pic}_{X/S}^\tau$（$Q\in\mathbf Q[t]$）在 $S$ 上都是射影的。

由 (4.1)，它们在 $S$ 上为有限型，因此可以应用 (2.4)。

\section*{参考文献}
\addcontentsline{toc}{section}{参考文献}

\begin{description}
\item[{[1]}] CHEVALLEY (Claude). -- Sur la théorie de Picard,
\emph{Amer. J. of Math.}, 第 82 卷，1960 年，第 435--490 页。

\item[{[2]}] GROTHENDIECK (A.) et DIEUDONNÉ (J.). -- Éléments de géométrie
algébrique, 第 I 章及以下各章。-- 巴黎，Presses universitaires de France，
1960--1961 年（Institut des Hautes Études Scientifiques，4, 8, 11,
\ldots）。

\item[{[3]}] GROTHENDIECK (Alexander). -- Technique de descente et théorèmes
d'existence en géométrie algébrique, I--V, \emph{Séminaire Bourbaki}, 第 12 卷，
1959/60 年，n\textsuperscript{o} 190，29 页及 n\textsuperscript{o} 195，22 页；
第 13 卷，1960/61 年，n\textsuperscript{o} 212，20 页及
n\textsuperscript{o} 221，28 页；第 14 卷，1961/62 年，
n\textsuperscript{o} 232，19 页。

\item[{[4]}] GROTHENDIECK (Alexander). -- Séminaire de géométrie algébrique,
1960/61. -- 巴黎，Institut des Hautes Études Scientifiques，1961 年。

\item[{[5]}] GROTHENDIECK (Alexander). -- Séminaire de géométrie algébrique,
1961/62, 由一组听课者整理。-- 巴黎，Institut des Hautes Études Scientifiques
（待刊）。

\item[{[6]}] MATSUSAKA (T.). -- The criteria for algebraic equivalence and the
torsion group, \emph{Amer. J. Math.}, 第 79 卷，1957 年，第 53--66 页。

\item[{[7]}] MUMFORD (D.) and TATE (J.). -- Séminaire de géométrie algébrique,
Harvard University, Spring Term 1962（待刊）。

\item[{[8]}] SERRE (Jean-Pierre). -- Sur la topologie des variétés algébriques
en caractéristique $p$, \emph{Symposium internacional de topología algebraica}
[1956, Mexico], 第 24--53 页。-- Mexico, Universidad nacional autónoma，1958 年。
\end{description}

% FGA Canon R30: scan restoration and explicit editorial apparatus.
\par\smallskip
\begingroup\footnotesize
\noindent\textit{FGA-CANON-DISP-0056 至 0058。3.2 的印刷本假设 $X$ 在 $k$ 上正规且固有，未要求几何正规（C122）。3.10 的两处印刷本固有性前提均指整个 Picard 概形，而非 $\operatorname{Pic}_{X/S}^{\tau}$（C124）。C123 恢复 2.10(i) 与 3.8 中印刷的 $\sigma$；圆圈与 $\tau$ 是转录错误，而非印刷读法。}\par
\endgroup

% END 236.tex

% BEGIN com.tex
\unit{1962年评注}
\sid{com}
\src{297}
\seminarzh{第14年度，1961/62，补编\\1962年5月}
\paperhead{代数几何基础\\评注}%
  {亚历山大·格罗滕迪克 著}

\section*{说明}
\addcontentsline{toc}{section}{说明}

1957年至1962年，我们在布尔巴基讨论班作了八次涉及代数几何基础的
报告。除第一篇外，这些报告均在概形的框架中撰写。所陈述的一切结果都将
收入让·迪厄多内与亚历山大·格罗滕迪克的《代数几何原本》
（\emph{Éléments de géométrie algébrique}）。然而，目前这些报告中
任何一篇的实质内容，都尚未由已经出版或正在编写的任何一章涵盖，也未由
任何其他书籍或论文涵盖；而且这种情形大概还会持续数年。这正是促使我们
汇集这些报告的主要原因：在正式定稿问世以前，先向使用者提供概形理论的
若干关键概念和结果。此外，阅读这些文字可以使人迅速熟悉上述结果与概念，
而不必受系统论著中必然繁复的细节所累，并将为研读这样的论著提供不可或缺
的动机。

为方便读者，我们在此汇集了若干按报告分组的评注与勘误，其中尤其会顺带
说明自原文写成以来取得的进展，并给出一些补充参考文献。

这些报告中的若干结果，已在 IHES 代数几何讨论班以及1961/62年度哈佛大学
连续举行的两个讨论班中得到详细论述；前一个由我主讲，后一个由
MUMFORD--TATE 主讲。相应的油印讲义（由 LICHTENBAUM 先生撰写）目前正在
编写之中。

\section*{所用缩写}
\addcontentsline{toc}{section}{所用缩写}

\begin{description}
\item[SGA :] 代数几何讨论班，法国高等科学研究所，1960--1962年。

\src{298}
\item[SHG :] GROTHENDIECK 讨论班，哈佛大学，1961/62年度。

\item[SHMT :] MUMFORD--TATE 讨论班，哈佛大学，1961/62年度。

\item[EGA :] GROTHENDIECK (A.) 与 DIEUDONNÉ (J.)，
《代数几何原本》（\emph{Éléments de géométrie algébrique}）。
巴黎：法国大学出版社，1960、1961、……（法国高等科学研究所，
《数学出版物》第4、8、11、……号）。

\item[TDTE :] GROTHENDIECK (Alexander)，《代数几何中的下降技术与
存在性定理》（\emph{Technique de descente et théorèmes d'existence
en géométrie algébrique}）I--VI，\emph{Séminaire Bourbaki}，第12卷，
1959/60年度，第190、195号；第13卷，1960/61年度，第212、221号；第14卷，
1961/62年度，第232、236号。
\end{description}

% END com.tex

% BEGIN e149.tex
\clearpage
\chapter*{第149次报告勘误}
\addcontentsline{toc}{chapter}{第149次报告勘误}
\sid{e149}
\setcounter{footnote}{0}
\src{298}

\noindent\textbf{1. 凝聚代数层的对偶定理。}

\begin{center}
[布尔巴基讨论班，第9卷，1956/57年度，第149号，25页。]
\end{center}

\noindent\textbf{第14页评注。} --- 正如报告人在1958年国际数学家大会上的
报告中指出的，\footnote{Grothendieck (Alexander). --- The cohomology
theory of abstract algebraic varieties, \emph{Proceedings of the
International Congress of Mathematicians} [1958, Edinburgh], p. 103--118.
--- Cambridge, at the University Press, 1960.}这里提出的问题现已全部解决。

关于凝聚层对偶性的更多资料，读者可参阅\emph{前引文献}\footnotemark[1]
第112--115页、EGA III 第2部分（编写中）；在此期间，还可参阅 SGA 1962。
更系统的论述将收入 EGA 后续的一章（在原定纲要中为第IX章）。

% END e149.tex

% BEGIN e182.tex
\clearpage
\chapter*{第182次报告勘误}
\addcontentsline{toc}{chapter}{第182次报告勘误}
\sid{e182}
\setcounter{footnote}{1}
\src{298}

\noindent\textbf{2. 形式几何与代数几何。}

\begin{center}
[布尔巴基讨论班，第11卷，1958/59年度，第182号，28页。]
\end{center}

\noindent\textbf{第8页，定理6。} --- 假设“$f$ 拟射影”可以换成较弱的假设
“$f$ 分离”；这是由下述结果（参见 SGA VIII，6.2）所得：每个拟有限且分离的
态射 $f:X\longrightarrow Y$ 都是拟射影的。

\src{299}
\noindent\textbf{第12页及第14页评注1。} --- 须指出，J.-P. Serre 构造了
代数闭域 $k$ 上的一个三维非奇异射影簇，该域的特征为 $p>0$；它并非由一个
整局部环上的固有概形约化而来，而该环的剩余域为 $k$、分式域的特征为零。
\footnote{Serre (Jean-Pierre). --- Exemples de variétés projectives en
caractéristique $p$ non relevables en caractéristique zéro, \emph{Proc. Nat.
Acad. Sc. U.S.A.}, t. 47, 1961, p. 108--109.}据说 Mumford 已找到一个以
非奇异射影曲面为例的类似结果。

\noindent\textbf{第15页，评注2和3。} --- 报告人目前对这里所猜想的结果已不
那么乐观。不过，评注3末尾提到的射影空间结构变形问题，已由 Hironaka 肯定地
解决；阿贝尔簇的相应问题则由 Koizumi 解决。

\noindent\textbf{第16页，第3行。} --- 须指出，Mumford 最近构造了亏格为
$g$ 的曲线的模概形（参见 SHMI）。此外，定理10证明，模理论中的“$n$ 级
Jacobi 概形”是非奇异的（甚至在 $\mathbf Z$ 上光滑）。

\noindent\textbf{第23页，第12行，推论5的第1行。} --- 补入：“$X$ 或 $Y$
在 $k$ 上固有”。

第1至5节的实质包含在 EGA III 已出版的部分，第6和7节的实质包含在
SGA III。关于基本群的研究，见 SGA V、IX、X、XI；关于 Lefschetz 型定理
及许多悬而未决的问题，另见 SGA 1962（报告 X、XII、XIII）。只有驯分歧
覆盖理论（参见定理14）尚未正式写成。定理14的推论给出了对如下对象的完全
分类：代数闭域上代数曲线的 Galois 覆盖，其阶与该域的特征互素。该推论曾在
三个场合发挥关键作用：

1）在 Igusa 对任意特征下非奇异射影曲面的 Picard 不等式的证明中；

2）在对下述群的研究中（Ogg 与 Šafarevič 各自独立开展了这一研究）：
任意特征下一元函数域上定义的阿贝尔簇作用下主齐性丛所成的群；

3）在 Artin 最近对代数簇“Weil 上同调”的若干关键定理的证明中。

% END e182.tex

% BEGIN e190.tex
\clearpage
\chapter*{第190次报告勘误}
\addcontentsline{toc}{chapter}{第190次报告勘误}
\sid{e190}
\setcounter{footnote}{0}
\src{300}

\noindent\textbf{3. 下降技术与代数几何中的存在性定理 [TDTE I]。
I：一般理论。沿忠实平坦态射的下降。}

\begin{center}
[布尔巴基讨论班，第12卷，1959/60年度，第190号，29页。]
\end{center}

\noindent\textbf{第1页，第10行。} --- 现在看来，说下降技术“是代数几何中
大多数存在性定理的基础”似乎过甚。对于 TDTE I 至 VI 系列前两篇报告所处理
的非射影技术，这在很大程度上成立；但对于射影技术（TDTE IV、V、VI）则不
成立。

\noindent\textbf{第5页，第16行。} --- 无需假设 $\alpha$ 是一个
$\mathcal F$-下降态射。

\noindent\textbf{第20页，评注。} --- 拟有限、平展且满射的态射
$S'\longrightarrow S$，或者态射
$\operatorname{Spec}(\widehat A)\longrightarrow\operatorname{Spec}(A)$，
未必是严格下降态射；即使 $A$ 是代数闭域 $k$ 上一条代数曲线的局部环，且
$S=\operatorname{Spec}(A)$，情形仍然如此。确实可以找到固有光滑态射
$f:X\longrightarrow S$，使 $X$ 成为一个以 $S$ 为基、以椭圆曲线 $E$ 为
结构群的主丛，并且 $f':X'\longrightarrow S'$ 射影而 $f$ 不射影。因此，这同时
也是阿贝尔概形作用下非等平凡主齐性丛的一个例子。

\noindent\textbf{第26页，第9行。} --- 应读作“Chow--Lang”，而非
“Igusa--Lang”。

关于下降理论的各项细节，见 SGA VI、VII、VIII。

% END e190.tex

% BEGIN e195.tex
\clearpage
\chapter*{第195次报告勘误}
\addcontentsline{toc}{chapter}{第195次报告勘误}
\sid{e195}
\setcounter{footnote}{0}
\src{300}

\noindent\textbf{4. TDTE II：形式模理论中的存在性定理。}

\begin{center}
[布尔巴基讨论班，第12卷，1959/60年度，第195号，22页。]
\end{center}

\noindent\textbf{第8页。} --- 命题5.1中所写的公式只有在 $\Lambda$ 是域时
才正确；一般情形下，必须把 $\mathfrak m_\xi/\mathfrak m_\xi^2$ 换成这个
空间对 $\mathfrak n_\xi/\mathfrak n_\xi^2$ 的像所取的商，其中
$\mathfrak n$ 是 $\Lambda$ 的极大理想。此外，关于 $O$ 在 $\Lambda$ 上
光滑所给出的定义，只有在剩余域扩张 $k'/k$ 可分时才正确。一般情形参见
SGA III，1.1。

\noindent\textbf{第14页，评注。} --- 这里提出的问题在射影情形下已由
“Hilbert 概形”彻底解决（参见 TDTE IV）。Nagata 和 Hironaka 的例子反而
表明，若不作射影性假设，所考察的函子未必可表；即便只限于分类一个三维完备
非奇异簇的零维子簇，情形也是如此。这与这类簇的对称平方可能不存在这一事实
有关。

\src{301}
\noindent\textbf{第15页，第3项。} --- 更完整的研究见 TDTE V。

\noindent\textbf{第16页，命题3.1。} --- 把“若 $f$ 固有”改为
“若 $f$ 固有且可分”。

\noindent\textbf{第18、19页，评注，第2项。} --- 我最近证明，域上
阿贝尔簇的形式模概形确实在 Witt 环上光滑；换言之，若一个阿廷局部环是
另一环的商，则该阿廷局部环上的任一阿贝尔概形，都是由后一环上的某个阿贝尔
概形约化而来。证明只用到第182次报告第12页所引入的提升障碍类的变换性质。
还应回顾，亏格为 $g$ 的曲线以及极化阿贝尔概形的模概形已由 Mumford 构造
（参见 SHMT）。

\noindent\textbf{第19页，第5段。} --- 为定义 $U$ 所考察的截面，必须假设
它不仅在 $\mathfrak X$ 上，而且在每个 $X_n$ 上都是非零因子。

\noindent\textbf{第20页，第12、13行。} --- 把
\[
  G_x=G\bigl(\operatorname{Spec}(\widehat{\mathcal O}_{\mathfrak X,x})\bigr),
\]
改为
\[
  G_x=G\bigl(\operatorname{Spec}(\mathcal O_{\mathfrak X,x})\bigr)
  \subset
  G\bigl(\operatorname{Spec}(\widehat{\mathcal O}_{\mathfrak X,x})\bigr),
\]
并把“$\mathcal O'_{0x}$”改为“$\widehat{\mathcal O}'_{0x}$”。

% END e195.tex

% BEGIN e212.tex
\clearpage
\chapter*{第212次报告勘误}
\addcontentsline{toc}{chapter}{第212次报告勘误}
\sid{e212}
\setcounter{footnote}{2}
\src{301}

\noindent\textbf{5. TDTE III：商预概形。}

\begin{center}
[布尔巴基讨论班，第13卷，1960/61年度，第212号，20页。]
\end{center}

\noindent\textbf{第15页，断言 (i)。} --- 关于 $Y=X/R$ 在 $S$ 上拟射影
这一事实，目前看来只在 $R$ 来自一个等价关系时得到证明。

\noindent\textbf{第18、19页。} --- 正如我们在下一篇报告末尾所指出的，
这里所考察的猜想终究是错误的。评注8.1中指出的“肯定性事实”，似乎已由多位
作者同时证明（Nagata、Rosenlicht、Grothendieck，\ldots）。

还应指出，TDTE V 把这里发展的取商理论用于构造 Picard 概形；这一做法也可
由适当地使用 Hilbert 概形来替代（参见 SHMT）。正如第8节所说，这里所述
\src{302}
理论最重要的缺口，是缺少一个判定按非固有等价关系所取之商存在的判据，例如
由射影群的某些作用产生的等价关系。一个沿此方向的重要定理由 Mumford
得到\footnote{MUMFORD (David). --- An elementary theorem in geometric
invariant theory, \emph{Bull. Amer. Math. Soc.}, t. 67, 1961,
p. 483--486.}。关于其结果的一项改进以及对该理论的多种应用，见 SHMT。

% END e212.tex

% BEGIN e221.tex
\clearpage
\chapter*{第221次报告勘误}
\addcontentsline{toc}{chapter}{第221次报告勘误}
\sid{e221}
\setcounter{footnote}{3}
\src{302}

\noindent\textbf{6. TDTE IV：Hilbert 概形。}

\begin{center}
[布尔巴基讨论班，第13卷，1960/61年度，第221号，28页。]
\end{center}

\noindent\textbf{第6页，定理2.1，第3行。} --- 把“必要”改为“必要且充分”。

\noindent\textbf{第15页，第13行。} --- 把“$P\neq Q$”改为“$P<Q$”。

\noindent\textbf{第15行。} --- 把“若相反，$P=Q$”改为
“在相反情形下，将有 $P(n)\geq Q(n)$（当 $n$ 足够大时）”。

\noindent\textbf{第18行，公式 (*)。} --- 把“$<$”和“$>$”改为
“$\leq$”和“$\geq$”。

\noindent\textbf{倒数第5行。} --- 把“$r$”改为“$Z$”。

\noindent\textbf{第18页，评注3.9。} --- 应指出，Hartshorne 已研究代数闭域
上的 Hilbert 概形的连通分支。他证明 $\operatorname{Hilb}^{P}_{r}$ 均连通，
并确定了各对 $(r,P)$ 中哪些满足
$\operatorname{Hilb}^{P}_{r}\neq\varnothing$\footnote{HARTSHORNE (R.). --- Connectedness of the Hilbert
scheme (Thesis, Harvard, 1962).}。

\noindent\textbf{第23页，倒数第5行。} --- 把法文“une composante”改为
“d'une composante”（即把“一个分支”改为“一个分支的”）。

\noindent\textbf{第23、24页，评注5.5。} --- 关于 Zappa 的例子，应指出，
Mumford 刚刚构造了 $\mathbf P^3_{\mathbf C}$ 的 Hilbert 概形的一个不可约
分支；其一般点表示次数为14、亏格为24的非奇异曲线，而该分支在其一般点处
非约化。沿所得曲线中的一条作爆破，他还得到一个 $\mathbf C$ 上的三维正则
射影概形；其形式模概形在一般点处非约化，或者等价地，其在 $\mathbf C$ 上
解析几何意义下的局部模簇（参见 Cartan 讨论班，第13卷，1960/61年度）在
每一点处均非约化。

% END e221.tex

% BEGIN e232.tex
\clearpage
\chapter*{第232次报告勘误}
\addcontentsline{toc}{chapter}{第232次报告勘误}
\sid{e232}
\src{303}

\noindent\textbf{7. TDTE V：Picard 概形。存在性定理。}

\begin{center}
[布尔巴基讨论班，第14卷，1961/62年度，第232号，19页。]
\end{center}

\noindent\textbf{第7页，评注3.3。} --- 这里提出的问题已由 Mumford 作出
肯定解答（见下一篇报告的评注）。

\noindent\textbf{第13页，评注5.2。} --- 正如我在下一篇报告开头所指出的，
这里提出的存在性猜想是错误的；不过 Mumford 用他的方法证明了一个稍弱的
定理。

\noindent\textbf{第14页，第7行。} --- 在“完备”之后加上“具有代数闭剩余
域”。刚才提到的 Mumford 反例还表明，这一限制不可缺少。

\noindent\textbf{第17页，评注6.6。} --- 正如我们在下一篇报告开头所指出的，
这里提出的问题刚刚由 Murre 作出肯定解答。

\noindent\textbf{第18页，评注6.7。} --- 这里提出的问题有肯定解答（参见下一
篇报告评注的最后一段）。

\noindent\textbf{第18页，评注6.8，第4行。} --- 把“正则”改为“在 $k$ 上
光滑”。

% END e232.tex

% BEGIN e236.tex
\clearpage
\chapter*{第236次报告勘误与补遗}
\addcontentsline{toc}{chapter}{第236次报告勘误与补遗}
\sid{e236}
\src{303}

\noindent\textbf{8. TDTE VI：Picard 概形。一般性质。}

\begin{center}
[布尔巴基讨论班，第14卷，1961/62年度，第236号，23页。]
\end{center}

\noindent\textbf{第12页，第16行。} --- 把“仿射直线”改为“去掉原点的仿射
直线”。

\noindent\textbf{第18行。} --- 把“$X[t]$”改为“$X[t,t^{-1}]$”。

\noindent\textbf{第16页，命题3.1，第6行。} --- 把“然而，假设……”改为
“然而，由假设……”。

\noindent\textbf{第21页，第4节：Picard 概形的有限性定理。} --- 自本篇报告
写成以来，本节所提出的这一类有限性问题已几乎全部解决。下面指出目前在这一
方向上已知的主要事实。为简化陈述，我们默认其中出现的所有 Picard 预概形均
存在，尽管对这些陈述作显然的修改便可去掉任何存在性假设。以下用 $S$ 表示
Noether 概形，用 $X,Y$ 表示 $S$ 上的固有概形。

\medskip
\noindent(i) 设 $f:X\longrightarrow Y$ 是满射态射，则
\[
f^*: \operatorname{Pic}_{Y/S}\longrightarrow\operatorname{Pic}_{X/S}
\]
是有限型态射。

\src{304}
\noindent(ii) 假设 $Y$ 在 $S$ 上射影，并配有一个相对于 $S$ 丰沛的可逆模；
设 $X$ 是该模某一任意截面的零点概形，$f:X\longrightarrow Y$ 是典范浸入。
再假设 $Y/S$ 各纤维的不可约分支维数均为 $\geq3$。则
\[
f^*: \operatorname{Pic}_{Y/S}\longrightarrow\operatorname{Pic}_{X/S}
\]
是有限型态射。

\noindent(iii) 假设 $X$ 在 $S$ 上射影，且其所有几何纤维均为维数 $n$ 的
整概形。设 $\mathcal O_X(1)$ 是 $X$ 上的丰沛可逆模，可用来定义 Hilbert
多项式。设 $M$ 是 $\operatorname{Pic}_{X/S}$ 的一个子集；则 $M$ 拟紧，
当且仅当 $M$ 中各元素的 Hilbert 多项式
\[
a_0x^n+a_1x^{n-1}+a_2x^{n-2}+\cdots
\]
中的系数 $a_1$ 和 $a_2$ 保持有界，其中要略去单项式次数为负的系数；因此当 $n=1$ 时只出现 $a_1$，当 $n=0$ 时二者都不出现。

\noindent(iv) 对每个整数 $n\neq0$，$n$ 次幂同态在群预概形
$\operatorname{Pic}_{X/S}$ 上都是有限型态射。

注意，(i) 和 (ii) 还意味着：在这两个陈述各自的假设下，一个子集 $M$ 在
$\operatorname{Pic}_{Y/S}$ 中拟紧，当且仅当它在
$\operatorname{Pic}_{X/S}$ 中的像拟紧。由此可知，一个可逆模
$\mathcal L$ 若定义在 $Y$ 上，则它与零 $\tau$-等价，当且仅当它在 $X$ 上
的逆像具有这一性质；
换言之，$\operatorname{Pic}_{Y/S}^\tau$ 是
$\operatorname{Pic}_{X/S}^\tau$ 的逆像。特别地，要证明前一个预概形在
$S$ 上为有限型，只须对后一个证明这一点，因为
$\operatorname{Pic}_{Y/S}\longrightarrow
\operatorname{Pic}_{X/S}$ 是有限型态射。于是利用 (i)、Chow 引理和 (iii)，
得到：

\medskip
\noindent(v) $\operatorname{Pic}_{X/S}^\tau$ 在 $S$ 上为有限型。

一般地，把有限态射情形的 (i) 与 (ii) 结合起来，可把大多数有限性问题归约到
$X/S$ 的几何纤维均为维数 $\leq2$ 的正规整概形的情形；往往还可把 (i) 用于
不必有限的满射态射，再利用代数曲面的奇点消解（任意特征下由 Abhyankar
证明），将问题归约到 $X/S$ 甚至光滑、从而几何纤维为二维非奇异概形的情形。
结合 (v) 以及控制非奇异射影曲面的 Néron--Severi 群之秩的 Picard--Igusa
不等式，这例如可证明 Néron 有限性定理的如下推广：

\medskip
\noindent(vi) 设 $X/S$ 在 $S$ 上固有，除此之外任意。则 $X_i/k_i$ 是
$X/S$ 的几何纤维，其 Néron--Severi 群
\[
\operatorname{Pic}_{X_i/k_i}/\operatorname{Pic}_{X_i/k_i}^\circ
\]

\src{305}
均为有限型，而且它们的秩以及挠子群的阶均保持有界。

同样归约到非奇异曲面情形，并利用该情形下的已知定理（即 Néron 有限性定理，
以及 Néron--Severi 群上的相交形式非退化这一事实），可推出：

\medskip
\noindent(vii) 设 $X$ 是代数闭域上的固有概形。则存在有限条闭整曲线
$C_i$（$1\leq i\leq r$）位于 $X$ 中，具有如下性质：对任一子集 $M$
（取自 $\operatorname{Pic}_{X/k}$），$M$ 拟紧，当且仅当整数
\[
\deg_{C_i'}\mathcal L\qquad(\mathcal L\in M)
\]
保持有界（其中 $C_i'$ 表示 $C_i$ 的正规化）。

这里可取 $r$ 为 Néron--Severi 群的秩。一旦知道该群为有限型，(vii) 就归结为
如下事实：由曲线 $C$（位于 $X$ 中）在 Néron--Severi 群上定义的线性形式，
只在该群
的挠元素上同时为零。当 $X$ 射影且非奇异时，这一结果以及 (v) 均由
Matsusaka 得到。利用 (vii)，容易得到与零 $\tau$-等价的可逆模在 $X$ 上的
如下刻画：

\medskip
\noindent(viii) 设 $X/k$ 是域上的固有概形，$\mathcal L$ 是 $X$ 上的可逆
模。下列条件等价：

\noindent a. $\mathcal L$ 与零 $\tau$-等价；

\noindent b. 对任意凝聚模 $F$，都有
\[
\chi(F\otimes\mathcal L)=\chi(F),
\]
其中 $\chi$ 表示 Euler--Poincaré 特征；

\noindent b'. 与 (b) 相同，但取 $F=\mathcal O_Y$，其中 $Y$ 是 $X$ 中的
一维闭整子概形；

\noindent c. 对每个这样的 $Y$，以 $Y'$ 表示其正规化曲线，则
\[
\deg\mathcal L_{Y'}=0.
\]

若 $X/k$ 射影，并配有丰沛可逆模 $\mathcal O_X(1)$，则上述条件还等价于：

\noindent d. 对每个整数 $m$，$\mathcal L^{\otimes m}(1)$ 均丰沛。

\noindent e.（若 $X$ 为整概形。）对每一对整数 $m,n$，都有
\[
\chi(\mathcal L^{\otimes m}(n))=\chi(\mathcal O_X(n))
\]

\src{306}
（即在 (b) 中取 $F=\mathcal L^{\otimes m}(n)$）。

为说明最后一个条件的充分性，应注意它意味着各
$\mathcal L^{\otimes m}$ 的 Hilbert 多项式全都相等；因此根据 Mumford
判据 (iii)，各 $\mathcal L^{\otimes m}$ 留在
$\operatorname{Pic}_{X/k}$ 的一个拟紧子集中，也就是说 (a) 成立。条件
(b)、(b')、(c) 应看作通常在非奇异射影簇上定义的数值等价概念在任意固有
概形上的变体。对于非奇异射影簇，(a) 与 (c) 的等价性当然早已为人所知
（Matsusaka）。

判据 (e) 还推出如下结果：

\medskip
\noindent(ix) 设 $f:X\longrightarrow S$ 是射影平坦态射，且其几何纤维均为
整概形。则 $\operatorname{Pic}_{X/S}^\tau$ 在
$\operatorname{Pic}_{X/S}$ 中既开又闭。

下面只对关键结果 (i)、(ii)、(iii) 的证明作几点说明。（结果 (iv) 与其他结果
稍有不同，它只需对有限、满射且纯不可分的态射使用 (i) 即可证明；更确切地
说，对 Frobenius 态射使用 (i)。）证明 (i) 时，本质上使用非平坦下降的思想
（尤其参见 TDTE I，第9页）。由于这里只追求有限性结果，缺少下降数据的有效
性判据并无妨碍。另一方面，Mumford 最近证明了 (iii) 的一个稍弱形式，即涉及
Hilbert 多项式所有系数的判据。他的论证极其简单，其思路来自 Nakai 丰沛性
判据的证明（Nakai 为非奇异曲面陈述了这一判据，Mumford 将其推广到任意射影
态射）。不过在我看来，这一论证只有在对 $X/S$ 的纤维再加上一项轻微限制时
才成立，即 Serre 条件 $(S_2)$；若几何纤维正规，则该条件满足。随后在 (ii)
的证明中使用这一受限判据：(i) 确实允许归约到 $Y/S$ 平坦且几何纤维均为正规
整概形的情形；再用 Mumford 判据，很容易归约到 $X/S$ 也满足同样条件的情形。
于是由维数假设，$Y$ 和 $X$ 的几何纤维在其闭点处的深度均为 $\geq2$，因而
可以应用 SGA 1962（第XII篇报告）所给形式的“等价判据”，完成 (ii) 的证明。
一旦得到 (i) 和 (ii)，便不难在 Mumford 判据中去掉对纤维的任何正规性假设，
并证明 (iii) 中陈述的更强形式。

\src{307}
最后还应指出，(i) 和 (ii) 的证明也表明：当 $S$ 是域 $k$ 的谱时，态射
\[
\operatorname{Pic}_{Y/k}\longrightarrow\operatorname{Pic}_{X/k}
\]
是仿射的（而不只是有限型的）。

% FGA Canon R30: scan restoration and explicit editorial apparatus.
\par\smallskip
\begingroup\footnotesize
\noindent\textit{FGA-CANON-DISP-0059（C125）。印刷本要求系数 $a_1$ 和 $a_2$ 保持有界，未说明 $n<2$ 的情形。现行读法仅明确低维时不存在的系数的约定，不宣称印刷定理为假。}\par
\endgroup

% END e236.tex

\end{document}
